% 高一67_华师大一附中_10月月考.tex
%   —— 高一上数学“2026-2027 年华师大一附中高一上 10 月月考”（讲评版）
% 试卷版：xelatex 高一67_华师大一附中_10月月考.tex
% 答案版：xelatex "\def\isanswer{1}% 高一67_华师大一附中_10月月考.tex
%   —— 高一上数学“2026-2027 年华师大一附中高一上 10 月月考”（讲评版）
% 试卷版：xelatex 高一67_华师大一附中_10月月考.tex
% 答案版：xelatex "\def\isanswer{1}% 高一67_华师大一附中_10月月考.tex
%   —— 高一上数学“2026-2027 年华师大一附中高一上 10 月月考”（讲评版）
% 试卷版：xelatex 高一67_华师大一附中_10月月考.tex
% 答案版：xelatex "\def\isanswer{1}% 高一67_华师大一附中_10月月考.tex
%   —— 高一上数学“2026-2027 年华师大一附中高一上 10 月月考”（讲评版）
% 试卷版：xelatex 高一67_华师大一附中_10月月考.tex
% 答案版：xelatex "\def\isanswer{1}\input{高一67_华师大一附中_10月月考.tex}"
% 紧凑版：xelatex "\def\iscompact{1}\input{高一67_华师大一附中_10月月考.tex}"
%
% 原始材料：微信公众号“魔都数学加油站”文章，作者破晓，连载“27学年高一67”；
% 原文标题“27学年高一67——2026-2027年上海市华师大一附中高一上10月月考”；
% 原文链接：https://mp.weixin.qq.com/s/Kdhx07TVjLoGTcIlnRzgcQ。
% 页面标注发布时间 2026-10-18 20:33（上海）。
% 本次转录只据 `原图/高一67_华师大一附中_10月月考/` 的 13 张 PNG 逐页读图。
% 正文为 13 张 PNG 页（`原图/高一67_华师大一附中_10月月考/01..13.png`）：
% 第 01、04、13 页 2560×3620，第 02、03、05—12 页 4762×6735
% （**同一篇各页分辨率不同**），均按页序读取。13 页全为试卷正文页（卷首标题在第 1 页页顶，
% 卷末解答题第 21 题解析止于第 13 页），无封面页、目录页、二维码或宣传页。
% 粉色斜排“魔都数学加油站”为卷面水印，与公众号页脚一并未录入。
%
% 卷首标题：照卷面印作“2026--2027 年华师大一附中高一上 10 月月考”（卷面为黑体加粗一行），
% 年份区间按全稿惯例排 en dash。原卷只有一行标题、无副标题。
%
% 篇幅与题号：原卷自第 3 题起编号（第 1、2 题不在本卷内）。填空题 3—12（10 题）、
% 选择题 13—16（4 题）、解答题 17—21（5 题），共 19 题。本稿沿用原节与原题号，
% 只改排为全稿统一的大节样式（原卷小节标题为“一、填空题”“二、选择题”“三、解答题”）。
%
% 讲评版：原卷每题下直接印【解析】。**全卷 19 题中 18 题都印【解析】（无一处只印
% 【答案】）**，按全稿口径这类题不另提 \ans{}；各题解析末句的结论（第 13、14、16 题
% 的“故选 B／D／B”、各填空与解答题的“故/综上…”）照录在解析内。
% **另有 1 题答案印在题干括号内**——第 15 题（“(  D  )”，且该题无【答案】/【解析】段），
% 本稿按 高一37 第 13 题、高一63 第 14 题成例把括号留空、答案另提为 \ans{$D$}。
%
% 副标题：原卷未印副标题。本稿按“知识模块”口径标作“集合与常用逻辑用语、不等式”：
% 19 题中集合与常用逻辑用语 12 题（第 3、4、5、7、10、11、12、13、15、16、18、21 题：
% 补集、并集、子集、充分条件、容斥计数、子集计数、新定义二元运算、条件判断、
% 反证法假设、新定义生成集、集合关系与二次方程、新定义 X_k 与 X^2），
% 不等式 7 题（第 6、8、9、14、17、19、20 题：存在量词不等式、解集反求系数、
% 解集为空、不等式性质、解不等式（组）、二次不等式与基本不等式、二次方程根的分布），
% 两类均超过三分之一，故并标。
%
% 与原卷的排版差异：
%   ① 原卷每题下直接印【解析】，本稿按全稿惯例将解析置于答案版，试卷版保持干净；
%   ② 原卷未印副标题，本稿按“知识模块”口径补出；
%   ③ 原卷填空横线约四个字宽，本稿按全稿惯例统一排 \blankline[4.4em]；
%   ④ 原卷实数集排斜体 $R$（第 14 题题干、第 17 题 (1) 题干、第 18 题解析、第 20 题题干），
%      整数集排斜体 $Z$（第 20 题 (3) 题干与解析、第 21 题题干），正整数集排粗体
%      $N^*$（第 11 题解析、第 19 题解析），本稿按全稿惯例统一改用 $\R$、$\Z$、$\N^*$
%      （记号归一，粗体／斜体不作区分）；空集原卷作 $\varnothing$，照录；
%   ⑤ 原卷在数学式内多处用半角逗号（如第 8 题 $(1,2)$ 与各 cases、第 11 题解析
%      $P_n=\{1,2,\cdots,n\}$,$n\in N^*$、第 12 题解析 $u=(\cdots)$,$v=(\cdots)$ 与
%      “2020,个 1”、第 19 题解析 $\{x|40<x<80,x\in N^*\}$），均照录未强改；
%      描述法集合的竖线照原卷排 \mid；句末句点按全稿惯例排西文句点；
%   ⑥ 原卷补集记号为上划线 $\bar A$（第 3 题题干与解析、第 11 题题干），照录
%      （同 高一64/65 的 $\bar B$ 成例）；
%   ⑦ 第 12 题新定义的运算记号 $*$ 照录（$\alpha*\beta$、$u*u$、$u*v$）；
%   ⑧ 第 17—21 题的子问原卷各自另起一行，本稿按全稿惯例用 enumerate 的 (1)(2)(3)；
%      解析中的子问标号一律用 \stepm{(1)}{…}；第 5 题解析的①②、第 21 题 (3) 解析的①②
%      按全稿惯例作 \stepm{①}{…}、\stepm{②}{…}（第 21 题 (3) 首行原卷“(3)”与“①”连排，
%      本稿合并记作 \stepm{(3)①}{…}）；
%   ⑨ 选择题作答括号原卷排半角“(　　)”（第 13、14、16 题括号内空白、第 15 题括号内印 D），
%      本稿按全稿惯例排 \mbox{（\quad）}；第 15 题括号内的答案移入 \ans{}，见上；
%   ⑩ 选项行排照原卷：第 13、15 题两行两条，第 14、16 题一行四条；
%   ⑪ 每道解答题原卷未留作答空白，本稿按全稿惯例加 \ansspace（末题用 \ansspace*）；
%      第 20、21 题题干长，各用一处 \ansneed 让题干与作答区不拆页（\ansneed 只作用于试卷版）。
%
% 原卷笔误/存疑（均照录未改，正文对应处另有 % 注）：
%   ① 第 3 题解析末句印作“则 $\bar A=(-\infty,0)\cup[2,3)$.”：按题干
%      $U=\{x\mid 1/x\leqslant 1/2\}$（即 $x<0$ 或 $x\geqslant2$）与 $A=[3,+\infty)$，
%      $\bar A$ 应为 $(-\infty,0)\cup\{2\}$；且与本解析首句自求之 $U$ 相矛盾（$3\notin U$），
%      疑为笔误；照录未改。
%   ② 第 12 题题干 $\alpha*\beta$ 定义式括号不配对：第四组“$+\cdots+x_{2020}+y_{2020}
%      -x_{2020}y_{2020}$”缺左括号、末尾连排两个右括号“))”；照录未改。
%   ③ 第 12 题解析“所以…中有 2020,个 1”一行，“$\cdots$”后脱逗号（作
%      “$x_2,\cdots x_{2020}$”，第 12 行“而…”同式样），且“2020”后衍一半角逗号
%      （作“2020,个 1”）；照录未改。
%   ④ 第 16 题解析乘积链印作“$a_1a_2<a_1a_3<a_1a_4<a_1a_5<a_2a_3<a_3a_3<a_{14}a_5$”：
%      “$a_3a_3$”同足标自乘、与单调链不合，“$a_{14}a_5$”下标 14 与 $a_1$—$a_5$ 记号
%      不合（疑为 $a_2a_4$、$a_4a_5$ 之误）；照录未改。
%   ⑤ 第 5 题题干作“$Q\subset P$”（真子集记号），解析两处均作“$Q\subseteq P$”，
%      同一题前后记号不一；照录（本稿分别排 \subset、\subseteq）。
%   ⑥ 第 17 题 (1) 解析“当 $a=1$ 时，原不等式为 $(x+1)^2<0$ 此时…”与
%      “当 $0<a<1$ 时，$-\frac1a<-1$ 此时…”两处“此时”前均脱标点；照录未改。
%   ⑦ 第 20 题 (2) 题干“若 $a=1$，$b=2k-1$，$c=k^2$，若方程 $f(x)=0$ 有两个大于 1 的根”
%      两个“若”连用，疑衍其一；照录未改。
%
% 插图：无（全卷 13 页无任何几何图形、无“如图/右图/图中”字样）。
%
% 行数核对：逐题数 \step/\stepm 条数与原卷【解析】印刷行数，18 题逐一相符
% （第 15 题无解析段，只提 \ans{}），合计 169 条 / 169 行
% （Q3—Q12：2、3、4、3、1、3、5、3、13、15；Q13—Q16：2、3、0、25；
% Q17—Q21：19、16、9、20、23）。
%
\documentclass[a4paper,11pt]{article}
\usepackage{common}

\begin{document}

\examtitlest[3]{2026--2027 年华师大一附中高一上 10 月月考}{集合与常用逻辑用语、不等式}

\bigsec{一、填空题}

\begin{enumerate}
\setcounter{enumi}{2}

\item 设全集 $U=\left\{x\mid\dfrac1x\leqslant\dfrac12\right\}$，$A=[3,+\infty)$，则 $\bar A=$\blankline[4.4em].

\ifanswer
\solbegin
\step{$\dfrac1x\leqslant\dfrac12$，解得 $x<0$ 或 $x\geqslant2$，故 $U=(-\infty,0)\cup[2,+\infty)$，}
% 原卷末句作 $\bar A=(-\infty,0)\cup[2,3)$，按题干 $U$ 应为 $(-\infty,0)\cup\{2\}$，与本解析首句矛盾，疑笔误；照录（见文件头“原卷笔误/存疑”①）
\step{$A=[3,+\infty)$，则 $\bar A=(-\infty,0)\cup[2,3)$.}
\solend
\fi

\item 已知集合 $M=\{y\mid y=x^2-2x-1\}$，$N=\{y\mid y=x^2-4x+1\}$，则 $M\cup N=$\blankline[4.4em].

\ifanswer
\solbegin
\step{$M=\{y\mid y=x^2-2x-1\}=\{y\mid y=(x-1)^2-2\}=\{y\mid y\geqslant-2\}$，}
\step{$N=\{y\mid y=x^2-4x+1\}=\{y\mid y=(x-2)^2-3\}=\{y\mid y\geqslant-3\}$，}
\step{则 $M\cup N=\{y\mid y\geqslant-3\}$.}
\solend
\fi

\item 已知 $P=\{x\mid x^2+x-6=0\}$，$Q=\{x\mid mx-1=0\}$，且 $Q\subset P$，则实数 $m$ 的值\blankline[4.4em].

\ifanswer
\solbegin
\step{$P=\{2,-3\}$，$Q=\{x\mid mx=1\}$；}
\stepm{①}{$m=0$ 时，$Q=\varnothing$，满足 $Q\subseteq P$；}
\stepm{②}{$m\neq0$ 时，要使 $Q\subseteq P$，则 $\dfrac1m=2$ 或 $-3$，所以 $m=\dfrac12$ 或 $-\dfrac13$；}
\step{所以 $m=0$ 或 $\dfrac12$ 或 $-\dfrac13$.}
\solend
\fi

\item 已知“存在正实数 $x$，使得关于 $x$ 的不等式 $ax-1>0$ 成立”为真命题，则实数 $a$ 的取值范围是\blankline[4.4em].

\ifanswer
\solbegin
\step{存在正实数 $x$，使得关于 $x$ 的不等式 $ax-1>0$ 成立为真命题，}
\step{即存在正实数 $x$，使得 $a>\dfrac1x$，因为 $\dfrac1x>0$，所以 $a>0$，}
\step{即 $a$ 的范围为 $(0,+\infty)$.}
\solend
\fi

\item 设 $\alpha:3\leqslant x\leqslant5$，$\beta:2m-1\leqslant x<m+4$，若 $\alpha$ 是 $\beta$ 的充分条件，则实数 $m$ 的取值范围是\blankline[4.4em].

\ifanswer
\solbegin
\step{若 $\alpha$ 是 $\beta$ 的充分条件，则 $\begin{cases}2m-1\leqslant3\\ m+4>5\end{cases}$，解得 $1<m\leqslant2$.}
\solend
\fi

\item 对于实数 $a$、$b$、$c$，已知关于 $x$ 的不等式 $ax^2+bx+c>0$ 的解集是 $(1,2)$，则关于 $x$ 的不等式 $cx^2-bx+a\leqslant0$ 的解集是\blankline[4.4em].

\ifanswer
\solbegin
\step{关于 $x$ 的不等式 $ax^2+bx+c>0$ 的解集是 $(1,2)$，则 $\begin{cases}a<0\\ -\dfrac ba=3\\ \dfrac ca=2\end{cases}$，则 $\begin{cases}b=-3a\\ c=2a\end{cases}$，}
\step{则关于 $x$ 的不等式 $cx^2-bx+a=2ax^2+3ax+a\leqslant0$，则 $2x^2+3x+1\geqslant0$，}
\step{则 $x\leqslant-1$ 或 $x\geqslant-\dfrac12$，则解集为 $(-\infty,-1]\cup\left[-\dfrac12,+\infty\right)$.}
\solend
\fi

\item 已知关于 $x$ 的不等式 $kx^2+kx+k+2\leqslant0$ 的解集为 $\varnothing$，则实数 $k$ 的取值范围为\blankline[4.4em].

\ifanswer
\solbegin
\step{当 $k=0$ 时，不等式化为 $2\leqslant0$，解集为 $\varnothing$，符合题意；}
\step{当 $k\neq0$ 时，若关于 $x$ 的不等式 $kx^2+kx+k+2\leqslant0$ 的解集为 $\varnothing$，}
\step{则 $\begin{cases}k>0\\ \Delta=k^2-4k(k+2)<0\end{cases}$，化简得 $\begin{cases}k>0\\ -3k^2-8k<0\end{cases}$，即 $\begin{cases}k>0\\ k(3k+8)>0\end{cases}$，}
\step{解得 $k>0$.}
\step{综上所述，$k\geqslant0$，即实数 $k$ 的取值范围为 $[0,+\infty)$.}
\solend
\fi

\item 对班级若干名学生调查对 $A$、$B$ 两事件的态度，有如下结果：赞成 $A$ 的人数是全体的五分之三，其余的不赞成；赞成 $B$ 的比赞成 $A$ 多 3 人，其余的不赞成. 已知对 $A$、$B$ 都赞成的学生有 18 人，对 $A$、$B$ 都不赞成的学生数比对 $A$、$B$ 都赞成的学生数的三分之一多 1 人，则班级总人数是\blankline[4.4em]人.

\ifanswer
\solbegin
\step{设班级总人数是 $x$ 人，则赞成 $A$ 的人数是 $\dfrac35x$，赞成 $B$ 的人数是 $\dfrac35x+3$，}
\step{$A$、$B$ 都不赞成的人数为 $18\times\dfrac13+1=7$ 人，}
\step{由容斥原理得 $x=\left(\dfrac35x+\dfrac35x+3-18\right)+7$，解得 $x=40$，即班级总人数是 40 人.}
\solend
\fi

\item 设全集 $U=\{1,2,3,4,5,6,7,8,9,10,11,12\}$，若集合 $A$ 满足：① $A\subseteq U$；②若 $x\in A$，必有 $2x\notin A$；③若 $x\in\bar A$，必有 $2x\notin\bar A$，则这样的集合 $A$ 有\blankline[4.4em]个.

\ifanswer
\solbegin
\step{设集合 $P_n=\{1,2,\cdots,n\}$,$n\in\N^*$，}
\step{任取 $x\in P_n$，将 $x$ 除以 2，若商仍为偶数，再除以 2，$\cdots$，}
\step{经过 $k$ 次以后，商必为奇数，此时记商为 $m$，}
\step{于是 $x=m\cdot2^k$，其中 $m$ 为奇数，$k\in\N^*$.}
\step{由条件得，若 $m\in A$，则 $x\in A$ 等价于 $k$ 为偶数；}
\step{若 $m\notin A$，则 $x\in A$ 等价于 $k$ 为奇数；}
\step{于是 $x$ 是否属于 $A$ 由 $m$ 是否属于 $A$ 确定；}
\step{设 $Q_n$ 是 $P_n$ 中所有奇数的集合，}
\step{设 $f(n)$ 为集合 $P_n$ 中满足题干条件的集合 $A$ 的个数，}
\step{因此 $f(n)$ 等于 $Q_n$ 的子集个数.}
\step{当 $n$ 为偶数（或奇数）时，$P_n$ 中奇数的个数是 $\dfrac n2$（或 $\dfrac{n+1}2$），}
\step{所以 $f(n)=\begin{cases}2^{\frac n2},n\text{为偶数}\\[2pt] 2^{\frac{n+1}2},n\text{为奇数}\end{cases}$，所以 $P_{12}=\{1,2,3,4,5,6,7,8,9,10,11,12\}$，}
\step{即同时满足三个条件的集合 $A$ 的个数为 $2^{\frac{12}2}=2^6=64$.}
\solend
\fi

\item 设集合 $\Omega=\{(x_1,x_2,\cdots,x_{2020})\mid x_i\in\{0,1\},i=1,2,\cdots,2020\}$，对于 $\Omega$ 中的任意两个元素 $\alpha=(x_1,x_2,\cdots,x_{2020})$，$\beta=(y_1,y_2,\cdots,y_{2020})$，记
% 原卷此式第四组缺左括号、末尾连排两个右括号“))”，括号不配对；照录（见文件头“原卷笔误/存疑”②）
$\alpha*\beta=(x_1+y_1-x_1y_1)+(x_2+y_2-x_2y_2)+\cdots+x_{2020}+y_{2020}-x_{2020}y_{2020}))$，设 $u$、$v\in\Omega$，若 $u*u+v*v=2020$，则 $u*v$ 的最大值和最小值之和为\blankline[4.4em].

\ifanswer
\solbegin
\step{设 $u=(x_1,x_2,\cdots,x_{2020})$,$v=(y_1,y_2,\cdots,y_{2020})$，}
\step{因为 $x_i\in\{0,1\}$，所以 $x_i(x_i-1)=0$，}
\step{所以 $u*u=(x_1+x_1-x_1^2)+(x_2+x_2-x_2^2)+\cdots+(x_{2020}+x_{2020}-x_{2020}^2)$}
\step{$=x_1+x_2+\cdots+x_{2020}$，}
\step{$v*v=(y_1+y_1-y_1^2)+(y_2+y_2-y_2^2)+\cdots+(y_{2020}+y_{2020}-y_{2020}^2)$}
\step{$=y_1+y_2+\cdots+y_{2020}$，}
\step{由 $u*u+v*v=2020$，得 $x_1+x_2+\cdots+x_{2020}+y_1+y_2+\cdots+y_{2020}=2020$，}
% 原卷此行“$\cdots$”后脱逗号、 “2020”后衍一半角逗号（“2020,个 1”）；照录（见文件头“原卷笔误/存疑”③）
\step{所以 $x_1,x_2,\cdots x_{2020},y_1,y_2,\cdots y_{2020}$ 中有 2020,个 1，2020 个 0，}
\step{$u*v=(x_1+y_1-x_1y_1)+(x_2+y_2-x_2y_2)+\cdots+(x_{2020}+y_{2020}-x_{2020}y_{2020})$}
\step{$=2020-(x_1y_1+x_2y_2+\cdots+x_{2020}y_{2020})$，}
\step{注意到只有 $x_i=y_i=1$ 时，$x_iy_i=1$，否则 $x_iy_i=0$，}
% 此行“$\cdots$”后同样脱逗号，与上一处同式样；照录（见文件头“原卷笔误/存疑”③）
\step{而 $x_1,x_2,\cdots x_{2020},y_1,y_2,\cdots y_{2020}$ 中有 2020 个 1，2020 个 0，}
\step{所以满足 $x_iy_i=1$ 的最多有 1010 个，所以 $u*v\geqslant2020-1010=1010$，}
\step{满足 $x_iy_i=0$ 的最多有 2020 个，所以 $u*v\leqslant2020-0=2020$，}
\step{则 $u*v$ 的最大值和最小值之和为 $2020+1010=3030$.}
\solend
\fi

\end{enumerate}

\bigsec{二、选择题}

\begin{enumerate}
\setcounter{enumi}{12}

\item 实数 $x$ 与 $y$，“$x+y>2$”的\mbox{（\quad）}条件是“$x>1$ 或 $y>1$”.

\keepwithprev
A. 充分非必要\qquad B. 必要非充分

C. 充要\qquad D. 既非充分又非必要

\ifanswer
\solbegin
\step{当 $x>1$ 或 $y>1$ 时，$x+y>2$ 不一定成立，充分性不成立；}
\step{当 $x+y>2$ 时，$x>1$ 或 $y>1$ 一定成立，必要性成立. 故选 $B$.}
\solend
\fi

\item 设 $a,b\in\R$，且 $a<0$，$-1<b<0$，那么下列不等式成立的是\mbox{（\quad）}

\keepwithprev
A. $a>ab>ab^2$\qquad B. $ab^2>ab>a$\qquad C. $ab>a>ab^2$\qquad D. $ab>ab^2>a$

\ifanswer
\solbegin
\step{因为 $-1<b<0$，$a<0$，所以 $ab>0$，$0<b^2<1$.}
\step{所以 $ab-ab^2=ab(1-b)>0$，$ab^2-a=a(b^2-1)>0$，}
\step{所以 $ab>ab^2>a$. 故选 $D$.}
\solend
\fi

% 第 15 题答案原卷印在题干括号内（“(  D  )”），且无【答案】/【解析】段；
% 按 高一37 第 13 题、高一63 第 14 题成例括号留空、答案另提 \ans{$D$}
\item 用反证法证明：“平面四边形中至少有一个内角不超过 90°”时，下列“假设”中正确的是\mbox{（\quad）}

\keepwithprev
A. 假设有两个内角超过 90°\qquad B. 假设有三个内角超过 90°

C. 假设至多有两个内角超过 90°\qquad D. 假设四个内角都超过 90°

\ifanswer
\ans{$D$}
\fi

\item 设 $A$ 为非空实数集，集合 $B=\{xy\mid x\in A,y\in A,x\neq y\}$，则对于以下两个结论说法正确的是\mbox{（\quad）}

结论①：若 $A$ 是五元集，则 $B$ 中的元素最多 10 个；

结论②：存在四元集 $A$，使得 $B=\{2,4,5,8,10,16\}$.

\keepwithprev
A. ①对②对\qquad B. ①对②错\qquad C. ①错②对\qquad D. ①错②错

\ifanswer
\solbegin
\step{对于①，想让元素个数达到最大值 10，则须避免产生符号不同的乘积，}
\step{否则正负分开计数容易导致总数受限，且容易产生相等值，}
\step{也最好避免 0 的出现，}
\step{故设 $A=\{a_1,a_2,a_3,a_4,a_5\}$，$0<a_1<a_2<a_3<a_4<a_5$.}
% 原卷此链“$a_3a_3$”同足标自乘、“$a_{14}a_5$”下标 14 均与文意不合，疑为 $a_2a_4$、$a_4a_5$ 之误；照录（见文件头“原卷笔误/存疑”④）
\step{因为 $a_1a_2<a_1a_3<a_1a_4<a_1a_5<a_2a_3<a_3a_3<a_{14}a_5$，$a_2a_3<a_2a_4<a_3a_4$，}
\step{所以 $B$ 中元素个数小于等于 10 个，}
\step{如设 $A=\{1,2,3,5,7\}$，则 $B=\{2,3,5,6,7,10,14,15,21,35\}$，}
\step{所以集合 $B$ 中元素个数的最大值为 10 个，故①正确；}
\step{对于②，由于集合 $B$ 中没有 0，所以集合 $A$ 中没有 0；}
\step{设集合 $A=\{a,b,c,d\}$，且 $a<b<c<d$，}
\step{当 $a<0<b<c<d$ 时，$a\times b<0$，$a\times c<0$，$a\times d<0$，}
\step{故集合 $B$ 至少有 3 个负数，与 $B=\{2,4,5,8,10,16\}$ 矛盾；}
\step{当 $a<b<0<c<d$ 时，$a\times b>0$，$c\times d>0$，}
\step{$a\times c<0$，$a\times d<0$，$b\times c<0$，$b\times d<0$，}
\step{故集合 $B$ 至少有 4 个负数，与 $B=\{2,4,5,8,10,16\}$ 矛盾；}
\step{当 $a<b<c<0<d$ 时，$a\times b>0$，$a\times c>0$，$b\times c>0$，}
\step{$a\times d<0$，$b\times d<0$，$c\times d<0$，}
\step{故集合 $B$ 至少有 3 个负数，与 $B=\{2,4,5,8,10,16\}$ 矛盾；}
\step{当 $0<a<b<c<d$ 时，则集合 $B=\{ab,ac,ad,bc,bd,cd\}$，}
\step{则必有 $ab=2$，$cd=16$，其 4 个正实数的乘积 $abcd=32$，}
\step{又 $ad=5$，$bc=8$ 或 $ad=8$，$bc=5$，其 4 个正实数的乘积 $abcd=40$，矛盾；}
\step{当 $a<b<c<d<0$ 时，与 $0<a<b<c<d$ 时的证明方法相同，}
\step{得集合 $A$ 也不存在；}
\step{综上不存在四元集 $A$，使其生成集 $B=\{2,4,5,8,10,16\}$，故②错误.}
\step{故选 $B$.}
\solend
\fi

\end{enumerate}

\bigsec{三、解答题}

\begin{enumerate}
\setcounter{enumi}{16}

\item 解下列不等式（组）

\begin{enumerate}[label=(\arabic*)]
\item 解不等式：$(ax+1)(x+1)<0(a\in\R)$.
\item 解不等式组：$\begin{cases}\dfrac{2x+4}{x+4}\geqslant1\\[4pt] (x+1)(2-|x|)\geqslant0\end{cases}$.
\end{enumerate}

\ifanswer
\solbegin
\stepm{(1)}{当 $a=0$ 时，原不等式为 $x+1<0$，即 $x<-1$，}
\step{所以此时不等式的解集为 $\{x\mid x<-1\}$；}
\step{当 $a\neq0$ 时，令 $(ax+1)(x+1)=0$，解得 $x_1=-\dfrac1a$,$x_2=-1$，}
\step{当 $a>1$ 时，$-\dfrac1a>-1$，此时不等式的解集为 $\left\{x\mid-1<x<-\dfrac1a\right\}$；}
% 以下两行“此时”前原卷均无标点；照录（见文件头“原卷笔误/存疑”⑥）
\step{当 $a=1$ 时，原不等式为 $(x+1)^2<0$ 此时不等式的解集为 $\varnothing$；}
\step{当 $0<a<1$ 时，$-\dfrac1a<-1$ 此时不等式的解集为 $\left\{x\mid-\dfrac1a<x<-1\right\}$；}
\step{当 $a<0$ 时，$-\dfrac1a>-1$，此时不等式的解集为 $\left\{x\mid x<-1\ \text{或}\ x>-\dfrac1a\right\}$.}
\step{综上所述，当 $a=0$ 时，$\{x\mid x<-1\}$；}
\step{当 $a>1$ 时，$\left\{x\mid-1<x<-\dfrac1a\right\}$；}
\step{当 $a=1$ 时，$\varnothing$；}
\step{当 $0<a<1$ 时，$\left\{x\mid-\dfrac1a<x<-1\right\}$；}
\step{当 $a<0$ 时，$\left\{x\mid x<-1\ \text{或}\ x>-\dfrac1a\right\}$.}
\stepm{(2)}{不等式 $\dfrac{2x+4}{x+4}\geqslant1$ 可化为 $\dfrac{2x+4}{x+4}-1\geqslant0$，即 $\dfrac x{x+4}\geqslant0$，}
\step{等价于 $\begin{cases}x(x+4)\geqslant0\\ x+4\neq0\end{cases}$，解得 $x<-4$ 或 $x\geqslant0$，}
\step{不等式 $(x+1)(2-|x|)\geqslant0$ 等价于 $\begin{cases}(x+1)(2-x)\geqslant0,x\geqslant0\\ (x+1)(2+x)\geqslant0,x<0\end{cases}$，}
\step{解不等式 $(x+1)(2-x)\geqslant0$，$x\geqslant0$，得 $0\leqslant x\leqslant2$，}
\step{解不等式 $(x+1)(2+x)\geqslant0$，$x<0$，得 $x\leqslant-2$ 或 $-1\leqslant x<0$，}
\step{所以 $(x+1)(2-|x|)\geqslant0$ 的解集为 $\{x\mid x\leqslant-2\ \text{或}\ -1\leqslant x\leqslant2\}$，}
\step{综上所述，$\begin{cases}\dfrac{2x+4}{x+4}\geqslant1\\[4pt] (x+1)(2-|x|)\geqslant0\end{cases}$ 的解集为 $\{x\mid x<-4\ \text{或}\ 0\leqslant x\leqslant2\}$.}
\solend
\fi

\ansspace[8]

% 第 18 题题干＋两小问与作答区 6 行须同页
\ansneed{20cm}
\item 已知 $a$ 为实数，设 $f(x)=x^2+2(a+1)x+1$.

\begin{enumerate}[label=(\arabic*)]
\item 已知对于任意实数 $x$，$\sqrt{f(x)}$ 均有意义，求实数 $a$ 的取值范围；
\item 设 $A=\{x\mid f(x)+a^2-2=0\}$，$B=\{-4,0\}$，$A\cup B=B$，求实数 $a$ 的取值范围.
\end{enumerate}

\ifanswer
\solbegin
\stepm{(1)}{$\sqrt{f(x)}$ 对任意实数 $x$ 都有意义，}
\step{即 $f(x)=x^2+2(a+1)x+1\geqslant0$ 对任意 $x\in\R$ 恒成立，}
\step{所以 $\Delta=4(a+1)^2-4\times1\times1\leqslant0$，即 $(a+1)^2\leqslant1$，解得 $-2\leqslant a\leqslant0$，}
\step{故 $a$ 的取值范围是 $[-2,0]$.}
\stepm{(2)}{由 $f(x)+a^2-2=0$ 得 $x^2+2(a+1)x+a^2-1=0$.}
\step{由 $A\cup B=B$ 得 $A\subseteq B=\{-4,0\}$，即该方程无实根，}
\step{或它的实根都在 $\{-4,0\}$ 中.}
\step{当方程无实根时，$\Delta=4(a+1)^2-4(a^2-1)=8a+8<0$，解得 $a<-1$；}
\step{当方程有实根时，$a\geqslant-1$，设两根为 $x_1$、$x_2$，}
\step{由根与系数的关系得 $x_1+x_2=-2(a+1)$、$x_1x_2=a^2-1$，}
\step{此时两根只能在 $\{-4,0\}$ 中取值；}
\step{若两根为 0 与 0，由 $x_1x_2=0$ 得 $a=\pm1$，}
\step{再由 $x_1+x_2=0$ 得 $a=-1$；}
\step{若两根为 0 与 $-4$，由 $x_1x_2=0$ 得 $a=\pm1$，再由 $x_1+x_2=-4$ 得 $a=1$.}
\step{经检验，$a=-1$ 时 $A=\{0\}$，$a=1$ 时 $A=\{-4,0\}$，都满足 $A\subseteq B$.}
\step{综上，实数 $a$ 的取值范围为 $(-\infty,-1]\cup\{1\}$.}
\solend
\fi

\ansspace[8]

\ansneed{18cm}
\item 共享单车给市民出行带来了诸多便利，某公司购买了一批单车投放到某地给居民使用，据市场分析，这批单车的运营累计利润 $y$（单位：元）与运营天数 $x$（单位：天，$x$ 为正整数）满足关系：$y=-\dfrac12x^2+60x-800$.

\begin{enumerate}[label=(\arabic*)]
\item 要使这批单车的运营累计利润高于 800 元，求运营天数的取值范围；
\item 判断这批单车的日平均运营利润能否超过 20 元，并说明理由.
\end{enumerate}

\ifanswer
\solbegin
\stepm{(1)}{由题意得 $y=-\dfrac12x^2+60x-800>800$，}
\step{即 $x^2-120x+3200<0$，解得 $40<x<80$，}
\step{即运营天数的取值范围为 $\{x\mid40<x<80,x\in\N^*\}$.}
\stepm{(2)}{不能，理由如下：}
\step{这批单车的日平均运营利润为 $\dfrac yx=-\dfrac12x-\dfrac{800}x+60=-\dfrac12\left(x+\dfrac{1600}x\right)+60$，}
\step{因 $x>0$，则 $x+\dfrac{1600}x\geqslant2\sqrt{x\cdot\dfrac{1600}x}=80$，}
\step{当且仅当 $x=\dfrac{1600}x$，即 $x=40$ 时取等号，}
\step{此时 $\dfrac yx=-\dfrac12\left(x+\dfrac{1600}x\right)+60\leqslant-\dfrac12\times80+60=20$，}
\step{故这批单车的日平均运营利润不能超过 20 元.}
\solend
\fi

\ansspace[8]

% 第 20 题题干＋三小问与作答区 8 行须同页
\ansneed{20cm}
\item 设 $f(x)=ax^2+bx+c$($a\neq0$，$b$、$c\in\R$).

\begin{enumerate}[label=(\arabic*)]
\item 若 $a=b=c=1$，设 $A=\{(x,y)\mid y=f(x)\}$，$B=\{(x,y)\mid y=1\}$，用列举法表示 $A\cap B$；
% 原卷两个“若”连用（“若 a=1，…，若方程 f(x)=0 有两个大于 1 的根”），疑衍其一；照录（见文件头“原卷笔误/存疑”⑦）
\item 若 $a=1$，$b=2k-1$，$c=k^2$，若方程 $f(x)=0$ 有两个大于 1 的根，求实数 $k$ 的取值范围；
\item 若 $a$、$b$、$c\in\Z$，且 $f(0)$、$f(1)$ 均为奇数，求证：方程 $f(x)=0$ 无整数解.
\end{enumerate}

\ifanswer
\solbegin
\stepm{(1)}{当 $a=b=c=1$ 时，$f(x)=x^2+x+1$，}
\step{所以 $A=\{(x,y)\mid y=x^2+x+1\}$，又 $y=1$，}
\step{由 $x^2+x+1=1$ 得 $x^2+x=0$，解得 $x=0$ 或 $x=-1$，}
\step{所以 $A\cap B=\{(0,1),(-1,1)\}$.}
\stepm{(2)}{当 $a=1$、$b=2k-1$、$c=k^2$ 时，$f(x)=x^2+(2k-1)x+k^2$.}
\step{设两根为 $x_1$、$x_2$，}
\step{则“两根都大于 1”等价于 $\Delta>0$、$f(1)>0$、$-\dfrac b{2a}>1$.}
\step{由 $\Delta=(2k-1)^2-4k^2=1-4k>0$ 得 $k<\dfrac14$；}
\step{由 $f(1)=1+(2k-1)+k^2=k^2+2k=k(k+2)>0$ 得 $k<-2$ 或 $k>0$；}
\step{由 $-\dfrac{2k-1}2=\dfrac{1-2k}2>1$ 得 $k<-\dfrac12$.}
\step{三者取交集得 $k<-2$，即 $k$ 的取值范围是 $(-\infty,-2)$.}
\stepm{(3)}{设 $x_0\in\Z$ 是方程 $f(x)=0$ 的整数解.}
\step{因为 $x_0(x_0-1)$ 是两个相邻整数之积，恒为偶数，所以 $ax_0(x_0-1)$ 为偶数.}
\step{而 $f(x_0)=ax_0^2+bx_0+c=ax_0(x_0-1)+(a+b)x_0+c$，}
\step{所以 $f(x_0)$ 与 $(a+b)x_0+c$ 同奇偶.}
\step{若 $x_0$ 为偶数，则 $(a+b)x_0$ 为偶数，$f(x_0)$ 与 $c=f(0)$ 同奇偶，为奇数；}
\step{若 $x_0$ 为奇数，则 $(a+b)x_0$ 与 $a+b$ 同奇偶，}
\step{$f(x_0)$ 与 $(a+b)+c=f(1)$ 同奇偶，也为奇数.}
\step{综上，$f(x_0)$ 恒为奇数，不可能等于 0，与 $f(x_0)=0$ 矛盾，}
\step{所以方程 $f(x)=0$ 无整数解.}
\solend
\fi

\ansspace[8]

% 压轴题（题干＋三小问）与作答区 10 行须同页
\ansneed{20cm}
\item 已知 $X$ 为有限集，且任取 $x\in X$，均有 $x\in\Z$，称 $X$ 为非空有限整数集. 定义：$X_k=\{x+k\mid x\in X\}$，$X^2=\{y\mid y=x^2,x\in X\}$.

\begin{enumerate}[label=(\arabic*)]
\item 若 $X=\{-2,2\}$，分别写出 $X_2$ 和 $X^2$（用列举法表示）；
\item 若 $X=\{-1,0,1\}$，$X_1\subseteq Y$，$(Y_{-1})^2\subseteq X$，求集合 $Y$ 中所有的元素；
\item 对于非空有限整数集 $A$、$B$，若 $A_k\subseteq B$，$(B^2)_{-2}\subseteq A$，$B$ 中的元素均为非负整数，$k$ 为整数，求证：

① $A$、$B$ 中元素的个数相同；

② 若集合 $A$ 中的元素个数为奇数，则 $k$ 为偶数.
\end{enumerate}

\ifanswer
\solbegin
\stepm{(1)}{由题意得 $X_2=\{0,4\}$，$X^2=\{4\}$；}
\stepm{(2)}{由题意得 $X_1=\{0,1,2\}$，即 $\{0,1,2\}\subseteq Y$，$Y_{-1}=\{y-1\mid y\in Y\}$，}
\step{$(Y_{-1})^2=\{(y-1)^2\mid y\in Y\}\subseteq X=\{-1,0,1\}$，}
\step{所以 $(y-1)^2=0$ 或 $(y-1)^2=1$，解得 $y=1$ 或 $y=2$ 或 $y=0$，}
\step{即 $y\in\{0,1,2\}$，即 $Y\subseteq\{0,1,2\}$，所以 $Y=\{0,1,2\}$；}
% 原卷 (3) 解析首行“(3)”与“①”连排，本稿合并记作 \stepm{(3)①}
\stepm{(3)①}{因为 $A_k=\{x+k\mid x\in A\}$，函数 $g(x)=x+k$ 在全体整数集上严格增，}
\step{所以 $\forall x_1$、$x_2\in A$，$x_1\neq x_2$，$g(x_1)\neq g(x_2)$，所以 $|A_k|=|A|$，}
\step{因为 $A_k\subseteq B$，所以 $|A|\leqslant|B|$，}
\step{设 $b\in B$，$B$ 中元素均为非负整数，$\forall b_1\neq b_2\in B$，$b_1^2\neq b_2^2$，}
\step{故函数 $f(b)=b^2-2$ 在 $B$ 上单射，所以 $|(B^2)_{-2}|=|B|$，}
\step{又因为 $(B^2)_{-2}\subseteq A$，所以 $|B|\leqslant|A|$，所以 $|A|=|B|$.}
\stepm{②}{由①得 $A_k\subseteq B$ 且 $|A_k|=|A|=|B|$，$(B^2)_{-2}\subseteq A$ 且 $|(B^2)_{-2}|=|B|=|A|$，}
\step{所以 $A_k=B$，$(B^2)_{-2}=A$，}
\step{所以 $B=A_k=\{a+k\mid a\in A\}$，$A=(B^2)_{-2}=\{b^2-2\mid b\in B\}$，}
\step{所以 $A=\{(a+k)^2-2\mid a\in A\}$，}
\step{所以映射 $h(a)=(a+k)^2-2$ 是集合 $A$ 到自身的一个双射.}
\step{假设 $k$ 是奇数，当 $a$ 为偶数时，$a+k$ 为奇数，$(a+k)^2-2$ 为奇数，}
\step{当 $a$ 为奇数时，$a+k$ 为偶数，$(a+k)^2-2$ 为偶数，}
\step{所以此时映射 $h(a)=(a+k)^2-2$，将集合 $A$ 中的奇数映射为偶数，}
\step{将集合 $A$ 中的偶数映射为奇数，即集合 $A$ 中的奇数元素和偶数元素必须配对，}
\step{此时集合 $A$ 的元素个数则为偶数，}
\step{与题设中“集合 $A$ 中的元素个数为奇数”矛盾.}
\step{故假设不成立，所以 $k$ 是偶数.}
\solend
\fi

\ansspace*[10]

\end{enumerate}

\end{document}
"
% 紧凑版：xelatex "\def\iscompact{1}% 高一67_华师大一附中_10月月考.tex
%   —— 高一上数学“2026-2027 年华师大一附中高一上 10 月月考”（讲评版）
% 试卷版：xelatex 高一67_华师大一附中_10月月考.tex
% 答案版：xelatex "\def\isanswer{1}\input{高一67_华师大一附中_10月月考.tex}"
% 紧凑版：xelatex "\def\iscompact{1}\input{高一67_华师大一附中_10月月考.tex}"
%
% 原始材料：微信公众号“魔都数学加油站”文章，作者破晓，连载“27学年高一67”；
% 原文标题“27学年高一67——2026-2027年上海市华师大一附中高一上10月月考”；
% 原文链接：https://mp.weixin.qq.com/s/Kdhx07TVjLoGTcIlnRzgcQ。
% 页面标注发布时间 2026-10-18 20:33（上海）。
% 本次转录只据 `原图/高一67_华师大一附中_10月月考/` 的 13 张 PNG 逐页读图。
% 正文为 13 张 PNG 页（`原图/高一67_华师大一附中_10月月考/01..13.png`）：
% 第 01、04、13 页 2560×3620，第 02、03、05—12 页 4762×6735
% （**同一篇各页分辨率不同**），均按页序读取。13 页全为试卷正文页（卷首标题在第 1 页页顶，
% 卷末解答题第 21 题解析止于第 13 页），无封面页、目录页、二维码或宣传页。
% 粉色斜排“魔都数学加油站”为卷面水印，与公众号页脚一并未录入。
%
% 卷首标题：照卷面印作“2026--2027 年华师大一附中高一上 10 月月考”（卷面为黑体加粗一行），
% 年份区间按全稿惯例排 en dash。原卷只有一行标题、无副标题。
%
% 篇幅与题号：原卷自第 3 题起编号（第 1、2 题不在本卷内）。填空题 3—12（10 题）、
% 选择题 13—16（4 题）、解答题 17—21（5 题），共 19 题。本稿沿用原节与原题号，
% 只改排为全稿统一的大节样式（原卷小节标题为“一、填空题”“二、选择题”“三、解答题”）。
%
% 讲评版：原卷每题下直接印【解析】。**全卷 19 题中 18 题都印【解析】（无一处只印
% 【答案】）**，按全稿口径这类题不另提 \ans{}；各题解析末句的结论（第 13、14、16 题
% 的“故选 B／D／B”、各填空与解答题的“故/综上…”）照录在解析内。
% **另有 1 题答案印在题干括号内**——第 15 题（“(  D  )”，且该题无【答案】/【解析】段），
% 本稿按 高一37 第 13 题、高一63 第 14 题成例把括号留空、答案另提为 \ans{$D$}。
%
% 副标题：原卷未印副标题。本稿按“知识模块”口径标作“集合与常用逻辑用语、不等式”：
% 19 题中集合与常用逻辑用语 12 题（第 3、4、5、7、10、11、12、13、15、16、18、21 题：
% 补集、并集、子集、充分条件、容斥计数、子集计数、新定义二元运算、条件判断、
% 反证法假设、新定义生成集、集合关系与二次方程、新定义 X_k 与 X^2），
% 不等式 7 题（第 6、8、9、14、17、19、20 题：存在量词不等式、解集反求系数、
% 解集为空、不等式性质、解不等式（组）、二次不等式与基本不等式、二次方程根的分布），
% 两类均超过三分之一，故并标。
%
% 与原卷的排版差异：
%   ① 原卷每题下直接印【解析】，本稿按全稿惯例将解析置于答案版，试卷版保持干净；
%   ② 原卷未印副标题，本稿按“知识模块”口径补出；
%   ③ 原卷填空横线约四个字宽，本稿按全稿惯例统一排 \blankline[4.4em]；
%   ④ 原卷实数集排斜体 $R$（第 14 题题干、第 17 题 (1) 题干、第 18 题解析、第 20 题题干），
%      整数集排斜体 $Z$（第 20 题 (3) 题干与解析、第 21 题题干），正整数集排粗体
%      $N^*$（第 11 题解析、第 19 题解析），本稿按全稿惯例统一改用 $\R$、$\Z$、$\N^*$
%      （记号归一，粗体／斜体不作区分）；空集原卷作 $\varnothing$，照录；
%   ⑤ 原卷在数学式内多处用半角逗号（如第 8 题 $(1,2)$ 与各 cases、第 11 题解析
%      $P_n=\{1,2,\cdots,n\}$,$n\in N^*$、第 12 题解析 $u=(\cdots)$,$v=(\cdots)$ 与
%      “2020,个 1”、第 19 题解析 $\{x|40<x<80,x\in N^*\}$），均照录未强改；
%      描述法集合的竖线照原卷排 \mid；句末句点按全稿惯例排西文句点；
%   ⑥ 原卷补集记号为上划线 $\bar A$（第 3 题题干与解析、第 11 题题干），照录
%      （同 高一64/65 的 $\bar B$ 成例）；
%   ⑦ 第 12 题新定义的运算记号 $*$ 照录（$\alpha*\beta$、$u*u$、$u*v$）；
%   ⑧ 第 17—21 题的子问原卷各自另起一行，本稿按全稿惯例用 enumerate 的 (1)(2)(3)；
%      解析中的子问标号一律用 \stepm{(1)}{…}；第 5 题解析的①②、第 21 题 (3) 解析的①②
%      按全稿惯例作 \stepm{①}{…}、\stepm{②}{…}（第 21 题 (3) 首行原卷“(3)”与“①”连排，
%      本稿合并记作 \stepm{(3)①}{…}）；
%   ⑨ 选择题作答括号原卷排半角“(　　)”（第 13、14、16 题括号内空白、第 15 题括号内印 D），
%      本稿按全稿惯例排 \mbox{（\quad）}；第 15 题括号内的答案移入 \ans{}，见上；
%   ⑩ 选项行排照原卷：第 13、15 题两行两条，第 14、16 题一行四条；
%   ⑪ 每道解答题原卷未留作答空白，本稿按全稿惯例加 \ansspace（末题用 \ansspace*）；
%      第 20、21 题题干长，各用一处 \ansneed 让题干与作答区不拆页（\ansneed 只作用于试卷版）。
%
% 原卷笔误/存疑（均照录未改，正文对应处另有 % 注）：
%   ① 第 3 题解析末句印作“则 $\bar A=(-\infty,0)\cup[2,3)$.”：按题干
%      $U=\{x\mid 1/x\leqslant 1/2\}$（即 $x<0$ 或 $x\geqslant2$）与 $A=[3,+\infty)$，
%      $\bar A$ 应为 $(-\infty,0)\cup\{2\}$；且与本解析首句自求之 $U$ 相矛盾（$3\notin U$），
%      疑为笔误；照录未改。
%   ② 第 12 题题干 $\alpha*\beta$ 定义式括号不配对：第四组“$+\cdots+x_{2020}+y_{2020}
%      -x_{2020}y_{2020}$”缺左括号、末尾连排两个右括号“))”；照录未改。
%   ③ 第 12 题解析“所以…中有 2020,个 1”一行，“$\cdots$”后脱逗号（作
%      “$x_2,\cdots x_{2020}$”，第 12 行“而…”同式样），且“2020”后衍一半角逗号
%      （作“2020,个 1”）；照录未改。
%   ④ 第 16 题解析乘积链印作“$a_1a_2<a_1a_3<a_1a_4<a_1a_5<a_2a_3<a_3a_3<a_{14}a_5$”：
%      “$a_3a_3$”同足标自乘、与单调链不合，“$a_{14}a_5$”下标 14 与 $a_1$—$a_5$ 记号
%      不合（疑为 $a_2a_4$、$a_4a_5$ 之误）；照录未改。
%   ⑤ 第 5 题题干作“$Q\subset P$”（真子集记号），解析两处均作“$Q\subseteq P$”，
%      同一题前后记号不一；照录（本稿分别排 \subset、\subseteq）。
%   ⑥ 第 17 题 (1) 解析“当 $a=1$ 时，原不等式为 $(x+1)^2<0$ 此时…”与
%      “当 $0<a<1$ 时，$-\frac1a<-1$ 此时…”两处“此时”前均脱标点；照录未改。
%   ⑦ 第 20 题 (2) 题干“若 $a=1$，$b=2k-1$，$c=k^2$，若方程 $f(x)=0$ 有两个大于 1 的根”
%      两个“若”连用，疑衍其一；照录未改。
%
% 插图：无（全卷 13 页无任何几何图形、无“如图/右图/图中”字样）。
%
% 行数核对：逐题数 \step/\stepm 条数与原卷【解析】印刷行数，18 题逐一相符
% （第 15 题无解析段，只提 \ans{}），合计 169 条 / 169 行
% （Q3—Q12：2、3、4、3、1、3、5、3、13、15；Q13—Q16：2、3、0、25；
% Q17—Q21：19、16、9、20、23）。
%
\documentclass[a4paper,11pt]{article}
\usepackage{common}

\begin{document}

\examtitlest[3]{2026--2027 年华师大一附中高一上 10 月月考}{集合与常用逻辑用语、不等式}

\bigsec{一、填空题}

\begin{enumerate}
\setcounter{enumi}{2}

\item 设全集 $U=\left\{x\mid\dfrac1x\leqslant\dfrac12\right\}$，$A=[3,+\infty)$，则 $\bar A=$\blankline[4.4em].

\ifanswer
\solbegin
\step{$\dfrac1x\leqslant\dfrac12$，解得 $x<0$ 或 $x\geqslant2$，故 $U=(-\infty,0)\cup[2,+\infty)$，}
% 原卷末句作 $\bar A=(-\infty,0)\cup[2,3)$，按题干 $U$ 应为 $(-\infty,0)\cup\{2\}$，与本解析首句矛盾，疑笔误；照录（见文件头“原卷笔误/存疑”①）
\step{$A=[3,+\infty)$，则 $\bar A=(-\infty,0)\cup[2,3)$.}
\solend
\fi

\item 已知集合 $M=\{y\mid y=x^2-2x-1\}$，$N=\{y\mid y=x^2-4x+1\}$，则 $M\cup N=$\blankline[4.4em].

\ifanswer
\solbegin
\step{$M=\{y\mid y=x^2-2x-1\}=\{y\mid y=(x-1)^2-2\}=\{y\mid y\geqslant-2\}$，}
\step{$N=\{y\mid y=x^2-4x+1\}=\{y\mid y=(x-2)^2-3\}=\{y\mid y\geqslant-3\}$，}
\step{则 $M\cup N=\{y\mid y\geqslant-3\}$.}
\solend
\fi

\item 已知 $P=\{x\mid x^2+x-6=0\}$，$Q=\{x\mid mx-1=0\}$，且 $Q\subset P$，则实数 $m$ 的值\blankline[4.4em].

\ifanswer
\solbegin
\step{$P=\{2,-3\}$，$Q=\{x\mid mx=1\}$；}
\stepm{①}{$m=0$ 时，$Q=\varnothing$，满足 $Q\subseteq P$；}
\stepm{②}{$m\neq0$ 时，要使 $Q\subseteq P$，则 $\dfrac1m=2$ 或 $-3$，所以 $m=\dfrac12$ 或 $-\dfrac13$；}
\step{所以 $m=0$ 或 $\dfrac12$ 或 $-\dfrac13$.}
\solend
\fi

\item 已知“存在正实数 $x$，使得关于 $x$ 的不等式 $ax-1>0$ 成立”为真命题，则实数 $a$ 的取值范围是\blankline[4.4em].

\ifanswer
\solbegin
\step{存在正实数 $x$，使得关于 $x$ 的不等式 $ax-1>0$ 成立为真命题，}
\step{即存在正实数 $x$，使得 $a>\dfrac1x$，因为 $\dfrac1x>0$，所以 $a>0$，}
\step{即 $a$ 的范围为 $(0,+\infty)$.}
\solend
\fi

\item 设 $\alpha:3\leqslant x\leqslant5$，$\beta:2m-1\leqslant x<m+4$，若 $\alpha$ 是 $\beta$ 的充分条件，则实数 $m$ 的取值范围是\blankline[4.4em].

\ifanswer
\solbegin
\step{若 $\alpha$ 是 $\beta$ 的充分条件，则 $\begin{cases}2m-1\leqslant3\\ m+4>5\end{cases}$，解得 $1<m\leqslant2$.}
\solend
\fi

\item 对于实数 $a$、$b$、$c$，已知关于 $x$ 的不等式 $ax^2+bx+c>0$ 的解集是 $(1,2)$，则关于 $x$ 的不等式 $cx^2-bx+a\leqslant0$ 的解集是\blankline[4.4em].

\ifanswer
\solbegin
\step{关于 $x$ 的不等式 $ax^2+bx+c>0$ 的解集是 $(1,2)$，则 $\begin{cases}a<0\\ -\dfrac ba=3\\ \dfrac ca=2\end{cases}$，则 $\begin{cases}b=-3a\\ c=2a\end{cases}$，}
\step{则关于 $x$ 的不等式 $cx^2-bx+a=2ax^2+3ax+a\leqslant0$，则 $2x^2+3x+1\geqslant0$，}
\step{则 $x\leqslant-1$ 或 $x\geqslant-\dfrac12$，则解集为 $(-\infty,-1]\cup\left[-\dfrac12,+\infty\right)$.}
\solend
\fi

\item 已知关于 $x$ 的不等式 $kx^2+kx+k+2\leqslant0$ 的解集为 $\varnothing$，则实数 $k$ 的取值范围为\blankline[4.4em].

\ifanswer
\solbegin
\step{当 $k=0$ 时，不等式化为 $2\leqslant0$，解集为 $\varnothing$，符合题意；}
\step{当 $k\neq0$ 时，若关于 $x$ 的不等式 $kx^2+kx+k+2\leqslant0$ 的解集为 $\varnothing$，}
\step{则 $\begin{cases}k>0\\ \Delta=k^2-4k(k+2)<0\end{cases}$，化简得 $\begin{cases}k>0\\ -3k^2-8k<0\end{cases}$，即 $\begin{cases}k>0\\ k(3k+8)>0\end{cases}$，}
\step{解得 $k>0$.}
\step{综上所述，$k\geqslant0$，即实数 $k$ 的取值范围为 $[0,+\infty)$.}
\solend
\fi

\item 对班级若干名学生调查对 $A$、$B$ 两事件的态度，有如下结果：赞成 $A$ 的人数是全体的五分之三，其余的不赞成；赞成 $B$ 的比赞成 $A$ 多 3 人，其余的不赞成. 已知对 $A$、$B$ 都赞成的学生有 18 人，对 $A$、$B$ 都不赞成的学生数比对 $A$、$B$ 都赞成的学生数的三分之一多 1 人，则班级总人数是\blankline[4.4em]人.

\ifanswer
\solbegin
\step{设班级总人数是 $x$ 人，则赞成 $A$ 的人数是 $\dfrac35x$，赞成 $B$ 的人数是 $\dfrac35x+3$，}
\step{$A$、$B$ 都不赞成的人数为 $18\times\dfrac13+1=7$ 人，}
\step{由容斥原理得 $x=\left(\dfrac35x+\dfrac35x+3-18\right)+7$，解得 $x=40$，即班级总人数是 40 人.}
\solend
\fi

\item 设全集 $U=\{1,2,3,4,5,6,7,8,9,10,11,12\}$，若集合 $A$ 满足：① $A\subseteq U$；②若 $x\in A$，必有 $2x\notin A$；③若 $x\in\bar A$，必有 $2x\notin\bar A$，则这样的集合 $A$ 有\blankline[4.4em]个.

\ifanswer
\solbegin
\step{设集合 $P_n=\{1,2,\cdots,n\}$,$n\in\N^*$，}
\step{任取 $x\in P_n$，将 $x$ 除以 2，若商仍为偶数，再除以 2，$\cdots$，}
\step{经过 $k$ 次以后，商必为奇数，此时记商为 $m$，}
\step{于是 $x=m\cdot2^k$，其中 $m$ 为奇数，$k\in\N^*$.}
\step{由条件得，若 $m\in A$，则 $x\in A$ 等价于 $k$ 为偶数；}
\step{若 $m\notin A$，则 $x\in A$ 等价于 $k$ 为奇数；}
\step{于是 $x$ 是否属于 $A$ 由 $m$ 是否属于 $A$ 确定；}
\step{设 $Q_n$ 是 $P_n$ 中所有奇数的集合，}
\step{设 $f(n)$ 为集合 $P_n$ 中满足题干条件的集合 $A$ 的个数，}
\step{因此 $f(n)$ 等于 $Q_n$ 的子集个数.}
\step{当 $n$ 为偶数（或奇数）时，$P_n$ 中奇数的个数是 $\dfrac n2$（或 $\dfrac{n+1}2$），}
\step{所以 $f(n)=\begin{cases}2^{\frac n2},n\text{为偶数}\\[2pt] 2^{\frac{n+1}2},n\text{为奇数}\end{cases}$，所以 $P_{12}=\{1,2,3,4,5,6,7,8,9,10,11,12\}$，}
\step{即同时满足三个条件的集合 $A$ 的个数为 $2^{\frac{12}2}=2^6=64$.}
\solend
\fi

\item 设集合 $\Omega=\{(x_1,x_2,\cdots,x_{2020})\mid x_i\in\{0,1\},i=1,2,\cdots,2020\}$，对于 $\Omega$ 中的任意两个元素 $\alpha=(x_1,x_2,\cdots,x_{2020})$，$\beta=(y_1,y_2,\cdots,y_{2020})$，记
% 原卷此式第四组缺左括号、末尾连排两个右括号“))”，括号不配对；照录（见文件头“原卷笔误/存疑”②）
$\alpha*\beta=(x_1+y_1-x_1y_1)+(x_2+y_2-x_2y_2)+\cdots+x_{2020}+y_{2020}-x_{2020}y_{2020}))$，设 $u$、$v\in\Omega$，若 $u*u+v*v=2020$，则 $u*v$ 的最大值和最小值之和为\blankline[4.4em].

\ifanswer
\solbegin
\step{设 $u=(x_1,x_2,\cdots,x_{2020})$,$v=(y_1,y_2,\cdots,y_{2020})$，}
\step{因为 $x_i\in\{0,1\}$，所以 $x_i(x_i-1)=0$，}
\step{所以 $u*u=(x_1+x_1-x_1^2)+(x_2+x_2-x_2^2)+\cdots+(x_{2020}+x_{2020}-x_{2020}^2)$}
\step{$=x_1+x_2+\cdots+x_{2020}$，}
\step{$v*v=(y_1+y_1-y_1^2)+(y_2+y_2-y_2^2)+\cdots+(y_{2020}+y_{2020}-y_{2020}^2)$}
\step{$=y_1+y_2+\cdots+y_{2020}$，}
\step{由 $u*u+v*v=2020$，得 $x_1+x_2+\cdots+x_{2020}+y_1+y_2+\cdots+y_{2020}=2020$，}
% 原卷此行“$\cdots$”后脱逗号、 “2020”后衍一半角逗号（“2020,个 1”）；照录（见文件头“原卷笔误/存疑”③）
\step{所以 $x_1,x_2,\cdots x_{2020},y_1,y_2,\cdots y_{2020}$ 中有 2020,个 1，2020 个 0，}
\step{$u*v=(x_1+y_1-x_1y_1)+(x_2+y_2-x_2y_2)+\cdots+(x_{2020}+y_{2020}-x_{2020}y_{2020})$}
\step{$=2020-(x_1y_1+x_2y_2+\cdots+x_{2020}y_{2020})$，}
\step{注意到只有 $x_i=y_i=1$ 时，$x_iy_i=1$，否则 $x_iy_i=0$，}
% 此行“$\cdots$”后同样脱逗号，与上一处同式样；照录（见文件头“原卷笔误/存疑”③）
\step{而 $x_1,x_2,\cdots x_{2020},y_1,y_2,\cdots y_{2020}$ 中有 2020 个 1，2020 个 0，}
\step{所以满足 $x_iy_i=1$ 的最多有 1010 个，所以 $u*v\geqslant2020-1010=1010$，}
\step{满足 $x_iy_i=0$ 的最多有 2020 个，所以 $u*v\leqslant2020-0=2020$，}
\step{则 $u*v$ 的最大值和最小值之和为 $2020+1010=3030$.}
\solend
\fi

\end{enumerate}

\bigsec{二、选择题}

\begin{enumerate}
\setcounter{enumi}{12}

\item 实数 $x$ 与 $y$，“$x+y>2$”的\mbox{（\quad）}条件是“$x>1$ 或 $y>1$”.

\keepwithprev
A. 充分非必要\qquad B. 必要非充分

C. 充要\qquad D. 既非充分又非必要

\ifanswer
\solbegin
\step{当 $x>1$ 或 $y>1$ 时，$x+y>2$ 不一定成立，充分性不成立；}
\step{当 $x+y>2$ 时，$x>1$ 或 $y>1$ 一定成立，必要性成立. 故选 $B$.}
\solend
\fi

\item 设 $a,b\in\R$，且 $a<0$，$-1<b<0$，那么下列不等式成立的是\mbox{（\quad）}

\keepwithprev
A. $a>ab>ab^2$\qquad B. $ab^2>ab>a$\qquad C. $ab>a>ab^2$\qquad D. $ab>ab^2>a$

\ifanswer
\solbegin
\step{因为 $-1<b<0$，$a<0$，所以 $ab>0$，$0<b^2<1$.}
\step{所以 $ab-ab^2=ab(1-b)>0$，$ab^2-a=a(b^2-1)>0$，}
\step{所以 $ab>ab^2>a$. 故选 $D$.}
\solend
\fi

% 第 15 题答案原卷印在题干括号内（“(  D  )”），且无【答案】/【解析】段；
% 按 高一37 第 13 题、高一63 第 14 题成例括号留空、答案另提 \ans{$D$}
\item 用反证法证明：“平面四边形中至少有一个内角不超过 90°”时，下列“假设”中正确的是\mbox{（\quad）}

\keepwithprev
A. 假设有两个内角超过 90°\qquad B. 假设有三个内角超过 90°

C. 假设至多有两个内角超过 90°\qquad D. 假设四个内角都超过 90°

\ifanswer
\ans{$D$}
\fi

\item 设 $A$ 为非空实数集，集合 $B=\{xy\mid x\in A,y\in A,x\neq y\}$，则对于以下两个结论说法正确的是\mbox{（\quad）}

结论①：若 $A$ 是五元集，则 $B$ 中的元素最多 10 个；

结论②：存在四元集 $A$，使得 $B=\{2,4,5,8,10,16\}$.

\keepwithprev
A. ①对②对\qquad B. ①对②错\qquad C. ①错②对\qquad D. ①错②错

\ifanswer
\solbegin
\step{对于①，想让元素个数达到最大值 10，则须避免产生符号不同的乘积，}
\step{否则正负分开计数容易导致总数受限，且容易产生相等值，}
\step{也最好避免 0 的出现，}
\step{故设 $A=\{a_1,a_2,a_3,a_4,a_5\}$，$0<a_1<a_2<a_3<a_4<a_5$.}
% 原卷此链“$a_3a_3$”同足标自乘、“$a_{14}a_5$”下标 14 均与文意不合，疑为 $a_2a_4$、$a_4a_5$ 之误；照录（见文件头“原卷笔误/存疑”④）
\step{因为 $a_1a_2<a_1a_3<a_1a_4<a_1a_5<a_2a_3<a_3a_3<a_{14}a_5$，$a_2a_3<a_2a_4<a_3a_4$，}
\step{所以 $B$ 中元素个数小于等于 10 个，}
\step{如设 $A=\{1,2,3,5,7\}$，则 $B=\{2,3,5,6,7,10,14,15,21,35\}$，}
\step{所以集合 $B$ 中元素个数的最大值为 10 个，故①正确；}
\step{对于②，由于集合 $B$ 中没有 0，所以集合 $A$ 中没有 0；}
\step{设集合 $A=\{a,b,c,d\}$，且 $a<b<c<d$，}
\step{当 $a<0<b<c<d$ 时，$a\times b<0$，$a\times c<0$，$a\times d<0$，}
\step{故集合 $B$ 至少有 3 个负数，与 $B=\{2,4,5,8,10,16\}$ 矛盾；}
\step{当 $a<b<0<c<d$ 时，$a\times b>0$，$c\times d>0$，}
\step{$a\times c<0$，$a\times d<0$，$b\times c<0$，$b\times d<0$，}
\step{故集合 $B$ 至少有 4 个负数，与 $B=\{2,4,5,8,10,16\}$ 矛盾；}
\step{当 $a<b<c<0<d$ 时，$a\times b>0$，$a\times c>0$，$b\times c>0$，}
\step{$a\times d<0$，$b\times d<0$，$c\times d<0$，}
\step{故集合 $B$ 至少有 3 个负数，与 $B=\{2,4,5,8,10,16\}$ 矛盾；}
\step{当 $0<a<b<c<d$ 时，则集合 $B=\{ab,ac,ad,bc,bd,cd\}$，}
\step{则必有 $ab=2$，$cd=16$，其 4 个正实数的乘积 $abcd=32$，}
\step{又 $ad=5$，$bc=8$ 或 $ad=8$，$bc=5$，其 4 个正实数的乘积 $abcd=40$，矛盾；}
\step{当 $a<b<c<d<0$ 时，与 $0<a<b<c<d$ 时的证明方法相同，}
\step{得集合 $A$ 也不存在；}
\step{综上不存在四元集 $A$，使其生成集 $B=\{2,4,5,8,10,16\}$，故②错误.}
\step{故选 $B$.}
\solend
\fi

\end{enumerate}

\bigsec{三、解答题}

\begin{enumerate}
\setcounter{enumi}{16}

\item 解下列不等式（组）

\begin{enumerate}[label=(\arabic*)]
\item 解不等式：$(ax+1)(x+1)<0(a\in\R)$.
\item 解不等式组：$\begin{cases}\dfrac{2x+4}{x+4}\geqslant1\\[4pt] (x+1)(2-|x|)\geqslant0\end{cases}$.
\end{enumerate}

\ifanswer
\solbegin
\stepm{(1)}{当 $a=0$ 时，原不等式为 $x+1<0$，即 $x<-1$，}
\step{所以此时不等式的解集为 $\{x\mid x<-1\}$；}
\step{当 $a\neq0$ 时，令 $(ax+1)(x+1)=0$，解得 $x_1=-\dfrac1a$,$x_2=-1$，}
\step{当 $a>1$ 时，$-\dfrac1a>-1$，此时不等式的解集为 $\left\{x\mid-1<x<-\dfrac1a\right\}$；}
% 以下两行“此时”前原卷均无标点；照录（见文件头“原卷笔误/存疑”⑥）
\step{当 $a=1$ 时，原不等式为 $(x+1)^2<0$ 此时不等式的解集为 $\varnothing$；}
\step{当 $0<a<1$ 时，$-\dfrac1a<-1$ 此时不等式的解集为 $\left\{x\mid-\dfrac1a<x<-1\right\}$；}
\step{当 $a<0$ 时，$-\dfrac1a>-1$，此时不等式的解集为 $\left\{x\mid x<-1\ \text{或}\ x>-\dfrac1a\right\}$.}
\step{综上所述，当 $a=0$ 时，$\{x\mid x<-1\}$；}
\step{当 $a>1$ 时，$\left\{x\mid-1<x<-\dfrac1a\right\}$；}
\step{当 $a=1$ 时，$\varnothing$；}
\step{当 $0<a<1$ 时，$\left\{x\mid-\dfrac1a<x<-1\right\}$；}
\step{当 $a<0$ 时，$\left\{x\mid x<-1\ \text{或}\ x>-\dfrac1a\right\}$.}
\stepm{(2)}{不等式 $\dfrac{2x+4}{x+4}\geqslant1$ 可化为 $\dfrac{2x+4}{x+4}-1\geqslant0$，即 $\dfrac x{x+4}\geqslant0$，}
\step{等价于 $\begin{cases}x(x+4)\geqslant0\\ x+4\neq0\end{cases}$，解得 $x<-4$ 或 $x\geqslant0$，}
\step{不等式 $(x+1)(2-|x|)\geqslant0$ 等价于 $\begin{cases}(x+1)(2-x)\geqslant0,x\geqslant0\\ (x+1)(2+x)\geqslant0,x<0\end{cases}$，}
\step{解不等式 $(x+1)(2-x)\geqslant0$，$x\geqslant0$，得 $0\leqslant x\leqslant2$，}
\step{解不等式 $(x+1)(2+x)\geqslant0$，$x<0$，得 $x\leqslant-2$ 或 $-1\leqslant x<0$，}
\step{所以 $(x+1)(2-|x|)\geqslant0$ 的解集为 $\{x\mid x\leqslant-2\ \text{或}\ -1\leqslant x\leqslant2\}$，}
\step{综上所述，$\begin{cases}\dfrac{2x+4}{x+4}\geqslant1\\[4pt] (x+1)(2-|x|)\geqslant0\end{cases}$ 的解集为 $\{x\mid x<-4\ \text{或}\ 0\leqslant x\leqslant2\}$.}
\solend
\fi

\ansspace[8]

% 第 18 题题干＋两小问与作答区 6 行须同页
\ansneed{20cm}
\item 已知 $a$ 为实数，设 $f(x)=x^2+2(a+1)x+1$.

\begin{enumerate}[label=(\arabic*)]
\item 已知对于任意实数 $x$，$\sqrt{f(x)}$ 均有意义，求实数 $a$ 的取值范围；
\item 设 $A=\{x\mid f(x)+a^2-2=0\}$，$B=\{-4,0\}$，$A\cup B=B$，求实数 $a$ 的取值范围.
\end{enumerate}

\ifanswer
\solbegin
\stepm{(1)}{$\sqrt{f(x)}$ 对任意实数 $x$ 都有意义，}
\step{即 $f(x)=x^2+2(a+1)x+1\geqslant0$ 对任意 $x\in\R$ 恒成立，}
\step{所以 $\Delta=4(a+1)^2-4\times1\times1\leqslant0$，即 $(a+1)^2\leqslant1$，解得 $-2\leqslant a\leqslant0$，}
\step{故 $a$ 的取值范围是 $[-2,0]$.}
\stepm{(2)}{由 $f(x)+a^2-2=0$ 得 $x^2+2(a+1)x+a^2-1=0$.}
\step{由 $A\cup B=B$ 得 $A\subseteq B=\{-4,0\}$，即该方程无实根，}
\step{或它的实根都在 $\{-4,0\}$ 中.}
\step{当方程无实根时，$\Delta=4(a+1)^2-4(a^2-1)=8a+8<0$，解得 $a<-1$；}
\step{当方程有实根时，$a\geqslant-1$，设两根为 $x_1$、$x_2$，}
\step{由根与系数的关系得 $x_1+x_2=-2(a+1)$、$x_1x_2=a^2-1$，}
\step{此时两根只能在 $\{-4,0\}$ 中取值；}
\step{若两根为 0 与 0，由 $x_1x_2=0$ 得 $a=\pm1$，}
\step{再由 $x_1+x_2=0$ 得 $a=-1$；}
\step{若两根为 0 与 $-4$，由 $x_1x_2=0$ 得 $a=\pm1$，再由 $x_1+x_2=-4$ 得 $a=1$.}
\step{经检验，$a=-1$ 时 $A=\{0\}$，$a=1$ 时 $A=\{-4,0\}$，都满足 $A\subseteq B$.}
\step{综上，实数 $a$ 的取值范围为 $(-\infty,-1]\cup\{1\}$.}
\solend
\fi

\ansspace[8]

\ansneed{18cm}
\item 共享单车给市民出行带来了诸多便利，某公司购买了一批单车投放到某地给居民使用，据市场分析，这批单车的运营累计利润 $y$（单位：元）与运营天数 $x$（单位：天，$x$ 为正整数）满足关系：$y=-\dfrac12x^2+60x-800$.

\begin{enumerate}[label=(\arabic*)]
\item 要使这批单车的运营累计利润高于 800 元，求运营天数的取值范围；
\item 判断这批单车的日平均运营利润能否超过 20 元，并说明理由.
\end{enumerate}

\ifanswer
\solbegin
\stepm{(1)}{由题意得 $y=-\dfrac12x^2+60x-800>800$，}
\step{即 $x^2-120x+3200<0$，解得 $40<x<80$，}
\step{即运营天数的取值范围为 $\{x\mid40<x<80,x\in\N^*\}$.}
\stepm{(2)}{不能，理由如下：}
\step{这批单车的日平均运营利润为 $\dfrac yx=-\dfrac12x-\dfrac{800}x+60=-\dfrac12\left(x+\dfrac{1600}x\right)+60$，}
\step{因 $x>0$，则 $x+\dfrac{1600}x\geqslant2\sqrt{x\cdot\dfrac{1600}x}=80$，}
\step{当且仅当 $x=\dfrac{1600}x$，即 $x=40$ 时取等号，}
\step{此时 $\dfrac yx=-\dfrac12\left(x+\dfrac{1600}x\right)+60\leqslant-\dfrac12\times80+60=20$，}
\step{故这批单车的日平均运营利润不能超过 20 元.}
\solend
\fi

\ansspace[8]

% 第 20 题题干＋三小问与作答区 8 行须同页
\ansneed{20cm}
\item 设 $f(x)=ax^2+bx+c$($a\neq0$，$b$、$c\in\R$).

\begin{enumerate}[label=(\arabic*)]
\item 若 $a=b=c=1$，设 $A=\{(x,y)\mid y=f(x)\}$，$B=\{(x,y)\mid y=1\}$，用列举法表示 $A\cap B$；
% 原卷两个“若”连用（“若 a=1，…，若方程 f(x)=0 有两个大于 1 的根”），疑衍其一；照录（见文件头“原卷笔误/存疑”⑦）
\item 若 $a=1$，$b=2k-1$，$c=k^2$，若方程 $f(x)=0$ 有两个大于 1 的根，求实数 $k$ 的取值范围；
\item 若 $a$、$b$、$c\in\Z$，且 $f(0)$、$f(1)$ 均为奇数，求证：方程 $f(x)=0$ 无整数解.
\end{enumerate}

\ifanswer
\solbegin
\stepm{(1)}{当 $a=b=c=1$ 时，$f(x)=x^2+x+1$，}
\step{所以 $A=\{(x,y)\mid y=x^2+x+1\}$，又 $y=1$，}
\step{由 $x^2+x+1=1$ 得 $x^2+x=0$，解得 $x=0$ 或 $x=-1$，}
\step{所以 $A\cap B=\{(0,1),(-1,1)\}$.}
\stepm{(2)}{当 $a=1$、$b=2k-1$、$c=k^2$ 时，$f(x)=x^2+(2k-1)x+k^2$.}
\step{设两根为 $x_1$、$x_2$，}
\step{则“两根都大于 1”等价于 $\Delta>0$、$f(1)>0$、$-\dfrac b{2a}>1$.}
\step{由 $\Delta=(2k-1)^2-4k^2=1-4k>0$ 得 $k<\dfrac14$；}
\step{由 $f(1)=1+(2k-1)+k^2=k^2+2k=k(k+2)>0$ 得 $k<-2$ 或 $k>0$；}
\step{由 $-\dfrac{2k-1}2=\dfrac{1-2k}2>1$ 得 $k<-\dfrac12$.}
\step{三者取交集得 $k<-2$，即 $k$ 的取值范围是 $(-\infty,-2)$.}
\stepm{(3)}{设 $x_0\in\Z$ 是方程 $f(x)=0$ 的整数解.}
\step{因为 $x_0(x_0-1)$ 是两个相邻整数之积，恒为偶数，所以 $ax_0(x_0-1)$ 为偶数.}
\step{而 $f(x_0)=ax_0^2+bx_0+c=ax_0(x_0-1)+(a+b)x_0+c$，}
\step{所以 $f(x_0)$ 与 $(a+b)x_0+c$ 同奇偶.}
\step{若 $x_0$ 为偶数，则 $(a+b)x_0$ 为偶数，$f(x_0)$ 与 $c=f(0)$ 同奇偶，为奇数；}
\step{若 $x_0$ 为奇数，则 $(a+b)x_0$ 与 $a+b$ 同奇偶，}
\step{$f(x_0)$ 与 $(a+b)+c=f(1)$ 同奇偶，也为奇数.}
\step{综上，$f(x_0)$ 恒为奇数，不可能等于 0，与 $f(x_0)=0$ 矛盾，}
\step{所以方程 $f(x)=0$ 无整数解.}
\solend
\fi

\ansspace[8]

% 压轴题（题干＋三小问）与作答区 10 行须同页
\ansneed{20cm}
\item 已知 $X$ 为有限集，且任取 $x\in X$，均有 $x\in\Z$，称 $X$ 为非空有限整数集. 定义：$X_k=\{x+k\mid x\in X\}$，$X^2=\{y\mid y=x^2,x\in X\}$.

\begin{enumerate}[label=(\arabic*)]
\item 若 $X=\{-2,2\}$，分别写出 $X_2$ 和 $X^2$（用列举法表示）；
\item 若 $X=\{-1,0,1\}$，$X_1\subseteq Y$，$(Y_{-1})^2\subseteq X$，求集合 $Y$ 中所有的元素；
\item 对于非空有限整数集 $A$、$B$，若 $A_k\subseteq B$，$(B^2)_{-2}\subseteq A$，$B$ 中的元素均为非负整数，$k$ 为整数，求证：

① $A$、$B$ 中元素的个数相同；

② 若集合 $A$ 中的元素个数为奇数，则 $k$ 为偶数.
\end{enumerate}

\ifanswer
\solbegin
\stepm{(1)}{由题意得 $X_2=\{0,4\}$，$X^2=\{4\}$；}
\stepm{(2)}{由题意得 $X_1=\{0,1,2\}$，即 $\{0,1,2\}\subseteq Y$，$Y_{-1}=\{y-1\mid y\in Y\}$，}
\step{$(Y_{-1})^2=\{(y-1)^2\mid y\in Y\}\subseteq X=\{-1,0,1\}$，}
\step{所以 $(y-1)^2=0$ 或 $(y-1)^2=1$，解得 $y=1$ 或 $y=2$ 或 $y=0$，}
\step{即 $y\in\{0,1,2\}$，即 $Y\subseteq\{0,1,2\}$，所以 $Y=\{0,1,2\}$；}
% 原卷 (3) 解析首行“(3)”与“①”连排，本稿合并记作 \stepm{(3)①}
\stepm{(3)①}{因为 $A_k=\{x+k\mid x\in A\}$，函数 $g(x)=x+k$ 在全体整数集上严格增，}
\step{所以 $\forall x_1$、$x_2\in A$，$x_1\neq x_2$，$g(x_1)\neq g(x_2)$，所以 $|A_k|=|A|$，}
\step{因为 $A_k\subseteq B$，所以 $|A|\leqslant|B|$，}
\step{设 $b\in B$，$B$ 中元素均为非负整数，$\forall b_1\neq b_2\in B$，$b_1^2\neq b_2^2$，}
\step{故函数 $f(b)=b^2-2$ 在 $B$ 上单射，所以 $|(B^2)_{-2}|=|B|$，}
\step{又因为 $(B^2)_{-2}\subseteq A$，所以 $|B|\leqslant|A|$，所以 $|A|=|B|$.}
\stepm{②}{由①得 $A_k\subseteq B$ 且 $|A_k|=|A|=|B|$，$(B^2)_{-2}\subseteq A$ 且 $|(B^2)_{-2}|=|B|=|A|$，}
\step{所以 $A_k=B$，$(B^2)_{-2}=A$，}
\step{所以 $B=A_k=\{a+k\mid a\in A\}$，$A=(B^2)_{-2}=\{b^2-2\mid b\in B\}$，}
\step{所以 $A=\{(a+k)^2-2\mid a\in A\}$，}
\step{所以映射 $h(a)=(a+k)^2-2$ 是集合 $A$ 到自身的一个双射.}
\step{假设 $k$ 是奇数，当 $a$ 为偶数时，$a+k$ 为奇数，$(a+k)^2-2$ 为奇数，}
\step{当 $a$ 为奇数时，$a+k$ 为偶数，$(a+k)^2-2$ 为偶数，}
\step{所以此时映射 $h(a)=(a+k)^2-2$，将集合 $A$ 中的奇数映射为偶数，}
\step{将集合 $A$ 中的偶数映射为奇数，即集合 $A$ 中的奇数元素和偶数元素必须配对，}
\step{此时集合 $A$ 的元素个数则为偶数，}
\step{与题设中“集合 $A$ 中的元素个数为奇数”矛盾.}
\step{故假设不成立，所以 $k$ 是偶数.}
\solend
\fi

\ansspace*[10]

\end{enumerate}

\end{document}
"
%
% 原始材料：微信公众号“魔都数学加油站”文章，作者破晓，连载“27学年高一67”；
% 原文标题“27学年高一67——2026-2027年上海市华师大一附中高一上10月月考”；
% 原文链接：https://mp.weixin.qq.com/s/Kdhx07TVjLoGTcIlnRzgcQ。
% 页面标注发布时间 2026-10-18 20:33（上海）。
% 本次转录只据 `原图/高一67_华师大一附中_10月月考/` 的 13 张 PNG 逐页读图。
% 正文为 13 张 PNG 页（`原图/高一67_华师大一附中_10月月考/01..13.png`）：
% 第 01、04、13 页 2560×3620，第 02、03、05—12 页 4762×6735
% （**同一篇各页分辨率不同**），均按页序读取。13 页全为试卷正文页（卷首标题在第 1 页页顶，
% 卷末解答题第 21 题解析止于第 13 页），无封面页、目录页、二维码或宣传页。
% 粉色斜排“魔都数学加油站”为卷面水印，与公众号页脚一并未录入。
%
% 卷首标题：照卷面印作“2026--2027 年华师大一附中高一上 10 月月考”（卷面为黑体加粗一行），
% 年份区间按全稿惯例排 en dash。原卷只有一行标题、无副标题。
%
% 篇幅与题号：原卷自第 3 题起编号（第 1、2 题不在本卷内）。填空题 3—12（10 题）、
% 选择题 13—16（4 题）、解答题 17—21（5 题），共 19 题。本稿沿用原节与原题号，
% 只改排为全稿统一的大节样式（原卷小节标题为“一、填空题”“二、选择题”“三、解答题”）。
%
% 讲评版：原卷每题下直接印【解析】。**全卷 19 题中 18 题都印【解析】（无一处只印
% 【答案】）**，按全稿口径这类题不另提 \ans{}；各题解析末句的结论（第 13、14、16 题
% 的“故选 B／D／B”、各填空与解答题的“故/综上…”）照录在解析内。
% **另有 1 题答案印在题干括号内**——第 15 题（“(  D  )”，且该题无【答案】/【解析】段），
% 本稿按 高一37 第 13 题、高一63 第 14 题成例把括号留空、答案另提为 \ans{$D$}。
%
% 副标题：原卷未印副标题。本稿按“知识模块”口径标作“集合与常用逻辑用语、不等式”：
% 19 题中集合与常用逻辑用语 12 题（第 3、4、5、7、10、11、12、13、15、16、18、21 题：
% 补集、并集、子集、充分条件、容斥计数、子集计数、新定义二元运算、条件判断、
% 反证法假设、新定义生成集、集合关系与二次方程、新定义 X_k 与 X^2），
% 不等式 7 题（第 6、8、9、14、17、19、20 题：存在量词不等式、解集反求系数、
% 解集为空、不等式性质、解不等式（组）、二次不等式与基本不等式、二次方程根的分布），
% 两类均超过三分之一，故并标。
%
% 与原卷的排版差异：
%   ① 原卷每题下直接印【解析】，本稿按全稿惯例将解析置于答案版，试卷版保持干净；
%   ② 原卷未印副标题，本稿按“知识模块”口径补出；
%   ③ 原卷填空横线约四个字宽，本稿按全稿惯例统一排 \blankline[4.4em]；
%   ④ 原卷实数集排斜体 $R$（第 14 题题干、第 17 题 (1) 题干、第 18 题解析、第 20 题题干），
%      整数集排斜体 $Z$（第 20 题 (3) 题干与解析、第 21 题题干），正整数集排粗体
%      $N^*$（第 11 题解析、第 19 题解析），本稿按全稿惯例统一改用 $\R$、$\Z$、$\N^*$
%      （记号归一，粗体／斜体不作区分）；空集原卷作 $\varnothing$，照录；
%   ⑤ 原卷在数学式内多处用半角逗号（如第 8 题 $(1,2)$ 与各 cases、第 11 题解析
%      $P_n=\{1,2,\cdots,n\}$,$n\in N^*$、第 12 题解析 $u=(\cdots)$,$v=(\cdots)$ 与
%      “2020,个 1”、第 19 题解析 $\{x|40<x<80,x\in N^*\}$），均照录未强改；
%      描述法集合的竖线照原卷排 \mid；句末句点按全稿惯例排西文句点；
%   ⑥ 原卷补集记号为上划线 $\bar A$（第 3 题题干与解析、第 11 题题干），照录
%      （同 高一64/65 的 $\bar B$ 成例）；
%   ⑦ 第 12 题新定义的运算记号 $*$ 照录（$\alpha*\beta$、$u*u$、$u*v$）；
%   ⑧ 第 17—21 题的子问原卷各自另起一行，本稿按全稿惯例用 enumerate 的 (1)(2)(3)；
%      解析中的子问标号一律用 \stepm{(1)}{…}；第 5 题解析的①②、第 21 题 (3) 解析的①②
%      按全稿惯例作 \stepm{①}{…}、\stepm{②}{…}（第 21 题 (3) 首行原卷“(3)”与“①”连排，
%      本稿合并记作 \stepm{(3)①}{…}）；
%   ⑨ 选择题作答括号原卷排半角“(　　)”（第 13、14、16 题括号内空白、第 15 题括号内印 D），
%      本稿按全稿惯例排 \mbox{（\quad）}；第 15 题括号内的答案移入 \ans{}，见上；
%   ⑩ 选项行排照原卷：第 13、15 题两行两条，第 14、16 题一行四条；
%   ⑪ 每道解答题原卷未留作答空白，本稿按全稿惯例加 \ansspace（末题用 \ansspace*）；
%      第 20、21 题题干长，各用一处 \ansneed 让题干与作答区不拆页（\ansneed 只作用于试卷版）。
%
% 原卷笔误/存疑（均照录未改，正文对应处另有 % 注）：
%   ① 第 3 题解析末句印作“则 $\bar A=(-\infty,0)\cup[2,3)$.”：按题干
%      $U=\{x\mid 1/x\leqslant 1/2\}$（即 $x<0$ 或 $x\geqslant2$）与 $A=[3,+\infty)$，
%      $\bar A$ 应为 $(-\infty,0)\cup\{2\}$；且与本解析首句自求之 $U$ 相矛盾（$3\notin U$），
%      疑为笔误；照录未改。
%   ② 第 12 题题干 $\alpha*\beta$ 定义式括号不配对：第四组“$+\cdots+x_{2020}+y_{2020}
%      -x_{2020}y_{2020}$”缺左括号、末尾连排两个右括号“))”；照录未改。
%   ③ 第 12 题解析“所以…中有 2020,个 1”一行，“$\cdots$”后脱逗号（作
%      “$x_2,\cdots x_{2020}$”，第 12 行“而…”同式样），且“2020”后衍一半角逗号
%      （作“2020,个 1”）；照录未改。
%   ④ 第 16 题解析乘积链印作“$a_1a_2<a_1a_3<a_1a_4<a_1a_5<a_2a_3<a_3a_3<a_{14}a_5$”：
%      “$a_3a_3$”同足标自乘、与单调链不合，“$a_{14}a_5$”下标 14 与 $a_1$—$a_5$ 记号
%      不合（疑为 $a_2a_4$、$a_4a_5$ 之误）；照录未改。
%   ⑤ 第 5 题题干作“$Q\subset P$”（真子集记号），解析两处均作“$Q\subseteq P$”，
%      同一题前后记号不一；照录（本稿分别排 \subset、\subseteq）。
%   ⑥ 第 17 题 (1) 解析“当 $a=1$ 时，原不等式为 $(x+1)^2<0$ 此时…”与
%      “当 $0<a<1$ 时，$-\frac1a<-1$ 此时…”两处“此时”前均脱标点；照录未改。
%   ⑦ 第 20 题 (2) 题干“若 $a=1$，$b=2k-1$，$c=k^2$，若方程 $f(x)=0$ 有两个大于 1 的根”
%      两个“若”连用，疑衍其一；照录未改。
%
% 插图：无（全卷 13 页无任何几何图形、无“如图/右图/图中”字样）。
%
% 行数核对：逐题数 \step/\stepm 条数与原卷【解析】印刷行数，18 题逐一相符
% （第 15 题无解析段，只提 \ans{}），合计 169 条 / 169 行
% （Q3—Q12：2、3、4、3、1、3、5、3、13、15；Q13—Q16：2、3、0、25；
% Q17—Q21：19、16、9、20、23）。
%
\documentclass[a4paper,11pt]{article}
\usepackage{common}

\begin{document}

\examtitlest[3]{2026--2027 年华师大一附中高一上 10 月月考}{集合与常用逻辑用语、不等式}

\bigsec{一、填空题}

\begin{enumerate}
\setcounter{enumi}{2}

\item 设全集 $U=\left\{x\mid\dfrac1x\leqslant\dfrac12\right\}$，$A=[3,+\infty)$，则 $\bar A=$\blankline[4.4em].

\ifanswer
\solbegin
\step{$\dfrac1x\leqslant\dfrac12$，解得 $x<0$ 或 $x\geqslant2$，故 $U=(-\infty,0)\cup[2,+\infty)$，}
% 原卷末句作 $\bar A=(-\infty,0)\cup[2,3)$，按题干 $U$ 应为 $(-\infty,0)\cup\{2\}$，与本解析首句矛盾，疑笔误；照录（见文件头“原卷笔误/存疑”①）
\step{$A=[3,+\infty)$，则 $\bar A=(-\infty,0)\cup[2,3)$.}
\solend
\fi

\item 已知集合 $M=\{y\mid y=x^2-2x-1\}$，$N=\{y\mid y=x^2-4x+1\}$，则 $M\cup N=$\blankline[4.4em].

\ifanswer
\solbegin
\step{$M=\{y\mid y=x^2-2x-1\}=\{y\mid y=(x-1)^2-2\}=\{y\mid y\geqslant-2\}$，}
\step{$N=\{y\mid y=x^2-4x+1\}=\{y\mid y=(x-2)^2-3\}=\{y\mid y\geqslant-3\}$，}
\step{则 $M\cup N=\{y\mid y\geqslant-3\}$.}
\solend
\fi

\item 已知 $P=\{x\mid x^2+x-6=0\}$，$Q=\{x\mid mx-1=0\}$，且 $Q\subset P$，则实数 $m$ 的值\blankline[4.4em].

\ifanswer
\solbegin
\step{$P=\{2,-3\}$，$Q=\{x\mid mx=1\}$；}
\stepm{①}{$m=0$ 时，$Q=\varnothing$，满足 $Q\subseteq P$；}
\stepm{②}{$m\neq0$ 时，要使 $Q\subseteq P$，则 $\dfrac1m=2$ 或 $-3$，所以 $m=\dfrac12$ 或 $-\dfrac13$；}
\step{所以 $m=0$ 或 $\dfrac12$ 或 $-\dfrac13$.}
\solend
\fi

\item 已知“存在正实数 $x$，使得关于 $x$ 的不等式 $ax-1>0$ 成立”为真命题，则实数 $a$ 的取值范围是\blankline[4.4em].

\ifanswer
\solbegin
\step{存在正实数 $x$，使得关于 $x$ 的不等式 $ax-1>0$ 成立为真命题，}
\step{即存在正实数 $x$，使得 $a>\dfrac1x$，因为 $\dfrac1x>0$，所以 $a>0$，}
\step{即 $a$ 的范围为 $(0,+\infty)$.}
\solend
\fi

\item 设 $\alpha:3\leqslant x\leqslant5$，$\beta:2m-1\leqslant x<m+4$，若 $\alpha$ 是 $\beta$ 的充分条件，则实数 $m$ 的取值范围是\blankline[4.4em].

\ifanswer
\solbegin
\step{若 $\alpha$ 是 $\beta$ 的充分条件，则 $\begin{cases}2m-1\leqslant3\\ m+4>5\end{cases}$，解得 $1<m\leqslant2$.}
\solend
\fi

\item 对于实数 $a$、$b$、$c$，已知关于 $x$ 的不等式 $ax^2+bx+c>0$ 的解集是 $(1,2)$，则关于 $x$ 的不等式 $cx^2-bx+a\leqslant0$ 的解集是\blankline[4.4em].

\ifanswer
\solbegin
\step{关于 $x$ 的不等式 $ax^2+bx+c>0$ 的解集是 $(1,2)$，则 $\begin{cases}a<0\\ -\dfrac ba=3\\ \dfrac ca=2\end{cases}$，则 $\begin{cases}b=-3a\\ c=2a\end{cases}$，}
\step{则关于 $x$ 的不等式 $cx^2-bx+a=2ax^2+3ax+a\leqslant0$，则 $2x^2+3x+1\geqslant0$，}
\step{则 $x\leqslant-1$ 或 $x\geqslant-\dfrac12$，则解集为 $(-\infty,-1]\cup\left[-\dfrac12,+\infty\right)$.}
\solend
\fi

\item 已知关于 $x$ 的不等式 $kx^2+kx+k+2\leqslant0$ 的解集为 $\varnothing$，则实数 $k$ 的取值范围为\blankline[4.4em].

\ifanswer
\solbegin
\step{当 $k=0$ 时，不等式化为 $2\leqslant0$，解集为 $\varnothing$，符合题意；}
\step{当 $k\neq0$ 时，若关于 $x$ 的不等式 $kx^2+kx+k+2\leqslant0$ 的解集为 $\varnothing$，}
\step{则 $\begin{cases}k>0\\ \Delta=k^2-4k(k+2)<0\end{cases}$，化简得 $\begin{cases}k>0\\ -3k^2-8k<0\end{cases}$，即 $\begin{cases}k>0\\ k(3k+8)>0\end{cases}$，}
\step{解得 $k>0$.}
\step{综上所述，$k\geqslant0$，即实数 $k$ 的取值范围为 $[0,+\infty)$.}
\solend
\fi

\item 对班级若干名学生调查对 $A$、$B$ 两事件的态度，有如下结果：赞成 $A$ 的人数是全体的五分之三，其余的不赞成；赞成 $B$ 的比赞成 $A$ 多 3 人，其余的不赞成. 已知对 $A$、$B$ 都赞成的学生有 18 人，对 $A$、$B$ 都不赞成的学生数比对 $A$、$B$ 都赞成的学生数的三分之一多 1 人，则班级总人数是\blankline[4.4em]人.

\ifanswer
\solbegin
\step{设班级总人数是 $x$ 人，则赞成 $A$ 的人数是 $\dfrac35x$，赞成 $B$ 的人数是 $\dfrac35x+3$，}
\step{$A$、$B$ 都不赞成的人数为 $18\times\dfrac13+1=7$ 人，}
\step{由容斥原理得 $x=\left(\dfrac35x+\dfrac35x+3-18\right)+7$，解得 $x=40$，即班级总人数是 40 人.}
\solend
\fi

\item 设全集 $U=\{1,2,3,4,5,6,7,8,9,10,11,12\}$，若集合 $A$ 满足：① $A\subseteq U$；②若 $x\in A$，必有 $2x\notin A$；③若 $x\in\bar A$，必有 $2x\notin\bar A$，则这样的集合 $A$ 有\blankline[4.4em]个.

\ifanswer
\solbegin
\step{设集合 $P_n=\{1,2,\cdots,n\}$,$n\in\N^*$，}
\step{任取 $x\in P_n$，将 $x$ 除以 2，若商仍为偶数，再除以 2，$\cdots$，}
\step{经过 $k$ 次以后，商必为奇数，此时记商为 $m$，}
\step{于是 $x=m\cdot2^k$，其中 $m$ 为奇数，$k\in\N^*$.}
\step{由条件得，若 $m\in A$，则 $x\in A$ 等价于 $k$ 为偶数；}
\step{若 $m\notin A$，则 $x\in A$ 等价于 $k$ 为奇数；}
\step{于是 $x$ 是否属于 $A$ 由 $m$ 是否属于 $A$ 确定；}
\step{设 $Q_n$ 是 $P_n$ 中所有奇数的集合，}
\step{设 $f(n)$ 为集合 $P_n$ 中满足题干条件的集合 $A$ 的个数，}
\step{因此 $f(n)$ 等于 $Q_n$ 的子集个数.}
\step{当 $n$ 为偶数（或奇数）时，$P_n$ 中奇数的个数是 $\dfrac n2$（或 $\dfrac{n+1}2$），}
\step{所以 $f(n)=\begin{cases}2^{\frac n2},n\text{为偶数}\\[2pt] 2^{\frac{n+1}2},n\text{为奇数}\end{cases}$，所以 $P_{12}=\{1,2,3,4,5,6,7,8,9,10,11,12\}$，}
\step{即同时满足三个条件的集合 $A$ 的个数为 $2^{\frac{12}2}=2^6=64$.}
\solend
\fi

\item 设集合 $\Omega=\{(x_1,x_2,\cdots,x_{2020})\mid x_i\in\{0,1\},i=1,2,\cdots,2020\}$，对于 $\Omega$ 中的任意两个元素 $\alpha=(x_1,x_2,\cdots,x_{2020})$，$\beta=(y_1,y_2,\cdots,y_{2020})$，记
% 原卷此式第四组缺左括号、末尾连排两个右括号“))”，括号不配对；照录（见文件头“原卷笔误/存疑”②）
$\alpha*\beta=(x_1+y_1-x_1y_1)+(x_2+y_2-x_2y_2)+\cdots+x_{2020}+y_{2020}-x_{2020}y_{2020}))$，设 $u$、$v\in\Omega$，若 $u*u+v*v=2020$，则 $u*v$ 的最大值和最小值之和为\blankline[4.4em].

\ifanswer
\solbegin
\step{设 $u=(x_1,x_2,\cdots,x_{2020})$,$v=(y_1,y_2,\cdots,y_{2020})$，}
\step{因为 $x_i\in\{0,1\}$，所以 $x_i(x_i-1)=0$，}
\step{所以 $u*u=(x_1+x_1-x_1^2)+(x_2+x_2-x_2^2)+\cdots+(x_{2020}+x_{2020}-x_{2020}^2)$}
\step{$=x_1+x_2+\cdots+x_{2020}$，}
\step{$v*v=(y_1+y_1-y_1^2)+(y_2+y_2-y_2^2)+\cdots+(y_{2020}+y_{2020}-y_{2020}^2)$}
\step{$=y_1+y_2+\cdots+y_{2020}$，}
\step{由 $u*u+v*v=2020$，得 $x_1+x_2+\cdots+x_{2020}+y_1+y_2+\cdots+y_{2020}=2020$，}
% 原卷此行“$\cdots$”后脱逗号、 “2020”后衍一半角逗号（“2020,个 1”）；照录（见文件头“原卷笔误/存疑”③）
\step{所以 $x_1,x_2,\cdots x_{2020},y_1,y_2,\cdots y_{2020}$ 中有 2020,个 1，2020 个 0，}
\step{$u*v=(x_1+y_1-x_1y_1)+(x_2+y_2-x_2y_2)+\cdots+(x_{2020}+y_{2020}-x_{2020}y_{2020})$}
\step{$=2020-(x_1y_1+x_2y_2+\cdots+x_{2020}y_{2020})$，}
\step{注意到只有 $x_i=y_i=1$ 时，$x_iy_i=1$，否则 $x_iy_i=0$，}
% 此行“$\cdots$”后同样脱逗号，与上一处同式样；照录（见文件头“原卷笔误/存疑”③）
\step{而 $x_1,x_2,\cdots x_{2020},y_1,y_2,\cdots y_{2020}$ 中有 2020 个 1，2020 个 0，}
\step{所以满足 $x_iy_i=1$ 的最多有 1010 个，所以 $u*v\geqslant2020-1010=1010$，}
\step{满足 $x_iy_i=0$ 的最多有 2020 个，所以 $u*v\leqslant2020-0=2020$，}
\step{则 $u*v$ 的最大值和最小值之和为 $2020+1010=3030$.}
\solend
\fi

\end{enumerate}

\bigsec{二、选择题}

\begin{enumerate}
\setcounter{enumi}{12}

\item 实数 $x$ 与 $y$，“$x+y>2$”的\mbox{（\quad）}条件是“$x>1$ 或 $y>1$”.

\keepwithprev
A. 充分非必要\qquad B. 必要非充分

C. 充要\qquad D. 既非充分又非必要

\ifanswer
\solbegin
\step{当 $x>1$ 或 $y>1$ 时，$x+y>2$ 不一定成立，充分性不成立；}
\step{当 $x+y>2$ 时，$x>1$ 或 $y>1$ 一定成立，必要性成立. 故选 $B$.}
\solend
\fi

\item 设 $a,b\in\R$，且 $a<0$，$-1<b<0$，那么下列不等式成立的是\mbox{（\quad）}

\keepwithprev
A. $a>ab>ab^2$\qquad B. $ab^2>ab>a$\qquad C. $ab>a>ab^2$\qquad D. $ab>ab^2>a$

\ifanswer
\solbegin
\step{因为 $-1<b<0$，$a<0$，所以 $ab>0$，$0<b^2<1$.}
\step{所以 $ab-ab^2=ab(1-b)>0$，$ab^2-a=a(b^2-1)>0$，}
\step{所以 $ab>ab^2>a$. 故选 $D$.}
\solend
\fi

% 第 15 题答案原卷印在题干括号内（“(  D  )”），且无【答案】/【解析】段；
% 按 高一37 第 13 题、高一63 第 14 题成例括号留空、答案另提 \ans{$D$}
\item 用反证法证明：“平面四边形中至少有一个内角不超过 90°”时，下列“假设”中正确的是\mbox{（\quad）}

\keepwithprev
A. 假设有两个内角超过 90°\qquad B. 假设有三个内角超过 90°

C. 假设至多有两个内角超过 90°\qquad D. 假设四个内角都超过 90°

\ifanswer
\ans{$D$}
\fi

\item 设 $A$ 为非空实数集，集合 $B=\{xy\mid x\in A,y\in A,x\neq y\}$，则对于以下两个结论说法正确的是\mbox{（\quad）}

结论①：若 $A$ 是五元集，则 $B$ 中的元素最多 10 个；

结论②：存在四元集 $A$，使得 $B=\{2,4,5,8,10,16\}$.

\keepwithprev
A. ①对②对\qquad B. ①对②错\qquad C. ①错②对\qquad D. ①错②错

\ifanswer
\solbegin
\step{对于①，想让元素个数达到最大值 10，则须避免产生符号不同的乘积，}
\step{否则正负分开计数容易导致总数受限，且容易产生相等值，}
\step{也最好避免 0 的出现，}
\step{故设 $A=\{a_1,a_2,a_3,a_4,a_5\}$，$0<a_1<a_2<a_3<a_4<a_5$.}
% 原卷此链“$a_3a_3$”同足标自乘、“$a_{14}a_5$”下标 14 均与文意不合，疑为 $a_2a_4$、$a_4a_5$ 之误；照录（见文件头“原卷笔误/存疑”④）
\step{因为 $a_1a_2<a_1a_3<a_1a_4<a_1a_5<a_2a_3<a_3a_3<a_{14}a_5$，$a_2a_3<a_2a_4<a_3a_4$，}
\step{所以 $B$ 中元素个数小于等于 10 个，}
\step{如设 $A=\{1,2,3,5,7\}$，则 $B=\{2,3,5,6,7,10,14,15,21,35\}$，}
\step{所以集合 $B$ 中元素个数的最大值为 10 个，故①正确；}
\step{对于②，由于集合 $B$ 中没有 0，所以集合 $A$ 中没有 0；}
\step{设集合 $A=\{a,b,c,d\}$，且 $a<b<c<d$，}
\step{当 $a<0<b<c<d$ 时，$a\times b<0$，$a\times c<0$，$a\times d<0$，}
\step{故集合 $B$ 至少有 3 个负数，与 $B=\{2,4,5,8,10,16\}$ 矛盾；}
\step{当 $a<b<0<c<d$ 时，$a\times b>0$，$c\times d>0$，}
\step{$a\times c<0$，$a\times d<0$，$b\times c<0$，$b\times d<0$，}
\step{故集合 $B$ 至少有 4 个负数，与 $B=\{2,4,5,8,10,16\}$ 矛盾；}
\step{当 $a<b<c<0<d$ 时，$a\times b>0$，$a\times c>0$，$b\times c>0$，}
\step{$a\times d<0$，$b\times d<0$，$c\times d<0$，}
\step{故集合 $B$ 至少有 3 个负数，与 $B=\{2,4,5,8,10,16\}$ 矛盾；}
\step{当 $0<a<b<c<d$ 时，则集合 $B=\{ab,ac,ad,bc,bd,cd\}$，}
\step{则必有 $ab=2$，$cd=16$，其 4 个正实数的乘积 $abcd=32$，}
\step{又 $ad=5$，$bc=8$ 或 $ad=8$，$bc=5$，其 4 个正实数的乘积 $abcd=40$，矛盾；}
\step{当 $a<b<c<d<0$ 时，与 $0<a<b<c<d$ 时的证明方法相同，}
\step{得集合 $A$ 也不存在；}
\step{综上不存在四元集 $A$，使其生成集 $B=\{2,4,5,8,10,16\}$，故②错误.}
\step{故选 $B$.}
\solend
\fi

\end{enumerate}

\bigsec{三、解答题}

\begin{enumerate}
\setcounter{enumi}{16}

\item 解下列不等式（组）

\begin{enumerate}[label=(\arabic*)]
\item 解不等式：$(ax+1)(x+1)<0(a\in\R)$.
\item 解不等式组：$\begin{cases}\dfrac{2x+4}{x+4}\geqslant1\\[4pt] (x+1)(2-|x|)\geqslant0\end{cases}$.
\end{enumerate}

\ifanswer
\solbegin
\stepm{(1)}{当 $a=0$ 时，原不等式为 $x+1<0$，即 $x<-1$，}
\step{所以此时不等式的解集为 $\{x\mid x<-1\}$；}
\step{当 $a\neq0$ 时，令 $(ax+1)(x+1)=0$，解得 $x_1=-\dfrac1a$,$x_2=-1$，}
\step{当 $a>1$ 时，$-\dfrac1a>-1$，此时不等式的解集为 $\left\{x\mid-1<x<-\dfrac1a\right\}$；}
% 以下两行“此时”前原卷均无标点；照录（见文件头“原卷笔误/存疑”⑥）
\step{当 $a=1$ 时，原不等式为 $(x+1)^2<0$ 此时不等式的解集为 $\varnothing$；}
\step{当 $0<a<1$ 时，$-\dfrac1a<-1$ 此时不等式的解集为 $\left\{x\mid-\dfrac1a<x<-1\right\}$；}
\step{当 $a<0$ 时，$-\dfrac1a>-1$，此时不等式的解集为 $\left\{x\mid x<-1\ \text{或}\ x>-\dfrac1a\right\}$.}
\step{综上所述，当 $a=0$ 时，$\{x\mid x<-1\}$；}
\step{当 $a>1$ 时，$\left\{x\mid-1<x<-\dfrac1a\right\}$；}
\step{当 $a=1$ 时，$\varnothing$；}
\step{当 $0<a<1$ 时，$\left\{x\mid-\dfrac1a<x<-1\right\}$；}
\step{当 $a<0$ 时，$\left\{x\mid x<-1\ \text{或}\ x>-\dfrac1a\right\}$.}
\stepm{(2)}{不等式 $\dfrac{2x+4}{x+4}\geqslant1$ 可化为 $\dfrac{2x+4}{x+4}-1\geqslant0$，即 $\dfrac x{x+4}\geqslant0$，}
\step{等价于 $\begin{cases}x(x+4)\geqslant0\\ x+4\neq0\end{cases}$，解得 $x<-4$ 或 $x\geqslant0$，}
\step{不等式 $(x+1)(2-|x|)\geqslant0$ 等价于 $\begin{cases}(x+1)(2-x)\geqslant0,x\geqslant0\\ (x+1)(2+x)\geqslant0,x<0\end{cases}$，}
\step{解不等式 $(x+1)(2-x)\geqslant0$，$x\geqslant0$，得 $0\leqslant x\leqslant2$，}
\step{解不等式 $(x+1)(2+x)\geqslant0$，$x<0$，得 $x\leqslant-2$ 或 $-1\leqslant x<0$，}
\step{所以 $(x+1)(2-|x|)\geqslant0$ 的解集为 $\{x\mid x\leqslant-2\ \text{或}\ -1\leqslant x\leqslant2\}$，}
\step{综上所述，$\begin{cases}\dfrac{2x+4}{x+4}\geqslant1\\[4pt] (x+1)(2-|x|)\geqslant0\end{cases}$ 的解集为 $\{x\mid x<-4\ \text{或}\ 0\leqslant x\leqslant2\}$.}
\solend
\fi

\ansspace[8]

% 第 18 题题干＋两小问与作答区 6 行须同页
\ansneed{20cm}
\item 已知 $a$ 为实数，设 $f(x)=x^2+2(a+1)x+1$.

\begin{enumerate}[label=(\arabic*)]
\item 已知对于任意实数 $x$，$\sqrt{f(x)}$ 均有意义，求实数 $a$ 的取值范围；
\item 设 $A=\{x\mid f(x)+a^2-2=0\}$，$B=\{-4,0\}$，$A\cup B=B$，求实数 $a$ 的取值范围.
\end{enumerate}

\ifanswer
\solbegin
\stepm{(1)}{$\sqrt{f(x)}$ 对任意实数 $x$ 都有意义，}
\step{即 $f(x)=x^2+2(a+1)x+1\geqslant0$ 对任意 $x\in\R$ 恒成立，}
\step{所以 $\Delta=4(a+1)^2-4\times1\times1\leqslant0$，即 $(a+1)^2\leqslant1$，解得 $-2\leqslant a\leqslant0$，}
\step{故 $a$ 的取值范围是 $[-2,0]$.}
\stepm{(2)}{由 $f(x)+a^2-2=0$ 得 $x^2+2(a+1)x+a^2-1=0$.}
\step{由 $A\cup B=B$ 得 $A\subseteq B=\{-4,0\}$，即该方程无实根，}
\step{或它的实根都在 $\{-4,0\}$ 中.}
\step{当方程无实根时，$\Delta=4(a+1)^2-4(a^2-1)=8a+8<0$，解得 $a<-1$；}
\step{当方程有实根时，$a\geqslant-1$，设两根为 $x_1$、$x_2$，}
\step{由根与系数的关系得 $x_1+x_2=-2(a+1)$、$x_1x_2=a^2-1$，}
\step{此时两根只能在 $\{-4,0\}$ 中取值；}
\step{若两根为 0 与 0，由 $x_1x_2=0$ 得 $a=\pm1$，}
\step{再由 $x_1+x_2=0$ 得 $a=-1$；}
\step{若两根为 0 与 $-4$，由 $x_1x_2=0$ 得 $a=\pm1$，再由 $x_1+x_2=-4$ 得 $a=1$.}
\step{经检验，$a=-1$ 时 $A=\{0\}$，$a=1$ 时 $A=\{-4,0\}$，都满足 $A\subseteq B$.}
\step{综上，实数 $a$ 的取值范围为 $(-\infty,-1]\cup\{1\}$.}
\solend
\fi

\ansspace[8]

\ansneed{18cm}
\item 共享单车给市民出行带来了诸多便利，某公司购买了一批单车投放到某地给居民使用，据市场分析，这批单车的运营累计利润 $y$（单位：元）与运营天数 $x$（单位：天，$x$ 为正整数）满足关系：$y=-\dfrac12x^2+60x-800$.

\begin{enumerate}[label=(\arabic*)]
\item 要使这批单车的运营累计利润高于 800 元，求运营天数的取值范围；
\item 判断这批单车的日平均运营利润能否超过 20 元，并说明理由.
\end{enumerate}

\ifanswer
\solbegin
\stepm{(1)}{由题意得 $y=-\dfrac12x^2+60x-800>800$，}
\step{即 $x^2-120x+3200<0$，解得 $40<x<80$，}
\step{即运营天数的取值范围为 $\{x\mid40<x<80,x\in\N^*\}$.}
\stepm{(2)}{不能，理由如下：}
\step{这批单车的日平均运营利润为 $\dfrac yx=-\dfrac12x-\dfrac{800}x+60=-\dfrac12\left(x+\dfrac{1600}x\right)+60$，}
\step{因 $x>0$，则 $x+\dfrac{1600}x\geqslant2\sqrt{x\cdot\dfrac{1600}x}=80$，}
\step{当且仅当 $x=\dfrac{1600}x$，即 $x=40$ 时取等号，}
\step{此时 $\dfrac yx=-\dfrac12\left(x+\dfrac{1600}x\right)+60\leqslant-\dfrac12\times80+60=20$，}
\step{故这批单车的日平均运营利润不能超过 20 元.}
\solend
\fi

\ansspace[8]

% 第 20 题题干＋三小问与作答区 8 行须同页
\ansneed{20cm}
\item 设 $f(x)=ax^2+bx+c$($a\neq0$，$b$、$c\in\R$).

\begin{enumerate}[label=(\arabic*)]
\item 若 $a=b=c=1$，设 $A=\{(x,y)\mid y=f(x)\}$，$B=\{(x,y)\mid y=1\}$，用列举法表示 $A\cap B$；
% 原卷两个“若”连用（“若 a=1，…，若方程 f(x)=0 有两个大于 1 的根”），疑衍其一；照录（见文件头“原卷笔误/存疑”⑦）
\item 若 $a=1$，$b=2k-1$，$c=k^2$，若方程 $f(x)=0$ 有两个大于 1 的根，求实数 $k$ 的取值范围；
\item 若 $a$、$b$、$c\in\Z$，且 $f(0)$、$f(1)$ 均为奇数，求证：方程 $f(x)=0$ 无整数解.
\end{enumerate}

\ifanswer
\solbegin
\stepm{(1)}{当 $a=b=c=1$ 时，$f(x)=x^2+x+1$，}
\step{所以 $A=\{(x,y)\mid y=x^2+x+1\}$，又 $y=1$，}
\step{由 $x^2+x+1=1$ 得 $x^2+x=0$，解得 $x=0$ 或 $x=-1$，}
\step{所以 $A\cap B=\{(0,1),(-1,1)\}$.}
\stepm{(2)}{当 $a=1$、$b=2k-1$、$c=k^2$ 时，$f(x)=x^2+(2k-1)x+k^2$.}
\step{设两根为 $x_1$、$x_2$，}
\step{则“两根都大于 1”等价于 $\Delta>0$、$f(1)>0$、$-\dfrac b{2a}>1$.}
\step{由 $\Delta=(2k-1)^2-4k^2=1-4k>0$ 得 $k<\dfrac14$；}
\step{由 $f(1)=1+(2k-1)+k^2=k^2+2k=k(k+2)>0$ 得 $k<-2$ 或 $k>0$；}
\step{由 $-\dfrac{2k-1}2=\dfrac{1-2k}2>1$ 得 $k<-\dfrac12$.}
\step{三者取交集得 $k<-2$，即 $k$ 的取值范围是 $(-\infty,-2)$.}
\stepm{(3)}{设 $x_0\in\Z$ 是方程 $f(x)=0$ 的整数解.}
\step{因为 $x_0(x_0-1)$ 是两个相邻整数之积，恒为偶数，所以 $ax_0(x_0-1)$ 为偶数.}
\step{而 $f(x_0)=ax_0^2+bx_0+c=ax_0(x_0-1)+(a+b)x_0+c$，}
\step{所以 $f(x_0)$ 与 $(a+b)x_0+c$ 同奇偶.}
\step{若 $x_0$ 为偶数，则 $(a+b)x_0$ 为偶数，$f(x_0)$ 与 $c=f(0)$ 同奇偶，为奇数；}
\step{若 $x_0$ 为奇数，则 $(a+b)x_0$ 与 $a+b$ 同奇偶，}
\step{$f(x_0)$ 与 $(a+b)+c=f(1)$ 同奇偶，也为奇数.}
\step{综上，$f(x_0)$ 恒为奇数，不可能等于 0，与 $f(x_0)=0$ 矛盾，}
\step{所以方程 $f(x)=0$ 无整数解.}
\solend
\fi

\ansspace[8]

% 压轴题（题干＋三小问）与作答区 10 行须同页
\ansneed{20cm}
\item 已知 $X$ 为有限集，且任取 $x\in X$，均有 $x\in\Z$，称 $X$ 为非空有限整数集. 定义：$X_k=\{x+k\mid x\in X\}$，$X^2=\{y\mid y=x^2,x\in X\}$.

\begin{enumerate}[label=(\arabic*)]
\item 若 $X=\{-2,2\}$，分别写出 $X_2$ 和 $X^2$（用列举法表示）；
\item 若 $X=\{-1,0,1\}$，$X_1\subseteq Y$，$(Y_{-1})^2\subseteq X$，求集合 $Y$ 中所有的元素；
\item 对于非空有限整数集 $A$、$B$，若 $A_k\subseteq B$，$(B^2)_{-2}\subseteq A$，$B$ 中的元素均为非负整数，$k$ 为整数，求证：

① $A$、$B$ 中元素的个数相同；

② 若集合 $A$ 中的元素个数为奇数，则 $k$ 为偶数.
\end{enumerate}

\ifanswer
\solbegin
\stepm{(1)}{由题意得 $X_2=\{0,4\}$，$X^2=\{4\}$；}
\stepm{(2)}{由题意得 $X_1=\{0,1,2\}$，即 $\{0,1,2\}\subseteq Y$，$Y_{-1}=\{y-1\mid y\in Y\}$，}
\step{$(Y_{-1})^2=\{(y-1)^2\mid y\in Y\}\subseteq X=\{-1,0,1\}$，}
\step{所以 $(y-1)^2=0$ 或 $(y-1)^2=1$，解得 $y=1$ 或 $y=2$ 或 $y=0$，}
\step{即 $y\in\{0,1,2\}$，即 $Y\subseteq\{0,1,2\}$，所以 $Y=\{0,1,2\}$；}
% 原卷 (3) 解析首行“(3)”与“①”连排，本稿合并记作 \stepm{(3)①}
\stepm{(3)①}{因为 $A_k=\{x+k\mid x\in A\}$，函数 $g(x)=x+k$ 在全体整数集上严格增，}
\step{所以 $\forall x_1$、$x_2\in A$，$x_1\neq x_2$，$g(x_1)\neq g(x_2)$，所以 $|A_k|=|A|$，}
\step{因为 $A_k\subseteq B$，所以 $|A|\leqslant|B|$，}
\step{设 $b\in B$，$B$ 中元素均为非负整数，$\forall b_1\neq b_2\in B$，$b_1^2\neq b_2^2$，}
\step{故函数 $f(b)=b^2-2$ 在 $B$ 上单射，所以 $|(B^2)_{-2}|=|B|$，}
\step{又因为 $(B^2)_{-2}\subseteq A$，所以 $|B|\leqslant|A|$，所以 $|A|=|B|$.}
\stepm{②}{由①得 $A_k\subseteq B$ 且 $|A_k|=|A|=|B|$，$(B^2)_{-2}\subseteq A$ 且 $|(B^2)_{-2}|=|B|=|A|$，}
\step{所以 $A_k=B$，$(B^2)_{-2}=A$，}
\step{所以 $B=A_k=\{a+k\mid a\in A\}$，$A=(B^2)_{-2}=\{b^2-2\mid b\in B\}$，}
\step{所以 $A=\{(a+k)^2-2\mid a\in A\}$，}
\step{所以映射 $h(a)=(a+k)^2-2$ 是集合 $A$ 到自身的一个双射.}
\step{假设 $k$ 是奇数，当 $a$ 为偶数时，$a+k$ 为奇数，$(a+k)^2-2$ 为奇数，}
\step{当 $a$ 为奇数时，$a+k$ 为偶数，$(a+k)^2-2$ 为偶数，}
\step{所以此时映射 $h(a)=(a+k)^2-2$，将集合 $A$ 中的奇数映射为偶数，}
\step{将集合 $A$ 中的偶数映射为奇数，即集合 $A$ 中的奇数元素和偶数元素必须配对，}
\step{此时集合 $A$ 的元素个数则为偶数，}
\step{与题设中“集合 $A$ 中的元素个数为奇数”矛盾.}
\step{故假设不成立，所以 $k$ 是偶数.}
\solend
\fi

\ansspace*[10]

\end{enumerate}

\end{document}
"
% 紧凑版：xelatex "\def\iscompact{1}% 高一67_华师大一附中_10月月考.tex
%   —— 高一上数学“2026-2027 年华师大一附中高一上 10 月月考”（讲评版）
% 试卷版：xelatex 高一67_华师大一附中_10月月考.tex
% 答案版：xelatex "\def\isanswer{1}% 高一67_华师大一附中_10月月考.tex
%   —— 高一上数学“2026-2027 年华师大一附中高一上 10 月月考”（讲评版）
% 试卷版：xelatex 高一67_华师大一附中_10月月考.tex
% 答案版：xelatex "\def\isanswer{1}\input{高一67_华师大一附中_10月月考.tex}"
% 紧凑版：xelatex "\def\iscompact{1}\input{高一67_华师大一附中_10月月考.tex}"
%
% 原始材料：微信公众号“魔都数学加油站”文章，作者破晓，连载“27学年高一67”；
% 原文标题“27学年高一67——2026-2027年上海市华师大一附中高一上10月月考”；
% 原文链接：https://mp.weixin.qq.com/s/Kdhx07TVjLoGTcIlnRzgcQ。
% 页面标注发布时间 2026-10-18 20:33（上海）。
% 本次转录只据 `原图/高一67_华师大一附中_10月月考/` 的 13 张 PNG 逐页读图。
% 正文为 13 张 PNG 页（`原图/高一67_华师大一附中_10月月考/01..13.png`）：
% 第 01、04、13 页 2560×3620，第 02、03、05—12 页 4762×6735
% （**同一篇各页分辨率不同**），均按页序读取。13 页全为试卷正文页（卷首标题在第 1 页页顶，
% 卷末解答题第 21 题解析止于第 13 页），无封面页、目录页、二维码或宣传页。
% 粉色斜排“魔都数学加油站”为卷面水印，与公众号页脚一并未录入。
%
% 卷首标题：照卷面印作“2026--2027 年华师大一附中高一上 10 月月考”（卷面为黑体加粗一行），
% 年份区间按全稿惯例排 en dash。原卷只有一行标题、无副标题。
%
% 篇幅与题号：原卷自第 3 题起编号（第 1、2 题不在本卷内）。填空题 3—12（10 题）、
% 选择题 13—16（4 题）、解答题 17—21（5 题），共 19 题。本稿沿用原节与原题号，
% 只改排为全稿统一的大节样式（原卷小节标题为“一、填空题”“二、选择题”“三、解答题”）。
%
% 讲评版：原卷每题下直接印【解析】。**全卷 19 题中 18 题都印【解析】（无一处只印
% 【答案】）**，按全稿口径这类题不另提 \ans{}；各题解析末句的结论（第 13、14、16 题
% 的“故选 B／D／B”、各填空与解答题的“故/综上…”）照录在解析内。
% **另有 1 题答案印在题干括号内**——第 15 题（“(  D  )”，且该题无【答案】/【解析】段），
% 本稿按 高一37 第 13 题、高一63 第 14 题成例把括号留空、答案另提为 \ans{$D$}。
%
% 副标题：原卷未印副标题。本稿按“知识模块”口径标作“集合与常用逻辑用语、不等式”：
% 19 题中集合与常用逻辑用语 12 题（第 3、4、5、7、10、11、12、13、15、16、18、21 题：
% 补集、并集、子集、充分条件、容斥计数、子集计数、新定义二元运算、条件判断、
% 反证法假设、新定义生成集、集合关系与二次方程、新定义 X_k 与 X^2），
% 不等式 7 题（第 6、8、9、14、17、19、20 题：存在量词不等式、解集反求系数、
% 解集为空、不等式性质、解不等式（组）、二次不等式与基本不等式、二次方程根的分布），
% 两类均超过三分之一，故并标。
%
% 与原卷的排版差异：
%   ① 原卷每题下直接印【解析】，本稿按全稿惯例将解析置于答案版，试卷版保持干净；
%   ② 原卷未印副标题，本稿按“知识模块”口径补出；
%   ③ 原卷填空横线约四个字宽，本稿按全稿惯例统一排 \blankline[4.4em]；
%   ④ 原卷实数集排斜体 $R$（第 14 题题干、第 17 题 (1) 题干、第 18 题解析、第 20 题题干），
%      整数集排斜体 $Z$（第 20 题 (3) 题干与解析、第 21 题题干），正整数集排粗体
%      $N^*$（第 11 题解析、第 19 题解析），本稿按全稿惯例统一改用 $\R$、$\Z$、$\N^*$
%      （记号归一，粗体／斜体不作区分）；空集原卷作 $\varnothing$，照录；
%   ⑤ 原卷在数学式内多处用半角逗号（如第 8 题 $(1,2)$ 与各 cases、第 11 题解析
%      $P_n=\{1,2,\cdots,n\}$,$n\in N^*$、第 12 题解析 $u=(\cdots)$,$v=(\cdots)$ 与
%      “2020,个 1”、第 19 题解析 $\{x|40<x<80,x\in N^*\}$），均照录未强改；
%      描述法集合的竖线照原卷排 \mid；句末句点按全稿惯例排西文句点；
%   ⑥ 原卷补集记号为上划线 $\bar A$（第 3 题题干与解析、第 11 题题干），照录
%      （同 高一64/65 的 $\bar B$ 成例）；
%   ⑦ 第 12 题新定义的运算记号 $*$ 照录（$\alpha*\beta$、$u*u$、$u*v$）；
%   ⑧ 第 17—21 题的子问原卷各自另起一行，本稿按全稿惯例用 enumerate 的 (1)(2)(3)；
%      解析中的子问标号一律用 \stepm{(1)}{…}；第 5 题解析的①②、第 21 题 (3) 解析的①②
%      按全稿惯例作 \stepm{①}{…}、\stepm{②}{…}（第 21 题 (3) 首行原卷“(3)”与“①”连排，
%      本稿合并记作 \stepm{(3)①}{…}）；
%   ⑨ 选择题作答括号原卷排半角“(　　)”（第 13、14、16 题括号内空白、第 15 题括号内印 D），
%      本稿按全稿惯例排 \mbox{（\quad）}；第 15 题括号内的答案移入 \ans{}，见上；
%   ⑩ 选项行排照原卷：第 13、15 题两行两条，第 14、16 题一行四条；
%   ⑪ 每道解答题原卷未留作答空白，本稿按全稿惯例加 \ansspace（末题用 \ansspace*）；
%      第 20、21 题题干长，各用一处 \ansneed 让题干与作答区不拆页（\ansneed 只作用于试卷版）。
%
% 原卷笔误/存疑（均照录未改，正文对应处另有 % 注）：
%   ① 第 3 题解析末句印作“则 $\bar A=(-\infty,0)\cup[2,3)$.”：按题干
%      $U=\{x\mid 1/x\leqslant 1/2\}$（即 $x<0$ 或 $x\geqslant2$）与 $A=[3,+\infty)$，
%      $\bar A$ 应为 $(-\infty,0)\cup\{2\}$；且与本解析首句自求之 $U$ 相矛盾（$3\notin U$），
%      疑为笔误；照录未改。
%   ② 第 12 题题干 $\alpha*\beta$ 定义式括号不配对：第四组“$+\cdots+x_{2020}+y_{2020}
%      -x_{2020}y_{2020}$”缺左括号、末尾连排两个右括号“))”；照录未改。
%   ③ 第 12 题解析“所以…中有 2020,个 1”一行，“$\cdots$”后脱逗号（作
%      “$x_2,\cdots x_{2020}$”，第 12 行“而…”同式样），且“2020”后衍一半角逗号
%      （作“2020,个 1”）；照录未改。
%   ④ 第 16 题解析乘积链印作“$a_1a_2<a_1a_3<a_1a_4<a_1a_5<a_2a_3<a_3a_3<a_{14}a_5$”：
%      “$a_3a_3$”同足标自乘、与单调链不合，“$a_{14}a_5$”下标 14 与 $a_1$—$a_5$ 记号
%      不合（疑为 $a_2a_4$、$a_4a_5$ 之误）；照录未改。
%   ⑤ 第 5 题题干作“$Q\subset P$”（真子集记号），解析两处均作“$Q\subseteq P$”，
%      同一题前后记号不一；照录（本稿分别排 \subset、\subseteq）。
%   ⑥ 第 17 题 (1) 解析“当 $a=1$ 时，原不等式为 $(x+1)^2<0$ 此时…”与
%      “当 $0<a<1$ 时，$-\frac1a<-1$ 此时…”两处“此时”前均脱标点；照录未改。
%   ⑦ 第 20 题 (2) 题干“若 $a=1$，$b=2k-1$，$c=k^2$，若方程 $f(x)=0$ 有两个大于 1 的根”
%      两个“若”连用，疑衍其一；照录未改。
%
% 插图：无（全卷 13 页无任何几何图形、无“如图/右图/图中”字样）。
%
% 行数核对：逐题数 \step/\stepm 条数与原卷【解析】印刷行数，18 题逐一相符
% （第 15 题无解析段，只提 \ans{}），合计 169 条 / 169 行
% （Q3—Q12：2、3、4、3、1、3、5、3、13、15；Q13—Q16：2、3、0、25；
% Q17—Q21：19、16、9、20、23）。
%
\documentclass[a4paper,11pt]{article}
\usepackage{common}

\begin{document}

\examtitlest[3]{2026--2027 年华师大一附中高一上 10 月月考}{集合与常用逻辑用语、不等式}

\bigsec{一、填空题}

\begin{enumerate}
\setcounter{enumi}{2}

\item 设全集 $U=\left\{x\mid\dfrac1x\leqslant\dfrac12\right\}$，$A=[3,+\infty)$，则 $\bar A=$\blankline[4.4em].

\ifanswer
\solbegin
\step{$\dfrac1x\leqslant\dfrac12$，解得 $x<0$ 或 $x\geqslant2$，故 $U=(-\infty,0)\cup[2,+\infty)$，}
% 原卷末句作 $\bar A=(-\infty,0)\cup[2,3)$，按题干 $U$ 应为 $(-\infty,0)\cup\{2\}$，与本解析首句矛盾，疑笔误；照录（见文件头“原卷笔误/存疑”①）
\step{$A=[3,+\infty)$，则 $\bar A=(-\infty,0)\cup[2,3)$.}
\solend
\fi

\item 已知集合 $M=\{y\mid y=x^2-2x-1\}$，$N=\{y\mid y=x^2-4x+1\}$，则 $M\cup N=$\blankline[4.4em].

\ifanswer
\solbegin
\step{$M=\{y\mid y=x^2-2x-1\}=\{y\mid y=(x-1)^2-2\}=\{y\mid y\geqslant-2\}$，}
\step{$N=\{y\mid y=x^2-4x+1\}=\{y\mid y=(x-2)^2-3\}=\{y\mid y\geqslant-3\}$，}
\step{则 $M\cup N=\{y\mid y\geqslant-3\}$.}
\solend
\fi

\item 已知 $P=\{x\mid x^2+x-6=0\}$，$Q=\{x\mid mx-1=0\}$，且 $Q\subset P$，则实数 $m$ 的值\blankline[4.4em].

\ifanswer
\solbegin
\step{$P=\{2,-3\}$，$Q=\{x\mid mx=1\}$；}
\stepm{①}{$m=0$ 时，$Q=\varnothing$，满足 $Q\subseteq P$；}
\stepm{②}{$m\neq0$ 时，要使 $Q\subseteq P$，则 $\dfrac1m=2$ 或 $-3$，所以 $m=\dfrac12$ 或 $-\dfrac13$；}
\step{所以 $m=0$ 或 $\dfrac12$ 或 $-\dfrac13$.}
\solend
\fi

\item 已知“存在正实数 $x$，使得关于 $x$ 的不等式 $ax-1>0$ 成立”为真命题，则实数 $a$ 的取值范围是\blankline[4.4em].

\ifanswer
\solbegin
\step{存在正实数 $x$，使得关于 $x$ 的不等式 $ax-1>0$ 成立为真命题，}
\step{即存在正实数 $x$，使得 $a>\dfrac1x$，因为 $\dfrac1x>0$，所以 $a>0$，}
\step{即 $a$ 的范围为 $(0,+\infty)$.}
\solend
\fi

\item 设 $\alpha:3\leqslant x\leqslant5$，$\beta:2m-1\leqslant x<m+4$，若 $\alpha$ 是 $\beta$ 的充分条件，则实数 $m$ 的取值范围是\blankline[4.4em].

\ifanswer
\solbegin
\step{若 $\alpha$ 是 $\beta$ 的充分条件，则 $\begin{cases}2m-1\leqslant3\\ m+4>5\end{cases}$，解得 $1<m\leqslant2$.}
\solend
\fi

\item 对于实数 $a$、$b$、$c$，已知关于 $x$ 的不等式 $ax^2+bx+c>0$ 的解集是 $(1,2)$，则关于 $x$ 的不等式 $cx^2-bx+a\leqslant0$ 的解集是\blankline[4.4em].

\ifanswer
\solbegin
\step{关于 $x$ 的不等式 $ax^2+bx+c>0$ 的解集是 $(1,2)$，则 $\begin{cases}a<0\\ -\dfrac ba=3\\ \dfrac ca=2\end{cases}$，则 $\begin{cases}b=-3a\\ c=2a\end{cases}$，}
\step{则关于 $x$ 的不等式 $cx^2-bx+a=2ax^2+3ax+a\leqslant0$，则 $2x^2+3x+1\geqslant0$，}
\step{则 $x\leqslant-1$ 或 $x\geqslant-\dfrac12$，则解集为 $(-\infty,-1]\cup\left[-\dfrac12,+\infty\right)$.}
\solend
\fi

\item 已知关于 $x$ 的不等式 $kx^2+kx+k+2\leqslant0$ 的解集为 $\varnothing$，则实数 $k$ 的取值范围为\blankline[4.4em].

\ifanswer
\solbegin
\step{当 $k=0$ 时，不等式化为 $2\leqslant0$，解集为 $\varnothing$，符合题意；}
\step{当 $k\neq0$ 时，若关于 $x$ 的不等式 $kx^2+kx+k+2\leqslant0$ 的解集为 $\varnothing$，}
\step{则 $\begin{cases}k>0\\ \Delta=k^2-4k(k+2)<0\end{cases}$，化简得 $\begin{cases}k>0\\ -3k^2-8k<0\end{cases}$，即 $\begin{cases}k>0\\ k(3k+8)>0\end{cases}$，}
\step{解得 $k>0$.}
\step{综上所述，$k\geqslant0$，即实数 $k$ 的取值范围为 $[0,+\infty)$.}
\solend
\fi

\item 对班级若干名学生调查对 $A$、$B$ 两事件的态度，有如下结果：赞成 $A$ 的人数是全体的五分之三，其余的不赞成；赞成 $B$ 的比赞成 $A$ 多 3 人，其余的不赞成. 已知对 $A$、$B$ 都赞成的学生有 18 人，对 $A$、$B$ 都不赞成的学生数比对 $A$、$B$ 都赞成的学生数的三分之一多 1 人，则班级总人数是\blankline[4.4em]人.

\ifanswer
\solbegin
\step{设班级总人数是 $x$ 人，则赞成 $A$ 的人数是 $\dfrac35x$，赞成 $B$ 的人数是 $\dfrac35x+3$，}
\step{$A$、$B$ 都不赞成的人数为 $18\times\dfrac13+1=7$ 人，}
\step{由容斥原理得 $x=\left(\dfrac35x+\dfrac35x+3-18\right)+7$，解得 $x=40$，即班级总人数是 40 人.}
\solend
\fi

\item 设全集 $U=\{1,2,3,4,5,6,7,8,9,10,11,12\}$，若集合 $A$ 满足：① $A\subseteq U$；②若 $x\in A$，必有 $2x\notin A$；③若 $x\in\bar A$，必有 $2x\notin\bar A$，则这样的集合 $A$ 有\blankline[4.4em]个.

\ifanswer
\solbegin
\step{设集合 $P_n=\{1,2,\cdots,n\}$,$n\in\N^*$，}
\step{任取 $x\in P_n$，将 $x$ 除以 2，若商仍为偶数，再除以 2，$\cdots$，}
\step{经过 $k$ 次以后，商必为奇数，此时记商为 $m$，}
\step{于是 $x=m\cdot2^k$，其中 $m$ 为奇数，$k\in\N^*$.}
\step{由条件得，若 $m\in A$，则 $x\in A$ 等价于 $k$ 为偶数；}
\step{若 $m\notin A$，则 $x\in A$ 等价于 $k$ 为奇数；}
\step{于是 $x$ 是否属于 $A$ 由 $m$ 是否属于 $A$ 确定；}
\step{设 $Q_n$ 是 $P_n$ 中所有奇数的集合，}
\step{设 $f(n)$ 为集合 $P_n$ 中满足题干条件的集合 $A$ 的个数，}
\step{因此 $f(n)$ 等于 $Q_n$ 的子集个数.}
\step{当 $n$ 为偶数（或奇数）时，$P_n$ 中奇数的个数是 $\dfrac n2$（或 $\dfrac{n+1}2$），}
\step{所以 $f(n)=\begin{cases}2^{\frac n2},n\text{为偶数}\\[2pt] 2^{\frac{n+1}2},n\text{为奇数}\end{cases}$，所以 $P_{12}=\{1,2,3,4,5,6,7,8,9,10,11,12\}$，}
\step{即同时满足三个条件的集合 $A$ 的个数为 $2^{\frac{12}2}=2^6=64$.}
\solend
\fi

\item 设集合 $\Omega=\{(x_1,x_2,\cdots,x_{2020})\mid x_i\in\{0,1\},i=1,2,\cdots,2020\}$，对于 $\Omega$ 中的任意两个元素 $\alpha=(x_1,x_2,\cdots,x_{2020})$，$\beta=(y_1,y_2,\cdots,y_{2020})$，记
% 原卷此式第四组缺左括号、末尾连排两个右括号“))”，括号不配对；照录（见文件头“原卷笔误/存疑”②）
$\alpha*\beta=(x_1+y_1-x_1y_1)+(x_2+y_2-x_2y_2)+\cdots+x_{2020}+y_{2020}-x_{2020}y_{2020}))$，设 $u$、$v\in\Omega$，若 $u*u+v*v=2020$，则 $u*v$ 的最大值和最小值之和为\blankline[4.4em].

\ifanswer
\solbegin
\step{设 $u=(x_1,x_2,\cdots,x_{2020})$,$v=(y_1,y_2,\cdots,y_{2020})$，}
\step{因为 $x_i\in\{0,1\}$，所以 $x_i(x_i-1)=0$，}
\step{所以 $u*u=(x_1+x_1-x_1^2)+(x_2+x_2-x_2^2)+\cdots+(x_{2020}+x_{2020}-x_{2020}^2)$}
\step{$=x_1+x_2+\cdots+x_{2020}$，}
\step{$v*v=(y_1+y_1-y_1^2)+(y_2+y_2-y_2^2)+\cdots+(y_{2020}+y_{2020}-y_{2020}^2)$}
\step{$=y_1+y_2+\cdots+y_{2020}$，}
\step{由 $u*u+v*v=2020$，得 $x_1+x_2+\cdots+x_{2020}+y_1+y_2+\cdots+y_{2020}=2020$，}
% 原卷此行“$\cdots$”后脱逗号、 “2020”后衍一半角逗号（“2020,个 1”）；照录（见文件头“原卷笔误/存疑”③）
\step{所以 $x_1,x_2,\cdots x_{2020},y_1,y_2,\cdots y_{2020}$ 中有 2020,个 1，2020 个 0，}
\step{$u*v=(x_1+y_1-x_1y_1)+(x_2+y_2-x_2y_2)+\cdots+(x_{2020}+y_{2020}-x_{2020}y_{2020})$}
\step{$=2020-(x_1y_1+x_2y_2+\cdots+x_{2020}y_{2020})$，}
\step{注意到只有 $x_i=y_i=1$ 时，$x_iy_i=1$，否则 $x_iy_i=0$，}
% 此行“$\cdots$”后同样脱逗号，与上一处同式样；照录（见文件头“原卷笔误/存疑”③）
\step{而 $x_1,x_2,\cdots x_{2020},y_1,y_2,\cdots y_{2020}$ 中有 2020 个 1，2020 个 0，}
\step{所以满足 $x_iy_i=1$ 的最多有 1010 个，所以 $u*v\geqslant2020-1010=1010$，}
\step{满足 $x_iy_i=0$ 的最多有 2020 个，所以 $u*v\leqslant2020-0=2020$，}
\step{则 $u*v$ 的最大值和最小值之和为 $2020+1010=3030$.}
\solend
\fi

\end{enumerate}

\bigsec{二、选择题}

\begin{enumerate}
\setcounter{enumi}{12}

\item 实数 $x$ 与 $y$，“$x+y>2$”的\mbox{（\quad）}条件是“$x>1$ 或 $y>1$”.

\keepwithprev
A. 充分非必要\qquad B. 必要非充分

C. 充要\qquad D. 既非充分又非必要

\ifanswer
\solbegin
\step{当 $x>1$ 或 $y>1$ 时，$x+y>2$ 不一定成立，充分性不成立；}
\step{当 $x+y>2$ 时，$x>1$ 或 $y>1$ 一定成立，必要性成立. 故选 $B$.}
\solend
\fi

\item 设 $a,b\in\R$，且 $a<0$，$-1<b<0$，那么下列不等式成立的是\mbox{（\quad）}

\keepwithprev
A. $a>ab>ab^2$\qquad B. $ab^2>ab>a$\qquad C. $ab>a>ab^2$\qquad D. $ab>ab^2>a$

\ifanswer
\solbegin
\step{因为 $-1<b<0$，$a<0$，所以 $ab>0$，$0<b^2<1$.}
\step{所以 $ab-ab^2=ab(1-b)>0$，$ab^2-a=a(b^2-1)>0$，}
\step{所以 $ab>ab^2>a$. 故选 $D$.}
\solend
\fi

% 第 15 题答案原卷印在题干括号内（“(  D  )”），且无【答案】/【解析】段；
% 按 高一37 第 13 题、高一63 第 14 题成例括号留空、答案另提 \ans{$D$}
\item 用反证法证明：“平面四边形中至少有一个内角不超过 90°”时，下列“假设”中正确的是\mbox{（\quad）}

\keepwithprev
A. 假设有两个内角超过 90°\qquad B. 假设有三个内角超过 90°

C. 假设至多有两个内角超过 90°\qquad D. 假设四个内角都超过 90°

\ifanswer
\ans{$D$}
\fi

\item 设 $A$ 为非空实数集，集合 $B=\{xy\mid x\in A,y\in A,x\neq y\}$，则对于以下两个结论说法正确的是\mbox{（\quad）}

结论①：若 $A$ 是五元集，则 $B$ 中的元素最多 10 个；

结论②：存在四元集 $A$，使得 $B=\{2,4,5,8,10,16\}$.

\keepwithprev
A. ①对②对\qquad B. ①对②错\qquad C. ①错②对\qquad D. ①错②错

\ifanswer
\solbegin
\step{对于①，想让元素个数达到最大值 10，则须避免产生符号不同的乘积，}
\step{否则正负分开计数容易导致总数受限，且容易产生相等值，}
\step{也最好避免 0 的出现，}
\step{故设 $A=\{a_1,a_2,a_3,a_4,a_5\}$，$0<a_1<a_2<a_3<a_4<a_5$.}
% 原卷此链“$a_3a_3$”同足标自乘、“$a_{14}a_5$”下标 14 均与文意不合，疑为 $a_2a_4$、$a_4a_5$ 之误；照录（见文件头“原卷笔误/存疑”④）
\step{因为 $a_1a_2<a_1a_3<a_1a_4<a_1a_5<a_2a_3<a_3a_3<a_{14}a_5$，$a_2a_3<a_2a_4<a_3a_4$，}
\step{所以 $B$ 中元素个数小于等于 10 个，}
\step{如设 $A=\{1,2,3,5,7\}$，则 $B=\{2,3,5,6,7,10,14,15,21,35\}$，}
\step{所以集合 $B$ 中元素个数的最大值为 10 个，故①正确；}
\step{对于②，由于集合 $B$ 中没有 0，所以集合 $A$ 中没有 0；}
\step{设集合 $A=\{a,b,c,d\}$，且 $a<b<c<d$，}
\step{当 $a<0<b<c<d$ 时，$a\times b<0$，$a\times c<0$，$a\times d<0$，}
\step{故集合 $B$ 至少有 3 个负数，与 $B=\{2,4,5,8,10,16\}$ 矛盾；}
\step{当 $a<b<0<c<d$ 时，$a\times b>0$，$c\times d>0$，}
\step{$a\times c<0$，$a\times d<0$，$b\times c<0$，$b\times d<0$，}
\step{故集合 $B$ 至少有 4 个负数，与 $B=\{2,4,5,8,10,16\}$ 矛盾；}
\step{当 $a<b<c<0<d$ 时，$a\times b>0$，$a\times c>0$，$b\times c>0$，}
\step{$a\times d<0$，$b\times d<0$，$c\times d<0$，}
\step{故集合 $B$ 至少有 3 个负数，与 $B=\{2,4,5,8,10,16\}$ 矛盾；}
\step{当 $0<a<b<c<d$ 时，则集合 $B=\{ab,ac,ad,bc,bd,cd\}$，}
\step{则必有 $ab=2$，$cd=16$，其 4 个正实数的乘积 $abcd=32$，}
\step{又 $ad=5$，$bc=8$ 或 $ad=8$，$bc=5$，其 4 个正实数的乘积 $abcd=40$，矛盾；}
\step{当 $a<b<c<d<0$ 时，与 $0<a<b<c<d$ 时的证明方法相同，}
\step{得集合 $A$ 也不存在；}
\step{综上不存在四元集 $A$，使其生成集 $B=\{2,4,5,8,10,16\}$，故②错误.}
\step{故选 $B$.}
\solend
\fi

\end{enumerate}

\bigsec{三、解答题}

\begin{enumerate}
\setcounter{enumi}{16}

\item 解下列不等式（组）

\begin{enumerate}[label=(\arabic*)]
\item 解不等式：$(ax+1)(x+1)<0(a\in\R)$.
\item 解不等式组：$\begin{cases}\dfrac{2x+4}{x+4}\geqslant1\\[4pt] (x+1)(2-|x|)\geqslant0\end{cases}$.
\end{enumerate}

\ifanswer
\solbegin
\stepm{(1)}{当 $a=0$ 时，原不等式为 $x+1<0$，即 $x<-1$，}
\step{所以此时不等式的解集为 $\{x\mid x<-1\}$；}
\step{当 $a\neq0$ 时，令 $(ax+1)(x+1)=0$，解得 $x_1=-\dfrac1a$,$x_2=-1$，}
\step{当 $a>1$ 时，$-\dfrac1a>-1$，此时不等式的解集为 $\left\{x\mid-1<x<-\dfrac1a\right\}$；}
% 以下两行“此时”前原卷均无标点；照录（见文件头“原卷笔误/存疑”⑥）
\step{当 $a=1$ 时，原不等式为 $(x+1)^2<0$ 此时不等式的解集为 $\varnothing$；}
\step{当 $0<a<1$ 时，$-\dfrac1a<-1$ 此时不等式的解集为 $\left\{x\mid-\dfrac1a<x<-1\right\}$；}
\step{当 $a<0$ 时，$-\dfrac1a>-1$，此时不等式的解集为 $\left\{x\mid x<-1\ \text{或}\ x>-\dfrac1a\right\}$.}
\step{综上所述，当 $a=0$ 时，$\{x\mid x<-1\}$；}
\step{当 $a>1$ 时，$\left\{x\mid-1<x<-\dfrac1a\right\}$；}
\step{当 $a=1$ 时，$\varnothing$；}
\step{当 $0<a<1$ 时，$\left\{x\mid-\dfrac1a<x<-1\right\}$；}
\step{当 $a<0$ 时，$\left\{x\mid x<-1\ \text{或}\ x>-\dfrac1a\right\}$.}
\stepm{(2)}{不等式 $\dfrac{2x+4}{x+4}\geqslant1$ 可化为 $\dfrac{2x+4}{x+4}-1\geqslant0$，即 $\dfrac x{x+4}\geqslant0$，}
\step{等价于 $\begin{cases}x(x+4)\geqslant0\\ x+4\neq0\end{cases}$，解得 $x<-4$ 或 $x\geqslant0$，}
\step{不等式 $(x+1)(2-|x|)\geqslant0$ 等价于 $\begin{cases}(x+1)(2-x)\geqslant0,x\geqslant0\\ (x+1)(2+x)\geqslant0,x<0\end{cases}$，}
\step{解不等式 $(x+1)(2-x)\geqslant0$，$x\geqslant0$，得 $0\leqslant x\leqslant2$，}
\step{解不等式 $(x+1)(2+x)\geqslant0$，$x<0$，得 $x\leqslant-2$ 或 $-1\leqslant x<0$，}
\step{所以 $(x+1)(2-|x|)\geqslant0$ 的解集为 $\{x\mid x\leqslant-2\ \text{或}\ -1\leqslant x\leqslant2\}$，}
\step{综上所述，$\begin{cases}\dfrac{2x+4}{x+4}\geqslant1\\[4pt] (x+1)(2-|x|)\geqslant0\end{cases}$ 的解集为 $\{x\mid x<-4\ \text{或}\ 0\leqslant x\leqslant2\}$.}
\solend
\fi

\ansspace[8]

% 第 18 题题干＋两小问与作答区 6 行须同页
\ansneed{20cm}
\item 已知 $a$ 为实数，设 $f(x)=x^2+2(a+1)x+1$.

\begin{enumerate}[label=(\arabic*)]
\item 已知对于任意实数 $x$，$\sqrt{f(x)}$ 均有意义，求实数 $a$ 的取值范围；
\item 设 $A=\{x\mid f(x)+a^2-2=0\}$，$B=\{-4,0\}$，$A\cup B=B$，求实数 $a$ 的取值范围.
\end{enumerate}

\ifanswer
\solbegin
\stepm{(1)}{$\sqrt{f(x)}$ 对任意实数 $x$ 都有意义，}
\step{即 $f(x)=x^2+2(a+1)x+1\geqslant0$ 对任意 $x\in\R$ 恒成立，}
\step{所以 $\Delta=4(a+1)^2-4\times1\times1\leqslant0$，即 $(a+1)^2\leqslant1$，解得 $-2\leqslant a\leqslant0$，}
\step{故 $a$ 的取值范围是 $[-2,0]$.}
\stepm{(2)}{由 $f(x)+a^2-2=0$ 得 $x^2+2(a+1)x+a^2-1=0$.}
\step{由 $A\cup B=B$ 得 $A\subseteq B=\{-4,0\}$，即该方程无实根，}
\step{或它的实根都在 $\{-4,0\}$ 中.}
\step{当方程无实根时，$\Delta=4(a+1)^2-4(a^2-1)=8a+8<0$，解得 $a<-1$；}
\step{当方程有实根时，$a\geqslant-1$，设两根为 $x_1$、$x_2$，}
\step{由根与系数的关系得 $x_1+x_2=-2(a+1)$、$x_1x_2=a^2-1$，}
\step{此时两根只能在 $\{-4,0\}$ 中取值；}
\step{若两根为 0 与 0，由 $x_1x_2=0$ 得 $a=\pm1$，}
\step{再由 $x_1+x_2=0$ 得 $a=-1$；}
\step{若两根为 0 与 $-4$，由 $x_1x_2=0$ 得 $a=\pm1$，再由 $x_1+x_2=-4$ 得 $a=1$.}
\step{经检验，$a=-1$ 时 $A=\{0\}$，$a=1$ 时 $A=\{-4,0\}$，都满足 $A\subseteq B$.}
\step{综上，实数 $a$ 的取值范围为 $(-\infty,-1]\cup\{1\}$.}
\solend
\fi

\ansspace[8]

\ansneed{18cm}
\item 共享单车给市民出行带来了诸多便利，某公司购买了一批单车投放到某地给居民使用，据市场分析，这批单车的运营累计利润 $y$（单位：元）与运营天数 $x$（单位：天，$x$ 为正整数）满足关系：$y=-\dfrac12x^2+60x-800$.

\begin{enumerate}[label=(\arabic*)]
\item 要使这批单车的运营累计利润高于 800 元，求运营天数的取值范围；
\item 判断这批单车的日平均运营利润能否超过 20 元，并说明理由.
\end{enumerate}

\ifanswer
\solbegin
\stepm{(1)}{由题意得 $y=-\dfrac12x^2+60x-800>800$，}
\step{即 $x^2-120x+3200<0$，解得 $40<x<80$，}
\step{即运营天数的取值范围为 $\{x\mid40<x<80,x\in\N^*\}$.}
\stepm{(2)}{不能，理由如下：}
\step{这批单车的日平均运营利润为 $\dfrac yx=-\dfrac12x-\dfrac{800}x+60=-\dfrac12\left(x+\dfrac{1600}x\right)+60$，}
\step{因 $x>0$，则 $x+\dfrac{1600}x\geqslant2\sqrt{x\cdot\dfrac{1600}x}=80$，}
\step{当且仅当 $x=\dfrac{1600}x$，即 $x=40$ 时取等号，}
\step{此时 $\dfrac yx=-\dfrac12\left(x+\dfrac{1600}x\right)+60\leqslant-\dfrac12\times80+60=20$，}
\step{故这批单车的日平均运营利润不能超过 20 元.}
\solend
\fi

\ansspace[8]

% 第 20 题题干＋三小问与作答区 8 行须同页
\ansneed{20cm}
\item 设 $f(x)=ax^2+bx+c$($a\neq0$，$b$、$c\in\R$).

\begin{enumerate}[label=(\arabic*)]
\item 若 $a=b=c=1$，设 $A=\{(x,y)\mid y=f(x)\}$，$B=\{(x,y)\mid y=1\}$，用列举法表示 $A\cap B$；
% 原卷两个“若”连用（“若 a=1，…，若方程 f(x)=0 有两个大于 1 的根”），疑衍其一；照录（见文件头“原卷笔误/存疑”⑦）
\item 若 $a=1$，$b=2k-1$，$c=k^2$，若方程 $f(x)=0$ 有两个大于 1 的根，求实数 $k$ 的取值范围；
\item 若 $a$、$b$、$c\in\Z$，且 $f(0)$、$f(1)$ 均为奇数，求证：方程 $f(x)=0$ 无整数解.
\end{enumerate}

\ifanswer
\solbegin
\stepm{(1)}{当 $a=b=c=1$ 时，$f(x)=x^2+x+1$，}
\step{所以 $A=\{(x,y)\mid y=x^2+x+1\}$，又 $y=1$，}
\step{由 $x^2+x+1=1$ 得 $x^2+x=0$，解得 $x=0$ 或 $x=-1$，}
\step{所以 $A\cap B=\{(0,1),(-1,1)\}$.}
\stepm{(2)}{当 $a=1$、$b=2k-1$、$c=k^2$ 时，$f(x)=x^2+(2k-1)x+k^2$.}
\step{设两根为 $x_1$、$x_2$，}
\step{则“两根都大于 1”等价于 $\Delta>0$、$f(1)>0$、$-\dfrac b{2a}>1$.}
\step{由 $\Delta=(2k-1)^2-4k^2=1-4k>0$ 得 $k<\dfrac14$；}
\step{由 $f(1)=1+(2k-1)+k^2=k^2+2k=k(k+2)>0$ 得 $k<-2$ 或 $k>0$；}
\step{由 $-\dfrac{2k-1}2=\dfrac{1-2k}2>1$ 得 $k<-\dfrac12$.}
\step{三者取交集得 $k<-2$，即 $k$ 的取值范围是 $(-\infty,-2)$.}
\stepm{(3)}{设 $x_0\in\Z$ 是方程 $f(x)=0$ 的整数解.}
\step{因为 $x_0(x_0-1)$ 是两个相邻整数之积，恒为偶数，所以 $ax_0(x_0-1)$ 为偶数.}
\step{而 $f(x_0)=ax_0^2+bx_0+c=ax_0(x_0-1)+(a+b)x_0+c$，}
\step{所以 $f(x_0)$ 与 $(a+b)x_0+c$ 同奇偶.}
\step{若 $x_0$ 为偶数，则 $(a+b)x_0$ 为偶数，$f(x_0)$ 与 $c=f(0)$ 同奇偶，为奇数；}
\step{若 $x_0$ 为奇数，则 $(a+b)x_0$ 与 $a+b$ 同奇偶，}
\step{$f(x_0)$ 与 $(a+b)+c=f(1)$ 同奇偶，也为奇数.}
\step{综上，$f(x_0)$ 恒为奇数，不可能等于 0，与 $f(x_0)=0$ 矛盾，}
\step{所以方程 $f(x)=0$ 无整数解.}
\solend
\fi

\ansspace[8]

% 压轴题（题干＋三小问）与作答区 10 行须同页
\ansneed{20cm}
\item 已知 $X$ 为有限集，且任取 $x\in X$，均有 $x\in\Z$，称 $X$ 为非空有限整数集. 定义：$X_k=\{x+k\mid x\in X\}$，$X^2=\{y\mid y=x^2,x\in X\}$.

\begin{enumerate}[label=(\arabic*)]
\item 若 $X=\{-2,2\}$，分别写出 $X_2$ 和 $X^2$（用列举法表示）；
\item 若 $X=\{-1,0,1\}$，$X_1\subseteq Y$，$(Y_{-1})^2\subseteq X$，求集合 $Y$ 中所有的元素；
\item 对于非空有限整数集 $A$、$B$，若 $A_k\subseteq B$，$(B^2)_{-2}\subseteq A$，$B$ 中的元素均为非负整数，$k$ 为整数，求证：

① $A$、$B$ 中元素的个数相同；

② 若集合 $A$ 中的元素个数为奇数，则 $k$ 为偶数.
\end{enumerate}

\ifanswer
\solbegin
\stepm{(1)}{由题意得 $X_2=\{0,4\}$，$X^2=\{4\}$；}
\stepm{(2)}{由题意得 $X_1=\{0,1,2\}$，即 $\{0,1,2\}\subseteq Y$，$Y_{-1}=\{y-1\mid y\in Y\}$，}
\step{$(Y_{-1})^2=\{(y-1)^2\mid y\in Y\}\subseteq X=\{-1,0,1\}$，}
\step{所以 $(y-1)^2=0$ 或 $(y-1)^2=1$，解得 $y=1$ 或 $y=2$ 或 $y=0$，}
\step{即 $y\in\{0,1,2\}$，即 $Y\subseteq\{0,1,2\}$，所以 $Y=\{0,1,2\}$；}
% 原卷 (3) 解析首行“(3)”与“①”连排，本稿合并记作 \stepm{(3)①}
\stepm{(3)①}{因为 $A_k=\{x+k\mid x\in A\}$，函数 $g(x)=x+k$ 在全体整数集上严格增，}
\step{所以 $\forall x_1$、$x_2\in A$，$x_1\neq x_2$，$g(x_1)\neq g(x_2)$，所以 $|A_k|=|A|$，}
\step{因为 $A_k\subseteq B$，所以 $|A|\leqslant|B|$，}
\step{设 $b\in B$，$B$ 中元素均为非负整数，$\forall b_1\neq b_2\in B$，$b_1^2\neq b_2^2$，}
\step{故函数 $f(b)=b^2-2$ 在 $B$ 上单射，所以 $|(B^2)_{-2}|=|B|$，}
\step{又因为 $(B^2)_{-2}\subseteq A$，所以 $|B|\leqslant|A|$，所以 $|A|=|B|$.}
\stepm{②}{由①得 $A_k\subseteq B$ 且 $|A_k|=|A|=|B|$，$(B^2)_{-2}\subseteq A$ 且 $|(B^2)_{-2}|=|B|=|A|$，}
\step{所以 $A_k=B$，$(B^2)_{-2}=A$，}
\step{所以 $B=A_k=\{a+k\mid a\in A\}$，$A=(B^2)_{-2}=\{b^2-2\mid b\in B\}$，}
\step{所以 $A=\{(a+k)^2-2\mid a\in A\}$，}
\step{所以映射 $h(a)=(a+k)^2-2$ 是集合 $A$ 到自身的一个双射.}
\step{假设 $k$ 是奇数，当 $a$ 为偶数时，$a+k$ 为奇数，$(a+k)^2-2$ 为奇数，}
\step{当 $a$ 为奇数时，$a+k$ 为偶数，$(a+k)^2-2$ 为偶数，}
\step{所以此时映射 $h(a)=(a+k)^2-2$，将集合 $A$ 中的奇数映射为偶数，}
\step{将集合 $A$ 中的偶数映射为奇数，即集合 $A$ 中的奇数元素和偶数元素必须配对，}
\step{此时集合 $A$ 的元素个数则为偶数，}
\step{与题设中“集合 $A$ 中的元素个数为奇数”矛盾.}
\step{故假设不成立，所以 $k$ 是偶数.}
\solend
\fi

\ansspace*[10]

\end{enumerate}

\end{document}
"
% 紧凑版：xelatex "\def\iscompact{1}% 高一67_华师大一附中_10月月考.tex
%   —— 高一上数学“2026-2027 年华师大一附中高一上 10 月月考”（讲评版）
% 试卷版：xelatex 高一67_华师大一附中_10月月考.tex
% 答案版：xelatex "\def\isanswer{1}\input{高一67_华师大一附中_10月月考.tex}"
% 紧凑版：xelatex "\def\iscompact{1}\input{高一67_华师大一附中_10月月考.tex}"
%
% 原始材料：微信公众号“魔都数学加油站”文章，作者破晓，连载“27学年高一67”；
% 原文标题“27学年高一67——2026-2027年上海市华师大一附中高一上10月月考”；
% 原文链接：https://mp.weixin.qq.com/s/Kdhx07TVjLoGTcIlnRzgcQ。
% 页面标注发布时间 2026-10-18 20:33（上海）。
% 本次转录只据 `原图/高一67_华师大一附中_10月月考/` 的 13 张 PNG 逐页读图。
% 正文为 13 张 PNG 页（`原图/高一67_华师大一附中_10月月考/01..13.png`）：
% 第 01、04、13 页 2560×3620，第 02、03、05—12 页 4762×6735
% （**同一篇各页分辨率不同**），均按页序读取。13 页全为试卷正文页（卷首标题在第 1 页页顶，
% 卷末解答题第 21 题解析止于第 13 页），无封面页、目录页、二维码或宣传页。
% 粉色斜排“魔都数学加油站”为卷面水印，与公众号页脚一并未录入。
%
% 卷首标题：照卷面印作“2026--2027 年华师大一附中高一上 10 月月考”（卷面为黑体加粗一行），
% 年份区间按全稿惯例排 en dash。原卷只有一行标题、无副标题。
%
% 篇幅与题号：原卷自第 3 题起编号（第 1、2 题不在本卷内）。填空题 3—12（10 题）、
% 选择题 13—16（4 题）、解答题 17—21（5 题），共 19 题。本稿沿用原节与原题号，
% 只改排为全稿统一的大节样式（原卷小节标题为“一、填空题”“二、选择题”“三、解答题”）。
%
% 讲评版：原卷每题下直接印【解析】。**全卷 19 题中 18 题都印【解析】（无一处只印
% 【答案】）**，按全稿口径这类题不另提 \ans{}；各题解析末句的结论（第 13、14、16 题
% 的“故选 B／D／B”、各填空与解答题的“故/综上…”）照录在解析内。
% **另有 1 题答案印在题干括号内**——第 15 题（“(  D  )”，且该题无【答案】/【解析】段），
% 本稿按 高一37 第 13 题、高一63 第 14 题成例把括号留空、答案另提为 \ans{$D$}。
%
% 副标题：原卷未印副标题。本稿按“知识模块”口径标作“集合与常用逻辑用语、不等式”：
% 19 题中集合与常用逻辑用语 12 题（第 3、4、5、7、10、11、12、13、15、16、18、21 题：
% 补集、并集、子集、充分条件、容斥计数、子集计数、新定义二元运算、条件判断、
% 反证法假设、新定义生成集、集合关系与二次方程、新定义 X_k 与 X^2），
% 不等式 7 题（第 6、8、9、14、17、19、20 题：存在量词不等式、解集反求系数、
% 解集为空、不等式性质、解不等式（组）、二次不等式与基本不等式、二次方程根的分布），
% 两类均超过三分之一，故并标。
%
% 与原卷的排版差异：
%   ① 原卷每题下直接印【解析】，本稿按全稿惯例将解析置于答案版，试卷版保持干净；
%   ② 原卷未印副标题，本稿按“知识模块”口径补出；
%   ③ 原卷填空横线约四个字宽，本稿按全稿惯例统一排 \blankline[4.4em]；
%   ④ 原卷实数集排斜体 $R$（第 14 题题干、第 17 题 (1) 题干、第 18 题解析、第 20 题题干），
%      整数集排斜体 $Z$（第 20 题 (3) 题干与解析、第 21 题题干），正整数集排粗体
%      $N^*$（第 11 题解析、第 19 题解析），本稿按全稿惯例统一改用 $\R$、$\Z$、$\N^*$
%      （记号归一，粗体／斜体不作区分）；空集原卷作 $\varnothing$，照录；
%   ⑤ 原卷在数学式内多处用半角逗号（如第 8 题 $(1,2)$ 与各 cases、第 11 题解析
%      $P_n=\{1,2,\cdots,n\}$,$n\in N^*$、第 12 题解析 $u=(\cdots)$,$v=(\cdots)$ 与
%      “2020,个 1”、第 19 题解析 $\{x|40<x<80,x\in N^*\}$），均照录未强改；
%      描述法集合的竖线照原卷排 \mid；句末句点按全稿惯例排西文句点；
%   ⑥ 原卷补集记号为上划线 $\bar A$（第 3 题题干与解析、第 11 题题干），照录
%      （同 高一64/65 的 $\bar B$ 成例）；
%   ⑦ 第 12 题新定义的运算记号 $*$ 照录（$\alpha*\beta$、$u*u$、$u*v$）；
%   ⑧ 第 17—21 题的子问原卷各自另起一行，本稿按全稿惯例用 enumerate 的 (1)(2)(3)；
%      解析中的子问标号一律用 \stepm{(1)}{…}；第 5 题解析的①②、第 21 题 (3) 解析的①②
%      按全稿惯例作 \stepm{①}{…}、\stepm{②}{…}（第 21 题 (3) 首行原卷“(3)”与“①”连排，
%      本稿合并记作 \stepm{(3)①}{…}）；
%   ⑨ 选择题作答括号原卷排半角“(　　)”（第 13、14、16 题括号内空白、第 15 题括号内印 D），
%      本稿按全稿惯例排 \mbox{（\quad）}；第 15 题括号内的答案移入 \ans{}，见上；
%   ⑩ 选项行排照原卷：第 13、15 题两行两条，第 14、16 题一行四条；
%   ⑪ 每道解答题原卷未留作答空白，本稿按全稿惯例加 \ansspace（末题用 \ansspace*）；
%      第 20、21 题题干长，各用一处 \ansneed 让题干与作答区不拆页（\ansneed 只作用于试卷版）。
%
% 原卷笔误/存疑（均照录未改，正文对应处另有 % 注）：
%   ① 第 3 题解析末句印作“则 $\bar A=(-\infty,0)\cup[2,3)$.”：按题干
%      $U=\{x\mid 1/x\leqslant 1/2\}$（即 $x<0$ 或 $x\geqslant2$）与 $A=[3,+\infty)$，
%      $\bar A$ 应为 $(-\infty,0)\cup\{2\}$；且与本解析首句自求之 $U$ 相矛盾（$3\notin U$），
%      疑为笔误；照录未改。
%   ② 第 12 题题干 $\alpha*\beta$ 定义式括号不配对：第四组“$+\cdots+x_{2020}+y_{2020}
%      -x_{2020}y_{2020}$”缺左括号、末尾连排两个右括号“))”；照录未改。
%   ③ 第 12 题解析“所以…中有 2020,个 1”一行，“$\cdots$”后脱逗号（作
%      “$x_2,\cdots x_{2020}$”，第 12 行“而…”同式样），且“2020”后衍一半角逗号
%      （作“2020,个 1”）；照录未改。
%   ④ 第 16 题解析乘积链印作“$a_1a_2<a_1a_3<a_1a_4<a_1a_5<a_2a_3<a_3a_3<a_{14}a_5$”：
%      “$a_3a_3$”同足标自乘、与单调链不合，“$a_{14}a_5$”下标 14 与 $a_1$—$a_5$ 记号
%      不合（疑为 $a_2a_4$、$a_4a_5$ 之误）；照录未改。
%   ⑤ 第 5 题题干作“$Q\subset P$”（真子集记号），解析两处均作“$Q\subseteq P$”，
%      同一题前后记号不一；照录（本稿分别排 \subset、\subseteq）。
%   ⑥ 第 17 题 (1) 解析“当 $a=1$ 时，原不等式为 $(x+1)^2<0$ 此时…”与
%      “当 $0<a<1$ 时，$-\frac1a<-1$ 此时…”两处“此时”前均脱标点；照录未改。
%   ⑦ 第 20 题 (2) 题干“若 $a=1$，$b=2k-1$，$c=k^2$，若方程 $f(x)=0$ 有两个大于 1 的根”
%      两个“若”连用，疑衍其一；照录未改。
%
% 插图：无（全卷 13 页无任何几何图形、无“如图/右图/图中”字样）。
%
% 行数核对：逐题数 \step/\stepm 条数与原卷【解析】印刷行数，18 题逐一相符
% （第 15 题无解析段，只提 \ans{}），合计 169 条 / 169 行
% （Q3—Q12：2、3、4、3、1、3、5、3、13、15；Q13—Q16：2、3、0、25；
% Q17—Q21：19、16、9、20、23）。
%
\documentclass[a4paper,11pt]{article}
\usepackage{common}

\begin{document}

\examtitlest[3]{2026--2027 年华师大一附中高一上 10 月月考}{集合与常用逻辑用语、不等式}

\bigsec{一、填空题}

\begin{enumerate}
\setcounter{enumi}{2}

\item 设全集 $U=\left\{x\mid\dfrac1x\leqslant\dfrac12\right\}$，$A=[3,+\infty)$，则 $\bar A=$\blankline[4.4em].

\ifanswer
\solbegin
\step{$\dfrac1x\leqslant\dfrac12$，解得 $x<0$ 或 $x\geqslant2$，故 $U=(-\infty,0)\cup[2,+\infty)$，}
% 原卷末句作 $\bar A=(-\infty,0)\cup[2,3)$，按题干 $U$ 应为 $(-\infty,0)\cup\{2\}$，与本解析首句矛盾，疑笔误；照录（见文件头“原卷笔误/存疑”①）
\step{$A=[3,+\infty)$，则 $\bar A=(-\infty,0)\cup[2,3)$.}
\solend
\fi

\item 已知集合 $M=\{y\mid y=x^2-2x-1\}$，$N=\{y\mid y=x^2-4x+1\}$，则 $M\cup N=$\blankline[4.4em].

\ifanswer
\solbegin
\step{$M=\{y\mid y=x^2-2x-1\}=\{y\mid y=(x-1)^2-2\}=\{y\mid y\geqslant-2\}$，}
\step{$N=\{y\mid y=x^2-4x+1\}=\{y\mid y=(x-2)^2-3\}=\{y\mid y\geqslant-3\}$，}
\step{则 $M\cup N=\{y\mid y\geqslant-3\}$.}
\solend
\fi

\item 已知 $P=\{x\mid x^2+x-6=0\}$，$Q=\{x\mid mx-1=0\}$，且 $Q\subset P$，则实数 $m$ 的值\blankline[4.4em].

\ifanswer
\solbegin
\step{$P=\{2,-3\}$，$Q=\{x\mid mx=1\}$；}
\stepm{①}{$m=0$ 时，$Q=\varnothing$，满足 $Q\subseteq P$；}
\stepm{②}{$m\neq0$ 时，要使 $Q\subseteq P$，则 $\dfrac1m=2$ 或 $-3$，所以 $m=\dfrac12$ 或 $-\dfrac13$；}
\step{所以 $m=0$ 或 $\dfrac12$ 或 $-\dfrac13$.}
\solend
\fi

\item 已知“存在正实数 $x$，使得关于 $x$ 的不等式 $ax-1>0$ 成立”为真命题，则实数 $a$ 的取值范围是\blankline[4.4em].

\ifanswer
\solbegin
\step{存在正实数 $x$，使得关于 $x$ 的不等式 $ax-1>0$ 成立为真命题，}
\step{即存在正实数 $x$，使得 $a>\dfrac1x$，因为 $\dfrac1x>0$，所以 $a>0$，}
\step{即 $a$ 的范围为 $(0,+\infty)$.}
\solend
\fi

\item 设 $\alpha:3\leqslant x\leqslant5$，$\beta:2m-1\leqslant x<m+4$，若 $\alpha$ 是 $\beta$ 的充分条件，则实数 $m$ 的取值范围是\blankline[4.4em].

\ifanswer
\solbegin
\step{若 $\alpha$ 是 $\beta$ 的充分条件，则 $\begin{cases}2m-1\leqslant3\\ m+4>5\end{cases}$，解得 $1<m\leqslant2$.}
\solend
\fi

\item 对于实数 $a$、$b$、$c$，已知关于 $x$ 的不等式 $ax^2+bx+c>0$ 的解集是 $(1,2)$，则关于 $x$ 的不等式 $cx^2-bx+a\leqslant0$ 的解集是\blankline[4.4em].

\ifanswer
\solbegin
\step{关于 $x$ 的不等式 $ax^2+bx+c>0$ 的解集是 $(1,2)$，则 $\begin{cases}a<0\\ -\dfrac ba=3\\ \dfrac ca=2\end{cases}$，则 $\begin{cases}b=-3a\\ c=2a\end{cases}$，}
\step{则关于 $x$ 的不等式 $cx^2-bx+a=2ax^2+3ax+a\leqslant0$，则 $2x^2+3x+1\geqslant0$，}
\step{则 $x\leqslant-1$ 或 $x\geqslant-\dfrac12$，则解集为 $(-\infty,-1]\cup\left[-\dfrac12,+\infty\right)$.}
\solend
\fi

\item 已知关于 $x$ 的不等式 $kx^2+kx+k+2\leqslant0$ 的解集为 $\varnothing$，则实数 $k$ 的取值范围为\blankline[4.4em].

\ifanswer
\solbegin
\step{当 $k=0$ 时，不等式化为 $2\leqslant0$，解集为 $\varnothing$，符合题意；}
\step{当 $k\neq0$ 时，若关于 $x$ 的不等式 $kx^2+kx+k+2\leqslant0$ 的解集为 $\varnothing$，}
\step{则 $\begin{cases}k>0\\ \Delta=k^2-4k(k+2)<0\end{cases}$，化简得 $\begin{cases}k>0\\ -3k^2-8k<0\end{cases}$，即 $\begin{cases}k>0\\ k(3k+8)>0\end{cases}$，}
\step{解得 $k>0$.}
\step{综上所述，$k\geqslant0$，即实数 $k$ 的取值范围为 $[0,+\infty)$.}
\solend
\fi

\item 对班级若干名学生调查对 $A$、$B$ 两事件的态度，有如下结果：赞成 $A$ 的人数是全体的五分之三，其余的不赞成；赞成 $B$ 的比赞成 $A$ 多 3 人，其余的不赞成. 已知对 $A$、$B$ 都赞成的学生有 18 人，对 $A$、$B$ 都不赞成的学生数比对 $A$、$B$ 都赞成的学生数的三分之一多 1 人，则班级总人数是\blankline[4.4em]人.

\ifanswer
\solbegin
\step{设班级总人数是 $x$ 人，则赞成 $A$ 的人数是 $\dfrac35x$，赞成 $B$ 的人数是 $\dfrac35x+3$，}
\step{$A$、$B$ 都不赞成的人数为 $18\times\dfrac13+1=7$ 人，}
\step{由容斥原理得 $x=\left(\dfrac35x+\dfrac35x+3-18\right)+7$，解得 $x=40$，即班级总人数是 40 人.}
\solend
\fi

\item 设全集 $U=\{1,2,3,4,5,6,7,8,9,10,11,12\}$，若集合 $A$ 满足：① $A\subseteq U$；②若 $x\in A$，必有 $2x\notin A$；③若 $x\in\bar A$，必有 $2x\notin\bar A$，则这样的集合 $A$ 有\blankline[4.4em]个.

\ifanswer
\solbegin
\step{设集合 $P_n=\{1,2,\cdots,n\}$,$n\in\N^*$，}
\step{任取 $x\in P_n$，将 $x$ 除以 2，若商仍为偶数，再除以 2，$\cdots$，}
\step{经过 $k$ 次以后，商必为奇数，此时记商为 $m$，}
\step{于是 $x=m\cdot2^k$，其中 $m$ 为奇数，$k\in\N^*$.}
\step{由条件得，若 $m\in A$，则 $x\in A$ 等价于 $k$ 为偶数；}
\step{若 $m\notin A$，则 $x\in A$ 等价于 $k$ 为奇数；}
\step{于是 $x$ 是否属于 $A$ 由 $m$ 是否属于 $A$ 确定；}
\step{设 $Q_n$ 是 $P_n$ 中所有奇数的集合，}
\step{设 $f(n)$ 为集合 $P_n$ 中满足题干条件的集合 $A$ 的个数，}
\step{因此 $f(n)$ 等于 $Q_n$ 的子集个数.}
\step{当 $n$ 为偶数（或奇数）时，$P_n$ 中奇数的个数是 $\dfrac n2$（或 $\dfrac{n+1}2$），}
\step{所以 $f(n)=\begin{cases}2^{\frac n2},n\text{为偶数}\\[2pt] 2^{\frac{n+1}2},n\text{为奇数}\end{cases}$，所以 $P_{12}=\{1,2,3,4,5,6,7,8,9,10,11,12\}$，}
\step{即同时满足三个条件的集合 $A$ 的个数为 $2^{\frac{12}2}=2^6=64$.}
\solend
\fi

\item 设集合 $\Omega=\{(x_1,x_2,\cdots,x_{2020})\mid x_i\in\{0,1\},i=1,2,\cdots,2020\}$，对于 $\Omega$ 中的任意两个元素 $\alpha=(x_1,x_2,\cdots,x_{2020})$，$\beta=(y_1,y_2,\cdots,y_{2020})$，记
% 原卷此式第四组缺左括号、末尾连排两个右括号“))”，括号不配对；照录（见文件头“原卷笔误/存疑”②）
$\alpha*\beta=(x_1+y_1-x_1y_1)+(x_2+y_2-x_2y_2)+\cdots+x_{2020}+y_{2020}-x_{2020}y_{2020}))$，设 $u$、$v\in\Omega$，若 $u*u+v*v=2020$，则 $u*v$ 的最大值和最小值之和为\blankline[4.4em].

\ifanswer
\solbegin
\step{设 $u=(x_1,x_2,\cdots,x_{2020})$,$v=(y_1,y_2,\cdots,y_{2020})$，}
\step{因为 $x_i\in\{0,1\}$，所以 $x_i(x_i-1)=0$，}
\step{所以 $u*u=(x_1+x_1-x_1^2)+(x_2+x_2-x_2^2)+\cdots+(x_{2020}+x_{2020}-x_{2020}^2)$}
\step{$=x_1+x_2+\cdots+x_{2020}$，}
\step{$v*v=(y_1+y_1-y_1^2)+(y_2+y_2-y_2^2)+\cdots+(y_{2020}+y_{2020}-y_{2020}^2)$}
\step{$=y_1+y_2+\cdots+y_{2020}$，}
\step{由 $u*u+v*v=2020$，得 $x_1+x_2+\cdots+x_{2020}+y_1+y_2+\cdots+y_{2020}=2020$，}
% 原卷此行“$\cdots$”后脱逗号、 “2020”后衍一半角逗号（“2020,个 1”）；照录（见文件头“原卷笔误/存疑”③）
\step{所以 $x_1,x_2,\cdots x_{2020},y_1,y_2,\cdots y_{2020}$ 中有 2020,个 1，2020 个 0，}
\step{$u*v=(x_1+y_1-x_1y_1)+(x_2+y_2-x_2y_2)+\cdots+(x_{2020}+y_{2020}-x_{2020}y_{2020})$}
\step{$=2020-(x_1y_1+x_2y_2+\cdots+x_{2020}y_{2020})$，}
\step{注意到只有 $x_i=y_i=1$ 时，$x_iy_i=1$，否则 $x_iy_i=0$，}
% 此行“$\cdots$”后同样脱逗号，与上一处同式样；照录（见文件头“原卷笔误/存疑”③）
\step{而 $x_1,x_2,\cdots x_{2020},y_1,y_2,\cdots y_{2020}$ 中有 2020 个 1，2020 个 0，}
\step{所以满足 $x_iy_i=1$ 的最多有 1010 个，所以 $u*v\geqslant2020-1010=1010$，}
\step{满足 $x_iy_i=0$ 的最多有 2020 个，所以 $u*v\leqslant2020-0=2020$，}
\step{则 $u*v$ 的最大值和最小值之和为 $2020+1010=3030$.}
\solend
\fi

\end{enumerate}

\bigsec{二、选择题}

\begin{enumerate}
\setcounter{enumi}{12}

\item 实数 $x$ 与 $y$，“$x+y>2$”的\mbox{（\quad）}条件是“$x>1$ 或 $y>1$”.

\keepwithprev
A. 充分非必要\qquad B. 必要非充分

C. 充要\qquad D. 既非充分又非必要

\ifanswer
\solbegin
\step{当 $x>1$ 或 $y>1$ 时，$x+y>2$ 不一定成立，充分性不成立；}
\step{当 $x+y>2$ 时，$x>1$ 或 $y>1$ 一定成立，必要性成立. 故选 $B$.}
\solend
\fi

\item 设 $a,b\in\R$，且 $a<0$，$-1<b<0$，那么下列不等式成立的是\mbox{（\quad）}

\keepwithprev
A. $a>ab>ab^2$\qquad B. $ab^2>ab>a$\qquad C. $ab>a>ab^2$\qquad D. $ab>ab^2>a$

\ifanswer
\solbegin
\step{因为 $-1<b<0$，$a<0$，所以 $ab>0$，$0<b^2<1$.}
\step{所以 $ab-ab^2=ab(1-b)>0$，$ab^2-a=a(b^2-1)>0$，}
\step{所以 $ab>ab^2>a$. 故选 $D$.}
\solend
\fi

% 第 15 题答案原卷印在题干括号内（“(  D  )”），且无【答案】/【解析】段；
% 按 高一37 第 13 题、高一63 第 14 题成例括号留空、答案另提 \ans{$D$}
\item 用反证法证明：“平面四边形中至少有一个内角不超过 90°”时，下列“假设”中正确的是\mbox{（\quad）}

\keepwithprev
A. 假设有两个内角超过 90°\qquad B. 假设有三个内角超过 90°

C. 假设至多有两个内角超过 90°\qquad D. 假设四个内角都超过 90°

\ifanswer
\ans{$D$}
\fi

\item 设 $A$ 为非空实数集，集合 $B=\{xy\mid x\in A,y\in A,x\neq y\}$，则对于以下两个结论说法正确的是\mbox{（\quad）}

结论①：若 $A$ 是五元集，则 $B$ 中的元素最多 10 个；

结论②：存在四元集 $A$，使得 $B=\{2,4,5,8,10,16\}$.

\keepwithprev
A. ①对②对\qquad B. ①对②错\qquad C. ①错②对\qquad D. ①错②错

\ifanswer
\solbegin
\step{对于①，想让元素个数达到最大值 10，则须避免产生符号不同的乘积，}
\step{否则正负分开计数容易导致总数受限，且容易产生相等值，}
\step{也最好避免 0 的出现，}
\step{故设 $A=\{a_1,a_2,a_3,a_4,a_5\}$，$0<a_1<a_2<a_3<a_4<a_5$.}
% 原卷此链“$a_3a_3$”同足标自乘、“$a_{14}a_5$”下标 14 均与文意不合，疑为 $a_2a_4$、$a_4a_5$ 之误；照录（见文件头“原卷笔误/存疑”④）
\step{因为 $a_1a_2<a_1a_3<a_1a_4<a_1a_5<a_2a_3<a_3a_3<a_{14}a_5$，$a_2a_3<a_2a_4<a_3a_4$，}
\step{所以 $B$ 中元素个数小于等于 10 个，}
\step{如设 $A=\{1,2,3,5,7\}$，则 $B=\{2,3,5,6,7,10,14,15,21,35\}$，}
\step{所以集合 $B$ 中元素个数的最大值为 10 个，故①正确；}
\step{对于②，由于集合 $B$ 中没有 0，所以集合 $A$ 中没有 0；}
\step{设集合 $A=\{a,b,c,d\}$，且 $a<b<c<d$，}
\step{当 $a<0<b<c<d$ 时，$a\times b<0$，$a\times c<0$，$a\times d<0$，}
\step{故集合 $B$ 至少有 3 个负数，与 $B=\{2,4,5,8,10,16\}$ 矛盾；}
\step{当 $a<b<0<c<d$ 时，$a\times b>0$，$c\times d>0$，}
\step{$a\times c<0$，$a\times d<0$，$b\times c<0$，$b\times d<0$，}
\step{故集合 $B$ 至少有 4 个负数，与 $B=\{2,4,5,8,10,16\}$ 矛盾；}
\step{当 $a<b<c<0<d$ 时，$a\times b>0$，$a\times c>0$，$b\times c>0$，}
\step{$a\times d<0$，$b\times d<0$，$c\times d<0$，}
\step{故集合 $B$ 至少有 3 个负数，与 $B=\{2,4,5,8,10,16\}$ 矛盾；}
\step{当 $0<a<b<c<d$ 时，则集合 $B=\{ab,ac,ad,bc,bd,cd\}$，}
\step{则必有 $ab=2$，$cd=16$，其 4 个正实数的乘积 $abcd=32$，}
\step{又 $ad=5$，$bc=8$ 或 $ad=8$，$bc=5$，其 4 个正实数的乘积 $abcd=40$，矛盾；}
\step{当 $a<b<c<d<0$ 时，与 $0<a<b<c<d$ 时的证明方法相同，}
\step{得集合 $A$ 也不存在；}
\step{综上不存在四元集 $A$，使其生成集 $B=\{2,4,5,8,10,16\}$，故②错误.}
\step{故选 $B$.}
\solend
\fi

\end{enumerate}

\bigsec{三、解答题}

\begin{enumerate}
\setcounter{enumi}{16}

\item 解下列不等式（组）

\begin{enumerate}[label=(\arabic*)]
\item 解不等式：$(ax+1)(x+1)<0(a\in\R)$.
\item 解不等式组：$\begin{cases}\dfrac{2x+4}{x+4}\geqslant1\\[4pt] (x+1)(2-|x|)\geqslant0\end{cases}$.
\end{enumerate}

\ifanswer
\solbegin
\stepm{(1)}{当 $a=0$ 时，原不等式为 $x+1<0$，即 $x<-1$，}
\step{所以此时不等式的解集为 $\{x\mid x<-1\}$；}
\step{当 $a\neq0$ 时，令 $(ax+1)(x+1)=0$，解得 $x_1=-\dfrac1a$,$x_2=-1$，}
\step{当 $a>1$ 时，$-\dfrac1a>-1$，此时不等式的解集为 $\left\{x\mid-1<x<-\dfrac1a\right\}$；}
% 以下两行“此时”前原卷均无标点；照录（见文件头“原卷笔误/存疑”⑥）
\step{当 $a=1$ 时，原不等式为 $(x+1)^2<0$ 此时不等式的解集为 $\varnothing$；}
\step{当 $0<a<1$ 时，$-\dfrac1a<-1$ 此时不等式的解集为 $\left\{x\mid-\dfrac1a<x<-1\right\}$；}
\step{当 $a<0$ 时，$-\dfrac1a>-1$，此时不等式的解集为 $\left\{x\mid x<-1\ \text{或}\ x>-\dfrac1a\right\}$.}
\step{综上所述，当 $a=0$ 时，$\{x\mid x<-1\}$；}
\step{当 $a>1$ 时，$\left\{x\mid-1<x<-\dfrac1a\right\}$；}
\step{当 $a=1$ 时，$\varnothing$；}
\step{当 $0<a<1$ 时，$\left\{x\mid-\dfrac1a<x<-1\right\}$；}
\step{当 $a<0$ 时，$\left\{x\mid x<-1\ \text{或}\ x>-\dfrac1a\right\}$.}
\stepm{(2)}{不等式 $\dfrac{2x+4}{x+4}\geqslant1$ 可化为 $\dfrac{2x+4}{x+4}-1\geqslant0$，即 $\dfrac x{x+4}\geqslant0$，}
\step{等价于 $\begin{cases}x(x+4)\geqslant0\\ x+4\neq0\end{cases}$，解得 $x<-4$ 或 $x\geqslant0$，}
\step{不等式 $(x+1)(2-|x|)\geqslant0$ 等价于 $\begin{cases}(x+1)(2-x)\geqslant0,x\geqslant0\\ (x+1)(2+x)\geqslant0,x<0\end{cases}$，}
\step{解不等式 $(x+1)(2-x)\geqslant0$，$x\geqslant0$，得 $0\leqslant x\leqslant2$，}
\step{解不等式 $(x+1)(2+x)\geqslant0$，$x<0$，得 $x\leqslant-2$ 或 $-1\leqslant x<0$，}
\step{所以 $(x+1)(2-|x|)\geqslant0$ 的解集为 $\{x\mid x\leqslant-2\ \text{或}\ -1\leqslant x\leqslant2\}$，}
\step{综上所述，$\begin{cases}\dfrac{2x+4}{x+4}\geqslant1\\[4pt] (x+1)(2-|x|)\geqslant0\end{cases}$ 的解集为 $\{x\mid x<-4\ \text{或}\ 0\leqslant x\leqslant2\}$.}
\solend
\fi

\ansspace[8]

% 第 18 题题干＋两小问与作答区 6 行须同页
\ansneed{20cm}
\item 已知 $a$ 为实数，设 $f(x)=x^2+2(a+1)x+1$.

\begin{enumerate}[label=(\arabic*)]
\item 已知对于任意实数 $x$，$\sqrt{f(x)}$ 均有意义，求实数 $a$ 的取值范围；
\item 设 $A=\{x\mid f(x)+a^2-2=0\}$，$B=\{-4,0\}$，$A\cup B=B$，求实数 $a$ 的取值范围.
\end{enumerate}

\ifanswer
\solbegin
\stepm{(1)}{$\sqrt{f(x)}$ 对任意实数 $x$ 都有意义，}
\step{即 $f(x)=x^2+2(a+1)x+1\geqslant0$ 对任意 $x\in\R$ 恒成立，}
\step{所以 $\Delta=4(a+1)^2-4\times1\times1\leqslant0$，即 $(a+1)^2\leqslant1$，解得 $-2\leqslant a\leqslant0$，}
\step{故 $a$ 的取值范围是 $[-2,0]$.}
\stepm{(2)}{由 $f(x)+a^2-2=0$ 得 $x^2+2(a+1)x+a^2-1=0$.}
\step{由 $A\cup B=B$ 得 $A\subseteq B=\{-4,0\}$，即该方程无实根，}
\step{或它的实根都在 $\{-4,0\}$ 中.}
\step{当方程无实根时，$\Delta=4(a+1)^2-4(a^2-1)=8a+8<0$，解得 $a<-1$；}
\step{当方程有实根时，$a\geqslant-1$，设两根为 $x_1$、$x_2$，}
\step{由根与系数的关系得 $x_1+x_2=-2(a+1)$、$x_1x_2=a^2-1$，}
\step{此时两根只能在 $\{-4,0\}$ 中取值；}
\step{若两根为 0 与 0，由 $x_1x_2=0$ 得 $a=\pm1$，}
\step{再由 $x_1+x_2=0$ 得 $a=-1$；}
\step{若两根为 0 与 $-4$，由 $x_1x_2=0$ 得 $a=\pm1$，再由 $x_1+x_2=-4$ 得 $a=1$.}
\step{经检验，$a=-1$ 时 $A=\{0\}$，$a=1$ 时 $A=\{-4,0\}$，都满足 $A\subseteq B$.}
\step{综上，实数 $a$ 的取值范围为 $(-\infty,-1]\cup\{1\}$.}
\solend
\fi

\ansspace[8]

\ansneed{18cm}
\item 共享单车给市民出行带来了诸多便利，某公司购买了一批单车投放到某地给居民使用，据市场分析，这批单车的运营累计利润 $y$（单位：元）与运营天数 $x$（单位：天，$x$ 为正整数）满足关系：$y=-\dfrac12x^2+60x-800$.

\begin{enumerate}[label=(\arabic*)]
\item 要使这批单车的运营累计利润高于 800 元，求运营天数的取值范围；
\item 判断这批单车的日平均运营利润能否超过 20 元，并说明理由.
\end{enumerate}

\ifanswer
\solbegin
\stepm{(1)}{由题意得 $y=-\dfrac12x^2+60x-800>800$，}
\step{即 $x^2-120x+3200<0$，解得 $40<x<80$，}
\step{即运营天数的取值范围为 $\{x\mid40<x<80,x\in\N^*\}$.}
\stepm{(2)}{不能，理由如下：}
\step{这批单车的日平均运营利润为 $\dfrac yx=-\dfrac12x-\dfrac{800}x+60=-\dfrac12\left(x+\dfrac{1600}x\right)+60$，}
\step{因 $x>0$，则 $x+\dfrac{1600}x\geqslant2\sqrt{x\cdot\dfrac{1600}x}=80$，}
\step{当且仅当 $x=\dfrac{1600}x$，即 $x=40$ 时取等号，}
\step{此时 $\dfrac yx=-\dfrac12\left(x+\dfrac{1600}x\right)+60\leqslant-\dfrac12\times80+60=20$，}
\step{故这批单车的日平均运营利润不能超过 20 元.}
\solend
\fi

\ansspace[8]

% 第 20 题题干＋三小问与作答区 8 行须同页
\ansneed{20cm}
\item 设 $f(x)=ax^2+bx+c$($a\neq0$，$b$、$c\in\R$).

\begin{enumerate}[label=(\arabic*)]
\item 若 $a=b=c=1$，设 $A=\{(x,y)\mid y=f(x)\}$，$B=\{(x,y)\mid y=1\}$，用列举法表示 $A\cap B$；
% 原卷两个“若”连用（“若 a=1，…，若方程 f(x)=0 有两个大于 1 的根”），疑衍其一；照录（见文件头“原卷笔误/存疑”⑦）
\item 若 $a=1$，$b=2k-1$，$c=k^2$，若方程 $f(x)=0$ 有两个大于 1 的根，求实数 $k$ 的取值范围；
\item 若 $a$、$b$、$c\in\Z$，且 $f(0)$、$f(1)$ 均为奇数，求证：方程 $f(x)=0$ 无整数解.
\end{enumerate}

\ifanswer
\solbegin
\stepm{(1)}{当 $a=b=c=1$ 时，$f(x)=x^2+x+1$，}
\step{所以 $A=\{(x,y)\mid y=x^2+x+1\}$，又 $y=1$，}
\step{由 $x^2+x+1=1$ 得 $x^2+x=0$，解得 $x=0$ 或 $x=-1$，}
\step{所以 $A\cap B=\{(0,1),(-1,1)\}$.}
\stepm{(2)}{当 $a=1$、$b=2k-1$、$c=k^2$ 时，$f(x)=x^2+(2k-1)x+k^2$.}
\step{设两根为 $x_1$、$x_2$，}
\step{则“两根都大于 1”等价于 $\Delta>0$、$f(1)>0$、$-\dfrac b{2a}>1$.}
\step{由 $\Delta=(2k-1)^2-4k^2=1-4k>0$ 得 $k<\dfrac14$；}
\step{由 $f(1)=1+(2k-1)+k^2=k^2+2k=k(k+2)>0$ 得 $k<-2$ 或 $k>0$；}
\step{由 $-\dfrac{2k-1}2=\dfrac{1-2k}2>1$ 得 $k<-\dfrac12$.}
\step{三者取交集得 $k<-2$，即 $k$ 的取值范围是 $(-\infty,-2)$.}
\stepm{(3)}{设 $x_0\in\Z$ 是方程 $f(x)=0$ 的整数解.}
\step{因为 $x_0(x_0-1)$ 是两个相邻整数之积，恒为偶数，所以 $ax_0(x_0-1)$ 为偶数.}
\step{而 $f(x_0)=ax_0^2+bx_0+c=ax_0(x_0-1)+(a+b)x_0+c$，}
\step{所以 $f(x_0)$ 与 $(a+b)x_0+c$ 同奇偶.}
\step{若 $x_0$ 为偶数，则 $(a+b)x_0$ 为偶数，$f(x_0)$ 与 $c=f(0)$ 同奇偶，为奇数；}
\step{若 $x_0$ 为奇数，则 $(a+b)x_0$ 与 $a+b$ 同奇偶，}
\step{$f(x_0)$ 与 $(a+b)+c=f(1)$ 同奇偶，也为奇数.}
\step{综上，$f(x_0)$ 恒为奇数，不可能等于 0，与 $f(x_0)=0$ 矛盾，}
\step{所以方程 $f(x)=0$ 无整数解.}
\solend
\fi

\ansspace[8]

% 压轴题（题干＋三小问）与作答区 10 行须同页
\ansneed{20cm}
\item 已知 $X$ 为有限集，且任取 $x\in X$，均有 $x\in\Z$，称 $X$ 为非空有限整数集. 定义：$X_k=\{x+k\mid x\in X\}$，$X^2=\{y\mid y=x^2,x\in X\}$.

\begin{enumerate}[label=(\arabic*)]
\item 若 $X=\{-2,2\}$，分别写出 $X_2$ 和 $X^2$（用列举法表示）；
\item 若 $X=\{-1,0,1\}$，$X_1\subseteq Y$，$(Y_{-1})^2\subseteq X$，求集合 $Y$ 中所有的元素；
\item 对于非空有限整数集 $A$、$B$，若 $A_k\subseteq B$，$(B^2)_{-2}\subseteq A$，$B$ 中的元素均为非负整数，$k$ 为整数，求证：

① $A$、$B$ 中元素的个数相同；

② 若集合 $A$ 中的元素个数为奇数，则 $k$ 为偶数.
\end{enumerate}

\ifanswer
\solbegin
\stepm{(1)}{由题意得 $X_2=\{0,4\}$，$X^2=\{4\}$；}
\stepm{(2)}{由题意得 $X_1=\{0,1,2\}$，即 $\{0,1,2\}\subseteq Y$，$Y_{-1}=\{y-1\mid y\in Y\}$，}
\step{$(Y_{-1})^2=\{(y-1)^2\mid y\in Y\}\subseteq X=\{-1,0,1\}$，}
\step{所以 $(y-1)^2=0$ 或 $(y-1)^2=1$，解得 $y=1$ 或 $y=2$ 或 $y=0$，}
\step{即 $y\in\{0,1,2\}$，即 $Y\subseteq\{0,1,2\}$，所以 $Y=\{0,1,2\}$；}
% 原卷 (3) 解析首行“(3)”与“①”连排，本稿合并记作 \stepm{(3)①}
\stepm{(3)①}{因为 $A_k=\{x+k\mid x\in A\}$，函数 $g(x)=x+k$ 在全体整数集上严格增，}
\step{所以 $\forall x_1$、$x_2\in A$，$x_1\neq x_2$，$g(x_1)\neq g(x_2)$，所以 $|A_k|=|A|$，}
\step{因为 $A_k\subseteq B$，所以 $|A|\leqslant|B|$，}
\step{设 $b\in B$，$B$ 中元素均为非负整数，$\forall b_1\neq b_2\in B$，$b_1^2\neq b_2^2$，}
\step{故函数 $f(b)=b^2-2$ 在 $B$ 上单射，所以 $|(B^2)_{-2}|=|B|$，}
\step{又因为 $(B^2)_{-2}\subseteq A$，所以 $|B|\leqslant|A|$，所以 $|A|=|B|$.}
\stepm{②}{由①得 $A_k\subseteq B$ 且 $|A_k|=|A|=|B|$，$(B^2)_{-2}\subseteq A$ 且 $|(B^2)_{-2}|=|B|=|A|$，}
\step{所以 $A_k=B$，$(B^2)_{-2}=A$，}
\step{所以 $B=A_k=\{a+k\mid a\in A\}$，$A=(B^2)_{-2}=\{b^2-2\mid b\in B\}$，}
\step{所以 $A=\{(a+k)^2-2\mid a\in A\}$，}
\step{所以映射 $h(a)=(a+k)^2-2$ 是集合 $A$ 到自身的一个双射.}
\step{假设 $k$ 是奇数，当 $a$ 为偶数时，$a+k$ 为奇数，$(a+k)^2-2$ 为奇数，}
\step{当 $a$ 为奇数时，$a+k$ 为偶数，$(a+k)^2-2$ 为偶数，}
\step{所以此时映射 $h(a)=(a+k)^2-2$，将集合 $A$ 中的奇数映射为偶数，}
\step{将集合 $A$ 中的偶数映射为奇数，即集合 $A$ 中的奇数元素和偶数元素必须配对，}
\step{此时集合 $A$ 的元素个数则为偶数，}
\step{与题设中“集合 $A$ 中的元素个数为奇数”矛盾.}
\step{故假设不成立，所以 $k$ 是偶数.}
\solend
\fi

\ansspace*[10]

\end{enumerate}

\end{document}
"
%
% 原始材料：微信公众号“魔都数学加油站”文章，作者破晓，连载“27学年高一67”；
% 原文标题“27学年高一67——2026-2027年上海市华师大一附中高一上10月月考”；
% 原文链接：https://mp.weixin.qq.com/s/Kdhx07TVjLoGTcIlnRzgcQ。
% 页面标注发布时间 2026-10-18 20:33（上海）。
% 本次转录只据 `原图/高一67_华师大一附中_10月月考/` 的 13 张 PNG 逐页读图。
% 正文为 13 张 PNG 页（`原图/高一67_华师大一附中_10月月考/01..13.png`）：
% 第 01、04、13 页 2560×3620，第 02、03、05—12 页 4762×6735
% （**同一篇各页分辨率不同**），均按页序读取。13 页全为试卷正文页（卷首标题在第 1 页页顶，
% 卷末解答题第 21 题解析止于第 13 页），无封面页、目录页、二维码或宣传页。
% 粉色斜排“魔都数学加油站”为卷面水印，与公众号页脚一并未录入。
%
% 卷首标题：照卷面印作“2026--2027 年华师大一附中高一上 10 月月考”（卷面为黑体加粗一行），
% 年份区间按全稿惯例排 en dash。原卷只有一行标题、无副标题。
%
% 篇幅与题号：原卷自第 3 题起编号（第 1、2 题不在本卷内）。填空题 3—12（10 题）、
% 选择题 13—16（4 题）、解答题 17—21（5 题），共 19 题。本稿沿用原节与原题号，
% 只改排为全稿统一的大节样式（原卷小节标题为“一、填空题”“二、选择题”“三、解答题”）。
%
% 讲评版：原卷每题下直接印【解析】。**全卷 19 题中 18 题都印【解析】（无一处只印
% 【答案】）**，按全稿口径这类题不另提 \ans{}；各题解析末句的结论（第 13、14、16 题
% 的“故选 B／D／B”、各填空与解答题的“故/综上…”）照录在解析内。
% **另有 1 题答案印在题干括号内**——第 15 题（“(  D  )”，且该题无【答案】/【解析】段），
% 本稿按 高一37 第 13 题、高一63 第 14 题成例把括号留空、答案另提为 \ans{$D$}。
%
% 副标题：原卷未印副标题。本稿按“知识模块”口径标作“集合与常用逻辑用语、不等式”：
% 19 题中集合与常用逻辑用语 12 题（第 3、4、5、7、10、11、12、13、15、16、18、21 题：
% 补集、并集、子集、充分条件、容斥计数、子集计数、新定义二元运算、条件判断、
% 反证法假设、新定义生成集、集合关系与二次方程、新定义 X_k 与 X^2），
% 不等式 7 题（第 6、8、9、14、17、19、20 题：存在量词不等式、解集反求系数、
% 解集为空、不等式性质、解不等式（组）、二次不等式与基本不等式、二次方程根的分布），
% 两类均超过三分之一，故并标。
%
% 与原卷的排版差异：
%   ① 原卷每题下直接印【解析】，本稿按全稿惯例将解析置于答案版，试卷版保持干净；
%   ② 原卷未印副标题，本稿按“知识模块”口径补出；
%   ③ 原卷填空横线约四个字宽，本稿按全稿惯例统一排 \blankline[4.4em]；
%   ④ 原卷实数集排斜体 $R$（第 14 题题干、第 17 题 (1) 题干、第 18 题解析、第 20 题题干），
%      整数集排斜体 $Z$（第 20 题 (3) 题干与解析、第 21 题题干），正整数集排粗体
%      $N^*$（第 11 题解析、第 19 题解析），本稿按全稿惯例统一改用 $\R$、$\Z$、$\N^*$
%      （记号归一，粗体／斜体不作区分）；空集原卷作 $\varnothing$，照录；
%   ⑤ 原卷在数学式内多处用半角逗号（如第 8 题 $(1,2)$ 与各 cases、第 11 题解析
%      $P_n=\{1,2,\cdots,n\}$,$n\in N^*$、第 12 题解析 $u=(\cdots)$,$v=(\cdots)$ 与
%      “2020,个 1”、第 19 题解析 $\{x|40<x<80,x\in N^*\}$），均照录未强改；
%      描述法集合的竖线照原卷排 \mid；句末句点按全稿惯例排西文句点；
%   ⑥ 原卷补集记号为上划线 $\bar A$（第 3 题题干与解析、第 11 题题干），照录
%      （同 高一64/65 的 $\bar B$ 成例）；
%   ⑦ 第 12 题新定义的运算记号 $*$ 照录（$\alpha*\beta$、$u*u$、$u*v$）；
%   ⑧ 第 17—21 题的子问原卷各自另起一行，本稿按全稿惯例用 enumerate 的 (1)(2)(3)；
%      解析中的子问标号一律用 \stepm{(1)}{…}；第 5 题解析的①②、第 21 题 (3) 解析的①②
%      按全稿惯例作 \stepm{①}{…}、\stepm{②}{…}（第 21 题 (3) 首行原卷“(3)”与“①”连排，
%      本稿合并记作 \stepm{(3)①}{…}）；
%   ⑨ 选择题作答括号原卷排半角“(　　)”（第 13、14、16 题括号内空白、第 15 题括号内印 D），
%      本稿按全稿惯例排 \mbox{（\quad）}；第 15 题括号内的答案移入 \ans{}，见上；
%   ⑩ 选项行排照原卷：第 13、15 题两行两条，第 14、16 题一行四条；
%   ⑪ 每道解答题原卷未留作答空白，本稿按全稿惯例加 \ansspace（末题用 \ansspace*）；
%      第 20、21 题题干长，各用一处 \ansneed 让题干与作答区不拆页（\ansneed 只作用于试卷版）。
%
% 原卷笔误/存疑（均照录未改，正文对应处另有 % 注）：
%   ① 第 3 题解析末句印作“则 $\bar A=(-\infty,0)\cup[2,3)$.”：按题干
%      $U=\{x\mid 1/x\leqslant 1/2\}$（即 $x<0$ 或 $x\geqslant2$）与 $A=[3,+\infty)$，
%      $\bar A$ 应为 $(-\infty,0)\cup\{2\}$；且与本解析首句自求之 $U$ 相矛盾（$3\notin U$），
%      疑为笔误；照录未改。
%   ② 第 12 题题干 $\alpha*\beta$ 定义式括号不配对：第四组“$+\cdots+x_{2020}+y_{2020}
%      -x_{2020}y_{2020}$”缺左括号、末尾连排两个右括号“))”；照录未改。
%   ③ 第 12 题解析“所以…中有 2020,个 1”一行，“$\cdots$”后脱逗号（作
%      “$x_2,\cdots x_{2020}$”，第 12 行“而…”同式样），且“2020”后衍一半角逗号
%      （作“2020,个 1”）；照录未改。
%   ④ 第 16 题解析乘积链印作“$a_1a_2<a_1a_3<a_1a_4<a_1a_5<a_2a_3<a_3a_3<a_{14}a_5$”：
%      “$a_3a_3$”同足标自乘、与单调链不合，“$a_{14}a_5$”下标 14 与 $a_1$—$a_5$ 记号
%      不合（疑为 $a_2a_4$、$a_4a_5$ 之误）；照录未改。
%   ⑤ 第 5 题题干作“$Q\subset P$”（真子集记号），解析两处均作“$Q\subseteq P$”，
%      同一题前后记号不一；照录（本稿分别排 \subset、\subseteq）。
%   ⑥ 第 17 题 (1) 解析“当 $a=1$ 时，原不等式为 $(x+1)^2<0$ 此时…”与
%      “当 $0<a<1$ 时，$-\frac1a<-1$ 此时…”两处“此时”前均脱标点；照录未改。
%   ⑦ 第 20 题 (2) 题干“若 $a=1$，$b=2k-1$，$c=k^2$，若方程 $f(x)=0$ 有两个大于 1 的根”
%      两个“若”连用，疑衍其一；照录未改。
%
% 插图：无（全卷 13 页无任何几何图形、无“如图/右图/图中”字样）。
%
% 行数核对：逐题数 \step/\stepm 条数与原卷【解析】印刷行数，18 题逐一相符
% （第 15 题无解析段，只提 \ans{}），合计 169 条 / 169 行
% （Q3—Q12：2、3、4、3、1、3、5、3、13、15；Q13—Q16：2、3、0、25；
% Q17—Q21：19、16、9、20、23）。
%
\documentclass[a4paper,11pt]{article}
\usepackage{common}

\begin{document}

\examtitlest[3]{2026--2027 年华师大一附中高一上 10 月月考}{集合与常用逻辑用语、不等式}

\bigsec{一、填空题}

\begin{enumerate}
\setcounter{enumi}{2}

\item 设全集 $U=\left\{x\mid\dfrac1x\leqslant\dfrac12\right\}$，$A=[3,+\infty)$，则 $\bar A=$\blankline[4.4em].

\ifanswer
\solbegin
\step{$\dfrac1x\leqslant\dfrac12$，解得 $x<0$ 或 $x\geqslant2$，故 $U=(-\infty,0)\cup[2,+\infty)$，}
% 原卷末句作 $\bar A=(-\infty,0)\cup[2,3)$，按题干 $U$ 应为 $(-\infty,0)\cup\{2\}$，与本解析首句矛盾，疑笔误；照录（见文件头“原卷笔误/存疑”①）
\step{$A=[3,+\infty)$，则 $\bar A=(-\infty,0)\cup[2,3)$.}
\solend
\fi

\item 已知集合 $M=\{y\mid y=x^2-2x-1\}$，$N=\{y\mid y=x^2-4x+1\}$，则 $M\cup N=$\blankline[4.4em].

\ifanswer
\solbegin
\step{$M=\{y\mid y=x^2-2x-1\}=\{y\mid y=(x-1)^2-2\}=\{y\mid y\geqslant-2\}$，}
\step{$N=\{y\mid y=x^2-4x+1\}=\{y\mid y=(x-2)^2-3\}=\{y\mid y\geqslant-3\}$，}
\step{则 $M\cup N=\{y\mid y\geqslant-3\}$.}
\solend
\fi

\item 已知 $P=\{x\mid x^2+x-6=0\}$，$Q=\{x\mid mx-1=0\}$，且 $Q\subset P$，则实数 $m$ 的值\blankline[4.4em].

\ifanswer
\solbegin
\step{$P=\{2,-3\}$，$Q=\{x\mid mx=1\}$；}
\stepm{①}{$m=0$ 时，$Q=\varnothing$，满足 $Q\subseteq P$；}
\stepm{②}{$m\neq0$ 时，要使 $Q\subseteq P$，则 $\dfrac1m=2$ 或 $-3$，所以 $m=\dfrac12$ 或 $-\dfrac13$；}
\step{所以 $m=0$ 或 $\dfrac12$ 或 $-\dfrac13$.}
\solend
\fi

\item 已知“存在正实数 $x$，使得关于 $x$ 的不等式 $ax-1>0$ 成立”为真命题，则实数 $a$ 的取值范围是\blankline[4.4em].

\ifanswer
\solbegin
\step{存在正实数 $x$，使得关于 $x$ 的不等式 $ax-1>0$ 成立为真命题，}
\step{即存在正实数 $x$，使得 $a>\dfrac1x$，因为 $\dfrac1x>0$，所以 $a>0$，}
\step{即 $a$ 的范围为 $(0,+\infty)$.}
\solend
\fi

\item 设 $\alpha:3\leqslant x\leqslant5$，$\beta:2m-1\leqslant x<m+4$，若 $\alpha$ 是 $\beta$ 的充分条件，则实数 $m$ 的取值范围是\blankline[4.4em].

\ifanswer
\solbegin
\step{若 $\alpha$ 是 $\beta$ 的充分条件，则 $\begin{cases}2m-1\leqslant3\\ m+4>5\end{cases}$，解得 $1<m\leqslant2$.}
\solend
\fi

\item 对于实数 $a$、$b$、$c$，已知关于 $x$ 的不等式 $ax^2+bx+c>0$ 的解集是 $(1,2)$，则关于 $x$ 的不等式 $cx^2-bx+a\leqslant0$ 的解集是\blankline[4.4em].

\ifanswer
\solbegin
\step{关于 $x$ 的不等式 $ax^2+bx+c>0$ 的解集是 $(1,2)$，则 $\begin{cases}a<0\\ -\dfrac ba=3\\ \dfrac ca=2\end{cases}$，则 $\begin{cases}b=-3a\\ c=2a\end{cases}$，}
\step{则关于 $x$ 的不等式 $cx^2-bx+a=2ax^2+3ax+a\leqslant0$，则 $2x^2+3x+1\geqslant0$，}
\step{则 $x\leqslant-1$ 或 $x\geqslant-\dfrac12$，则解集为 $(-\infty,-1]\cup\left[-\dfrac12,+\infty\right)$.}
\solend
\fi

\item 已知关于 $x$ 的不等式 $kx^2+kx+k+2\leqslant0$ 的解集为 $\varnothing$，则实数 $k$ 的取值范围为\blankline[4.4em].

\ifanswer
\solbegin
\step{当 $k=0$ 时，不等式化为 $2\leqslant0$，解集为 $\varnothing$，符合题意；}
\step{当 $k\neq0$ 时，若关于 $x$ 的不等式 $kx^2+kx+k+2\leqslant0$ 的解集为 $\varnothing$，}
\step{则 $\begin{cases}k>0\\ \Delta=k^2-4k(k+2)<0\end{cases}$，化简得 $\begin{cases}k>0\\ -3k^2-8k<0\end{cases}$，即 $\begin{cases}k>0\\ k(3k+8)>0\end{cases}$，}
\step{解得 $k>0$.}
\step{综上所述，$k\geqslant0$，即实数 $k$ 的取值范围为 $[0,+\infty)$.}
\solend
\fi

\item 对班级若干名学生调查对 $A$、$B$ 两事件的态度，有如下结果：赞成 $A$ 的人数是全体的五分之三，其余的不赞成；赞成 $B$ 的比赞成 $A$ 多 3 人，其余的不赞成. 已知对 $A$、$B$ 都赞成的学生有 18 人，对 $A$、$B$ 都不赞成的学生数比对 $A$、$B$ 都赞成的学生数的三分之一多 1 人，则班级总人数是\blankline[4.4em]人.

\ifanswer
\solbegin
\step{设班级总人数是 $x$ 人，则赞成 $A$ 的人数是 $\dfrac35x$，赞成 $B$ 的人数是 $\dfrac35x+3$，}
\step{$A$、$B$ 都不赞成的人数为 $18\times\dfrac13+1=7$ 人，}
\step{由容斥原理得 $x=\left(\dfrac35x+\dfrac35x+3-18\right)+7$，解得 $x=40$，即班级总人数是 40 人.}
\solend
\fi

\item 设全集 $U=\{1,2,3,4,5,6,7,8,9,10,11,12\}$，若集合 $A$ 满足：① $A\subseteq U$；②若 $x\in A$，必有 $2x\notin A$；③若 $x\in\bar A$，必有 $2x\notin\bar A$，则这样的集合 $A$ 有\blankline[4.4em]个.

\ifanswer
\solbegin
\step{设集合 $P_n=\{1,2,\cdots,n\}$,$n\in\N^*$，}
\step{任取 $x\in P_n$，将 $x$ 除以 2，若商仍为偶数，再除以 2，$\cdots$，}
\step{经过 $k$ 次以后，商必为奇数，此时记商为 $m$，}
\step{于是 $x=m\cdot2^k$，其中 $m$ 为奇数，$k\in\N^*$.}
\step{由条件得，若 $m\in A$，则 $x\in A$ 等价于 $k$ 为偶数；}
\step{若 $m\notin A$，则 $x\in A$ 等价于 $k$ 为奇数；}
\step{于是 $x$ 是否属于 $A$ 由 $m$ 是否属于 $A$ 确定；}
\step{设 $Q_n$ 是 $P_n$ 中所有奇数的集合，}
\step{设 $f(n)$ 为集合 $P_n$ 中满足题干条件的集合 $A$ 的个数，}
\step{因此 $f(n)$ 等于 $Q_n$ 的子集个数.}
\step{当 $n$ 为偶数（或奇数）时，$P_n$ 中奇数的个数是 $\dfrac n2$（或 $\dfrac{n+1}2$），}
\step{所以 $f(n)=\begin{cases}2^{\frac n2},n\text{为偶数}\\[2pt] 2^{\frac{n+1}2},n\text{为奇数}\end{cases}$，所以 $P_{12}=\{1,2,3,4,5,6,7,8,9,10,11,12\}$，}
\step{即同时满足三个条件的集合 $A$ 的个数为 $2^{\frac{12}2}=2^6=64$.}
\solend
\fi

\item 设集合 $\Omega=\{(x_1,x_2,\cdots,x_{2020})\mid x_i\in\{0,1\},i=1,2,\cdots,2020\}$，对于 $\Omega$ 中的任意两个元素 $\alpha=(x_1,x_2,\cdots,x_{2020})$，$\beta=(y_1,y_2,\cdots,y_{2020})$，记
% 原卷此式第四组缺左括号、末尾连排两个右括号“))”，括号不配对；照录（见文件头“原卷笔误/存疑”②）
$\alpha*\beta=(x_1+y_1-x_1y_1)+(x_2+y_2-x_2y_2)+\cdots+x_{2020}+y_{2020}-x_{2020}y_{2020}))$，设 $u$、$v\in\Omega$，若 $u*u+v*v=2020$，则 $u*v$ 的最大值和最小值之和为\blankline[4.4em].

\ifanswer
\solbegin
\step{设 $u=(x_1,x_2,\cdots,x_{2020})$,$v=(y_1,y_2,\cdots,y_{2020})$，}
\step{因为 $x_i\in\{0,1\}$，所以 $x_i(x_i-1)=0$，}
\step{所以 $u*u=(x_1+x_1-x_1^2)+(x_2+x_2-x_2^2)+\cdots+(x_{2020}+x_{2020}-x_{2020}^2)$}
\step{$=x_1+x_2+\cdots+x_{2020}$，}
\step{$v*v=(y_1+y_1-y_1^2)+(y_2+y_2-y_2^2)+\cdots+(y_{2020}+y_{2020}-y_{2020}^2)$}
\step{$=y_1+y_2+\cdots+y_{2020}$，}
\step{由 $u*u+v*v=2020$，得 $x_1+x_2+\cdots+x_{2020}+y_1+y_2+\cdots+y_{2020}=2020$，}
% 原卷此行“$\cdots$”后脱逗号、 “2020”后衍一半角逗号（“2020,个 1”）；照录（见文件头“原卷笔误/存疑”③）
\step{所以 $x_1,x_2,\cdots x_{2020},y_1,y_2,\cdots y_{2020}$ 中有 2020,个 1，2020 个 0，}
\step{$u*v=(x_1+y_1-x_1y_1)+(x_2+y_2-x_2y_2)+\cdots+(x_{2020}+y_{2020}-x_{2020}y_{2020})$}
\step{$=2020-(x_1y_1+x_2y_2+\cdots+x_{2020}y_{2020})$，}
\step{注意到只有 $x_i=y_i=1$ 时，$x_iy_i=1$，否则 $x_iy_i=0$，}
% 此行“$\cdots$”后同样脱逗号，与上一处同式样；照录（见文件头“原卷笔误/存疑”③）
\step{而 $x_1,x_2,\cdots x_{2020},y_1,y_2,\cdots y_{2020}$ 中有 2020 个 1，2020 个 0，}
\step{所以满足 $x_iy_i=1$ 的最多有 1010 个，所以 $u*v\geqslant2020-1010=1010$，}
\step{满足 $x_iy_i=0$ 的最多有 2020 个，所以 $u*v\leqslant2020-0=2020$，}
\step{则 $u*v$ 的最大值和最小值之和为 $2020+1010=3030$.}
\solend
\fi

\end{enumerate}

\bigsec{二、选择题}

\begin{enumerate}
\setcounter{enumi}{12}

\item 实数 $x$ 与 $y$，“$x+y>2$”的\mbox{（\quad）}条件是“$x>1$ 或 $y>1$”.

\keepwithprev
A. 充分非必要\qquad B. 必要非充分

C. 充要\qquad D. 既非充分又非必要

\ifanswer
\solbegin
\step{当 $x>1$ 或 $y>1$ 时，$x+y>2$ 不一定成立，充分性不成立；}
\step{当 $x+y>2$ 时，$x>1$ 或 $y>1$ 一定成立，必要性成立. 故选 $B$.}
\solend
\fi

\item 设 $a,b\in\R$，且 $a<0$，$-1<b<0$，那么下列不等式成立的是\mbox{（\quad）}

\keepwithprev
A. $a>ab>ab^2$\qquad B. $ab^2>ab>a$\qquad C. $ab>a>ab^2$\qquad D. $ab>ab^2>a$

\ifanswer
\solbegin
\step{因为 $-1<b<0$，$a<0$，所以 $ab>0$，$0<b^2<1$.}
\step{所以 $ab-ab^2=ab(1-b)>0$，$ab^2-a=a(b^2-1)>0$，}
\step{所以 $ab>ab^2>a$. 故选 $D$.}
\solend
\fi

% 第 15 题答案原卷印在题干括号内（“(  D  )”），且无【答案】/【解析】段；
% 按 高一37 第 13 题、高一63 第 14 题成例括号留空、答案另提 \ans{$D$}
\item 用反证法证明：“平面四边形中至少有一个内角不超过 90°”时，下列“假设”中正确的是\mbox{（\quad）}

\keepwithprev
A. 假设有两个内角超过 90°\qquad B. 假设有三个内角超过 90°

C. 假设至多有两个内角超过 90°\qquad D. 假设四个内角都超过 90°

\ifanswer
\ans{$D$}
\fi

\item 设 $A$ 为非空实数集，集合 $B=\{xy\mid x\in A,y\in A,x\neq y\}$，则对于以下两个结论说法正确的是\mbox{（\quad）}

结论①：若 $A$ 是五元集，则 $B$ 中的元素最多 10 个；

结论②：存在四元集 $A$，使得 $B=\{2,4,5,8,10,16\}$.

\keepwithprev
A. ①对②对\qquad B. ①对②错\qquad C. ①错②对\qquad D. ①错②错

\ifanswer
\solbegin
\step{对于①，想让元素个数达到最大值 10，则须避免产生符号不同的乘积，}
\step{否则正负分开计数容易导致总数受限，且容易产生相等值，}
\step{也最好避免 0 的出现，}
\step{故设 $A=\{a_1,a_2,a_3,a_4,a_5\}$，$0<a_1<a_2<a_3<a_4<a_5$.}
% 原卷此链“$a_3a_3$”同足标自乘、“$a_{14}a_5$”下标 14 均与文意不合，疑为 $a_2a_4$、$a_4a_5$ 之误；照录（见文件头“原卷笔误/存疑”④）
\step{因为 $a_1a_2<a_1a_3<a_1a_4<a_1a_5<a_2a_3<a_3a_3<a_{14}a_5$，$a_2a_3<a_2a_4<a_3a_4$，}
\step{所以 $B$ 中元素个数小于等于 10 个，}
\step{如设 $A=\{1,2,3,5,7\}$，则 $B=\{2,3,5,6,7,10,14,15,21,35\}$，}
\step{所以集合 $B$ 中元素个数的最大值为 10 个，故①正确；}
\step{对于②，由于集合 $B$ 中没有 0，所以集合 $A$ 中没有 0；}
\step{设集合 $A=\{a,b,c,d\}$，且 $a<b<c<d$，}
\step{当 $a<0<b<c<d$ 时，$a\times b<0$，$a\times c<0$，$a\times d<0$，}
\step{故集合 $B$ 至少有 3 个负数，与 $B=\{2,4,5,8,10,16\}$ 矛盾；}
\step{当 $a<b<0<c<d$ 时，$a\times b>0$，$c\times d>0$，}
\step{$a\times c<0$，$a\times d<0$，$b\times c<0$，$b\times d<0$，}
\step{故集合 $B$ 至少有 4 个负数，与 $B=\{2,4,5,8,10,16\}$ 矛盾；}
\step{当 $a<b<c<0<d$ 时，$a\times b>0$，$a\times c>0$，$b\times c>0$，}
\step{$a\times d<0$，$b\times d<0$，$c\times d<0$，}
\step{故集合 $B$ 至少有 3 个负数，与 $B=\{2,4,5,8,10,16\}$ 矛盾；}
\step{当 $0<a<b<c<d$ 时，则集合 $B=\{ab,ac,ad,bc,bd,cd\}$，}
\step{则必有 $ab=2$，$cd=16$，其 4 个正实数的乘积 $abcd=32$，}
\step{又 $ad=5$，$bc=8$ 或 $ad=8$，$bc=5$，其 4 个正实数的乘积 $abcd=40$，矛盾；}
\step{当 $a<b<c<d<0$ 时，与 $0<a<b<c<d$ 时的证明方法相同，}
\step{得集合 $A$ 也不存在；}
\step{综上不存在四元集 $A$，使其生成集 $B=\{2,4,5,8,10,16\}$，故②错误.}
\step{故选 $B$.}
\solend
\fi

\end{enumerate}

\bigsec{三、解答题}

\begin{enumerate}
\setcounter{enumi}{16}

\item 解下列不等式（组）

\begin{enumerate}[label=(\arabic*)]
\item 解不等式：$(ax+1)(x+1)<0(a\in\R)$.
\item 解不等式组：$\begin{cases}\dfrac{2x+4}{x+4}\geqslant1\\[4pt] (x+1)(2-|x|)\geqslant0\end{cases}$.
\end{enumerate}

\ifanswer
\solbegin
\stepm{(1)}{当 $a=0$ 时，原不等式为 $x+1<0$，即 $x<-1$，}
\step{所以此时不等式的解集为 $\{x\mid x<-1\}$；}
\step{当 $a\neq0$ 时，令 $(ax+1)(x+1)=0$，解得 $x_1=-\dfrac1a$,$x_2=-1$，}
\step{当 $a>1$ 时，$-\dfrac1a>-1$，此时不等式的解集为 $\left\{x\mid-1<x<-\dfrac1a\right\}$；}
% 以下两行“此时”前原卷均无标点；照录（见文件头“原卷笔误/存疑”⑥）
\step{当 $a=1$ 时，原不等式为 $(x+1)^2<0$ 此时不等式的解集为 $\varnothing$；}
\step{当 $0<a<1$ 时，$-\dfrac1a<-1$ 此时不等式的解集为 $\left\{x\mid-\dfrac1a<x<-1\right\}$；}
\step{当 $a<0$ 时，$-\dfrac1a>-1$，此时不等式的解集为 $\left\{x\mid x<-1\ \text{或}\ x>-\dfrac1a\right\}$.}
\step{综上所述，当 $a=0$ 时，$\{x\mid x<-1\}$；}
\step{当 $a>1$ 时，$\left\{x\mid-1<x<-\dfrac1a\right\}$；}
\step{当 $a=1$ 时，$\varnothing$；}
\step{当 $0<a<1$ 时，$\left\{x\mid-\dfrac1a<x<-1\right\}$；}
\step{当 $a<0$ 时，$\left\{x\mid x<-1\ \text{或}\ x>-\dfrac1a\right\}$.}
\stepm{(2)}{不等式 $\dfrac{2x+4}{x+4}\geqslant1$ 可化为 $\dfrac{2x+4}{x+4}-1\geqslant0$，即 $\dfrac x{x+4}\geqslant0$，}
\step{等价于 $\begin{cases}x(x+4)\geqslant0\\ x+4\neq0\end{cases}$，解得 $x<-4$ 或 $x\geqslant0$，}
\step{不等式 $(x+1)(2-|x|)\geqslant0$ 等价于 $\begin{cases}(x+1)(2-x)\geqslant0,x\geqslant0\\ (x+1)(2+x)\geqslant0,x<0\end{cases}$，}
\step{解不等式 $(x+1)(2-x)\geqslant0$，$x\geqslant0$，得 $0\leqslant x\leqslant2$，}
\step{解不等式 $(x+1)(2+x)\geqslant0$，$x<0$，得 $x\leqslant-2$ 或 $-1\leqslant x<0$，}
\step{所以 $(x+1)(2-|x|)\geqslant0$ 的解集为 $\{x\mid x\leqslant-2\ \text{或}\ -1\leqslant x\leqslant2\}$，}
\step{综上所述，$\begin{cases}\dfrac{2x+4}{x+4}\geqslant1\\[4pt] (x+1)(2-|x|)\geqslant0\end{cases}$ 的解集为 $\{x\mid x<-4\ \text{或}\ 0\leqslant x\leqslant2\}$.}
\solend
\fi

\ansspace[8]

% 第 18 题题干＋两小问与作答区 6 行须同页
\ansneed{20cm}
\item 已知 $a$ 为实数，设 $f(x)=x^2+2(a+1)x+1$.

\begin{enumerate}[label=(\arabic*)]
\item 已知对于任意实数 $x$，$\sqrt{f(x)}$ 均有意义，求实数 $a$ 的取值范围；
\item 设 $A=\{x\mid f(x)+a^2-2=0\}$，$B=\{-4,0\}$，$A\cup B=B$，求实数 $a$ 的取值范围.
\end{enumerate}

\ifanswer
\solbegin
\stepm{(1)}{$\sqrt{f(x)}$ 对任意实数 $x$ 都有意义，}
\step{即 $f(x)=x^2+2(a+1)x+1\geqslant0$ 对任意 $x\in\R$ 恒成立，}
\step{所以 $\Delta=4(a+1)^2-4\times1\times1\leqslant0$，即 $(a+1)^2\leqslant1$，解得 $-2\leqslant a\leqslant0$，}
\step{故 $a$ 的取值范围是 $[-2,0]$.}
\stepm{(2)}{由 $f(x)+a^2-2=0$ 得 $x^2+2(a+1)x+a^2-1=0$.}
\step{由 $A\cup B=B$ 得 $A\subseteq B=\{-4,0\}$，即该方程无实根，}
\step{或它的实根都在 $\{-4,0\}$ 中.}
\step{当方程无实根时，$\Delta=4(a+1)^2-4(a^2-1)=8a+8<0$，解得 $a<-1$；}
\step{当方程有实根时，$a\geqslant-1$，设两根为 $x_1$、$x_2$，}
\step{由根与系数的关系得 $x_1+x_2=-2(a+1)$、$x_1x_2=a^2-1$，}
\step{此时两根只能在 $\{-4,0\}$ 中取值；}
\step{若两根为 0 与 0，由 $x_1x_2=0$ 得 $a=\pm1$，}
\step{再由 $x_1+x_2=0$ 得 $a=-1$；}
\step{若两根为 0 与 $-4$，由 $x_1x_2=0$ 得 $a=\pm1$，再由 $x_1+x_2=-4$ 得 $a=1$.}
\step{经检验，$a=-1$ 时 $A=\{0\}$，$a=1$ 时 $A=\{-4,0\}$，都满足 $A\subseteq B$.}
\step{综上，实数 $a$ 的取值范围为 $(-\infty,-1]\cup\{1\}$.}
\solend
\fi

\ansspace[8]

\ansneed{18cm}
\item 共享单车给市民出行带来了诸多便利，某公司购买了一批单车投放到某地给居民使用，据市场分析，这批单车的运营累计利润 $y$（单位：元）与运营天数 $x$（单位：天，$x$ 为正整数）满足关系：$y=-\dfrac12x^2+60x-800$.

\begin{enumerate}[label=(\arabic*)]
\item 要使这批单车的运营累计利润高于 800 元，求运营天数的取值范围；
\item 判断这批单车的日平均运营利润能否超过 20 元，并说明理由.
\end{enumerate}

\ifanswer
\solbegin
\stepm{(1)}{由题意得 $y=-\dfrac12x^2+60x-800>800$，}
\step{即 $x^2-120x+3200<0$，解得 $40<x<80$，}
\step{即运营天数的取值范围为 $\{x\mid40<x<80,x\in\N^*\}$.}
\stepm{(2)}{不能，理由如下：}
\step{这批单车的日平均运营利润为 $\dfrac yx=-\dfrac12x-\dfrac{800}x+60=-\dfrac12\left(x+\dfrac{1600}x\right)+60$，}
\step{因 $x>0$，则 $x+\dfrac{1600}x\geqslant2\sqrt{x\cdot\dfrac{1600}x}=80$，}
\step{当且仅当 $x=\dfrac{1600}x$，即 $x=40$ 时取等号，}
\step{此时 $\dfrac yx=-\dfrac12\left(x+\dfrac{1600}x\right)+60\leqslant-\dfrac12\times80+60=20$，}
\step{故这批单车的日平均运营利润不能超过 20 元.}
\solend
\fi

\ansspace[8]

% 第 20 题题干＋三小问与作答区 8 行须同页
\ansneed{20cm}
\item 设 $f(x)=ax^2+bx+c$($a\neq0$，$b$、$c\in\R$).

\begin{enumerate}[label=(\arabic*)]
\item 若 $a=b=c=1$，设 $A=\{(x,y)\mid y=f(x)\}$，$B=\{(x,y)\mid y=1\}$，用列举法表示 $A\cap B$；
% 原卷两个“若”连用（“若 a=1，…，若方程 f(x)=0 有两个大于 1 的根”），疑衍其一；照录（见文件头“原卷笔误/存疑”⑦）
\item 若 $a=1$，$b=2k-1$，$c=k^2$，若方程 $f(x)=0$ 有两个大于 1 的根，求实数 $k$ 的取值范围；
\item 若 $a$、$b$、$c\in\Z$，且 $f(0)$、$f(1)$ 均为奇数，求证：方程 $f(x)=0$ 无整数解.
\end{enumerate}

\ifanswer
\solbegin
\stepm{(1)}{当 $a=b=c=1$ 时，$f(x)=x^2+x+1$，}
\step{所以 $A=\{(x,y)\mid y=x^2+x+1\}$，又 $y=1$，}
\step{由 $x^2+x+1=1$ 得 $x^2+x=0$，解得 $x=0$ 或 $x=-1$，}
\step{所以 $A\cap B=\{(0,1),(-1,1)\}$.}
\stepm{(2)}{当 $a=1$、$b=2k-1$、$c=k^2$ 时，$f(x)=x^2+(2k-1)x+k^2$.}
\step{设两根为 $x_1$、$x_2$，}
\step{则“两根都大于 1”等价于 $\Delta>0$、$f(1)>0$、$-\dfrac b{2a}>1$.}
\step{由 $\Delta=(2k-1)^2-4k^2=1-4k>0$ 得 $k<\dfrac14$；}
\step{由 $f(1)=1+(2k-1)+k^2=k^2+2k=k(k+2)>0$ 得 $k<-2$ 或 $k>0$；}
\step{由 $-\dfrac{2k-1}2=\dfrac{1-2k}2>1$ 得 $k<-\dfrac12$.}
\step{三者取交集得 $k<-2$，即 $k$ 的取值范围是 $(-\infty,-2)$.}
\stepm{(3)}{设 $x_0\in\Z$ 是方程 $f(x)=0$ 的整数解.}
\step{因为 $x_0(x_0-1)$ 是两个相邻整数之积，恒为偶数，所以 $ax_0(x_0-1)$ 为偶数.}
\step{而 $f(x_0)=ax_0^2+bx_0+c=ax_0(x_0-1)+(a+b)x_0+c$，}
\step{所以 $f(x_0)$ 与 $(a+b)x_0+c$ 同奇偶.}
\step{若 $x_0$ 为偶数，则 $(a+b)x_0$ 为偶数，$f(x_0)$ 与 $c=f(0)$ 同奇偶，为奇数；}
\step{若 $x_0$ 为奇数，则 $(a+b)x_0$ 与 $a+b$ 同奇偶，}
\step{$f(x_0)$ 与 $(a+b)+c=f(1)$ 同奇偶，也为奇数.}
\step{综上，$f(x_0)$ 恒为奇数，不可能等于 0，与 $f(x_0)=0$ 矛盾，}
\step{所以方程 $f(x)=0$ 无整数解.}
\solend
\fi

\ansspace[8]

% 压轴题（题干＋三小问）与作答区 10 行须同页
\ansneed{20cm}
\item 已知 $X$ 为有限集，且任取 $x\in X$，均有 $x\in\Z$，称 $X$ 为非空有限整数集. 定义：$X_k=\{x+k\mid x\in X\}$，$X^2=\{y\mid y=x^2,x\in X\}$.

\begin{enumerate}[label=(\arabic*)]
\item 若 $X=\{-2,2\}$，分别写出 $X_2$ 和 $X^2$（用列举法表示）；
\item 若 $X=\{-1,0,1\}$，$X_1\subseteq Y$，$(Y_{-1})^2\subseteq X$，求集合 $Y$ 中所有的元素；
\item 对于非空有限整数集 $A$、$B$，若 $A_k\subseteq B$，$(B^2)_{-2}\subseteq A$，$B$ 中的元素均为非负整数，$k$ 为整数，求证：

① $A$、$B$ 中元素的个数相同；

② 若集合 $A$ 中的元素个数为奇数，则 $k$ 为偶数.
\end{enumerate}

\ifanswer
\solbegin
\stepm{(1)}{由题意得 $X_2=\{0,4\}$，$X^2=\{4\}$；}
\stepm{(2)}{由题意得 $X_1=\{0,1,2\}$，即 $\{0,1,2\}\subseteq Y$，$Y_{-1}=\{y-1\mid y\in Y\}$，}
\step{$(Y_{-1})^2=\{(y-1)^2\mid y\in Y\}\subseteq X=\{-1,0,1\}$，}
\step{所以 $(y-1)^2=0$ 或 $(y-1)^2=1$，解得 $y=1$ 或 $y=2$ 或 $y=0$，}
\step{即 $y\in\{0,1,2\}$，即 $Y\subseteq\{0,1,2\}$，所以 $Y=\{0,1,2\}$；}
% 原卷 (3) 解析首行“(3)”与“①”连排，本稿合并记作 \stepm{(3)①}
\stepm{(3)①}{因为 $A_k=\{x+k\mid x\in A\}$，函数 $g(x)=x+k$ 在全体整数集上严格增，}
\step{所以 $\forall x_1$、$x_2\in A$，$x_1\neq x_2$，$g(x_1)\neq g(x_2)$，所以 $|A_k|=|A|$，}
\step{因为 $A_k\subseteq B$，所以 $|A|\leqslant|B|$，}
\step{设 $b\in B$，$B$ 中元素均为非负整数，$\forall b_1\neq b_2\in B$，$b_1^2\neq b_2^2$，}
\step{故函数 $f(b)=b^2-2$ 在 $B$ 上单射，所以 $|(B^2)_{-2}|=|B|$，}
\step{又因为 $(B^2)_{-2}\subseteq A$，所以 $|B|\leqslant|A|$，所以 $|A|=|B|$.}
\stepm{②}{由①得 $A_k\subseteq B$ 且 $|A_k|=|A|=|B|$，$(B^2)_{-2}\subseteq A$ 且 $|(B^2)_{-2}|=|B|=|A|$，}
\step{所以 $A_k=B$，$(B^2)_{-2}=A$，}
\step{所以 $B=A_k=\{a+k\mid a\in A\}$，$A=(B^2)_{-2}=\{b^2-2\mid b\in B\}$，}
\step{所以 $A=\{(a+k)^2-2\mid a\in A\}$，}
\step{所以映射 $h(a)=(a+k)^2-2$ 是集合 $A$ 到自身的一个双射.}
\step{假设 $k$ 是奇数，当 $a$ 为偶数时，$a+k$ 为奇数，$(a+k)^2-2$ 为奇数，}
\step{当 $a$ 为奇数时，$a+k$ 为偶数，$(a+k)^2-2$ 为偶数，}
\step{所以此时映射 $h(a)=(a+k)^2-2$，将集合 $A$ 中的奇数映射为偶数，}
\step{将集合 $A$ 中的偶数映射为奇数，即集合 $A$ 中的奇数元素和偶数元素必须配对，}
\step{此时集合 $A$ 的元素个数则为偶数，}
\step{与题设中“集合 $A$ 中的元素个数为奇数”矛盾.}
\step{故假设不成立，所以 $k$ 是偶数.}
\solend
\fi

\ansspace*[10]

\end{enumerate}

\end{document}
"
%
% 原始材料：微信公众号“魔都数学加油站”文章，作者破晓，连载“27学年高一67”；
% 原文标题“27学年高一67——2026-2027年上海市华师大一附中高一上10月月考”；
% 原文链接：https://mp.weixin.qq.com/s/Kdhx07TVjLoGTcIlnRzgcQ。
% 页面标注发布时间 2026-10-18 20:33（上海）。
% 本次转录只据 `原图/高一67_华师大一附中_10月月考/` 的 13 张 PNG 逐页读图。
% 正文为 13 张 PNG 页（`原图/高一67_华师大一附中_10月月考/01..13.png`）：
% 第 01、04、13 页 2560×3620，第 02、03、05—12 页 4762×6735
% （**同一篇各页分辨率不同**），均按页序读取。13 页全为试卷正文页（卷首标题在第 1 页页顶，
% 卷末解答题第 21 题解析止于第 13 页），无封面页、目录页、二维码或宣传页。
% 粉色斜排“魔都数学加油站”为卷面水印，与公众号页脚一并未录入。
%
% 卷首标题：照卷面印作“2026--2027 年华师大一附中高一上 10 月月考”（卷面为黑体加粗一行），
% 年份区间按全稿惯例排 en dash。原卷只有一行标题、无副标题。
%
% 篇幅与题号：原卷自第 3 题起编号（第 1、2 题不在本卷内）。填空题 3—12（10 题）、
% 选择题 13—16（4 题）、解答题 17—21（5 题），共 19 题。本稿沿用原节与原题号，
% 只改排为全稿统一的大节样式（原卷小节标题为“一、填空题”“二、选择题”“三、解答题”）。
%
% 讲评版：原卷每题下直接印【解析】。**全卷 19 题中 18 题都印【解析】（无一处只印
% 【答案】）**，按全稿口径这类题不另提 \ans{}；各题解析末句的结论（第 13、14、16 题
% 的“故选 B／D／B”、各填空与解答题的“故/综上…”）照录在解析内。
% **另有 1 题答案印在题干括号内**——第 15 题（“(  D  )”，且该题无【答案】/【解析】段），
% 本稿按 高一37 第 13 题、高一63 第 14 题成例把括号留空、答案另提为 \ans{$D$}。
%
% 副标题：原卷未印副标题。本稿按“知识模块”口径标作“集合与常用逻辑用语、不等式”：
% 19 题中集合与常用逻辑用语 12 题（第 3、4、5、7、10、11、12、13、15、16、18、21 题：
% 补集、并集、子集、充分条件、容斥计数、子集计数、新定义二元运算、条件判断、
% 反证法假设、新定义生成集、集合关系与二次方程、新定义 X_k 与 X^2），
% 不等式 7 题（第 6、8、9、14、17、19、20 题：存在量词不等式、解集反求系数、
% 解集为空、不等式性质、解不等式（组）、二次不等式与基本不等式、二次方程根的分布），
% 两类均超过三分之一，故并标。
%
% 与原卷的排版差异：
%   ① 原卷每题下直接印【解析】，本稿按全稿惯例将解析置于答案版，试卷版保持干净；
%   ② 原卷未印副标题，本稿按“知识模块”口径补出；
%   ③ 原卷填空横线约四个字宽，本稿按全稿惯例统一排 \blankline[4.4em]；
%   ④ 原卷实数集排斜体 $R$（第 14 题题干、第 17 题 (1) 题干、第 18 题解析、第 20 题题干），
%      整数集排斜体 $Z$（第 20 题 (3) 题干与解析、第 21 题题干），正整数集排粗体
%      $N^*$（第 11 题解析、第 19 题解析），本稿按全稿惯例统一改用 $\R$、$\Z$、$\N^*$
%      （记号归一，粗体／斜体不作区分）；空集原卷作 $\varnothing$，照录；
%   ⑤ 原卷在数学式内多处用半角逗号（如第 8 题 $(1,2)$ 与各 cases、第 11 题解析
%      $P_n=\{1,2,\cdots,n\}$,$n\in N^*$、第 12 题解析 $u=(\cdots)$,$v=(\cdots)$ 与
%      “2020,个 1”、第 19 题解析 $\{x|40<x<80,x\in N^*\}$），均照录未强改；
%      描述法集合的竖线照原卷排 \mid；句末句点按全稿惯例排西文句点；
%   ⑥ 原卷补集记号为上划线 $\bar A$（第 3 题题干与解析、第 11 题题干），照录
%      （同 高一64/65 的 $\bar B$ 成例）；
%   ⑦ 第 12 题新定义的运算记号 $*$ 照录（$\alpha*\beta$、$u*u$、$u*v$）；
%   ⑧ 第 17—21 题的子问原卷各自另起一行，本稿按全稿惯例用 enumerate 的 (1)(2)(3)；
%      解析中的子问标号一律用 \stepm{(1)}{…}；第 5 题解析的①②、第 21 题 (3) 解析的①②
%      按全稿惯例作 \stepm{①}{…}、\stepm{②}{…}（第 21 题 (3) 首行原卷“(3)”与“①”连排，
%      本稿合并记作 \stepm{(3)①}{…}）；
%   ⑨ 选择题作答括号原卷排半角“(　　)”（第 13、14、16 题括号内空白、第 15 题括号内印 D），
%      本稿按全稿惯例排 \mbox{（\quad）}；第 15 题括号内的答案移入 \ans{}，见上；
%   ⑩ 选项行排照原卷：第 13、15 题两行两条，第 14、16 题一行四条；
%   ⑪ 每道解答题原卷未留作答空白，本稿按全稿惯例加 \ansspace（末题用 \ansspace*）；
%      第 20、21 题题干长，各用一处 \ansneed 让题干与作答区不拆页（\ansneed 只作用于试卷版）。
%
% 原卷笔误/存疑（均照录未改，正文对应处另有 % 注）：
%   ① 第 3 题解析末句印作“则 $\bar A=(-\infty,0)\cup[2,3)$.”：按题干
%      $U=\{x\mid 1/x\leqslant 1/2\}$（即 $x<0$ 或 $x\geqslant2$）与 $A=[3,+\infty)$，
%      $\bar A$ 应为 $(-\infty,0)\cup\{2\}$；且与本解析首句自求之 $U$ 相矛盾（$3\notin U$），
%      疑为笔误；照录未改。
%   ② 第 12 题题干 $\alpha*\beta$ 定义式括号不配对：第四组“$+\cdots+x_{2020}+y_{2020}
%      -x_{2020}y_{2020}$”缺左括号、末尾连排两个右括号“))”；照录未改。
%   ③ 第 12 题解析“所以…中有 2020,个 1”一行，“$\cdots$”后脱逗号（作
%      “$x_2,\cdots x_{2020}$”，第 12 行“而…”同式样），且“2020”后衍一半角逗号
%      （作“2020,个 1”）；照录未改。
%   ④ 第 16 题解析乘积链印作“$a_1a_2<a_1a_3<a_1a_4<a_1a_5<a_2a_3<a_3a_3<a_{14}a_5$”：
%      “$a_3a_3$”同足标自乘、与单调链不合，“$a_{14}a_5$”下标 14 与 $a_1$—$a_5$ 记号
%      不合（疑为 $a_2a_4$、$a_4a_5$ 之误）；照录未改。
%   ⑤ 第 5 题题干作“$Q\subset P$”（真子集记号），解析两处均作“$Q\subseteq P$”，
%      同一题前后记号不一；照录（本稿分别排 \subset、\subseteq）。
%   ⑥ 第 17 题 (1) 解析“当 $a=1$ 时，原不等式为 $(x+1)^2<0$ 此时…”与
%      “当 $0<a<1$ 时，$-\frac1a<-1$ 此时…”两处“此时”前均脱标点；照录未改。
%   ⑦ 第 20 题 (2) 题干“若 $a=1$，$b=2k-1$，$c=k^2$，若方程 $f(x)=0$ 有两个大于 1 的根”
%      两个“若”连用，疑衍其一；照录未改。
%
% 插图：无（全卷 13 页无任何几何图形、无“如图/右图/图中”字样）。
%
% 行数核对：逐题数 \step/\stepm 条数与原卷【解析】印刷行数，18 题逐一相符
% （第 15 题无解析段，只提 \ans{}），合计 169 条 / 169 行
% （Q3—Q12：2、3、4、3、1、3、5、3、13、15；Q13—Q16：2、3、0、25；
% Q17—Q21：19、16、9、20、23）。
%
\documentclass[a4paper,11pt]{article}
\usepackage{common}

\begin{document}

\examtitlest[3]{2026--2027 年华师大一附中高一上 10 月月考}{集合与常用逻辑用语、不等式}

\bigsec{一、填空题}

\begin{enumerate}
\setcounter{enumi}{2}

\item 设全集 $U=\left\{x\mid\dfrac1x\leqslant\dfrac12\right\}$，$A=[3,+\infty)$，则 $\bar A=$\blankline[4.4em].

\ifanswer
\solbegin
\step{$\dfrac1x\leqslant\dfrac12$，解得 $x<0$ 或 $x\geqslant2$，故 $U=(-\infty,0)\cup[2,+\infty)$，}
% 原卷末句作 $\bar A=(-\infty,0)\cup[2,3)$，按题干 $U$ 应为 $(-\infty,0)\cup\{2\}$，与本解析首句矛盾，疑笔误；照录（见文件头“原卷笔误/存疑”①）
\step{$A=[3,+\infty)$，则 $\bar A=(-\infty,0)\cup[2,3)$.}
\solend
\fi

\item 已知集合 $M=\{y\mid y=x^2-2x-1\}$，$N=\{y\mid y=x^2-4x+1\}$，则 $M\cup N=$\blankline[4.4em].

\ifanswer
\solbegin
\step{$M=\{y\mid y=x^2-2x-1\}=\{y\mid y=(x-1)^2-2\}=\{y\mid y\geqslant-2\}$，}
\step{$N=\{y\mid y=x^2-4x+1\}=\{y\mid y=(x-2)^2-3\}=\{y\mid y\geqslant-3\}$，}
\step{则 $M\cup N=\{y\mid y\geqslant-3\}$.}
\solend
\fi

\item 已知 $P=\{x\mid x^2+x-6=0\}$，$Q=\{x\mid mx-1=0\}$，且 $Q\subset P$，则实数 $m$ 的值\blankline[4.4em].

\ifanswer
\solbegin
\step{$P=\{2,-3\}$，$Q=\{x\mid mx=1\}$；}
\stepm{①}{$m=0$ 时，$Q=\varnothing$，满足 $Q\subseteq P$；}
\stepm{②}{$m\neq0$ 时，要使 $Q\subseteq P$，则 $\dfrac1m=2$ 或 $-3$，所以 $m=\dfrac12$ 或 $-\dfrac13$；}
\step{所以 $m=0$ 或 $\dfrac12$ 或 $-\dfrac13$.}
\solend
\fi

\item 已知“存在正实数 $x$，使得关于 $x$ 的不等式 $ax-1>0$ 成立”为真命题，则实数 $a$ 的取值范围是\blankline[4.4em].

\ifanswer
\solbegin
\step{存在正实数 $x$，使得关于 $x$ 的不等式 $ax-1>0$ 成立为真命题，}
\step{即存在正实数 $x$，使得 $a>\dfrac1x$，因为 $\dfrac1x>0$，所以 $a>0$，}
\step{即 $a$ 的范围为 $(0,+\infty)$.}
\solend
\fi

\item 设 $\alpha:3\leqslant x\leqslant5$，$\beta:2m-1\leqslant x<m+4$，若 $\alpha$ 是 $\beta$ 的充分条件，则实数 $m$ 的取值范围是\blankline[4.4em].

\ifanswer
\solbegin
\step{若 $\alpha$ 是 $\beta$ 的充分条件，则 $\begin{cases}2m-1\leqslant3\\ m+4>5\end{cases}$，解得 $1<m\leqslant2$.}
\solend
\fi

\item 对于实数 $a$、$b$、$c$，已知关于 $x$ 的不等式 $ax^2+bx+c>0$ 的解集是 $(1,2)$，则关于 $x$ 的不等式 $cx^2-bx+a\leqslant0$ 的解集是\blankline[4.4em].

\ifanswer
\solbegin
\step{关于 $x$ 的不等式 $ax^2+bx+c>0$ 的解集是 $(1,2)$，则 $\begin{cases}a<0\\ -\dfrac ba=3\\ \dfrac ca=2\end{cases}$，则 $\begin{cases}b=-3a\\ c=2a\end{cases}$，}
\step{则关于 $x$ 的不等式 $cx^2-bx+a=2ax^2+3ax+a\leqslant0$，则 $2x^2+3x+1\geqslant0$，}
\step{则 $x\leqslant-1$ 或 $x\geqslant-\dfrac12$，则解集为 $(-\infty,-1]\cup\left[-\dfrac12,+\infty\right)$.}
\solend
\fi

\item 已知关于 $x$ 的不等式 $kx^2+kx+k+2\leqslant0$ 的解集为 $\varnothing$，则实数 $k$ 的取值范围为\blankline[4.4em].

\ifanswer
\solbegin
\step{当 $k=0$ 时，不等式化为 $2\leqslant0$，解集为 $\varnothing$，符合题意；}
\step{当 $k\neq0$ 时，若关于 $x$ 的不等式 $kx^2+kx+k+2\leqslant0$ 的解集为 $\varnothing$，}
\step{则 $\begin{cases}k>0\\ \Delta=k^2-4k(k+2)<0\end{cases}$，化简得 $\begin{cases}k>0\\ -3k^2-8k<0\end{cases}$，即 $\begin{cases}k>0\\ k(3k+8)>0\end{cases}$，}
\step{解得 $k>0$.}
\step{综上所述，$k\geqslant0$，即实数 $k$ 的取值范围为 $[0,+\infty)$.}
\solend
\fi

\item 对班级若干名学生调查对 $A$、$B$ 两事件的态度，有如下结果：赞成 $A$ 的人数是全体的五分之三，其余的不赞成；赞成 $B$ 的比赞成 $A$ 多 3 人，其余的不赞成. 已知对 $A$、$B$ 都赞成的学生有 18 人，对 $A$、$B$ 都不赞成的学生数比对 $A$、$B$ 都赞成的学生数的三分之一多 1 人，则班级总人数是\blankline[4.4em]人.

\ifanswer
\solbegin
\step{设班级总人数是 $x$ 人，则赞成 $A$ 的人数是 $\dfrac35x$，赞成 $B$ 的人数是 $\dfrac35x+3$，}
\step{$A$、$B$ 都不赞成的人数为 $18\times\dfrac13+1=7$ 人，}
\step{由容斥原理得 $x=\left(\dfrac35x+\dfrac35x+3-18\right)+7$，解得 $x=40$，即班级总人数是 40 人.}
\solend
\fi

\item 设全集 $U=\{1,2,3,4,5,6,7,8,9,10,11,12\}$，若集合 $A$ 满足：① $A\subseteq U$；②若 $x\in A$，必有 $2x\notin A$；③若 $x\in\bar A$，必有 $2x\notin\bar A$，则这样的集合 $A$ 有\blankline[4.4em]个.

\ifanswer
\solbegin
\step{设集合 $P_n=\{1,2,\cdots,n\}$,$n\in\N^*$，}
\step{任取 $x\in P_n$，将 $x$ 除以 2，若商仍为偶数，再除以 2，$\cdots$，}
\step{经过 $k$ 次以后，商必为奇数，此时记商为 $m$，}
\step{于是 $x=m\cdot2^k$，其中 $m$ 为奇数，$k\in\N^*$.}
\step{由条件得，若 $m\in A$，则 $x\in A$ 等价于 $k$ 为偶数；}
\step{若 $m\notin A$，则 $x\in A$ 等价于 $k$ 为奇数；}
\step{于是 $x$ 是否属于 $A$ 由 $m$ 是否属于 $A$ 确定；}
\step{设 $Q_n$ 是 $P_n$ 中所有奇数的集合，}
\step{设 $f(n)$ 为集合 $P_n$ 中满足题干条件的集合 $A$ 的个数，}
\step{因此 $f(n)$ 等于 $Q_n$ 的子集个数.}
\step{当 $n$ 为偶数（或奇数）时，$P_n$ 中奇数的个数是 $\dfrac n2$（或 $\dfrac{n+1}2$），}
\step{所以 $f(n)=\begin{cases}2^{\frac n2},n\text{为偶数}\\[2pt] 2^{\frac{n+1}2},n\text{为奇数}\end{cases}$，所以 $P_{12}=\{1,2,3,4,5,6,7,8,9,10,11,12\}$，}
\step{即同时满足三个条件的集合 $A$ 的个数为 $2^{\frac{12}2}=2^6=64$.}
\solend
\fi

\item 设集合 $\Omega=\{(x_1,x_2,\cdots,x_{2020})\mid x_i\in\{0,1\},i=1,2,\cdots,2020\}$，对于 $\Omega$ 中的任意两个元素 $\alpha=(x_1,x_2,\cdots,x_{2020})$，$\beta=(y_1,y_2,\cdots,y_{2020})$，记
% 原卷此式第四组缺左括号、末尾连排两个右括号“))”，括号不配对；照录（见文件头“原卷笔误/存疑”②）
$\alpha*\beta=(x_1+y_1-x_1y_1)+(x_2+y_2-x_2y_2)+\cdots+x_{2020}+y_{2020}-x_{2020}y_{2020}))$，设 $u$、$v\in\Omega$，若 $u*u+v*v=2020$，则 $u*v$ 的最大值和最小值之和为\blankline[4.4em].

\ifanswer
\solbegin
\step{设 $u=(x_1,x_2,\cdots,x_{2020})$,$v=(y_1,y_2,\cdots,y_{2020})$，}
\step{因为 $x_i\in\{0,1\}$，所以 $x_i(x_i-1)=0$，}
\step{所以 $u*u=(x_1+x_1-x_1^2)+(x_2+x_2-x_2^2)+\cdots+(x_{2020}+x_{2020}-x_{2020}^2)$}
\step{$=x_1+x_2+\cdots+x_{2020}$，}
\step{$v*v=(y_1+y_1-y_1^2)+(y_2+y_2-y_2^2)+\cdots+(y_{2020}+y_{2020}-y_{2020}^2)$}
\step{$=y_1+y_2+\cdots+y_{2020}$，}
\step{由 $u*u+v*v=2020$，得 $x_1+x_2+\cdots+x_{2020}+y_1+y_2+\cdots+y_{2020}=2020$，}
% 原卷此行“$\cdots$”后脱逗号、 “2020”后衍一半角逗号（“2020,个 1”）；照录（见文件头“原卷笔误/存疑”③）
\step{所以 $x_1,x_2,\cdots x_{2020},y_1,y_2,\cdots y_{2020}$ 中有 2020,个 1，2020 个 0，}
\step{$u*v=(x_1+y_1-x_1y_1)+(x_2+y_2-x_2y_2)+\cdots+(x_{2020}+y_{2020}-x_{2020}y_{2020})$}
\step{$=2020-(x_1y_1+x_2y_2+\cdots+x_{2020}y_{2020})$，}
\step{注意到只有 $x_i=y_i=1$ 时，$x_iy_i=1$，否则 $x_iy_i=0$，}
% 此行“$\cdots$”后同样脱逗号，与上一处同式样；照录（见文件头“原卷笔误/存疑”③）
\step{而 $x_1,x_2,\cdots x_{2020},y_1,y_2,\cdots y_{2020}$ 中有 2020 个 1，2020 个 0，}
\step{所以满足 $x_iy_i=1$ 的最多有 1010 个，所以 $u*v\geqslant2020-1010=1010$，}
\step{满足 $x_iy_i=0$ 的最多有 2020 个，所以 $u*v\leqslant2020-0=2020$，}
\step{则 $u*v$ 的最大值和最小值之和为 $2020+1010=3030$.}
\solend
\fi

\end{enumerate}

\bigsec{二、选择题}

\begin{enumerate}
\setcounter{enumi}{12}

\item 实数 $x$ 与 $y$，“$x+y>2$”的\mbox{（\quad）}条件是“$x>1$ 或 $y>1$”.

\keepwithprev
A. 充分非必要\qquad B. 必要非充分

C. 充要\qquad D. 既非充分又非必要

\ifanswer
\solbegin
\step{当 $x>1$ 或 $y>1$ 时，$x+y>2$ 不一定成立，充分性不成立；}
\step{当 $x+y>2$ 时，$x>1$ 或 $y>1$ 一定成立，必要性成立. 故选 $B$.}
\solend
\fi

\item 设 $a,b\in\R$，且 $a<0$，$-1<b<0$，那么下列不等式成立的是\mbox{（\quad）}

\keepwithprev
A. $a>ab>ab^2$\qquad B. $ab^2>ab>a$\qquad C. $ab>a>ab^2$\qquad D. $ab>ab^2>a$

\ifanswer
\solbegin
\step{因为 $-1<b<0$，$a<0$，所以 $ab>0$，$0<b^2<1$.}
\step{所以 $ab-ab^2=ab(1-b)>0$，$ab^2-a=a(b^2-1)>0$，}
\step{所以 $ab>ab^2>a$. 故选 $D$.}
\solend
\fi

% 第 15 题答案原卷印在题干括号内（“(  D  )”），且无【答案】/【解析】段；
% 按 高一37 第 13 题、高一63 第 14 题成例括号留空、答案另提 \ans{$D$}
\item 用反证法证明：“平面四边形中至少有一个内角不超过 90°”时，下列“假设”中正确的是\mbox{（\quad）}

\keepwithprev
A. 假设有两个内角超过 90°\qquad B. 假设有三个内角超过 90°

C. 假设至多有两个内角超过 90°\qquad D. 假设四个内角都超过 90°

\ifanswer
\ans{$D$}
\fi

\item 设 $A$ 为非空实数集，集合 $B=\{xy\mid x\in A,y\in A,x\neq y\}$，则对于以下两个结论说法正确的是\mbox{（\quad）}

结论①：若 $A$ 是五元集，则 $B$ 中的元素最多 10 个；

结论②：存在四元集 $A$，使得 $B=\{2,4,5,8,10,16\}$.

\keepwithprev
A. ①对②对\qquad B. ①对②错\qquad C. ①错②对\qquad D. ①错②错

\ifanswer
\solbegin
\step{对于①，想让元素个数达到最大值 10，则须避免产生符号不同的乘积，}
\step{否则正负分开计数容易导致总数受限，且容易产生相等值，}
\step{也最好避免 0 的出现，}
\step{故设 $A=\{a_1,a_2,a_3,a_4,a_5\}$，$0<a_1<a_2<a_3<a_4<a_5$.}
% 原卷此链“$a_3a_3$”同足标自乘、“$a_{14}a_5$”下标 14 均与文意不合，疑为 $a_2a_4$、$a_4a_5$ 之误；照录（见文件头“原卷笔误/存疑”④）
\step{因为 $a_1a_2<a_1a_3<a_1a_4<a_1a_5<a_2a_3<a_3a_3<a_{14}a_5$，$a_2a_3<a_2a_4<a_3a_4$，}
\step{所以 $B$ 中元素个数小于等于 10 个，}
\step{如设 $A=\{1,2,3,5,7\}$，则 $B=\{2,3,5,6,7,10,14,15,21,35\}$，}
\step{所以集合 $B$ 中元素个数的最大值为 10 个，故①正确；}
\step{对于②，由于集合 $B$ 中没有 0，所以集合 $A$ 中没有 0；}
\step{设集合 $A=\{a,b,c,d\}$，且 $a<b<c<d$，}
\step{当 $a<0<b<c<d$ 时，$a\times b<0$，$a\times c<0$，$a\times d<0$，}
\step{故集合 $B$ 至少有 3 个负数，与 $B=\{2,4,5,8,10,16\}$ 矛盾；}
\step{当 $a<b<0<c<d$ 时，$a\times b>0$，$c\times d>0$，}
\step{$a\times c<0$，$a\times d<0$，$b\times c<0$，$b\times d<0$，}
\step{故集合 $B$ 至少有 4 个负数，与 $B=\{2,4,5,8,10,16\}$ 矛盾；}
\step{当 $a<b<c<0<d$ 时，$a\times b>0$，$a\times c>0$，$b\times c>0$，}
\step{$a\times d<0$，$b\times d<0$，$c\times d<0$，}
\step{故集合 $B$ 至少有 3 个负数，与 $B=\{2,4,5,8,10,16\}$ 矛盾；}
\step{当 $0<a<b<c<d$ 时，则集合 $B=\{ab,ac,ad,bc,bd,cd\}$，}
\step{则必有 $ab=2$，$cd=16$，其 4 个正实数的乘积 $abcd=32$，}
\step{又 $ad=5$，$bc=8$ 或 $ad=8$，$bc=5$，其 4 个正实数的乘积 $abcd=40$，矛盾；}
\step{当 $a<b<c<d<0$ 时，与 $0<a<b<c<d$ 时的证明方法相同，}
\step{得集合 $A$ 也不存在；}
\step{综上不存在四元集 $A$，使其生成集 $B=\{2,4,5,8,10,16\}$，故②错误.}
\step{故选 $B$.}
\solend
\fi

\end{enumerate}

\bigsec{三、解答题}

\begin{enumerate}
\setcounter{enumi}{16}

\item 解下列不等式（组）

\begin{enumerate}[label=(\arabic*)]
\item 解不等式：$(ax+1)(x+1)<0(a\in\R)$.
\item 解不等式组：$\begin{cases}\dfrac{2x+4}{x+4}\geqslant1\\[4pt] (x+1)(2-|x|)\geqslant0\end{cases}$.
\end{enumerate}

\ifanswer
\solbegin
\stepm{(1)}{当 $a=0$ 时，原不等式为 $x+1<0$，即 $x<-1$，}
\step{所以此时不等式的解集为 $\{x\mid x<-1\}$；}
\step{当 $a\neq0$ 时，令 $(ax+1)(x+1)=0$，解得 $x_1=-\dfrac1a$,$x_2=-1$，}
\step{当 $a>1$ 时，$-\dfrac1a>-1$，此时不等式的解集为 $\left\{x\mid-1<x<-\dfrac1a\right\}$；}
% 以下两行“此时”前原卷均无标点；照录（见文件头“原卷笔误/存疑”⑥）
\step{当 $a=1$ 时，原不等式为 $(x+1)^2<0$ 此时不等式的解集为 $\varnothing$；}
\step{当 $0<a<1$ 时，$-\dfrac1a<-1$ 此时不等式的解集为 $\left\{x\mid-\dfrac1a<x<-1\right\}$；}
\step{当 $a<0$ 时，$-\dfrac1a>-1$，此时不等式的解集为 $\left\{x\mid x<-1\ \text{或}\ x>-\dfrac1a\right\}$.}
\step{综上所述，当 $a=0$ 时，$\{x\mid x<-1\}$；}
\step{当 $a>1$ 时，$\left\{x\mid-1<x<-\dfrac1a\right\}$；}
\step{当 $a=1$ 时，$\varnothing$；}
\step{当 $0<a<1$ 时，$\left\{x\mid-\dfrac1a<x<-1\right\}$；}
\step{当 $a<0$ 时，$\left\{x\mid x<-1\ \text{或}\ x>-\dfrac1a\right\}$.}
\stepm{(2)}{不等式 $\dfrac{2x+4}{x+4}\geqslant1$ 可化为 $\dfrac{2x+4}{x+4}-1\geqslant0$，即 $\dfrac x{x+4}\geqslant0$，}
\step{等价于 $\begin{cases}x(x+4)\geqslant0\\ x+4\neq0\end{cases}$，解得 $x<-4$ 或 $x\geqslant0$，}
\step{不等式 $(x+1)(2-|x|)\geqslant0$ 等价于 $\begin{cases}(x+1)(2-x)\geqslant0,x\geqslant0\\ (x+1)(2+x)\geqslant0,x<0\end{cases}$，}
\step{解不等式 $(x+1)(2-x)\geqslant0$，$x\geqslant0$，得 $0\leqslant x\leqslant2$，}
\step{解不等式 $(x+1)(2+x)\geqslant0$，$x<0$，得 $x\leqslant-2$ 或 $-1\leqslant x<0$，}
\step{所以 $(x+1)(2-|x|)\geqslant0$ 的解集为 $\{x\mid x\leqslant-2\ \text{或}\ -1\leqslant x\leqslant2\}$，}
\step{综上所述，$\begin{cases}\dfrac{2x+4}{x+4}\geqslant1\\[4pt] (x+1)(2-|x|)\geqslant0\end{cases}$ 的解集为 $\{x\mid x<-4\ \text{或}\ 0\leqslant x\leqslant2\}$.}
\solend
\fi

\ansspace[8]

% 第 18 题题干＋两小问与作答区 6 行须同页
\ansneed{20cm}
\item 已知 $a$ 为实数，设 $f(x)=x^2+2(a+1)x+1$.

\begin{enumerate}[label=(\arabic*)]
\item 已知对于任意实数 $x$，$\sqrt{f(x)}$ 均有意义，求实数 $a$ 的取值范围；
\item 设 $A=\{x\mid f(x)+a^2-2=0\}$，$B=\{-4,0\}$，$A\cup B=B$，求实数 $a$ 的取值范围.
\end{enumerate}

\ifanswer
\solbegin
\stepm{(1)}{$\sqrt{f(x)}$ 对任意实数 $x$ 都有意义，}
\step{即 $f(x)=x^2+2(a+1)x+1\geqslant0$ 对任意 $x\in\R$ 恒成立，}
\step{所以 $\Delta=4(a+1)^2-4\times1\times1\leqslant0$，即 $(a+1)^2\leqslant1$，解得 $-2\leqslant a\leqslant0$，}
\step{故 $a$ 的取值范围是 $[-2,0]$.}
\stepm{(2)}{由 $f(x)+a^2-2=0$ 得 $x^2+2(a+1)x+a^2-1=0$.}
\step{由 $A\cup B=B$ 得 $A\subseteq B=\{-4,0\}$，即该方程无实根，}
\step{或它的实根都在 $\{-4,0\}$ 中.}
\step{当方程无实根时，$\Delta=4(a+1)^2-4(a^2-1)=8a+8<0$，解得 $a<-1$；}
\step{当方程有实根时，$a\geqslant-1$，设两根为 $x_1$、$x_2$，}
\step{由根与系数的关系得 $x_1+x_2=-2(a+1)$、$x_1x_2=a^2-1$，}
\step{此时两根只能在 $\{-4,0\}$ 中取值；}
\step{若两根为 0 与 0，由 $x_1x_2=0$ 得 $a=\pm1$，}
\step{再由 $x_1+x_2=0$ 得 $a=-1$；}
\step{若两根为 0 与 $-4$，由 $x_1x_2=0$ 得 $a=\pm1$，再由 $x_1+x_2=-4$ 得 $a=1$.}
\step{经检验，$a=-1$ 时 $A=\{0\}$，$a=1$ 时 $A=\{-4,0\}$，都满足 $A\subseteq B$.}
\step{综上，实数 $a$ 的取值范围为 $(-\infty,-1]\cup\{1\}$.}
\solend
\fi

\ansspace[8]

\ansneed{18cm}
\item 共享单车给市民出行带来了诸多便利，某公司购买了一批单车投放到某地给居民使用，据市场分析，这批单车的运营累计利润 $y$（单位：元）与运营天数 $x$（单位：天，$x$ 为正整数）满足关系：$y=-\dfrac12x^2+60x-800$.

\begin{enumerate}[label=(\arabic*)]
\item 要使这批单车的运营累计利润高于 800 元，求运营天数的取值范围；
\item 判断这批单车的日平均运营利润能否超过 20 元，并说明理由.
\end{enumerate}

\ifanswer
\solbegin
\stepm{(1)}{由题意得 $y=-\dfrac12x^2+60x-800>800$，}
\step{即 $x^2-120x+3200<0$，解得 $40<x<80$，}
\step{即运营天数的取值范围为 $\{x\mid40<x<80,x\in\N^*\}$.}
\stepm{(2)}{不能，理由如下：}
\step{这批单车的日平均运营利润为 $\dfrac yx=-\dfrac12x-\dfrac{800}x+60=-\dfrac12\left(x+\dfrac{1600}x\right)+60$，}
\step{因 $x>0$，则 $x+\dfrac{1600}x\geqslant2\sqrt{x\cdot\dfrac{1600}x}=80$，}
\step{当且仅当 $x=\dfrac{1600}x$，即 $x=40$ 时取等号，}
\step{此时 $\dfrac yx=-\dfrac12\left(x+\dfrac{1600}x\right)+60\leqslant-\dfrac12\times80+60=20$，}
\step{故这批单车的日平均运营利润不能超过 20 元.}
\solend
\fi

\ansspace[8]

% 第 20 题题干＋三小问与作答区 8 行须同页
\ansneed{20cm}
\item 设 $f(x)=ax^2+bx+c$($a\neq0$，$b$、$c\in\R$).

\begin{enumerate}[label=(\arabic*)]
\item 若 $a=b=c=1$，设 $A=\{(x,y)\mid y=f(x)\}$，$B=\{(x,y)\mid y=1\}$，用列举法表示 $A\cap B$；
% 原卷两个“若”连用（“若 a=1，…，若方程 f(x)=0 有两个大于 1 的根”），疑衍其一；照录（见文件头“原卷笔误/存疑”⑦）
\item 若 $a=1$，$b=2k-1$，$c=k^2$，若方程 $f(x)=0$ 有两个大于 1 的根，求实数 $k$ 的取值范围；
\item 若 $a$、$b$、$c\in\Z$，且 $f(0)$、$f(1)$ 均为奇数，求证：方程 $f(x)=0$ 无整数解.
\end{enumerate}

\ifanswer
\solbegin
\stepm{(1)}{当 $a=b=c=1$ 时，$f(x)=x^2+x+1$，}
\step{所以 $A=\{(x,y)\mid y=x^2+x+1\}$，又 $y=1$，}
\step{由 $x^2+x+1=1$ 得 $x^2+x=0$，解得 $x=0$ 或 $x=-1$，}
\step{所以 $A\cap B=\{(0,1),(-1,1)\}$.}
\stepm{(2)}{当 $a=1$、$b=2k-1$、$c=k^2$ 时，$f(x)=x^2+(2k-1)x+k^2$.}
\step{设两根为 $x_1$、$x_2$，}
\step{则“两根都大于 1”等价于 $\Delta>0$、$f(1)>0$、$-\dfrac b{2a}>1$.}
\step{由 $\Delta=(2k-1)^2-4k^2=1-4k>0$ 得 $k<\dfrac14$；}
\step{由 $f(1)=1+(2k-1)+k^2=k^2+2k=k(k+2)>0$ 得 $k<-2$ 或 $k>0$；}
\step{由 $-\dfrac{2k-1}2=\dfrac{1-2k}2>1$ 得 $k<-\dfrac12$.}
\step{三者取交集得 $k<-2$，即 $k$ 的取值范围是 $(-\infty,-2)$.}
\stepm{(3)}{设 $x_0\in\Z$ 是方程 $f(x)=0$ 的整数解.}
\step{因为 $x_0(x_0-1)$ 是两个相邻整数之积，恒为偶数，所以 $ax_0(x_0-1)$ 为偶数.}
\step{而 $f(x_0)=ax_0^2+bx_0+c=ax_0(x_0-1)+(a+b)x_0+c$，}
\step{所以 $f(x_0)$ 与 $(a+b)x_0+c$ 同奇偶.}
\step{若 $x_0$ 为偶数，则 $(a+b)x_0$ 为偶数，$f(x_0)$ 与 $c=f(0)$ 同奇偶，为奇数；}
\step{若 $x_0$ 为奇数，则 $(a+b)x_0$ 与 $a+b$ 同奇偶，}
\step{$f(x_0)$ 与 $(a+b)+c=f(1)$ 同奇偶，也为奇数.}
\step{综上，$f(x_0)$ 恒为奇数，不可能等于 0，与 $f(x_0)=0$ 矛盾，}
\step{所以方程 $f(x)=0$ 无整数解.}
\solend
\fi

\ansspace[8]

% 压轴题（题干＋三小问）与作答区 10 行须同页
\ansneed{20cm}
\item 已知 $X$ 为有限集，且任取 $x\in X$，均有 $x\in\Z$，称 $X$ 为非空有限整数集. 定义：$X_k=\{x+k\mid x\in X\}$，$X^2=\{y\mid y=x^2,x\in X\}$.

\begin{enumerate}[label=(\arabic*)]
\item 若 $X=\{-2,2\}$，分别写出 $X_2$ 和 $X^2$（用列举法表示）；
\item 若 $X=\{-1,0,1\}$，$X_1\subseteq Y$，$(Y_{-1})^2\subseteq X$，求集合 $Y$ 中所有的元素；
\item 对于非空有限整数集 $A$、$B$，若 $A_k\subseteq B$，$(B^2)_{-2}\subseteq A$，$B$ 中的元素均为非负整数，$k$ 为整数，求证：

① $A$、$B$ 中元素的个数相同；

② 若集合 $A$ 中的元素个数为奇数，则 $k$ 为偶数.
\end{enumerate}

\ifanswer
\solbegin
\stepm{(1)}{由题意得 $X_2=\{0,4\}$，$X^2=\{4\}$；}
\stepm{(2)}{由题意得 $X_1=\{0,1,2\}$，即 $\{0,1,2\}\subseteq Y$，$Y_{-1}=\{y-1\mid y\in Y\}$，}
\step{$(Y_{-1})^2=\{(y-1)^2\mid y\in Y\}\subseteq X=\{-1,0,1\}$，}
\step{所以 $(y-1)^2=0$ 或 $(y-1)^2=1$，解得 $y=1$ 或 $y=2$ 或 $y=0$，}
\step{即 $y\in\{0,1,2\}$，即 $Y\subseteq\{0,1,2\}$，所以 $Y=\{0,1,2\}$；}
% 原卷 (3) 解析首行“(3)”与“①”连排，本稿合并记作 \stepm{(3)①}
\stepm{(3)①}{因为 $A_k=\{x+k\mid x\in A\}$，函数 $g(x)=x+k$ 在全体整数集上严格增，}
\step{所以 $\forall x_1$、$x_2\in A$，$x_1\neq x_2$，$g(x_1)\neq g(x_2)$，所以 $|A_k|=|A|$，}
\step{因为 $A_k\subseteq B$，所以 $|A|\leqslant|B|$，}
\step{设 $b\in B$，$B$ 中元素均为非负整数，$\forall b_1\neq b_2\in B$，$b_1^2\neq b_2^2$，}
\step{故函数 $f(b)=b^2-2$ 在 $B$ 上单射，所以 $|(B^2)_{-2}|=|B|$，}
\step{又因为 $(B^2)_{-2}\subseteq A$，所以 $|B|\leqslant|A|$，所以 $|A|=|B|$.}
\stepm{②}{由①得 $A_k\subseteq B$ 且 $|A_k|=|A|=|B|$，$(B^2)_{-2}\subseteq A$ 且 $|(B^2)_{-2}|=|B|=|A|$，}
\step{所以 $A_k=B$，$(B^2)_{-2}=A$，}
\step{所以 $B=A_k=\{a+k\mid a\in A\}$，$A=(B^2)_{-2}=\{b^2-2\mid b\in B\}$，}
\step{所以 $A=\{(a+k)^2-2\mid a\in A\}$，}
\step{所以映射 $h(a)=(a+k)^2-2$ 是集合 $A$ 到自身的一个双射.}
\step{假设 $k$ 是奇数，当 $a$ 为偶数时，$a+k$ 为奇数，$(a+k)^2-2$ 为奇数，}
\step{当 $a$ 为奇数时，$a+k$ 为偶数，$(a+k)^2-2$ 为偶数，}
\step{所以此时映射 $h(a)=(a+k)^2-2$，将集合 $A$ 中的奇数映射为偶数，}
\step{将集合 $A$ 中的偶数映射为奇数，即集合 $A$ 中的奇数元素和偶数元素必须配对，}
\step{此时集合 $A$ 的元素个数则为偶数，}
\step{与题设中“集合 $A$ 中的元素个数为奇数”矛盾.}
\step{故假设不成立，所以 $k$ 是偶数.}
\solend
\fi

\ansspace*[10]

\end{enumerate}

\end{document}
"
% 紧凑版：xelatex "\def\iscompact{1}% 高一67_华师大一附中_10月月考.tex
%   —— 高一上数学“2026-2027 年华师大一附中高一上 10 月月考”（讲评版）
% 试卷版：xelatex 高一67_华师大一附中_10月月考.tex
% 答案版：xelatex "\def\isanswer{1}% 高一67_华师大一附中_10月月考.tex
%   —— 高一上数学“2026-2027 年华师大一附中高一上 10 月月考”（讲评版）
% 试卷版：xelatex 高一67_华师大一附中_10月月考.tex
% 答案版：xelatex "\def\isanswer{1}% 高一67_华师大一附中_10月月考.tex
%   —— 高一上数学“2026-2027 年华师大一附中高一上 10 月月考”（讲评版）
% 试卷版：xelatex 高一67_华师大一附中_10月月考.tex
% 答案版：xelatex "\def\isanswer{1}\input{高一67_华师大一附中_10月月考.tex}"
% 紧凑版：xelatex "\def\iscompact{1}\input{高一67_华师大一附中_10月月考.tex}"
%
% 原始材料：微信公众号“魔都数学加油站”文章，作者破晓，连载“27学年高一67”；
% 原文标题“27学年高一67——2026-2027年上海市华师大一附中高一上10月月考”；
% 原文链接：https://mp.weixin.qq.com/s/Kdhx07TVjLoGTcIlnRzgcQ。
% 页面标注发布时间 2026-10-18 20:33（上海）。
% 本次转录只据 `原图/高一67_华师大一附中_10月月考/` 的 13 张 PNG 逐页读图。
% 正文为 13 张 PNG 页（`原图/高一67_华师大一附中_10月月考/01..13.png`）：
% 第 01、04、13 页 2560×3620，第 02、03、05—12 页 4762×6735
% （**同一篇各页分辨率不同**），均按页序读取。13 页全为试卷正文页（卷首标题在第 1 页页顶，
% 卷末解答题第 21 题解析止于第 13 页），无封面页、目录页、二维码或宣传页。
% 粉色斜排“魔都数学加油站”为卷面水印，与公众号页脚一并未录入。
%
% 卷首标题：照卷面印作“2026--2027 年华师大一附中高一上 10 月月考”（卷面为黑体加粗一行），
% 年份区间按全稿惯例排 en dash。原卷只有一行标题、无副标题。
%
% 篇幅与题号：原卷自第 3 题起编号（第 1、2 题不在本卷内）。填空题 3—12（10 题）、
% 选择题 13—16（4 题）、解答题 17—21（5 题），共 19 题。本稿沿用原节与原题号，
% 只改排为全稿统一的大节样式（原卷小节标题为“一、填空题”“二、选择题”“三、解答题”）。
%
% 讲评版：原卷每题下直接印【解析】。**全卷 19 题中 18 题都印【解析】（无一处只印
% 【答案】）**，按全稿口径这类题不另提 \ans{}；各题解析末句的结论（第 13、14、16 题
% 的“故选 B／D／B”、各填空与解答题的“故/综上…”）照录在解析内。
% **另有 1 题答案印在题干括号内**——第 15 题（“(  D  )”，且该题无【答案】/【解析】段），
% 本稿按 高一37 第 13 题、高一63 第 14 题成例把括号留空、答案另提为 \ans{$D$}。
%
% 副标题：原卷未印副标题。本稿按“知识模块”口径标作“集合与常用逻辑用语、不等式”：
% 19 题中集合与常用逻辑用语 12 题（第 3、4、5、7、10、11、12、13、15、16、18、21 题：
% 补集、并集、子集、充分条件、容斥计数、子集计数、新定义二元运算、条件判断、
% 反证法假设、新定义生成集、集合关系与二次方程、新定义 X_k 与 X^2），
% 不等式 7 题（第 6、8、9、14、17、19、20 题：存在量词不等式、解集反求系数、
% 解集为空、不等式性质、解不等式（组）、二次不等式与基本不等式、二次方程根的分布），
% 两类均超过三分之一，故并标。
%
% 与原卷的排版差异：
%   ① 原卷每题下直接印【解析】，本稿按全稿惯例将解析置于答案版，试卷版保持干净；
%   ② 原卷未印副标题，本稿按“知识模块”口径补出；
%   ③ 原卷填空横线约四个字宽，本稿按全稿惯例统一排 \blankline[4.4em]；
%   ④ 原卷实数集排斜体 $R$（第 14 题题干、第 17 题 (1) 题干、第 18 题解析、第 20 题题干），
%      整数集排斜体 $Z$（第 20 题 (3) 题干与解析、第 21 题题干），正整数集排粗体
%      $N^*$（第 11 题解析、第 19 题解析），本稿按全稿惯例统一改用 $\R$、$\Z$、$\N^*$
%      （记号归一，粗体／斜体不作区分）；空集原卷作 $\varnothing$，照录；
%   ⑤ 原卷在数学式内多处用半角逗号（如第 8 题 $(1,2)$ 与各 cases、第 11 题解析
%      $P_n=\{1,2,\cdots,n\}$,$n\in N^*$、第 12 题解析 $u=(\cdots)$,$v=(\cdots)$ 与
%      “2020,个 1”、第 19 题解析 $\{x|40<x<80,x\in N^*\}$），均照录未强改；
%      描述法集合的竖线照原卷排 \mid；句末句点按全稿惯例排西文句点；
%   ⑥ 原卷补集记号为上划线 $\bar A$（第 3 题题干与解析、第 11 题题干），照录
%      （同 高一64/65 的 $\bar B$ 成例）；
%   ⑦ 第 12 题新定义的运算记号 $*$ 照录（$\alpha*\beta$、$u*u$、$u*v$）；
%   ⑧ 第 17—21 题的子问原卷各自另起一行，本稿按全稿惯例用 enumerate 的 (1)(2)(3)；
%      解析中的子问标号一律用 \stepm{(1)}{…}；第 5 题解析的①②、第 21 题 (3) 解析的①②
%      按全稿惯例作 \stepm{①}{…}、\stepm{②}{…}（第 21 题 (3) 首行原卷“(3)”与“①”连排，
%      本稿合并记作 \stepm{(3)①}{…}）；
%   ⑨ 选择题作答括号原卷排半角“(　　)”（第 13、14、16 题括号内空白、第 15 题括号内印 D），
%      本稿按全稿惯例排 \mbox{（\quad）}；第 15 题括号内的答案移入 \ans{}，见上；
%   ⑩ 选项行排照原卷：第 13、15 题两行两条，第 14、16 题一行四条；
%   ⑪ 每道解答题原卷未留作答空白，本稿按全稿惯例加 \ansspace（末题用 \ansspace*）；
%      第 20、21 题题干长，各用一处 \ansneed 让题干与作答区不拆页（\ansneed 只作用于试卷版）。
%
% 原卷笔误/存疑（均照录未改，正文对应处另有 % 注）：
%   ① 第 3 题解析末句印作“则 $\bar A=(-\infty,0)\cup[2,3)$.”：按题干
%      $U=\{x\mid 1/x\leqslant 1/2\}$（即 $x<0$ 或 $x\geqslant2$）与 $A=[3,+\infty)$，
%      $\bar A$ 应为 $(-\infty,0)\cup\{2\}$；且与本解析首句自求之 $U$ 相矛盾（$3\notin U$），
%      疑为笔误；照录未改。
%   ② 第 12 题题干 $\alpha*\beta$ 定义式括号不配对：第四组“$+\cdots+x_{2020}+y_{2020}
%      -x_{2020}y_{2020}$”缺左括号、末尾连排两个右括号“))”；照录未改。
%   ③ 第 12 题解析“所以…中有 2020,个 1”一行，“$\cdots$”后脱逗号（作
%      “$x_2,\cdots x_{2020}$”，第 12 行“而…”同式样），且“2020”后衍一半角逗号
%      （作“2020,个 1”）；照录未改。
%   ④ 第 16 题解析乘积链印作“$a_1a_2<a_1a_3<a_1a_4<a_1a_5<a_2a_3<a_3a_3<a_{14}a_5$”：
%      “$a_3a_3$”同足标自乘、与单调链不合，“$a_{14}a_5$”下标 14 与 $a_1$—$a_5$ 记号
%      不合（疑为 $a_2a_4$、$a_4a_5$ 之误）；照录未改。
%   ⑤ 第 5 题题干作“$Q\subset P$”（真子集记号），解析两处均作“$Q\subseteq P$”，
%      同一题前后记号不一；照录（本稿分别排 \subset、\subseteq）。
%   ⑥ 第 17 题 (1) 解析“当 $a=1$ 时，原不等式为 $(x+1)^2<0$ 此时…”与
%      “当 $0<a<1$ 时，$-\frac1a<-1$ 此时…”两处“此时”前均脱标点；照录未改。
%   ⑦ 第 20 题 (2) 题干“若 $a=1$，$b=2k-1$，$c=k^2$，若方程 $f(x)=0$ 有两个大于 1 的根”
%      两个“若”连用，疑衍其一；照录未改。
%
% 插图：无（全卷 13 页无任何几何图形、无“如图/右图/图中”字样）。
%
% 行数核对：逐题数 \step/\stepm 条数与原卷【解析】印刷行数，18 题逐一相符
% （第 15 题无解析段，只提 \ans{}），合计 169 条 / 169 行
% （Q3—Q12：2、3、4、3、1、3、5、3、13、15；Q13—Q16：2、3、0、25；
% Q17—Q21：19、16、9、20、23）。
%
\documentclass[a4paper,11pt]{article}
\usepackage{common}

\begin{document}

\examtitlest[3]{2026--2027 年华师大一附中高一上 10 月月考}{集合与常用逻辑用语、不等式}

\bigsec{一、填空题}

\begin{enumerate}
\setcounter{enumi}{2}

\item 设全集 $U=\left\{x\mid\dfrac1x\leqslant\dfrac12\right\}$，$A=[3,+\infty)$，则 $\bar A=$\blankline[4.4em].

\ifanswer
\solbegin
\step{$\dfrac1x\leqslant\dfrac12$，解得 $x<0$ 或 $x\geqslant2$，故 $U=(-\infty,0)\cup[2,+\infty)$，}
% 原卷末句作 $\bar A=(-\infty,0)\cup[2,3)$，按题干 $U$ 应为 $(-\infty,0)\cup\{2\}$，与本解析首句矛盾，疑笔误；照录（见文件头“原卷笔误/存疑”①）
\step{$A=[3,+\infty)$，则 $\bar A=(-\infty,0)\cup[2,3)$.}
\solend
\fi

\item 已知集合 $M=\{y\mid y=x^2-2x-1\}$，$N=\{y\mid y=x^2-4x+1\}$，则 $M\cup N=$\blankline[4.4em].

\ifanswer
\solbegin
\step{$M=\{y\mid y=x^2-2x-1\}=\{y\mid y=(x-1)^2-2\}=\{y\mid y\geqslant-2\}$，}
\step{$N=\{y\mid y=x^2-4x+1\}=\{y\mid y=(x-2)^2-3\}=\{y\mid y\geqslant-3\}$，}
\step{则 $M\cup N=\{y\mid y\geqslant-3\}$.}
\solend
\fi

\item 已知 $P=\{x\mid x^2+x-6=0\}$，$Q=\{x\mid mx-1=0\}$，且 $Q\subset P$，则实数 $m$ 的值\blankline[4.4em].

\ifanswer
\solbegin
\step{$P=\{2,-3\}$，$Q=\{x\mid mx=1\}$；}
\stepm{①}{$m=0$ 时，$Q=\varnothing$，满足 $Q\subseteq P$；}
\stepm{②}{$m\neq0$ 时，要使 $Q\subseteq P$，则 $\dfrac1m=2$ 或 $-3$，所以 $m=\dfrac12$ 或 $-\dfrac13$；}
\step{所以 $m=0$ 或 $\dfrac12$ 或 $-\dfrac13$.}
\solend
\fi

\item 已知“存在正实数 $x$，使得关于 $x$ 的不等式 $ax-1>0$ 成立”为真命题，则实数 $a$ 的取值范围是\blankline[4.4em].

\ifanswer
\solbegin
\step{存在正实数 $x$，使得关于 $x$ 的不等式 $ax-1>0$ 成立为真命题，}
\step{即存在正实数 $x$，使得 $a>\dfrac1x$，因为 $\dfrac1x>0$，所以 $a>0$，}
\step{即 $a$ 的范围为 $(0,+\infty)$.}
\solend
\fi

\item 设 $\alpha:3\leqslant x\leqslant5$，$\beta:2m-1\leqslant x<m+4$，若 $\alpha$ 是 $\beta$ 的充分条件，则实数 $m$ 的取值范围是\blankline[4.4em].

\ifanswer
\solbegin
\step{若 $\alpha$ 是 $\beta$ 的充分条件，则 $\begin{cases}2m-1\leqslant3\\ m+4>5\end{cases}$，解得 $1<m\leqslant2$.}
\solend
\fi

\item 对于实数 $a$、$b$、$c$，已知关于 $x$ 的不等式 $ax^2+bx+c>0$ 的解集是 $(1,2)$，则关于 $x$ 的不等式 $cx^2-bx+a\leqslant0$ 的解集是\blankline[4.4em].

\ifanswer
\solbegin
\step{关于 $x$ 的不等式 $ax^2+bx+c>0$ 的解集是 $(1,2)$，则 $\begin{cases}a<0\\ -\dfrac ba=3\\ \dfrac ca=2\end{cases}$，则 $\begin{cases}b=-3a\\ c=2a\end{cases}$，}
\step{则关于 $x$ 的不等式 $cx^2-bx+a=2ax^2+3ax+a\leqslant0$，则 $2x^2+3x+1\geqslant0$，}
\step{则 $x\leqslant-1$ 或 $x\geqslant-\dfrac12$，则解集为 $(-\infty,-1]\cup\left[-\dfrac12,+\infty\right)$.}
\solend
\fi

\item 已知关于 $x$ 的不等式 $kx^2+kx+k+2\leqslant0$ 的解集为 $\varnothing$，则实数 $k$ 的取值范围为\blankline[4.4em].

\ifanswer
\solbegin
\step{当 $k=0$ 时，不等式化为 $2\leqslant0$，解集为 $\varnothing$，符合题意；}
\step{当 $k\neq0$ 时，若关于 $x$ 的不等式 $kx^2+kx+k+2\leqslant0$ 的解集为 $\varnothing$，}
\step{则 $\begin{cases}k>0\\ \Delta=k^2-4k(k+2)<0\end{cases}$，化简得 $\begin{cases}k>0\\ -3k^2-8k<0\end{cases}$，即 $\begin{cases}k>0\\ k(3k+8)>0\end{cases}$，}
\step{解得 $k>0$.}
\step{综上所述，$k\geqslant0$，即实数 $k$ 的取值范围为 $[0,+\infty)$.}
\solend
\fi

\item 对班级若干名学生调查对 $A$、$B$ 两事件的态度，有如下结果：赞成 $A$ 的人数是全体的五分之三，其余的不赞成；赞成 $B$ 的比赞成 $A$ 多 3 人，其余的不赞成. 已知对 $A$、$B$ 都赞成的学生有 18 人，对 $A$、$B$ 都不赞成的学生数比对 $A$、$B$ 都赞成的学生数的三分之一多 1 人，则班级总人数是\blankline[4.4em]人.

\ifanswer
\solbegin
\step{设班级总人数是 $x$ 人，则赞成 $A$ 的人数是 $\dfrac35x$，赞成 $B$ 的人数是 $\dfrac35x+3$，}
\step{$A$、$B$ 都不赞成的人数为 $18\times\dfrac13+1=7$ 人，}
\step{由容斥原理得 $x=\left(\dfrac35x+\dfrac35x+3-18\right)+7$，解得 $x=40$，即班级总人数是 40 人.}
\solend
\fi

\item 设全集 $U=\{1,2,3,4,5,6,7,8,9,10,11,12\}$，若集合 $A$ 满足：① $A\subseteq U$；②若 $x\in A$，必有 $2x\notin A$；③若 $x\in\bar A$，必有 $2x\notin\bar A$，则这样的集合 $A$ 有\blankline[4.4em]个.

\ifanswer
\solbegin
\step{设集合 $P_n=\{1,2,\cdots,n\}$,$n\in\N^*$，}
\step{任取 $x\in P_n$，将 $x$ 除以 2，若商仍为偶数，再除以 2，$\cdots$，}
\step{经过 $k$ 次以后，商必为奇数，此时记商为 $m$，}
\step{于是 $x=m\cdot2^k$，其中 $m$ 为奇数，$k\in\N^*$.}
\step{由条件得，若 $m\in A$，则 $x\in A$ 等价于 $k$ 为偶数；}
\step{若 $m\notin A$，则 $x\in A$ 等价于 $k$ 为奇数；}
\step{于是 $x$ 是否属于 $A$ 由 $m$ 是否属于 $A$ 确定；}
\step{设 $Q_n$ 是 $P_n$ 中所有奇数的集合，}
\step{设 $f(n)$ 为集合 $P_n$ 中满足题干条件的集合 $A$ 的个数，}
\step{因此 $f(n)$ 等于 $Q_n$ 的子集个数.}
\step{当 $n$ 为偶数（或奇数）时，$P_n$ 中奇数的个数是 $\dfrac n2$（或 $\dfrac{n+1}2$），}
\step{所以 $f(n)=\begin{cases}2^{\frac n2},n\text{为偶数}\\[2pt] 2^{\frac{n+1}2},n\text{为奇数}\end{cases}$，所以 $P_{12}=\{1,2,3,4,5,6,7,8,9,10,11,12\}$，}
\step{即同时满足三个条件的集合 $A$ 的个数为 $2^{\frac{12}2}=2^6=64$.}
\solend
\fi

\item 设集合 $\Omega=\{(x_1,x_2,\cdots,x_{2020})\mid x_i\in\{0,1\},i=1,2,\cdots,2020\}$，对于 $\Omega$ 中的任意两个元素 $\alpha=(x_1,x_2,\cdots,x_{2020})$，$\beta=(y_1,y_2,\cdots,y_{2020})$，记
% 原卷此式第四组缺左括号、末尾连排两个右括号“))”，括号不配对；照录（见文件头“原卷笔误/存疑”②）
$\alpha*\beta=(x_1+y_1-x_1y_1)+(x_2+y_2-x_2y_2)+\cdots+x_{2020}+y_{2020}-x_{2020}y_{2020}))$，设 $u$、$v\in\Omega$，若 $u*u+v*v=2020$，则 $u*v$ 的最大值和最小值之和为\blankline[4.4em].

\ifanswer
\solbegin
\step{设 $u=(x_1,x_2,\cdots,x_{2020})$,$v=(y_1,y_2,\cdots,y_{2020})$，}
\step{因为 $x_i\in\{0,1\}$，所以 $x_i(x_i-1)=0$，}
\step{所以 $u*u=(x_1+x_1-x_1^2)+(x_2+x_2-x_2^2)+\cdots+(x_{2020}+x_{2020}-x_{2020}^2)$}
\step{$=x_1+x_2+\cdots+x_{2020}$，}
\step{$v*v=(y_1+y_1-y_1^2)+(y_2+y_2-y_2^2)+\cdots+(y_{2020}+y_{2020}-y_{2020}^2)$}
\step{$=y_1+y_2+\cdots+y_{2020}$，}
\step{由 $u*u+v*v=2020$，得 $x_1+x_2+\cdots+x_{2020}+y_1+y_2+\cdots+y_{2020}=2020$，}
% 原卷此行“$\cdots$”后脱逗号、 “2020”后衍一半角逗号（“2020,个 1”）；照录（见文件头“原卷笔误/存疑”③）
\step{所以 $x_1,x_2,\cdots x_{2020},y_1,y_2,\cdots y_{2020}$ 中有 2020,个 1，2020 个 0，}
\step{$u*v=(x_1+y_1-x_1y_1)+(x_2+y_2-x_2y_2)+\cdots+(x_{2020}+y_{2020}-x_{2020}y_{2020})$}
\step{$=2020-(x_1y_1+x_2y_2+\cdots+x_{2020}y_{2020})$，}
\step{注意到只有 $x_i=y_i=1$ 时，$x_iy_i=1$，否则 $x_iy_i=0$，}
% 此行“$\cdots$”后同样脱逗号，与上一处同式样；照录（见文件头“原卷笔误/存疑”③）
\step{而 $x_1,x_2,\cdots x_{2020},y_1,y_2,\cdots y_{2020}$ 中有 2020 个 1，2020 个 0，}
\step{所以满足 $x_iy_i=1$ 的最多有 1010 个，所以 $u*v\geqslant2020-1010=1010$，}
\step{满足 $x_iy_i=0$ 的最多有 2020 个，所以 $u*v\leqslant2020-0=2020$，}
\step{则 $u*v$ 的最大值和最小值之和为 $2020+1010=3030$.}
\solend
\fi

\end{enumerate}

\bigsec{二、选择题}

\begin{enumerate}
\setcounter{enumi}{12}

\item 实数 $x$ 与 $y$，“$x+y>2$”的\mbox{（\quad）}条件是“$x>1$ 或 $y>1$”.

\keepwithprev
A. 充分非必要\qquad B. 必要非充分

C. 充要\qquad D. 既非充分又非必要

\ifanswer
\solbegin
\step{当 $x>1$ 或 $y>1$ 时，$x+y>2$ 不一定成立，充分性不成立；}
\step{当 $x+y>2$ 时，$x>1$ 或 $y>1$ 一定成立，必要性成立. 故选 $B$.}
\solend
\fi

\item 设 $a,b\in\R$，且 $a<0$，$-1<b<0$，那么下列不等式成立的是\mbox{（\quad）}

\keepwithprev
A. $a>ab>ab^2$\qquad B. $ab^2>ab>a$\qquad C. $ab>a>ab^2$\qquad D. $ab>ab^2>a$

\ifanswer
\solbegin
\step{因为 $-1<b<0$，$a<0$，所以 $ab>0$，$0<b^2<1$.}
\step{所以 $ab-ab^2=ab(1-b)>0$，$ab^2-a=a(b^2-1)>0$，}
\step{所以 $ab>ab^2>a$. 故选 $D$.}
\solend
\fi

% 第 15 题答案原卷印在题干括号内（“(  D  )”），且无【答案】/【解析】段；
% 按 高一37 第 13 题、高一63 第 14 题成例括号留空、答案另提 \ans{$D$}
\item 用反证法证明：“平面四边形中至少有一个内角不超过 90°”时，下列“假设”中正确的是\mbox{（\quad）}

\keepwithprev
A. 假设有两个内角超过 90°\qquad B. 假设有三个内角超过 90°

C. 假设至多有两个内角超过 90°\qquad D. 假设四个内角都超过 90°

\ifanswer
\ans{$D$}
\fi

\item 设 $A$ 为非空实数集，集合 $B=\{xy\mid x\in A,y\in A,x\neq y\}$，则对于以下两个结论说法正确的是\mbox{（\quad）}

结论①：若 $A$ 是五元集，则 $B$ 中的元素最多 10 个；

结论②：存在四元集 $A$，使得 $B=\{2,4,5,8,10,16\}$.

\keepwithprev
A. ①对②对\qquad B. ①对②错\qquad C. ①错②对\qquad D. ①错②错

\ifanswer
\solbegin
\step{对于①，想让元素个数达到最大值 10，则须避免产生符号不同的乘积，}
\step{否则正负分开计数容易导致总数受限，且容易产生相等值，}
\step{也最好避免 0 的出现，}
\step{故设 $A=\{a_1,a_2,a_3,a_4,a_5\}$，$0<a_1<a_2<a_3<a_4<a_5$.}
% 原卷此链“$a_3a_3$”同足标自乘、“$a_{14}a_5$”下标 14 均与文意不合，疑为 $a_2a_4$、$a_4a_5$ 之误；照录（见文件头“原卷笔误/存疑”④）
\step{因为 $a_1a_2<a_1a_3<a_1a_4<a_1a_5<a_2a_3<a_3a_3<a_{14}a_5$，$a_2a_3<a_2a_4<a_3a_4$，}
\step{所以 $B$ 中元素个数小于等于 10 个，}
\step{如设 $A=\{1,2,3,5,7\}$，则 $B=\{2,3,5,6,7,10,14,15,21,35\}$，}
\step{所以集合 $B$ 中元素个数的最大值为 10 个，故①正确；}
\step{对于②，由于集合 $B$ 中没有 0，所以集合 $A$ 中没有 0；}
\step{设集合 $A=\{a,b,c,d\}$，且 $a<b<c<d$，}
\step{当 $a<0<b<c<d$ 时，$a\times b<0$，$a\times c<0$，$a\times d<0$，}
\step{故集合 $B$ 至少有 3 个负数，与 $B=\{2,4,5,8,10,16\}$ 矛盾；}
\step{当 $a<b<0<c<d$ 时，$a\times b>0$，$c\times d>0$，}
\step{$a\times c<0$，$a\times d<0$，$b\times c<0$，$b\times d<0$，}
\step{故集合 $B$ 至少有 4 个负数，与 $B=\{2,4,5,8,10,16\}$ 矛盾；}
\step{当 $a<b<c<0<d$ 时，$a\times b>0$，$a\times c>0$，$b\times c>0$，}
\step{$a\times d<0$，$b\times d<0$，$c\times d<0$，}
\step{故集合 $B$ 至少有 3 个负数，与 $B=\{2,4,5,8,10,16\}$ 矛盾；}
\step{当 $0<a<b<c<d$ 时，则集合 $B=\{ab,ac,ad,bc,bd,cd\}$，}
\step{则必有 $ab=2$，$cd=16$，其 4 个正实数的乘积 $abcd=32$，}
\step{又 $ad=5$，$bc=8$ 或 $ad=8$，$bc=5$，其 4 个正实数的乘积 $abcd=40$，矛盾；}
\step{当 $a<b<c<d<0$ 时，与 $0<a<b<c<d$ 时的证明方法相同，}
\step{得集合 $A$ 也不存在；}
\step{综上不存在四元集 $A$，使其生成集 $B=\{2,4,5,8,10,16\}$，故②错误.}
\step{故选 $B$.}
\solend
\fi

\end{enumerate}

\bigsec{三、解答题}

\begin{enumerate}
\setcounter{enumi}{16}

\item 解下列不等式（组）

\begin{enumerate}[label=(\arabic*)]
\item 解不等式：$(ax+1)(x+1)<0(a\in\R)$.
\item 解不等式组：$\begin{cases}\dfrac{2x+4}{x+4}\geqslant1\\[4pt] (x+1)(2-|x|)\geqslant0\end{cases}$.
\end{enumerate}

\ifanswer
\solbegin
\stepm{(1)}{当 $a=0$ 时，原不等式为 $x+1<0$，即 $x<-1$，}
\step{所以此时不等式的解集为 $\{x\mid x<-1\}$；}
\step{当 $a\neq0$ 时，令 $(ax+1)(x+1)=0$，解得 $x_1=-\dfrac1a$,$x_2=-1$，}
\step{当 $a>1$ 时，$-\dfrac1a>-1$，此时不等式的解集为 $\left\{x\mid-1<x<-\dfrac1a\right\}$；}
% 以下两行“此时”前原卷均无标点；照录（见文件头“原卷笔误/存疑”⑥）
\step{当 $a=1$ 时，原不等式为 $(x+1)^2<0$ 此时不等式的解集为 $\varnothing$；}
\step{当 $0<a<1$ 时，$-\dfrac1a<-1$ 此时不等式的解集为 $\left\{x\mid-\dfrac1a<x<-1\right\}$；}
\step{当 $a<0$ 时，$-\dfrac1a>-1$，此时不等式的解集为 $\left\{x\mid x<-1\ \text{或}\ x>-\dfrac1a\right\}$.}
\step{综上所述，当 $a=0$ 时，$\{x\mid x<-1\}$；}
\step{当 $a>1$ 时，$\left\{x\mid-1<x<-\dfrac1a\right\}$；}
\step{当 $a=1$ 时，$\varnothing$；}
\step{当 $0<a<1$ 时，$\left\{x\mid-\dfrac1a<x<-1\right\}$；}
\step{当 $a<0$ 时，$\left\{x\mid x<-1\ \text{或}\ x>-\dfrac1a\right\}$.}
\stepm{(2)}{不等式 $\dfrac{2x+4}{x+4}\geqslant1$ 可化为 $\dfrac{2x+4}{x+4}-1\geqslant0$，即 $\dfrac x{x+4}\geqslant0$，}
\step{等价于 $\begin{cases}x(x+4)\geqslant0\\ x+4\neq0\end{cases}$，解得 $x<-4$ 或 $x\geqslant0$，}
\step{不等式 $(x+1)(2-|x|)\geqslant0$ 等价于 $\begin{cases}(x+1)(2-x)\geqslant0,x\geqslant0\\ (x+1)(2+x)\geqslant0,x<0\end{cases}$，}
\step{解不等式 $(x+1)(2-x)\geqslant0$，$x\geqslant0$，得 $0\leqslant x\leqslant2$，}
\step{解不等式 $(x+1)(2+x)\geqslant0$，$x<0$，得 $x\leqslant-2$ 或 $-1\leqslant x<0$，}
\step{所以 $(x+1)(2-|x|)\geqslant0$ 的解集为 $\{x\mid x\leqslant-2\ \text{或}\ -1\leqslant x\leqslant2\}$，}
\step{综上所述，$\begin{cases}\dfrac{2x+4}{x+4}\geqslant1\\[4pt] (x+1)(2-|x|)\geqslant0\end{cases}$ 的解集为 $\{x\mid x<-4\ \text{或}\ 0\leqslant x\leqslant2\}$.}
\solend
\fi

\ansspace[8]

% 第 18 题题干＋两小问与作答区 6 行须同页
\ansneed{20cm}
\item 已知 $a$ 为实数，设 $f(x)=x^2+2(a+1)x+1$.

\begin{enumerate}[label=(\arabic*)]
\item 已知对于任意实数 $x$，$\sqrt{f(x)}$ 均有意义，求实数 $a$ 的取值范围；
\item 设 $A=\{x\mid f(x)+a^2-2=0\}$，$B=\{-4,0\}$，$A\cup B=B$，求实数 $a$ 的取值范围.
\end{enumerate}

\ifanswer
\solbegin
\stepm{(1)}{$\sqrt{f(x)}$ 对任意实数 $x$ 都有意义，}
\step{即 $f(x)=x^2+2(a+1)x+1\geqslant0$ 对任意 $x\in\R$ 恒成立，}
\step{所以 $\Delta=4(a+1)^2-4\times1\times1\leqslant0$，即 $(a+1)^2\leqslant1$，解得 $-2\leqslant a\leqslant0$，}
\step{故 $a$ 的取值范围是 $[-2,0]$.}
\stepm{(2)}{由 $f(x)+a^2-2=0$ 得 $x^2+2(a+1)x+a^2-1=0$.}
\step{由 $A\cup B=B$ 得 $A\subseteq B=\{-4,0\}$，即该方程无实根，}
\step{或它的实根都在 $\{-4,0\}$ 中.}
\step{当方程无实根时，$\Delta=4(a+1)^2-4(a^2-1)=8a+8<0$，解得 $a<-1$；}
\step{当方程有实根时，$a\geqslant-1$，设两根为 $x_1$、$x_2$，}
\step{由根与系数的关系得 $x_1+x_2=-2(a+1)$、$x_1x_2=a^2-1$，}
\step{此时两根只能在 $\{-4,0\}$ 中取值；}
\step{若两根为 0 与 0，由 $x_1x_2=0$ 得 $a=\pm1$，}
\step{再由 $x_1+x_2=0$ 得 $a=-1$；}
\step{若两根为 0 与 $-4$，由 $x_1x_2=0$ 得 $a=\pm1$，再由 $x_1+x_2=-4$ 得 $a=1$.}
\step{经检验，$a=-1$ 时 $A=\{0\}$，$a=1$ 时 $A=\{-4,0\}$，都满足 $A\subseteq B$.}
\step{综上，实数 $a$ 的取值范围为 $(-\infty,-1]\cup\{1\}$.}
\solend
\fi

\ansspace[8]

\ansneed{18cm}
\item 共享单车给市民出行带来了诸多便利，某公司购买了一批单车投放到某地给居民使用，据市场分析，这批单车的运营累计利润 $y$（单位：元）与运营天数 $x$（单位：天，$x$ 为正整数）满足关系：$y=-\dfrac12x^2+60x-800$.

\begin{enumerate}[label=(\arabic*)]
\item 要使这批单车的运营累计利润高于 800 元，求运营天数的取值范围；
\item 判断这批单车的日平均运营利润能否超过 20 元，并说明理由.
\end{enumerate}

\ifanswer
\solbegin
\stepm{(1)}{由题意得 $y=-\dfrac12x^2+60x-800>800$，}
\step{即 $x^2-120x+3200<0$，解得 $40<x<80$，}
\step{即运营天数的取值范围为 $\{x\mid40<x<80,x\in\N^*\}$.}
\stepm{(2)}{不能，理由如下：}
\step{这批单车的日平均运营利润为 $\dfrac yx=-\dfrac12x-\dfrac{800}x+60=-\dfrac12\left(x+\dfrac{1600}x\right)+60$，}
\step{因 $x>0$，则 $x+\dfrac{1600}x\geqslant2\sqrt{x\cdot\dfrac{1600}x}=80$，}
\step{当且仅当 $x=\dfrac{1600}x$，即 $x=40$ 时取等号，}
\step{此时 $\dfrac yx=-\dfrac12\left(x+\dfrac{1600}x\right)+60\leqslant-\dfrac12\times80+60=20$，}
\step{故这批单车的日平均运营利润不能超过 20 元.}
\solend
\fi

\ansspace[8]

% 第 20 题题干＋三小问与作答区 8 行须同页
\ansneed{20cm}
\item 设 $f(x)=ax^2+bx+c$($a\neq0$，$b$、$c\in\R$).

\begin{enumerate}[label=(\arabic*)]
\item 若 $a=b=c=1$，设 $A=\{(x,y)\mid y=f(x)\}$，$B=\{(x,y)\mid y=1\}$，用列举法表示 $A\cap B$；
% 原卷两个“若”连用（“若 a=1，…，若方程 f(x)=0 有两个大于 1 的根”），疑衍其一；照录（见文件头“原卷笔误/存疑”⑦）
\item 若 $a=1$，$b=2k-1$，$c=k^2$，若方程 $f(x)=0$ 有两个大于 1 的根，求实数 $k$ 的取值范围；
\item 若 $a$、$b$、$c\in\Z$，且 $f(0)$、$f(1)$ 均为奇数，求证：方程 $f(x)=0$ 无整数解.
\end{enumerate}

\ifanswer
\solbegin
\stepm{(1)}{当 $a=b=c=1$ 时，$f(x)=x^2+x+1$，}
\step{所以 $A=\{(x,y)\mid y=x^2+x+1\}$，又 $y=1$，}
\step{由 $x^2+x+1=1$ 得 $x^2+x=0$，解得 $x=0$ 或 $x=-1$，}
\step{所以 $A\cap B=\{(0,1),(-1,1)\}$.}
\stepm{(2)}{当 $a=1$、$b=2k-1$、$c=k^2$ 时，$f(x)=x^2+(2k-1)x+k^2$.}
\step{设两根为 $x_1$、$x_2$，}
\step{则“两根都大于 1”等价于 $\Delta>0$、$f(1)>0$、$-\dfrac b{2a}>1$.}
\step{由 $\Delta=(2k-1)^2-4k^2=1-4k>0$ 得 $k<\dfrac14$；}
\step{由 $f(1)=1+(2k-1)+k^2=k^2+2k=k(k+2)>0$ 得 $k<-2$ 或 $k>0$；}
\step{由 $-\dfrac{2k-1}2=\dfrac{1-2k}2>1$ 得 $k<-\dfrac12$.}
\step{三者取交集得 $k<-2$，即 $k$ 的取值范围是 $(-\infty,-2)$.}
\stepm{(3)}{设 $x_0\in\Z$ 是方程 $f(x)=0$ 的整数解.}
\step{因为 $x_0(x_0-1)$ 是两个相邻整数之积，恒为偶数，所以 $ax_0(x_0-1)$ 为偶数.}
\step{而 $f(x_0)=ax_0^2+bx_0+c=ax_0(x_0-1)+(a+b)x_0+c$，}
\step{所以 $f(x_0)$ 与 $(a+b)x_0+c$ 同奇偶.}
\step{若 $x_0$ 为偶数，则 $(a+b)x_0$ 为偶数，$f(x_0)$ 与 $c=f(0)$ 同奇偶，为奇数；}
\step{若 $x_0$ 为奇数，则 $(a+b)x_0$ 与 $a+b$ 同奇偶，}
\step{$f(x_0)$ 与 $(a+b)+c=f(1)$ 同奇偶，也为奇数.}
\step{综上，$f(x_0)$ 恒为奇数，不可能等于 0，与 $f(x_0)=0$ 矛盾，}
\step{所以方程 $f(x)=0$ 无整数解.}
\solend
\fi

\ansspace[8]

% 压轴题（题干＋三小问）与作答区 10 行须同页
\ansneed{20cm}
\item 已知 $X$ 为有限集，且任取 $x\in X$，均有 $x\in\Z$，称 $X$ 为非空有限整数集. 定义：$X_k=\{x+k\mid x\in X\}$，$X^2=\{y\mid y=x^2,x\in X\}$.

\begin{enumerate}[label=(\arabic*)]
\item 若 $X=\{-2,2\}$，分别写出 $X_2$ 和 $X^2$（用列举法表示）；
\item 若 $X=\{-1,0,1\}$，$X_1\subseteq Y$，$(Y_{-1})^2\subseteq X$，求集合 $Y$ 中所有的元素；
\item 对于非空有限整数集 $A$、$B$，若 $A_k\subseteq B$，$(B^2)_{-2}\subseteq A$，$B$ 中的元素均为非负整数，$k$ 为整数，求证：

① $A$、$B$ 中元素的个数相同；

② 若集合 $A$ 中的元素个数为奇数，则 $k$ 为偶数.
\end{enumerate}

\ifanswer
\solbegin
\stepm{(1)}{由题意得 $X_2=\{0,4\}$，$X^2=\{4\}$；}
\stepm{(2)}{由题意得 $X_1=\{0,1,2\}$，即 $\{0,1,2\}\subseteq Y$，$Y_{-1}=\{y-1\mid y\in Y\}$，}
\step{$(Y_{-1})^2=\{(y-1)^2\mid y\in Y\}\subseteq X=\{-1,0,1\}$，}
\step{所以 $(y-1)^2=0$ 或 $(y-1)^2=1$，解得 $y=1$ 或 $y=2$ 或 $y=0$，}
\step{即 $y\in\{0,1,2\}$，即 $Y\subseteq\{0,1,2\}$，所以 $Y=\{0,1,2\}$；}
% 原卷 (3) 解析首行“(3)”与“①”连排，本稿合并记作 \stepm{(3)①}
\stepm{(3)①}{因为 $A_k=\{x+k\mid x\in A\}$，函数 $g(x)=x+k$ 在全体整数集上严格增，}
\step{所以 $\forall x_1$、$x_2\in A$，$x_1\neq x_2$，$g(x_1)\neq g(x_2)$，所以 $|A_k|=|A|$，}
\step{因为 $A_k\subseteq B$，所以 $|A|\leqslant|B|$，}
\step{设 $b\in B$，$B$ 中元素均为非负整数，$\forall b_1\neq b_2\in B$，$b_1^2\neq b_2^2$，}
\step{故函数 $f(b)=b^2-2$ 在 $B$ 上单射，所以 $|(B^2)_{-2}|=|B|$，}
\step{又因为 $(B^2)_{-2}\subseteq A$，所以 $|B|\leqslant|A|$，所以 $|A|=|B|$.}
\stepm{②}{由①得 $A_k\subseteq B$ 且 $|A_k|=|A|=|B|$，$(B^2)_{-2}\subseteq A$ 且 $|(B^2)_{-2}|=|B|=|A|$，}
\step{所以 $A_k=B$，$(B^2)_{-2}=A$，}
\step{所以 $B=A_k=\{a+k\mid a\in A\}$，$A=(B^2)_{-2}=\{b^2-2\mid b\in B\}$，}
\step{所以 $A=\{(a+k)^2-2\mid a\in A\}$，}
\step{所以映射 $h(a)=(a+k)^2-2$ 是集合 $A$ 到自身的一个双射.}
\step{假设 $k$ 是奇数，当 $a$ 为偶数时，$a+k$ 为奇数，$(a+k)^2-2$ 为奇数，}
\step{当 $a$ 为奇数时，$a+k$ 为偶数，$(a+k)^2-2$ 为偶数，}
\step{所以此时映射 $h(a)=(a+k)^2-2$，将集合 $A$ 中的奇数映射为偶数，}
\step{将集合 $A$ 中的偶数映射为奇数，即集合 $A$ 中的奇数元素和偶数元素必须配对，}
\step{此时集合 $A$ 的元素个数则为偶数，}
\step{与题设中“集合 $A$ 中的元素个数为奇数”矛盾.}
\step{故假设不成立，所以 $k$ 是偶数.}
\solend
\fi

\ansspace*[10]

\end{enumerate}

\end{document}
"
% 紧凑版：xelatex "\def\iscompact{1}% 高一67_华师大一附中_10月月考.tex
%   —— 高一上数学“2026-2027 年华师大一附中高一上 10 月月考”（讲评版）
% 试卷版：xelatex 高一67_华师大一附中_10月月考.tex
% 答案版：xelatex "\def\isanswer{1}\input{高一67_华师大一附中_10月月考.tex}"
% 紧凑版：xelatex "\def\iscompact{1}\input{高一67_华师大一附中_10月月考.tex}"
%
% 原始材料：微信公众号“魔都数学加油站”文章，作者破晓，连载“27学年高一67”；
% 原文标题“27学年高一67——2026-2027年上海市华师大一附中高一上10月月考”；
% 原文链接：https://mp.weixin.qq.com/s/Kdhx07TVjLoGTcIlnRzgcQ。
% 页面标注发布时间 2026-10-18 20:33（上海）。
% 本次转录只据 `原图/高一67_华师大一附中_10月月考/` 的 13 张 PNG 逐页读图。
% 正文为 13 张 PNG 页（`原图/高一67_华师大一附中_10月月考/01..13.png`）：
% 第 01、04、13 页 2560×3620，第 02、03、05—12 页 4762×6735
% （**同一篇各页分辨率不同**），均按页序读取。13 页全为试卷正文页（卷首标题在第 1 页页顶，
% 卷末解答题第 21 题解析止于第 13 页），无封面页、目录页、二维码或宣传页。
% 粉色斜排“魔都数学加油站”为卷面水印，与公众号页脚一并未录入。
%
% 卷首标题：照卷面印作“2026--2027 年华师大一附中高一上 10 月月考”（卷面为黑体加粗一行），
% 年份区间按全稿惯例排 en dash。原卷只有一行标题、无副标题。
%
% 篇幅与题号：原卷自第 3 题起编号（第 1、2 题不在本卷内）。填空题 3—12（10 题）、
% 选择题 13—16（4 题）、解答题 17—21（5 题），共 19 题。本稿沿用原节与原题号，
% 只改排为全稿统一的大节样式（原卷小节标题为“一、填空题”“二、选择题”“三、解答题”）。
%
% 讲评版：原卷每题下直接印【解析】。**全卷 19 题中 18 题都印【解析】（无一处只印
% 【答案】）**，按全稿口径这类题不另提 \ans{}；各题解析末句的结论（第 13、14、16 题
% 的“故选 B／D／B”、各填空与解答题的“故/综上…”）照录在解析内。
% **另有 1 题答案印在题干括号内**——第 15 题（“(  D  )”，且该题无【答案】/【解析】段），
% 本稿按 高一37 第 13 题、高一63 第 14 题成例把括号留空、答案另提为 \ans{$D$}。
%
% 副标题：原卷未印副标题。本稿按“知识模块”口径标作“集合与常用逻辑用语、不等式”：
% 19 题中集合与常用逻辑用语 12 题（第 3、4、5、7、10、11、12、13、15、16、18、21 题：
% 补集、并集、子集、充分条件、容斥计数、子集计数、新定义二元运算、条件判断、
% 反证法假设、新定义生成集、集合关系与二次方程、新定义 X_k 与 X^2），
% 不等式 7 题（第 6、8、9、14、17、19、20 题：存在量词不等式、解集反求系数、
% 解集为空、不等式性质、解不等式（组）、二次不等式与基本不等式、二次方程根的分布），
% 两类均超过三分之一，故并标。
%
% 与原卷的排版差异：
%   ① 原卷每题下直接印【解析】，本稿按全稿惯例将解析置于答案版，试卷版保持干净；
%   ② 原卷未印副标题，本稿按“知识模块”口径补出；
%   ③ 原卷填空横线约四个字宽，本稿按全稿惯例统一排 \blankline[4.4em]；
%   ④ 原卷实数集排斜体 $R$（第 14 题题干、第 17 题 (1) 题干、第 18 题解析、第 20 题题干），
%      整数集排斜体 $Z$（第 20 题 (3) 题干与解析、第 21 题题干），正整数集排粗体
%      $N^*$（第 11 题解析、第 19 题解析），本稿按全稿惯例统一改用 $\R$、$\Z$、$\N^*$
%      （记号归一，粗体／斜体不作区分）；空集原卷作 $\varnothing$，照录；
%   ⑤ 原卷在数学式内多处用半角逗号（如第 8 题 $(1,2)$ 与各 cases、第 11 题解析
%      $P_n=\{1,2,\cdots,n\}$,$n\in N^*$、第 12 题解析 $u=(\cdots)$,$v=(\cdots)$ 与
%      “2020,个 1”、第 19 题解析 $\{x|40<x<80,x\in N^*\}$），均照录未强改；
%      描述法集合的竖线照原卷排 \mid；句末句点按全稿惯例排西文句点；
%   ⑥ 原卷补集记号为上划线 $\bar A$（第 3 题题干与解析、第 11 题题干），照录
%      （同 高一64/65 的 $\bar B$ 成例）；
%   ⑦ 第 12 题新定义的运算记号 $*$ 照录（$\alpha*\beta$、$u*u$、$u*v$）；
%   ⑧ 第 17—21 题的子问原卷各自另起一行，本稿按全稿惯例用 enumerate 的 (1)(2)(3)；
%      解析中的子问标号一律用 \stepm{(1)}{…}；第 5 题解析的①②、第 21 题 (3) 解析的①②
%      按全稿惯例作 \stepm{①}{…}、\stepm{②}{…}（第 21 题 (3) 首行原卷“(3)”与“①”连排，
%      本稿合并记作 \stepm{(3)①}{…}）；
%   ⑨ 选择题作答括号原卷排半角“(　　)”（第 13、14、16 题括号内空白、第 15 题括号内印 D），
%      本稿按全稿惯例排 \mbox{（\quad）}；第 15 题括号内的答案移入 \ans{}，见上；
%   ⑩ 选项行排照原卷：第 13、15 题两行两条，第 14、16 题一行四条；
%   ⑪ 每道解答题原卷未留作答空白，本稿按全稿惯例加 \ansspace（末题用 \ansspace*）；
%      第 20、21 题题干长，各用一处 \ansneed 让题干与作答区不拆页（\ansneed 只作用于试卷版）。
%
% 原卷笔误/存疑（均照录未改，正文对应处另有 % 注）：
%   ① 第 3 题解析末句印作“则 $\bar A=(-\infty,0)\cup[2,3)$.”：按题干
%      $U=\{x\mid 1/x\leqslant 1/2\}$（即 $x<0$ 或 $x\geqslant2$）与 $A=[3,+\infty)$，
%      $\bar A$ 应为 $(-\infty,0)\cup\{2\}$；且与本解析首句自求之 $U$ 相矛盾（$3\notin U$），
%      疑为笔误；照录未改。
%   ② 第 12 题题干 $\alpha*\beta$ 定义式括号不配对：第四组“$+\cdots+x_{2020}+y_{2020}
%      -x_{2020}y_{2020}$”缺左括号、末尾连排两个右括号“))”；照录未改。
%   ③ 第 12 题解析“所以…中有 2020,个 1”一行，“$\cdots$”后脱逗号（作
%      “$x_2,\cdots x_{2020}$”，第 12 行“而…”同式样），且“2020”后衍一半角逗号
%      （作“2020,个 1”）；照录未改。
%   ④ 第 16 题解析乘积链印作“$a_1a_2<a_1a_3<a_1a_4<a_1a_5<a_2a_3<a_3a_3<a_{14}a_5$”：
%      “$a_3a_3$”同足标自乘、与单调链不合，“$a_{14}a_5$”下标 14 与 $a_1$—$a_5$ 记号
%      不合（疑为 $a_2a_4$、$a_4a_5$ 之误）；照录未改。
%   ⑤ 第 5 题题干作“$Q\subset P$”（真子集记号），解析两处均作“$Q\subseteq P$”，
%      同一题前后记号不一；照录（本稿分别排 \subset、\subseteq）。
%   ⑥ 第 17 题 (1) 解析“当 $a=1$ 时，原不等式为 $(x+1)^2<0$ 此时…”与
%      “当 $0<a<1$ 时，$-\frac1a<-1$ 此时…”两处“此时”前均脱标点；照录未改。
%   ⑦ 第 20 题 (2) 题干“若 $a=1$，$b=2k-1$，$c=k^2$，若方程 $f(x)=0$ 有两个大于 1 的根”
%      两个“若”连用，疑衍其一；照录未改。
%
% 插图：无（全卷 13 页无任何几何图形、无“如图/右图/图中”字样）。
%
% 行数核对：逐题数 \step/\stepm 条数与原卷【解析】印刷行数，18 题逐一相符
% （第 15 题无解析段，只提 \ans{}），合计 169 条 / 169 行
% （Q3—Q12：2、3、4、3、1、3、5、3、13、15；Q13—Q16：2、3、0、25；
% Q17—Q21：19、16、9、20、23）。
%
\documentclass[a4paper,11pt]{article}
\usepackage{common}

\begin{document}

\examtitlest[3]{2026--2027 年华师大一附中高一上 10 月月考}{集合与常用逻辑用语、不等式}

\bigsec{一、填空题}

\begin{enumerate}
\setcounter{enumi}{2}

\item 设全集 $U=\left\{x\mid\dfrac1x\leqslant\dfrac12\right\}$，$A=[3,+\infty)$，则 $\bar A=$\blankline[4.4em].

\ifanswer
\solbegin
\step{$\dfrac1x\leqslant\dfrac12$，解得 $x<0$ 或 $x\geqslant2$，故 $U=(-\infty,0)\cup[2,+\infty)$，}
% 原卷末句作 $\bar A=(-\infty,0)\cup[2,3)$，按题干 $U$ 应为 $(-\infty,0)\cup\{2\}$，与本解析首句矛盾，疑笔误；照录（见文件头“原卷笔误/存疑”①）
\step{$A=[3,+\infty)$，则 $\bar A=(-\infty,0)\cup[2,3)$.}
\solend
\fi

\item 已知集合 $M=\{y\mid y=x^2-2x-1\}$，$N=\{y\mid y=x^2-4x+1\}$，则 $M\cup N=$\blankline[4.4em].

\ifanswer
\solbegin
\step{$M=\{y\mid y=x^2-2x-1\}=\{y\mid y=(x-1)^2-2\}=\{y\mid y\geqslant-2\}$，}
\step{$N=\{y\mid y=x^2-4x+1\}=\{y\mid y=(x-2)^2-3\}=\{y\mid y\geqslant-3\}$，}
\step{则 $M\cup N=\{y\mid y\geqslant-3\}$.}
\solend
\fi

\item 已知 $P=\{x\mid x^2+x-6=0\}$，$Q=\{x\mid mx-1=0\}$，且 $Q\subset P$，则实数 $m$ 的值\blankline[4.4em].

\ifanswer
\solbegin
\step{$P=\{2,-3\}$，$Q=\{x\mid mx=1\}$；}
\stepm{①}{$m=0$ 时，$Q=\varnothing$，满足 $Q\subseteq P$；}
\stepm{②}{$m\neq0$ 时，要使 $Q\subseteq P$，则 $\dfrac1m=2$ 或 $-3$，所以 $m=\dfrac12$ 或 $-\dfrac13$；}
\step{所以 $m=0$ 或 $\dfrac12$ 或 $-\dfrac13$.}
\solend
\fi

\item 已知“存在正实数 $x$，使得关于 $x$ 的不等式 $ax-1>0$ 成立”为真命题，则实数 $a$ 的取值范围是\blankline[4.4em].

\ifanswer
\solbegin
\step{存在正实数 $x$，使得关于 $x$ 的不等式 $ax-1>0$ 成立为真命题，}
\step{即存在正实数 $x$，使得 $a>\dfrac1x$，因为 $\dfrac1x>0$，所以 $a>0$，}
\step{即 $a$ 的范围为 $(0,+\infty)$.}
\solend
\fi

\item 设 $\alpha:3\leqslant x\leqslant5$，$\beta:2m-1\leqslant x<m+4$，若 $\alpha$ 是 $\beta$ 的充分条件，则实数 $m$ 的取值范围是\blankline[4.4em].

\ifanswer
\solbegin
\step{若 $\alpha$ 是 $\beta$ 的充分条件，则 $\begin{cases}2m-1\leqslant3\\ m+4>5\end{cases}$，解得 $1<m\leqslant2$.}
\solend
\fi

\item 对于实数 $a$、$b$、$c$，已知关于 $x$ 的不等式 $ax^2+bx+c>0$ 的解集是 $(1,2)$，则关于 $x$ 的不等式 $cx^2-bx+a\leqslant0$ 的解集是\blankline[4.4em].

\ifanswer
\solbegin
\step{关于 $x$ 的不等式 $ax^2+bx+c>0$ 的解集是 $(1,2)$，则 $\begin{cases}a<0\\ -\dfrac ba=3\\ \dfrac ca=2\end{cases}$，则 $\begin{cases}b=-3a\\ c=2a\end{cases}$，}
\step{则关于 $x$ 的不等式 $cx^2-bx+a=2ax^2+3ax+a\leqslant0$，则 $2x^2+3x+1\geqslant0$，}
\step{则 $x\leqslant-1$ 或 $x\geqslant-\dfrac12$，则解集为 $(-\infty,-1]\cup\left[-\dfrac12,+\infty\right)$.}
\solend
\fi

\item 已知关于 $x$ 的不等式 $kx^2+kx+k+2\leqslant0$ 的解集为 $\varnothing$，则实数 $k$ 的取值范围为\blankline[4.4em].

\ifanswer
\solbegin
\step{当 $k=0$ 时，不等式化为 $2\leqslant0$，解集为 $\varnothing$，符合题意；}
\step{当 $k\neq0$ 时，若关于 $x$ 的不等式 $kx^2+kx+k+2\leqslant0$ 的解集为 $\varnothing$，}
\step{则 $\begin{cases}k>0\\ \Delta=k^2-4k(k+2)<0\end{cases}$，化简得 $\begin{cases}k>0\\ -3k^2-8k<0\end{cases}$，即 $\begin{cases}k>0\\ k(3k+8)>0\end{cases}$，}
\step{解得 $k>0$.}
\step{综上所述，$k\geqslant0$，即实数 $k$ 的取值范围为 $[0,+\infty)$.}
\solend
\fi

\item 对班级若干名学生调查对 $A$、$B$ 两事件的态度，有如下结果：赞成 $A$ 的人数是全体的五分之三，其余的不赞成；赞成 $B$ 的比赞成 $A$ 多 3 人，其余的不赞成. 已知对 $A$、$B$ 都赞成的学生有 18 人，对 $A$、$B$ 都不赞成的学生数比对 $A$、$B$ 都赞成的学生数的三分之一多 1 人，则班级总人数是\blankline[4.4em]人.

\ifanswer
\solbegin
\step{设班级总人数是 $x$ 人，则赞成 $A$ 的人数是 $\dfrac35x$，赞成 $B$ 的人数是 $\dfrac35x+3$，}
\step{$A$、$B$ 都不赞成的人数为 $18\times\dfrac13+1=7$ 人，}
\step{由容斥原理得 $x=\left(\dfrac35x+\dfrac35x+3-18\right)+7$，解得 $x=40$，即班级总人数是 40 人.}
\solend
\fi

\item 设全集 $U=\{1,2,3,4,5,6,7,8,9,10,11,12\}$，若集合 $A$ 满足：① $A\subseteq U$；②若 $x\in A$，必有 $2x\notin A$；③若 $x\in\bar A$，必有 $2x\notin\bar A$，则这样的集合 $A$ 有\blankline[4.4em]个.

\ifanswer
\solbegin
\step{设集合 $P_n=\{1,2,\cdots,n\}$,$n\in\N^*$，}
\step{任取 $x\in P_n$，将 $x$ 除以 2，若商仍为偶数，再除以 2，$\cdots$，}
\step{经过 $k$ 次以后，商必为奇数，此时记商为 $m$，}
\step{于是 $x=m\cdot2^k$，其中 $m$ 为奇数，$k\in\N^*$.}
\step{由条件得，若 $m\in A$，则 $x\in A$ 等价于 $k$ 为偶数；}
\step{若 $m\notin A$，则 $x\in A$ 等价于 $k$ 为奇数；}
\step{于是 $x$ 是否属于 $A$ 由 $m$ 是否属于 $A$ 确定；}
\step{设 $Q_n$ 是 $P_n$ 中所有奇数的集合，}
\step{设 $f(n)$ 为集合 $P_n$ 中满足题干条件的集合 $A$ 的个数，}
\step{因此 $f(n)$ 等于 $Q_n$ 的子集个数.}
\step{当 $n$ 为偶数（或奇数）时，$P_n$ 中奇数的个数是 $\dfrac n2$（或 $\dfrac{n+1}2$），}
\step{所以 $f(n)=\begin{cases}2^{\frac n2},n\text{为偶数}\\[2pt] 2^{\frac{n+1}2},n\text{为奇数}\end{cases}$，所以 $P_{12}=\{1,2,3,4,5,6,7,8,9,10,11,12\}$，}
\step{即同时满足三个条件的集合 $A$ 的个数为 $2^{\frac{12}2}=2^6=64$.}
\solend
\fi

\item 设集合 $\Omega=\{(x_1,x_2,\cdots,x_{2020})\mid x_i\in\{0,1\},i=1,2,\cdots,2020\}$，对于 $\Omega$ 中的任意两个元素 $\alpha=(x_1,x_2,\cdots,x_{2020})$，$\beta=(y_1,y_2,\cdots,y_{2020})$，记
% 原卷此式第四组缺左括号、末尾连排两个右括号“))”，括号不配对；照录（见文件头“原卷笔误/存疑”②）
$\alpha*\beta=(x_1+y_1-x_1y_1)+(x_2+y_2-x_2y_2)+\cdots+x_{2020}+y_{2020}-x_{2020}y_{2020}))$，设 $u$、$v\in\Omega$，若 $u*u+v*v=2020$，则 $u*v$ 的最大值和最小值之和为\blankline[4.4em].

\ifanswer
\solbegin
\step{设 $u=(x_1,x_2,\cdots,x_{2020})$,$v=(y_1,y_2,\cdots,y_{2020})$，}
\step{因为 $x_i\in\{0,1\}$，所以 $x_i(x_i-1)=0$，}
\step{所以 $u*u=(x_1+x_1-x_1^2)+(x_2+x_2-x_2^2)+\cdots+(x_{2020}+x_{2020}-x_{2020}^2)$}
\step{$=x_1+x_2+\cdots+x_{2020}$，}
\step{$v*v=(y_1+y_1-y_1^2)+(y_2+y_2-y_2^2)+\cdots+(y_{2020}+y_{2020}-y_{2020}^2)$}
\step{$=y_1+y_2+\cdots+y_{2020}$，}
\step{由 $u*u+v*v=2020$，得 $x_1+x_2+\cdots+x_{2020}+y_1+y_2+\cdots+y_{2020}=2020$，}
% 原卷此行“$\cdots$”后脱逗号、 “2020”后衍一半角逗号（“2020,个 1”）；照录（见文件头“原卷笔误/存疑”③）
\step{所以 $x_1,x_2,\cdots x_{2020},y_1,y_2,\cdots y_{2020}$ 中有 2020,个 1，2020 个 0，}
\step{$u*v=(x_1+y_1-x_1y_1)+(x_2+y_2-x_2y_2)+\cdots+(x_{2020}+y_{2020}-x_{2020}y_{2020})$}
\step{$=2020-(x_1y_1+x_2y_2+\cdots+x_{2020}y_{2020})$，}
\step{注意到只有 $x_i=y_i=1$ 时，$x_iy_i=1$，否则 $x_iy_i=0$，}
% 此行“$\cdots$”后同样脱逗号，与上一处同式样；照录（见文件头“原卷笔误/存疑”③）
\step{而 $x_1,x_2,\cdots x_{2020},y_1,y_2,\cdots y_{2020}$ 中有 2020 个 1，2020 个 0，}
\step{所以满足 $x_iy_i=1$ 的最多有 1010 个，所以 $u*v\geqslant2020-1010=1010$，}
\step{满足 $x_iy_i=0$ 的最多有 2020 个，所以 $u*v\leqslant2020-0=2020$，}
\step{则 $u*v$ 的最大值和最小值之和为 $2020+1010=3030$.}
\solend
\fi

\end{enumerate}

\bigsec{二、选择题}

\begin{enumerate}
\setcounter{enumi}{12}

\item 实数 $x$ 与 $y$，“$x+y>2$”的\mbox{（\quad）}条件是“$x>1$ 或 $y>1$”.

\keepwithprev
A. 充分非必要\qquad B. 必要非充分

C. 充要\qquad D. 既非充分又非必要

\ifanswer
\solbegin
\step{当 $x>1$ 或 $y>1$ 时，$x+y>2$ 不一定成立，充分性不成立；}
\step{当 $x+y>2$ 时，$x>1$ 或 $y>1$ 一定成立，必要性成立. 故选 $B$.}
\solend
\fi

\item 设 $a,b\in\R$，且 $a<0$，$-1<b<0$，那么下列不等式成立的是\mbox{（\quad）}

\keepwithprev
A. $a>ab>ab^2$\qquad B. $ab^2>ab>a$\qquad C. $ab>a>ab^2$\qquad D. $ab>ab^2>a$

\ifanswer
\solbegin
\step{因为 $-1<b<0$，$a<0$，所以 $ab>0$，$0<b^2<1$.}
\step{所以 $ab-ab^2=ab(1-b)>0$，$ab^2-a=a(b^2-1)>0$，}
\step{所以 $ab>ab^2>a$. 故选 $D$.}
\solend
\fi

% 第 15 题答案原卷印在题干括号内（“(  D  )”），且无【答案】/【解析】段；
% 按 高一37 第 13 题、高一63 第 14 题成例括号留空、答案另提 \ans{$D$}
\item 用反证法证明：“平面四边形中至少有一个内角不超过 90°”时，下列“假设”中正确的是\mbox{（\quad）}

\keepwithprev
A. 假设有两个内角超过 90°\qquad B. 假设有三个内角超过 90°

C. 假设至多有两个内角超过 90°\qquad D. 假设四个内角都超过 90°

\ifanswer
\ans{$D$}
\fi

\item 设 $A$ 为非空实数集，集合 $B=\{xy\mid x\in A,y\in A,x\neq y\}$，则对于以下两个结论说法正确的是\mbox{（\quad）}

结论①：若 $A$ 是五元集，则 $B$ 中的元素最多 10 个；

结论②：存在四元集 $A$，使得 $B=\{2,4,5,8,10,16\}$.

\keepwithprev
A. ①对②对\qquad B. ①对②错\qquad C. ①错②对\qquad D. ①错②错

\ifanswer
\solbegin
\step{对于①，想让元素个数达到最大值 10，则须避免产生符号不同的乘积，}
\step{否则正负分开计数容易导致总数受限，且容易产生相等值，}
\step{也最好避免 0 的出现，}
\step{故设 $A=\{a_1,a_2,a_3,a_4,a_5\}$，$0<a_1<a_2<a_3<a_4<a_5$.}
% 原卷此链“$a_3a_3$”同足标自乘、“$a_{14}a_5$”下标 14 均与文意不合，疑为 $a_2a_4$、$a_4a_5$ 之误；照录（见文件头“原卷笔误/存疑”④）
\step{因为 $a_1a_2<a_1a_3<a_1a_4<a_1a_5<a_2a_3<a_3a_3<a_{14}a_5$，$a_2a_3<a_2a_4<a_3a_4$，}
\step{所以 $B$ 中元素个数小于等于 10 个，}
\step{如设 $A=\{1,2,3,5,7\}$，则 $B=\{2,3,5,6,7,10,14,15,21,35\}$，}
\step{所以集合 $B$ 中元素个数的最大值为 10 个，故①正确；}
\step{对于②，由于集合 $B$ 中没有 0，所以集合 $A$ 中没有 0；}
\step{设集合 $A=\{a,b,c,d\}$，且 $a<b<c<d$，}
\step{当 $a<0<b<c<d$ 时，$a\times b<0$，$a\times c<0$，$a\times d<0$，}
\step{故集合 $B$ 至少有 3 个负数，与 $B=\{2,4,5,8,10,16\}$ 矛盾；}
\step{当 $a<b<0<c<d$ 时，$a\times b>0$，$c\times d>0$，}
\step{$a\times c<0$，$a\times d<0$，$b\times c<0$，$b\times d<0$，}
\step{故集合 $B$ 至少有 4 个负数，与 $B=\{2,4,5,8,10,16\}$ 矛盾；}
\step{当 $a<b<c<0<d$ 时，$a\times b>0$，$a\times c>0$，$b\times c>0$，}
\step{$a\times d<0$，$b\times d<0$，$c\times d<0$，}
\step{故集合 $B$ 至少有 3 个负数，与 $B=\{2,4,5,8,10,16\}$ 矛盾；}
\step{当 $0<a<b<c<d$ 时，则集合 $B=\{ab,ac,ad,bc,bd,cd\}$，}
\step{则必有 $ab=2$，$cd=16$，其 4 个正实数的乘积 $abcd=32$，}
\step{又 $ad=5$，$bc=8$ 或 $ad=8$，$bc=5$，其 4 个正实数的乘积 $abcd=40$，矛盾；}
\step{当 $a<b<c<d<0$ 时，与 $0<a<b<c<d$ 时的证明方法相同，}
\step{得集合 $A$ 也不存在；}
\step{综上不存在四元集 $A$，使其生成集 $B=\{2,4,5,8,10,16\}$，故②错误.}
\step{故选 $B$.}
\solend
\fi

\end{enumerate}

\bigsec{三、解答题}

\begin{enumerate}
\setcounter{enumi}{16}

\item 解下列不等式（组）

\begin{enumerate}[label=(\arabic*)]
\item 解不等式：$(ax+1)(x+1)<0(a\in\R)$.
\item 解不等式组：$\begin{cases}\dfrac{2x+4}{x+4}\geqslant1\\[4pt] (x+1)(2-|x|)\geqslant0\end{cases}$.
\end{enumerate}

\ifanswer
\solbegin
\stepm{(1)}{当 $a=0$ 时，原不等式为 $x+1<0$，即 $x<-1$，}
\step{所以此时不等式的解集为 $\{x\mid x<-1\}$；}
\step{当 $a\neq0$ 时，令 $(ax+1)(x+1)=0$，解得 $x_1=-\dfrac1a$,$x_2=-1$，}
\step{当 $a>1$ 时，$-\dfrac1a>-1$，此时不等式的解集为 $\left\{x\mid-1<x<-\dfrac1a\right\}$；}
% 以下两行“此时”前原卷均无标点；照录（见文件头“原卷笔误/存疑”⑥）
\step{当 $a=1$ 时，原不等式为 $(x+1)^2<0$ 此时不等式的解集为 $\varnothing$；}
\step{当 $0<a<1$ 时，$-\dfrac1a<-1$ 此时不等式的解集为 $\left\{x\mid-\dfrac1a<x<-1\right\}$；}
\step{当 $a<0$ 时，$-\dfrac1a>-1$，此时不等式的解集为 $\left\{x\mid x<-1\ \text{或}\ x>-\dfrac1a\right\}$.}
\step{综上所述，当 $a=0$ 时，$\{x\mid x<-1\}$；}
\step{当 $a>1$ 时，$\left\{x\mid-1<x<-\dfrac1a\right\}$；}
\step{当 $a=1$ 时，$\varnothing$；}
\step{当 $0<a<1$ 时，$\left\{x\mid-\dfrac1a<x<-1\right\}$；}
\step{当 $a<0$ 时，$\left\{x\mid x<-1\ \text{或}\ x>-\dfrac1a\right\}$.}
\stepm{(2)}{不等式 $\dfrac{2x+4}{x+4}\geqslant1$ 可化为 $\dfrac{2x+4}{x+4}-1\geqslant0$，即 $\dfrac x{x+4}\geqslant0$，}
\step{等价于 $\begin{cases}x(x+4)\geqslant0\\ x+4\neq0\end{cases}$，解得 $x<-4$ 或 $x\geqslant0$，}
\step{不等式 $(x+1)(2-|x|)\geqslant0$ 等价于 $\begin{cases}(x+1)(2-x)\geqslant0,x\geqslant0\\ (x+1)(2+x)\geqslant0,x<0\end{cases}$，}
\step{解不等式 $(x+1)(2-x)\geqslant0$，$x\geqslant0$，得 $0\leqslant x\leqslant2$，}
\step{解不等式 $(x+1)(2+x)\geqslant0$，$x<0$，得 $x\leqslant-2$ 或 $-1\leqslant x<0$，}
\step{所以 $(x+1)(2-|x|)\geqslant0$ 的解集为 $\{x\mid x\leqslant-2\ \text{或}\ -1\leqslant x\leqslant2\}$，}
\step{综上所述，$\begin{cases}\dfrac{2x+4}{x+4}\geqslant1\\[4pt] (x+1)(2-|x|)\geqslant0\end{cases}$ 的解集为 $\{x\mid x<-4\ \text{或}\ 0\leqslant x\leqslant2\}$.}
\solend
\fi

\ansspace[8]

% 第 18 题题干＋两小问与作答区 6 行须同页
\ansneed{20cm}
\item 已知 $a$ 为实数，设 $f(x)=x^2+2(a+1)x+1$.

\begin{enumerate}[label=(\arabic*)]
\item 已知对于任意实数 $x$，$\sqrt{f(x)}$ 均有意义，求实数 $a$ 的取值范围；
\item 设 $A=\{x\mid f(x)+a^2-2=0\}$，$B=\{-4,0\}$，$A\cup B=B$，求实数 $a$ 的取值范围.
\end{enumerate}

\ifanswer
\solbegin
\stepm{(1)}{$\sqrt{f(x)}$ 对任意实数 $x$ 都有意义，}
\step{即 $f(x)=x^2+2(a+1)x+1\geqslant0$ 对任意 $x\in\R$ 恒成立，}
\step{所以 $\Delta=4(a+1)^2-4\times1\times1\leqslant0$，即 $(a+1)^2\leqslant1$，解得 $-2\leqslant a\leqslant0$，}
\step{故 $a$ 的取值范围是 $[-2,0]$.}
\stepm{(2)}{由 $f(x)+a^2-2=0$ 得 $x^2+2(a+1)x+a^2-1=0$.}
\step{由 $A\cup B=B$ 得 $A\subseteq B=\{-4,0\}$，即该方程无实根，}
\step{或它的实根都在 $\{-4,0\}$ 中.}
\step{当方程无实根时，$\Delta=4(a+1)^2-4(a^2-1)=8a+8<0$，解得 $a<-1$；}
\step{当方程有实根时，$a\geqslant-1$，设两根为 $x_1$、$x_2$，}
\step{由根与系数的关系得 $x_1+x_2=-2(a+1)$、$x_1x_2=a^2-1$，}
\step{此时两根只能在 $\{-4,0\}$ 中取值；}
\step{若两根为 0 与 0，由 $x_1x_2=0$ 得 $a=\pm1$，}
\step{再由 $x_1+x_2=0$ 得 $a=-1$；}
\step{若两根为 0 与 $-4$，由 $x_1x_2=0$ 得 $a=\pm1$，再由 $x_1+x_2=-4$ 得 $a=1$.}
\step{经检验，$a=-1$ 时 $A=\{0\}$，$a=1$ 时 $A=\{-4,0\}$，都满足 $A\subseteq B$.}
\step{综上，实数 $a$ 的取值范围为 $(-\infty,-1]\cup\{1\}$.}
\solend
\fi

\ansspace[8]

\ansneed{18cm}
\item 共享单车给市民出行带来了诸多便利，某公司购买了一批单车投放到某地给居民使用，据市场分析，这批单车的运营累计利润 $y$（单位：元）与运营天数 $x$（单位：天，$x$ 为正整数）满足关系：$y=-\dfrac12x^2+60x-800$.

\begin{enumerate}[label=(\arabic*)]
\item 要使这批单车的运营累计利润高于 800 元，求运营天数的取值范围；
\item 判断这批单车的日平均运营利润能否超过 20 元，并说明理由.
\end{enumerate}

\ifanswer
\solbegin
\stepm{(1)}{由题意得 $y=-\dfrac12x^2+60x-800>800$，}
\step{即 $x^2-120x+3200<0$，解得 $40<x<80$，}
\step{即运营天数的取值范围为 $\{x\mid40<x<80,x\in\N^*\}$.}
\stepm{(2)}{不能，理由如下：}
\step{这批单车的日平均运营利润为 $\dfrac yx=-\dfrac12x-\dfrac{800}x+60=-\dfrac12\left(x+\dfrac{1600}x\right)+60$，}
\step{因 $x>0$，则 $x+\dfrac{1600}x\geqslant2\sqrt{x\cdot\dfrac{1600}x}=80$，}
\step{当且仅当 $x=\dfrac{1600}x$，即 $x=40$ 时取等号，}
\step{此时 $\dfrac yx=-\dfrac12\left(x+\dfrac{1600}x\right)+60\leqslant-\dfrac12\times80+60=20$，}
\step{故这批单车的日平均运营利润不能超过 20 元.}
\solend
\fi

\ansspace[8]

% 第 20 题题干＋三小问与作答区 8 行须同页
\ansneed{20cm}
\item 设 $f(x)=ax^2+bx+c$($a\neq0$，$b$、$c\in\R$).

\begin{enumerate}[label=(\arabic*)]
\item 若 $a=b=c=1$，设 $A=\{(x,y)\mid y=f(x)\}$，$B=\{(x,y)\mid y=1\}$，用列举法表示 $A\cap B$；
% 原卷两个“若”连用（“若 a=1，…，若方程 f(x)=0 有两个大于 1 的根”），疑衍其一；照录（见文件头“原卷笔误/存疑”⑦）
\item 若 $a=1$，$b=2k-1$，$c=k^2$，若方程 $f(x)=0$ 有两个大于 1 的根，求实数 $k$ 的取值范围；
\item 若 $a$、$b$、$c\in\Z$，且 $f(0)$、$f(1)$ 均为奇数，求证：方程 $f(x)=0$ 无整数解.
\end{enumerate}

\ifanswer
\solbegin
\stepm{(1)}{当 $a=b=c=1$ 时，$f(x)=x^2+x+1$，}
\step{所以 $A=\{(x,y)\mid y=x^2+x+1\}$，又 $y=1$，}
\step{由 $x^2+x+1=1$ 得 $x^2+x=0$，解得 $x=0$ 或 $x=-1$，}
\step{所以 $A\cap B=\{(0,1),(-1,1)\}$.}
\stepm{(2)}{当 $a=1$、$b=2k-1$、$c=k^2$ 时，$f(x)=x^2+(2k-1)x+k^2$.}
\step{设两根为 $x_1$、$x_2$，}
\step{则“两根都大于 1”等价于 $\Delta>0$、$f(1)>0$、$-\dfrac b{2a}>1$.}
\step{由 $\Delta=(2k-1)^2-4k^2=1-4k>0$ 得 $k<\dfrac14$；}
\step{由 $f(1)=1+(2k-1)+k^2=k^2+2k=k(k+2)>0$ 得 $k<-2$ 或 $k>0$；}
\step{由 $-\dfrac{2k-1}2=\dfrac{1-2k}2>1$ 得 $k<-\dfrac12$.}
\step{三者取交集得 $k<-2$，即 $k$ 的取值范围是 $(-\infty,-2)$.}
\stepm{(3)}{设 $x_0\in\Z$ 是方程 $f(x)=0$ 的整数解.}
\step{因为 $x_0(x_0-1)$ 是两个相邻整数之积，恒为偶数，所以 $ax_0(x_0-1)$ 为偶数.}
\step{而 $f(x_0)=ax_0^2+bx_0+c=ax_0(x_0-1)+(a+b)x_0+c$，}
\step{所以 $f(x_0)$ 与 $(a+b)x_0+c$ 同奇偶.}
\step{若 $x_0$ 为偶数，则 $(a+b)x_0$ 为偶数，$f(x_0)$ 与 $c=f(0)$ 同奇偶，为奇数；}
\step{若 $x_0$ 为奇数，则 $(a+b)x_0$ 与 $a+b$ 同奇偶，}
\step{$f(x_0)$ 与 $(a+b)+c=f(1)$ 同奇偶，也为奇数.}
\step{综上，$f(x_0)$ 恒为奇数，不可能等于 0，与 $f(x_0)=0$ 矛盾，}
\step{所以方程 $f(x)=0$ 无整数解.}
\solend
\fi

\ansspace[8]

% 压轴题（题干＋三小问）与作答区 10 行须同页
\ansneed{20cm}
\item 已知 $X$ 为有限集，且任取 $x\in X$，均有 $x\in\Z$，称 $X$ 为非空有限整数集. 定义：$X_k=\{x+k\mid x\in X\}$，$X^2=\{y\mid y=x^2,x\in X\}$.

\begin{enumerate}[label=(\arabic*)]
\item 若 $X=\{-2,2\}$，分别写出 $X_2$ 和 $X^2$（用列举法表示）；
\item 若 $X=\{-1,0,1\}$，$X_1\subseteq Y$，$(Y_{-1})^2\subseteq X$，求集合 $Y$ 中所有的元素；
\item 对于非空有限整数集 $A$、$B$，若 $A_k\subseteq B$，$(B^2)_{-2}\subseteq A$，$B$ 中的元素均为非负整数，$k$ 为整数，求证：

① $A$、$B$ 中元素的个数相同；

② 若集合 $A$ 中的元素个数为奇数，则 $k$ 为偶数.
\end{enumerate}

\ifanswer
\solbegin
\stepm{(1)}{由题意得 $X_2=\{0,4\}$，$X^2=\{4\}$；}
\stepm{(2)}{由题意得 $X_1=\{0,1,2\}$，即 $\{0,1,2\}\subseteq Y$，$Y_{-1}=\{y-1\mid y\in Y\}$，}
\step{$(Y_{-1})^2=\{(y-1)^2\mid y\in Y\}\subseteq X=\{-1,0,1\}$，}
\step{所以 $(y-1)^2=0$ 或 $(y-1)^2=1$，解得 $y=1$ 或 $y=2$ 或 $y=0$，}
\step{即 $y\in\{0,1,2\}$，即 $Y\subseteq\{0,1,2\}$，所以 $Y=\{0,1,2\}$；}
% 原卷 (3) 解析首行“(3)”与“①”连排，本稿合并记作 \stepm{(3)①}
\stepm{(3)①}{因为 $A_k=\{x+k\mid x\in A\}$，函数 $g(x)=x+k$ 在全体整数集上严格增，}
\step{所以 $\forall x_1$、$x_2\in A$，$x_1\neq x_2$，$g(x_1)\neq g(x_2)$，所以 $|A_k|=|A|$，}
\step{因为 $A_k\subseteq B$，所以 $|A|\leqslant|B|$，}
\step{设 $b\in B$，$B$ 中元素均为非负整数，$\forall b_1\neq b_2\in B$，$b_1^2\neq b_2^2$，}
\step{故函数 $f(b)=b^2-2$ 在 $B$ 上单射，所以 $|(B^2)_{-2}|=|B|$，}
\step{又因为 $(B^2)_{-2}\subseteq A$，所以 $|B|\leqslant|A|$，所以 $|A|=|B|$.}
\stepm{②}{由①得 $A_k\subseteq B$ 且 $|A_k|=|A|=|B|$，$(B^2)_{-2}\subseteq A$ 且 $|(B^2)_{-2}|=|B|=|A|$，}
\step{所以 $A_k=B$，$(B^2)_{-2}=A$，}
\step{所以 $B=A_k=\{a+k\mid a\in A\}$，$A=(B^2)_{-2}=\{b^2-2\mid b\in B\}$，}
\step{所以 $A=\{(a+k)^2-2\mid a\in A\}$，}
\step{所以映射 $h(a)=(a+k)^2-2$ 是集合 $A$ 到自身的一个双射.}
\step{假设 $k$ 是奇数，当 $a$ 为偶数时，$a+k$ 为奇数，$(a+k)^2-2$ 为奇数，}
\step{当 $a$ 为奇数时，$a+k$ 为偶数，$(a+k)^2-2$ 为偶数，}
\step{所以此时映射 $h(a)=(a+k)^2-2$，将集合 $A$ 中的奇数映射为偶数，}
\step{将集合 $A$ 中的偶数映射为奇数，即集合 $A$ 中的奇数元素和偶数元素必须配对，}
\step{此时集合 $A$ 的元素个数则为偶数，}
\step{与题设中“集合 $A$ 中的元素个数为奇数”矛盾.}
\step{故假设不成立，所以 $k$ 是偶数.}
\solend
\fi

\ansspace*[10]

\end{enumerate}

\end{document}
"
%
% 原始材料：微信公众号“魔都数学加油站”文章，作者破晓，连载“27学年高一67”；
% 原文标题“27学年高一67——2026-2027年上海市华师大一附中高一上10月月考”；
% 原文链接：https://mp.weixin.qq.com/s/Kdhx07TVjLoGTcIlnRzgcQ。
% 页面标注发布时间 2026-10-18 20:33（上海）。
% 本次转录只据 `原图/高一67_华师大一附中_10月月考/` 的 13 张 PNG 逐页读图。
% 正文为 13 张 PNG 页（`原图/高一67_华师大一附中_10月月考/01..13.png`）：
% 第 01、04、13 页 2560×3620，第 02、03、05—12 页 4762×6735
% （**同一篇各页分辨率不同**），均按页序读取。13 页全为试卷正文页（卷首标题在第 1 页页顶，
% 卷末解答题第 21 题解析止于第 13 页），无封面页、目录页、二维码或宣传页。
% 粉色斜排“魔都数学加油站”为卷面水印，与公众号页脚一并未录入。
%
% 卷首标题：照卷面印作“2026--2027 年华师大一附中高一上 10 月月考”（卷面为黑体加粗一行），
% 年份区间按全稿惯例排 en dash。原卷只有一行标题、无副标题。
%
% 篇幅与题号：原卷自第 3 题起编号（第 1、2 题不在本卷内）。填空题 3—12（10 题）、
% 选择题 13—16（4 题）、解答题 17—21（5 题），共 19 题。本稿沿用原节与原题号，
% 只改排为全稿统一的大节样式（原卷小节标题为“一、填空题”“二、选择题”“三、解答题”）。
%
% 讲评版：原卷每题下直接印【解析】。**全卷 19 题中 18 题都印【解析】（无一处只印
% 【答案】）**，按全稿口径这类题不另提 \ans{}；各题解析末句的结论（第 13、14、16 题
% 的“故选 B／D／B”、各填空与解答题的“故/综上…”）照录在解析内。
% **另有 1 题答案印在题干括号内**——第 15 题（“(  D  )”，且该题无【答案】/【解析】段），
% 本稿按 高一37 第 13 题、高一63 第 14 题成例把括号留空、答案另提为 \ans{$D$}。
%
% 副标题：原卷未印副标题。本稿按“知识模块”口径标作“集合与常用逻辑用语、不等式”：
% 19 题中集合与常用逻辑用语 12 题（第 3、4、5、7、10、11、12、13、15、16、18、21 题：
% 补集、并集、子集、充分条件、容斥计数、子集计数、新定义二元运算、条件判断、
% 反证法假设、新定义生成集、集合关系与二次方程、新定义 X_k 与 X^2），
% 不等式 7 题（第 6、8、9、14、17、19、20 题：存在量词不等式、解集反求系数、
% 解集为空、不等式性质、解不等式（组）、二次不等式与基本不等式、二次方程根的分布），
% 两类均超过三分之一，故并标。
%
% 与原卷的排版差异：
%   ① 原卷每题下直接印【解析】，本稿按全稿惯例将解析置于答案版，试卷版保持干净；
%   ② 原卷未印副标题，本稿按“知识模块”口径补出；
%   ③ 原卷填空横线约四个字宽，本稿按全稿惯例统一排 \blankline[4.4em]；
%   ④ 原卷实数集排斜体 $R$（第 14 题题干、第 17 题 (1) 题干、第 18 题解析、第 20 题题干），
%      整数集排斜体 $Z$（第 20 题 (3) 题干与解析、第 21 题题干），正整数集排粗体
%      $N^*$（第 11 题解析、第 19 题解析），本稿按全稿惯例统一改用 $\R$、$\Z$、$\N^*$
%      （记号归一，粗体／斜体不作区分）；空集原卷作 $\varnothing$，照录；
%   ⑤ 原卷在数学式内多处用半角逗号（如第 8 题 $(1,2)$ 与各 cases、第 11 题解析
%      $P_n=\{1,2,\cdots,n\}$,$n\in N^*$、第 12 题解析 $u=(\cdots)$,$v=(\cdots)$ 与
%      “2020,个 1”、第 19 题解析 $\{x|40<x<80,x\in N^*\}$），均照录未强改；
%      描述法集合的竖线照原卷排 \mid；句末句点按全稿惯例排西文句点；
%   ⑥ 原卷补集记号为上划线 $\bar A$（第 3 题题干与解析、第 11 题题干），照录
%      （同 高一64/65 的 $\bar B$ 成例）；
%   ⑦ 第 12 题新定义的运算记号 $*$ 照录（$\alpha*\beta$、$u*u$、$u*v$）；
%   ⑧ 第 17—21 题的子问原卷各自另起一行，本稿按全稿惯例用 enumerate 的 (1)(2)(3)；
%      解析中的子问标号一律用 \stepm{(1)}{…}；第 5 题解析的①②、第 21 题 (3) 解析的①②
%      按全稿惯例作 \stepm{①}{…}、\stepm{②}{…}（第 21 题 (3) 首行原卷“(3)”与“①”连排，
%      本稿合并记作 \stepm{(3)①}{…}）；
%   ⑨ 选择题作答括号原卷排半角“(　　)”（第 13、14、16 题括号内空白、第 15 题括号内印 D），
%      本稿按全稿惯例排 \mbox{（\quad）}；第 15 题括号内的答案移入 \ans{}，见上；
%   ⑩ 选项行排照原卷：第 13、15 题两行两条，第 14、16 题一行四条；
%   ⑪ 每道解答题原卷未留作答空白，本稿按全稿惯例加 \ansspace（末题用 \ansspace*）；
%      第 20、21 题题干长，各用一处 \ansneed 让题干与作答区不拆页（\ansneed 只作用于试卷版）。
%
% 原卷笔误/存疑（均照录未改，正文对应处另有 % 注）：
%   ① 第 3 题解析末句印作“则 $\bar A=(-\infty,0)\cup[2,3)$.”：按题干
%      $U=\{x\mid 1/x\leqslant 1/2\}$（即 $x<0$ 或 $x\geqslant2$）与 $A=[3,+\infty)$，
%      $\bar A$ 应为 $(-\infty,0)\cup\{2\}$；且与本解析首句自求之 $U$ 相矛盾（$3\notin U$），
%      疑为笔误；照录未改。
%   ② 第 12 题题干 $\alpha*\beta$ 定义式括号不配对：第四组“$+\cdots+x_{2020}+y_{2020}
%      -x_{2020}y_{2020}$”缺左括号、末尾连排两个右括号“))”；照录未改。
%   ③ 第 12 题解析“所以…中有 2020,个 1”一行，“$\cdots$”后脱逗号（作
%      “$x_2,\cdots x_{2020}$”，第 12 行“而…”同式样），且“2020”后衍一半角逗号
%      （作“2020,个 1”）；照录未改。
%   ④ 第 16 题解析乘积链印作“$a_1a_2<a_1a_3<a_1a_4<a_1a_5<a_2a_3<a_3a_3<a_{14}a_5$”：
%      “$a_3a_3$”同足标自乘、与单调链不合，“$a_{14}a_5$”下标 14 与 $a_1$—$a_5$ 记号
%      不合（疑为 $a_2a_4$、$a_4a_5$ 之误）；照录未改。
%   ⑤ 第 5 题题干作“$Q\subset P$”（真子集记号），解析两处均作“$Q\subseteq P$”，
%      同一题前后记号不一；照录（本稿分别排 \subset、\subseteq）。
%   ⑥ 第 17 题 (1) 解析“当 $a=1$ 时，原不等式为 $(x+1)^2<0$ 此时…”与
%      “当 $0<a<1$ 时，$-\frac1a<-1$ 此时…”两处“此时”前均脱标点；照录未改。
%   ⑦ 第 20 题 (2) 题干“若 $a=1$，$b=2k-1$，$c=k^2$，若方程 $f(x)=0$ 有两个大于 1 的根”
%      两个“若”连用，疑衍其一；照录未改。
%
% 插图：无（全卷 13 页无任何几何图形、无“如图/右图/图中”字样）。
%
% 行数核对：逐题数 \step/\stepm 条数与原卷【解析】印刷行数，18 题逐一相符
% （第 15 题无解析段，只提 \ans{}），合计 169 条 / 169 行
% （Q3—Q12：2、3、4、3、1、3、5、3、13、15；Q13—Q16：2、3、0、25；
% Q17—Q21：19、16、9、20、23）。
%
\documentclass[a4paper,11pt]{article}
\usepackage{common}

\begin{document}

\examtitlest[3]{2026--2027 年华师大一附中高一上 10 月月考}{集合与常用逻辑用语、不等式}

\bigsec{一、填空题}

\begin{enumerate}
\setcounter{enumi}{2}

\item 设全集 $U=\left\{x\mid\dfrac1x\leqslant\dfrac12\right\}$，$A=[3,+\infty)$，则 $\bar A=$\blankline[4.4em].

\ifanswer
\solbegin
\step{$\dfrac1x\leqslant\dfrac12$，解得 $x<0$ 或 $x\geqslant2$，故 $U=(-\infty,0)\cup[2,+\infty)$，}
% 原卷末句作 $\bar A=(-\infty,0)\cup[2,3)$，按题干 $U$ 应为 $(-\infty,0)\cup\{2\}$，与本解析首句矛盾，疑笔误；照录（见文件头“原卷笔误/存疑”①）
\step{$A=[3,+\infty)$，则 $\bar A=(-\infty,0)\cup[2,3)$.}
\solend
\fi

\item 已知集合 $M=\{y\mid y=x^2-2x-1\}$，$N=\{y\mid y=x^2-4x+1\}$，则 $M\cup N=$\blankline[4.4em].

\ifanswer
\solbegin
\step{$M=\{y\mid y=x^2-2x-1\}=\{y\mid y=(x-1)^2-2\}=\{y\mid y\geqslant-2\}$，}
\step{$N=\{y\mid y=x^2-4x+1\}=\{y\mid y=(x-2)^2-3\}=\{y\mid y\geqslant-3\}$，}
\step{则 $M\cup N=\{y\mid y\geqslant-3\}$.}
\solend
\fi

\item 已知 $P=\{x\mid x^2+x-6=0\}$，$Q=\{x\mid mx-1=0\}$，且 $Q\subset P$，则实数 $m$ 的值\blankline[4.4em].

\ifanswer
\solbegin
\step{$P=\{2,-3\}$，$Q=\{x\mid mx=1\}$；}
\stepm{①}{$m=0$ 时，$Q=\varnothing$，满足 $Q\subseteq P$；}
\stepm{②}{$m\neq0$ 时，要使 $Q\subseteq P$，则 $\dfrac1m=2$ 或 $-3$，所以 $m=\dfrac12$ 或 $-\dfrac13$；}
\step{所以 $m=0$ 或 $\dfrac12$ 或 $-\dfrac13$.}
\solend
\fi

\item 已知“存在正实数 $x$，使得关于 $x$ 的不等式 $ax-1>0$ 成立”为真命题，则实数 $a$ 的取值范围是\blankline[4.4em].

\ifanswer
\solbegin
\step{存在正实数 $x$，使得关于 $x$ 的不等式 $ax-1>0$ 成立为真命题，}
\step{即存在正实数 $x$，使得 $a>\dfrac1x$，因为 $\dfrac1x>0$，所以 $a>0$，}
\step{即 $a$ 的范围为 $(0,+\infty)$.}
\solend
\fi

\item 设 $\alpha:3\leqslant x\leqslant5$，$\beta:2m-1\leqslant x<m+4$，若 $\alpha$ 是 $\beta$ 的充分条件，则实数 $m$ 的取值范围是\blankline[4.4em].

\ifanswer
\solbegin
\step{若 $\alpha$ 是 $\beta$ 的充分条件，则 $\begin{cases}2m-1\leqslant3\\ m+4>5\end{cases}$，解得 $1<m\leqslant2$.}
\solend
\fi

\item 对于实数 $a$、$b$、$c$，已知关于 $x$ 的不等式 $ax^2+bx+c>0$ 的解集是 $(1,2)$，则关于 $x$ 的不等式 $cx^2-bx+a\leqslant0$ 的解集是\blankline[4.4em].

\ifanswer
\solbegin
\step{关于 $x$ 的不等式 $ax^2+bx+c>0$ 的解集是 $(1,2)$，则 $\begin{cases}a<0\\ -\dfrac ba=3\\ \dfrac ca=2\end{cases}$，则 $\begin{cases}b=-3a\\ c=2a\end{cases}$，}
\step{则关于 $x$ 的不等式 $cx^2-bx+a=2ax^2+3ax+a\leqslant0$，则 $2x^2+3x+1\geqslant0$，}
\step{则 $x\leqslant-1$ 或 $x\geqslant-\dfrac12$，则解集为 $(-\infty,-1]\cup\left[-\dfrac12,+\infty\right)$.}
\solend
\fi

\item 已知关于 $x$ 的不等式 $kx^2+kx+k+2\leqslant0$ 的解集为 $\varnothing$，则实数 $k$ 的取值范围为\blankline[4.4em].

\ifanswer
\solbegin
\step{当 $k=0$ 时，不等式化为 $2\leqslant0$，解集为 $\varnothing$，符合题意；}
\step{当 $k\neq0$ 时，若关于 $x$ 的不等式 $kx^2+kx+k+2\leqslant0$ 的解集为 $\varnothing$，}
\step{则 $\begin{cases}k>0\\ \Delta=k^2-4k(k+2)<0\end{cases}$，化简得 $\begin{cases}k>0\\ -3k^2-8k<0\end{cases}$，即 $\begin{cases}k>0\\ k(3k+8)>0\end{cases}$，}
\step{解得 $k>0$.}
\step{综上所述，$k\geqslant0$，即实数 $k$ 的取值范围为 $[0,+\infty)$.}
\solend
\fi

\item 对班级若干名学生调查对 $A$、$B$ 两事件的态度，有如下结果：赞成 $A$ 的人数是全体的五分之三，其余的不赞成；赞成 $B$ 的比赞成 $A$ 多 3 人，其余的不赞成. 已知对 $A$、$B$ 都赞成的学生有 18 人，对 $A$、$B$ 都不赞成的学生数比对 $A$、$B$ 都赞成的学生数的三分之一多 1 人，则班级总人数是\blankline[4.4em]人.

\ifanswer
\solbegin
\step{设班级总人数是 $x$ 人，则赞成 $A$ 的人数是 $\dfrac35x$，赞成 $B$ 的人数是 $\dfrac35x+3$，}
\step{$A$、$B$ 都不赞成的人数为 $18\times\dfrac13+1=7$ 人，}
\step{由容斥原理得 $x=\left(\dfrac35x+\dfrac35x+3-18\right)+7$，解得 $x=40$，即班级总人数是 40 人.}
\solend
\fi

\item 设全集 $U=\{1,2,3,4,5,6,7,8,9,10,11,12\}$，若集合 $A$ 满足：① $A\subseteq U$；②若 $x\in A$，必有 $2x\notin A$；③若 $x\in\bar A$，必有 $2x\notin\bar A$，则这样的集合 $A$ 有\blankline[4.4em]个.

\ifanswer
\solbegin
\step{设集合 $P_n=\{1,2,\cdots,n\}$,$n\in\N^*$，}
\step{任取 $x\in P_n$，将 $x$ 除以 2，若商仍为偶数，再除以 2，$\cdots$，}
\step{经过 $k$ 次以后，商必为奇数，此时记商为 $m$，}
\step{于是 $x=m\cdot2^k$，其中 $m$ 为奇数，$k\in\N^*$.}
\step{由条件得，若 $m\in A$，则 $x\in A$ 等价于 $k$ 为偶数；}
\step{若 $m\notin A$，则 $x\in A$ 等价于 $k$ 为奇数；}
\step{于是 $x$ 是否属于 $A$ 由 $m$ 是否属于 $A$ 确定；}
\step{设 $Q_n$ 是 $P_n$ 中所有奇数的集合，}
\step{设 $f(n)$ 为集合 $P_n$ 中满足题干条件的集合 $A$ 的个数，}
\step{因此 $f(n)$ 等于 $Q_n$ 的子集个数.}
\step{当 $n$ 为偶数（或奇数）时，$P_n$ 中奇数的个数是 $\dfrac n2$（或 $\dfrac{n+1}2$），}
\step{所以 $f(n)=\begin{cases}2^{\frac n2},n\text{为偶数}\\[2pt] 2^{\frac{n+1}2},n\text{为奇数}\end{cases}$，所以 $P_{12}=\{1,2,3,4,5,6,7,8,9,10,11,12\}$，}
\step{即同时满足三个条件的集合 $A$ 的个数为 $2^{\frac{12}2}=2^6=64$.}
\solend
\fi

\item 设集合 $\Omega=\{(x_1,x_2,\cdots,x_{2020})\mid x_i\in\{0,1\},i=1,2,\cdots,2020\}$，对于 $\Omega$ 中的任意两个元素 $\alpha=(x_1,x_2,\cdots,x_{2020})$，$\beta=(y_1,y_2,\cdots,y_{2020})$，记
% 原卷此式第四组缺左括号、末尾连排两个右括号“))”，括号不配对；照录（见文件头“原卷笔误/存疑”②）
$\alpha*\beta=(x_1+y_1-x_1y_1)+(x_2+y_2-x_2y_2)+\cdots+x_{2020}+y_{2020}-x_{2020}y_{2020}))$，设 $u$、$v\in\Omega$，若 $u*u+v*v=2020$，则 $u*v$ 的最大值和最小值之和为\blankline[4.4em].

\ifanswer
\solbegin
\step{设 $u=(x_1,x_2,\cdots,x_{2020})$,$v=(y_1,y_2,\cdots,y_{2020})$，}
\step{因为 $x_i\in\{0,1\}$，所以 $x_i(x_i-1)=0$，}
\step{所以 $u*u=(x_1+x_1-x_1^2)+(x_2+x_2-x_2^2)+\cdots+(x_{2020}+x_{2020}-x_{2020}^2)$}
\step{$=x_1+x_2+\cdots+x_{2020}$，}
\step{$v*v=(y_1+y_1-y_1^2)+(y_2+y_2-y_2^2)+\cdots+(y_{2020}+y_{2020}-y_{2020}^2)$}
\step{$=y_1+y_2+\cdots+y_{2020}$，}
\step{由 $u*u+v*v=2020$，得 $x_1+x_2+\cdots+x_{2020}+y_1+y_2+\cdots+y_{2020}=2020$，}
% 原卷此行“$\cdots$”后脱逗号、 “2020”后衍一半角逗号（“2020,个 1”）；照录（见文件头“原卷笔误/存疑”③）
\step{所以 $x_1,x_2,\cdots x_{2020},y_1,y_2,\cdots y_{2020}$ 中有 2020,个 1，2020 个 0，}
\step{$u*v=(x_1+y_1-x_1y_1)+(x_2+y_2-x_2y_2)+\cdots+(x_{2020}+y_{2020}-x_{2020}y_{2020})$}
\step{$=2020-(x_1y_1+x_2y_2+\cdots+x_{2020}y_{2020})$，}
\step{注意到只有 $x_i=y_i=1$ 时，$x_iy_i=1$，否则 $x_iy_i=0$，}
% 此行“$\cdots$”后同样脱逗号，与上一处同式样；照录（见文件头“原卷笔误/存疑”③）
\step{而 $x_1,x_2,\cdots x_{2020},y_1,y_2,\cdots y_{2020}$ 中有 2020 个 1，2020 个 0，}
\step{所以满足 $x_iy_i=1$ 的最多有 1010 个，所以 $u*v\geqslant2020-1010=1010$，}
\step{满足 $x_iy_i=0$ 的最多有 2020 个，所以 $u*v\leqslant2020-0=2020$，}
\step{则 $u*v$ 的最大值和最小值之和为 $2020+1010=3030$.}
\solend
\fi

\end{enumerate}

\bigsec{二、选择题}

\begin{enumerate}
\setcounter{enumi}{12}

\item 实数 $x$ 与 $y$，“$x+y>2$”的\mbox{（\quad）}条件是“$x>1$ 或 $y>1$”.

\keepwithprev
A. 充分非必要\qquad B. 必要非充分

C. 充要\qquad D. 既非充分又非必要

\ifanswer
\solbegin
\step{当 $x>1$ 或 $y>1$ 时，$x+y>2$ 不一定成立，充分性不成立；}
\step{当 $x+y>2$ 时，$x>1$ 或 $y>1$ 一定成立，必要性成立. 故选 $B$.}
\solend
\fi

\item 设 $a,b\in\R$，且 $a<0$，$-1<b<0$，那么下列不等式成立的是\mbox{（\quad）}

\keepwithprev
A. $a>ab>ab^2$\qquad B. $ab^2>ab>a$\qquad C. $ab>a>ab^2$\qquad D. $ab>ab^2>a$

\ifanswer
\solbegin
\step{因为 $-1<b<0$，$a<0$，所以 $ab>0$，$0<b^2<1$.}
\step{所以 $ab-ab^2=ab(1-b)>0$，$ab^2-a=a(b^2-1)>0$，}
\step{所以 $ab>ab^2>a$. 故选 $D$.}
\solend
\fi

% 第 15 题答案原卷印在题干括号内（“(  D  )”），且无【答案】/【解析】段；
% 按 高一37 第 13 题、高一63 第 14 题成例括号留空、答案另提 \ans{$D$}
\item 用反证法证明：“平面四边形中至少有一个内角不超过 90°”时，下列“假设”中正确的是\mbox{（\quad）}

\keepwithprev
A. 假设有两个内角超过 90°\qquad B. 假设有三个内角超过 90°

C. 假设至多有两个内角超过 90°\qquad D. 假设四个内角都超过 90°

\ifanswer
\ans{$D$}
\fi

\item 设 $A$ 为非空实数集，集合 $B=\{xy\mid x\in A,y\in A,x\neq y\}$，则对于以下两个结论说法正确的是\mbox{（\quad）}

结论①：若 $A$ 是五元集，则 $B$ 中的元素最多 10 个；

结论②：存在四元集 $A$，使得 $B=\{2,4,5,8,10,16\}$.

\keepwithprev
A. ①对②对\qquad B. ①对②错\qquad C. ①错②对\qquad D. ①错②错

\ifanswer
\solbegin
\step{对于①，想让元素个数达到最大值 10，则须避免产生符号不同的乘积，}
\step{否则正负分开计数容易导致总数受限，且容易产生相等值，}
\step{也最好避免 0 的出现，}
\step{故设 $A=\{a_1,a_2,a_3,a_4,a_5\}$，$0<a_1<a_2<a_3<a_4<a_5$.}
% 原卷此链“$a_3a_3$”同足标自乘、“$a_{14}a_5$”下标 14 均与文意不合，疑为 $a_2a_4$、$a_4a_5$ 之误；照录（见文件头“原卷笔误/存疑”④）
\step{因为 $a_1a_2<a_1a_3<a_1a_4<a_1a_5<a_2a_3<a_3a_3<a_{14}a_5$，$a_2a_3<a_2a_4<a_3a_4$，}
\step{所以 $B$ 中元素个数小于等于 10 个，}
\step{如设 $A=\{1,2,3,5,7\}$，则 $B=\{2,3,5,6,7,10,14,15,21,35\}$，}
\step{所以集合 $B$ 中元素个数的最大值为 10 个，故①正确；}
\step{对于②，由于集合 $B$ 中没有 0，所以集合 $A$ 中没有 0；}
\step{设集合 $A=\{a,b,c,d\}$，且 $a<b<c<d$，}
\step{当 $a<0<b<c<d$ 时，$a\times b<0$，$a\times c<0$，$a\times d<0$，}
\step{故集合 $B$ 至少有 3 个负数，与 $B=\{2,4,5,8,10,16\}$ 矛盾；}
\step{当 $a<b<0<c<d$ 时，$a\times b>0$，$c\times d>0$，}
\step{$a\times c<0$，$a\times d<0$，$b\times c<0$，$b\times d<0$，}
\step{故集合 $B$ 至少有 4 个负数，与 $B=\{2,4,5,8,10,16\}$ 矛盾；}
\step{当 $a<b<c<0<d$ 时，$a\times b>0$，$a\times c>0$，$b\times c>0$，}
\step{$a\times d<0$，$b\times d<0$，$c\times d<0$，}
\step{故集合 $B$ 至少有 3 个负数，与 $B=\{2,4,5,8,10,16\}$ 矛盾；}
\step{当 $0<a<b<c<d$ 时，则集合 $B=\{ab,ac,ad,bc,bd,cd\}$，}
\step{则必有 $ab=2$，$cd=16$，其 4 个正实数的乘积 $abcd=32$，}
\step{又 $ad=5$，$bc=8$ 或 $ad=8$，$bc=5$，其 4 个正实数的乘积 $abcd=40$，矛盾；}
\step{当 $a<b<c<d<0$ 时，与 $0<a<b<c<d$ 时的证明方法相同，}
\step{得集合 $A$ 也不存在；}
\step{综上不存在四元集 $A$，使其生成集 $B=\{2,4,5,8,10,16\}$，故②错误.}
\step{故选 $B$.}
\solend
\fi

\end{enumerate}

\bigsec{三、解答题}

\begin{enumerate}
\setcounter{enumi}{16}

\item 解下列不等式（组）

\begin{enumerate}[label=(\arabic*)]
\item 解不等式：$(ax+1)(x+1)<0(a\in\R)$.
\item 解不等式组：$\begin{cases}\dfrac{2x+4}{x+4}\geqslant1\\[4pt] (x+1)(2-|x|)\geqslant0\end{cases}$.
\end{enumerate}

\ifanswer
\solbegin
\stepm{(1)}{当 $a=0$ 时，原不等式为 $x+1<0$，即 $x<-1$，}
\step{所以此时不等式的解集为 $\{x\mid x<-1\}$；}
\step{当 $a\neq0$ 时，令 $(ax+1)(x+1)=0$，解得 $x_1=-\dfrac1a$,$x_2=-1$，}
\step{当 $a>1$ 时，$-\dfrac1a>-1$，此时不等式的解集为 $\left\{x\mid-1<x<-\dfrac1a\right\}$；}
% 以下两行“此时”前原卷均无标点；照录（见文件头“原卷笔误/存疑”⑥）
\step{当 $a=1$ 时，原不等式为 $(x+1)^2<0$ 此时不等式的解集为 $\varnothing$；}
\step{当 $0<a<1$ 时，$-\dfrac1a<-1$ 此时不等式的解集为 $\left\{x\mid-\dfrac1a<x<-1\right\}$；}
\step{当 $a<0$ 时，$-\dfrac1a>-1$，此时不等式的解集为 $\left\{x\mid x<-1\ \text{或}\ x>-\dfrac1a\right\}$.}
\step{综上所述，当 $a=0$ 时，$\{x\mid x<-1\}$；}
\step{当 $a>1$ 时，$\left\{x\mid-1<x<-\dfrac1a\right\}$；}
\step{当 $a=1$ 时，$\varnothing$；}
\step{当 $0<a<1$ 时，$\left\{x\mid-\dfrac1a<x<-1\right\}$；}
\step{当 $a<0$ 时，$\left\{x\mid x<-1\ \text{或}\ x>-\dfrac1a\right\}$.}
\stepm{(2)}{不等式 $\dfrac{2x+4}{x+4}\geqslant1$ 可化为 $\dfrac{2x+4}{x+4}-1\geqslant0$，即 $\dfrac x{x+4}\geqslant0$，}
\step{等价于 $\begin{cases}x(x+4)\geqslant0\\ x+4\neq0\end{cases}$，解得 $x<-4$ 或 $x\geqslant0$，}
\step{不等式 $(x+1)(2-|x|)\geqslant0$ 等价于 $\begin{cases}(x+1)(2-x)\geqslant0,x\geqslant0\\ (x+1)(2+x)\geqslant0,x<0\end{cases}$，}
\step{解不等式 $(x+1)(2-x)\geqslant0$，$x\geqslant0$，得 $0\leqslant x\leqslant2$，}
\step{解不等式 $(x+1)(2+x)\geqslant0$，$x<0$，得 $x\leqslant-2$ 或 $-1\leqslant x<0$，}
\step{所以 $(x+1)(2-|x|)\geqslant0$ 的解集为 $\{x\mid x\leqslant-2\ \text{或}\ -1\leqslant x\leqslant2\}$，}
\step{综上所述，$\begin{cases}\dfrac{2x+4}{x+4}\geqslant1\\[4pt] (x+1)(2-|x|)\geqslant0\end{cases}$ 的解集为 $\{x\mid x<-4\ \text{或}\ 0\leqslant x\leqslant2\}$.}
\solend
\fi

\ansspace[8]

% 第 18 题题干＋两小问与作答区 6 行须同页
\ansneed{20cm}
\item 已知 $a$ 为实数，设 $f(x)=x^2+2(a+1)x+1$.

\begin{enumerate}[label=(\arabic*)]
\item 已知对于任意实数 $x$，$\sqrt{f(x)}$ 均有意义，求实数 $a$ 的取值范围；
\item 设 $A=\{x\mid f(x)+a^2-2=0\}$，$B=\{-4,0\}$，$A\cup B=B$，求实数 $a$ 的取值范围.
\end{enumerate}

\ifanswer
\solbegin
\stepm{(1)}{$\sqrt{f(x)}$ 对任意实数 $x$ 都有意义，}
\step{即 $f(x)=x^2+2(a+1)x+1\geqslant0$ 对任意 $x\in\R$ 恒成立，}
\step{所以 $\Delta=4(a+1)^2-4\times1\times1\leqslant0$，即 $(a+1)^2\leqslant1$，解得 $-2\leqslant a\leqslant0$，}
\step{故 $a$ 的取值范围是 $[-2,0]$.}
\stepm{(2)}{由 $f(x)+a^2-2=0$ 得 $x^2+2(a+1)x+a^2-1=0$.}
\step{由 $A\cup B=B$ 得 $A\subseteq B=\{-4,0\}$，即该方程无实根，}
\step{或它的实根都在 $\{-4,0\}$ 中.}
\step{当方程无实根时，$\Delta=4(a+1)^2-4(a^2-1)=8a+8<0$，解得 $a<-1$；}
\step{当方程有实根时，$a\geqslant-1$，设两根为 $x_1$、$x_2$，}
\step{由根与系数的关系得 $x_1+x_2=-2(a+1)$、$x_1x_2=a^2-1$，}
\step{此时两根只能在 $\{-4,0\}$ 中取值；}
\step{若两根为 0 与 0，由 $x_1x_2=0$ 得 $a=\pm1$，}
\step{再由 $x_1+x_2=0$ 得 $a=-1$；}
\step{若两根为 0 与 $-4$，由 $x_1x_2=0$ 得 $a=\pm1$，再由 $x_1+x_2=-4$ 得 $a=1$.}
\step{经检验，$a=-1$ 时 $A=\{0\}$，$a=1$ 时 $A=\{-4,0\}$，都满足 $A\subseteq B$.}
\step{综上，实数 $a$ 的取值范围为 $(-\infty,-1]\cup\{1\}$.}
\solend
\fi

\ansspace[8]

\ansneed{18cm}
\item 共享单车给市民出行带来了诸多便利，某公司购买了一批单车投放到某地给居民使用，据市场分析，这批单车的运营累计利润 $y$（单位：元）与运营天数 $x$（单位：天，$x$ 为正整数）满足关系：$y=-\dfrac12x^2+60x-800$.

\begin{enumerate}[label=(\arabic*)]
\item 要使这批单车的运营累计利润高于 800 元，求运营天数的取值范围；
\item 判断这批单车的日平均运营利润能否超过 20 元，并说明理由.
\end{enumerate}

\ifanswer
\solbegin
\stepm{(1)}{由题意得 $y=-\dfrac12x^2+60x-800>800$，}
\step{即 $x^2-120x+3200<0$，解得 $40<x<80$，}
\step{即运营天数的取值范围为 $\{x\mid40<x<80,x\in\N^*\}$.}
\stepm{(2)}{不能，理由如下：}
\step{这批单车的日平均运营利润为 $\dfrac yx=-\dfrac12x-\dfrac{800}x+60=-\dfrac12\left(x+\dfrac{1600}x\right)+60$，}
\step{因 $x>0$，则 $x+\dfrac{1600}x\geqslant2\sqrt{x\cdot\dfrac{1600}x}=80$，}
\step{当且仅当 $x=\dfrac{1600}x$，即 $x=40$ 时取等号，}
\step{此时 $\dfrac yx=-\dfrac12\left(x+\dfrac{1600}x\right)+60\leqslant-\dfrac12\times80+60=20$，}
\step{故这批单车的日平均运营利润不能超过 20 元.}
\solend
\fi

\ansspace[8]

% 第 20 题题干＋三小问与作答区 8 行须同页
\ansneed{20cm}
\item 设 $f(x)=ax^2+bx+c$($a\neq0$，$b$、$c\in\R$).

\begin{enumerate}[label=(\arabic*)]
\item 若 $a=b=c=1$，设 $A=\{(x,y)\mid y=f(x)\}$，$B=\{(x,y)\mid y=1\}$，用列举法表示 $A\cap B$；
% 原卷两个“若”连用（“若 a=1，…，若方程 f(x)=0 有两个大于 1 的根”），疑衍其一；照录（见文件头“原卷笔误/存疑”⑦）
\item 若 $a=1$，$b=2k-1$，$c=k^2$，若方程 $f(x)=0$ 有两个大于 1 的根，求实数 $k$ 的取值范围；
\item 若 $a$、$b$、$c\in\Z$，且 $f(0)$、$f(1)$ 均为奇数，求证：方程 $f(x)=0$ 无整数解.
\end{enumerate}

\ifanswer
\solbegin
\stepm{(1)}{当 $a=b=c=1$ 时，$f(x)=x^2+x+1$，}
\step{所以 $A=\{(x,y)\mid y=x^2+x+1\}$，又 $y=1$，}
\step{由 $x^2+x+1=1$ 得 $x^2+x=0$，解得 $x=0$ 或 $x=-1$，}
\step{所以 $A\cap B=\{(0,1),(-1,1)\}$.}
\stepm{(2)}{当 $a=1$、$b=2k-1$、$c=k^2$ 时，$f(x)=x^2+(2k-1)x+k^2$.}
\step{设两根为 $x_1$、$x_2$，}
\step{则“两根都大于 1”等价于 $\Delta>0$、$f(1)>0$、$-\dfrac b{2a}>1$.}
\step{由 $\Delta=(2k-1)^2-4k^2=1-4k>0$ 得 $k<\dfrac14$；}
\step{由 $f(1)=1+(2k-1)+k^2=k^2+2k=k(k+2)>0$ 得 $k<-2$ 或 $k>0$；}
\step{由 $-\dfrac{2k-1}2=\dfrac{1-2k}2>1$ 得 $k<-\dfrac12$.}
\step{三者取交集得 $k<-2$，即 $k$ 的取值范围是 $(-\infty,-2)$.}
\stepm{(3)}{设 $x_0\in\Z$ 是方程 $f(x)=0$ 的整数解.}
\step{因为 $x_0(x_0-1)$ 是两个相邻整数之积，恒为偶数，所以 $ax_0(x_0-1)$ 为偶数.}
\step{而 $f(x_0)=ax_0^2+bx_0+c=ax_0(x_0-1)+(a+b)x_0+c$，}
\step{所以 $f(x_0)$ 与 $(a+b)x_0+c$ 同奇偶.}
\step{若 $x_0$ 为偶数，则 $(a+b)x_0$ 为偶数，$f(x_0)$ 与 $c=f(0)$ 同奇偶，为奇数；}
\step{若 $x_0$ 为奇数，则 $(a+b)x_0$ 与 $a+b$ 同奇偶，}
\step{$f(x_0)$ 与 $(a+b)+c=f(1)$ 同奇偶，也为奇数.}
\step{综上，$f(x_0)$ 恒为奇数，不可能等于 0，与 $f(x_0)=0$ 矛盾，}
\step{所以方程 $f(x)=0$ 无整数解.}
\solend
\fi

\ansspace[8]

% 压轴题（题干＋三小问）与作答区 10 行须同页
\ansneed{20cm}
\item 已知 $X$ 为有限集，且任取 $x\in X$，均有 $x\in\Z$，称 $X$ 为非空有限整数集. 定义：$X_k=\{x+k\mid x\in X\}$，$X^2=\{y\mid y=x^2,x\in X\}$.

\begin{enumerate}[label=(\arabic*)]
\item 若 $X=\{-2,2\}$，分别写出 $X_2$ 和 $X^2$（用列举法表示）；
\item 若 $X=\{-1,0,1\}$，$X_1\subseteq Y$，$(Y_{-1})^2\subseteq X$，求集合 $Y$ 中所有的元素；
\item 对于非空有限整数集 $A$、$B$，若 $A_k\subseteq B$，$(B^2)_{-2}\subseteq A$，$B$ 中的元素均为非负整数，$k$ 为整数，求证：

① $A$、$B$ 中元素的个数相同；

② 若集合 $A$ 中的元素个数为奇数，则 $k$ 为偶数.
\end{enumerate}

\ifanswer
\solbegin
\stepm{(1)}{由题意得 $X_2=\{0,4\}$，$X^2=\{4\}$；}
\stepm{(2)}{由题意得 $X_1=\{0,1,2\}$，即 $\{0,1,2\}\subseteq Y$，$Y_{-1}=\{y-1\mid y\in Y\}$，}
\step{$(Y_{-1})^2=\{(y-1)^2\mid y\in Y\}\subseteq X=\{-1,0,1\}$，}
\step{所以 $(y-1)^2=0$ 或 $(y-1)^2=1$，解得 $y=1$ 或 $y=2$ 或 $y=0$，}
\step{即 $y\in\{0,1,2\}$，即 $Y\subseteq\{0,1,2\}$，所以 $Y=\{0,1,2\}$；}
% 原卷 (3) 解析首行“(3)”与“①”连排，本稿合并记作 \stepm{(3)①}
\stepm{(3)①}{因为 $A_k=\{x+k\mid x\in A\}$，函数 $g(x)=x+k$ 在全体整数集上严格增，}
\step{所以 $\forall x_1$、$x_2\in A$，$x_1\neq x_2$，$g(x_1)\neq g(x_2)$，所以 $|A_k|=|A|$，}
\step{因为 $A_k\subseteq B$，所以 $|A|\leqslant|B|$，}
\step{设 $b\in B$，$B$ 中元素均为非负整数，$\forall b_1\neq b_2\in B$，$b_1^2\neq b_2^2$，}
\step{故函数 $f(b)=b^2-2$ 在 $B$ 上单射，所以 $|(B^2)_{-2}|=|B|$，}
\step{又因为 $(B^2)_{-2}\subseteq A$，所以 $|B|\leqslant|A|$，所以 $|A|=|B|$.}
\stepm{②}{由①得 $A_k\subseteq B$ 且 $|A_k|=|A|=|B|$，$(B^2)_{-2}\subseteq A$ 且 $|(B^2)_{-2}|=|B|=|A|$，}
\step{所以 $A_k=B$，$(B^2)_{-2}=A$，}
\step{所以 $B=A_k=\{a+k\mid a\in A\}$，$A=(B^2)_{-2}=\{b^2-2\mid b\in B\}$，}
\step{所以 $A=\{(a+k)^2-2\mid a\in A\}$，}
\step{所以映射 $h(a)=(a+k)^2-2$ 是集合 $A$ 到自身的一个双射.}
\step{假设 $k$ 是奇数，当 $a$ 为偶数时，$a+k$ 为奇数，$(a+k)^2-2$ 为奇数，}
\step{当 $a$ 为奇数时，$a+k$ 为偶数，$(a+k)^2-2$ 为偶数，}
\step{所以此时映射 $h(a)=(a+k)^2-2$，将集合 $A$ 中的奇数映射为偶数，}
\step{将集合 $A$ 中的偶数映射为奇数，即集合 $A$ 中的奇数元素和偶数元素必须配对，}
\step{此时集合 $A$ 的元素个数则为偶数，}
\step{与题设中“集合 $A$ 中的元素个数为奇数”矛盾.}
\step{故假设不成立，所以 $k$ 是偶数.}
\solend
\fi

\ansspace*[10]

\end{enumerate}

\end{document}
"
% 紧凑版：xelatex "\def\iscompact{1}% 高一67_华师大一附中_10月月考.tex
%   —— 高一上数学“2026-2027 年华师大一附中高一上 10 月月考”（讲评版）
% 试卷版：xelatex 高一67_华师大一附中_10月月考.tex
% 答案版：xelatex "\def\isanswer{1}% 高一67_华师大一附中_10月月考.tex
%   —— 高一上数学“2026-2027 年华师大一附中高一上 10 月月考”（讲评版）
% 试卷版：xelatex 高一67_华师大一附中_10月月考.tex
% 答案版：xelatex "\def\isanswer{1}\input{高一67_华师大一附中_10月月考.tex}"
% 紧凑版：xelatex "\def\iscompact{1}\input{高一67_华师大一附中_10月月考.tex}"
%
% 原始材料：微信公众号“魔都数学加油站”文章，作者破晓，连载“27学年高一67”；
% 原文标题“27学年高一67——2026-2027年上海市华师大一附中高一上10月月考”；
% 原文链接：https://mp.weixin.qq.com/s/Kdhx07TVjLoGTcIlnRzgcQ。
% 页面标注发布时间 2026-10-18 20:33（上海）。
% 本次转录只据 `原图/高一67_华师大一附中_10月月考/` 的 13 张 PNG 逐页读图。
% 正文为 13 张 PNG 页（`原图/高一67_华师大一附中_10月月考/01..13.png`）：
% 第 01、04、13 页 2560×3620，第 02、03、05—12 页 4762×6735
% （**同一篇各页分辨率不同**），均按页序读取。13 页全为试卷正文页（卷首标题在第 1 页页顶，
% 卷末解答题第 21 题解析止于第 13 页），无封面页、目录页、二维码或宣传页。
% 粉色斜排“魔都数学加油站”为卷面水印，与公众号页脚一并未录入。
%
% 卷首标题：照卷面印作“2026--2027 年华师大一附中高一上 10 月月考”（卷面为黑体加粗一行），
% 年份区间按全稿惯例排 en dash。原卷只有一行标题、无副标题。
%
% 篇幅与题号：原卷自第 3 题起编号（第 1、2 题不在本卷内）。填空题 3—12（10 题）、
% 选择题 13—16（4 题）、解答题 17—21（5 题），共 19 题。本稿沿用原节与原题号，
% 只改排为全稿统一的大节样式（原卷小节标题为“一、填空题”“二、选择题”“三、解答题”）。
%
% 讲评版：原卷每题下直接印【解析】。**全卷 19 题中 18 题都印【解析】（无一处只印
% 【答案】）**，按全稿口径这类题不另提 \ans{}；各题解析末句的结论（第 13、14、16 题
% 的“故选 B／D／B”、各填空与解答题的“故/综上…”）照录在解析内。
% **另有 1 题答案印在题干括号内**——第 15 题（“(  D  )”，且该题无【答案】/【解析】段），
% 本稿按 高一37 第 13 题、高一63 第 14 题成例把括号留空、答案另提为 \ans{$D$}。
%
% 副标题：原卷未印副标题。本稿按“知识模块”口径标作“集合与常用逻辑用语、不等式”：
% 19 题中集合与常用逻辑用语 12 题（第 3、4、5、7、10、11、12、13、15、16、18、21 题：
% 补集、并集、子集、充分条件、容斥计数、子集计数、新定义二元运算、条件判断、
% 反证法假设、新定义生成集、集合关系与二次方程、新定义 X_k 与 X^2），
% 不等式 7 题（第 6、8、9、14、17、19、20 题：存在量词不等式、解集反求系数、
% 解集为空、不等式性质、解不等式（组）、二次不等式与基本不等式、二次方程根的分布），
% 两类均超过三分之一，故并标。
%
% 与原卷的排版差异：
%   ① 原卷每题下直接印【解析】，本稿按全稿惯例将解析置于答案版，试卷版保持干净；
%   ② 原卷未印副标题，本稿按“知识模块”口径补出；
%   ③ 原卷填空横线约四个字宽，本稿按全稿惯例统一排 \blankline[4.4em]；
%   ④ 原卷实数集排斜体 $R$（第 14 题题干、第 17 题 (1) 题干、第 18 题解析、第 20 题题干），
%      整数集排斜体 $Z$（第 20 题 (3) 题干与解析、第 21 题题干），正整数集排粗体
%      $N^*$（第 11 题解析、第 19 题解析），本稿按全稿惯例统一改用 $\R$、$\Z$、$\N^*$
%      （记号归一，粗体／斜体不作区分）；空集原卷作 $\varnothing$，照录；
%   ⑤ 原卷在数学式内多处用半角逗号（如第 8 题 $(1,2)$ 与各 cases、第 11 题解析
%      $P_n=\{1,2,\cdots,n\}$,$n\in N^*$、第 12 题解析 $u=(\cdots)$,$v=(\cdots)$ 与
%      “2020,个 1”、第 19 题解析 $\{x|40<x<80,x\in N^*\}$），均照录未强改；
%      描述法集合的竖线照原卷排 \mid；句末句点按全稿惯例排西文句点；
%   ⑥ 原卷补集记号为上划线 $\bar A$（第 3 题题干与解析、第 11 题题干），照录
%      （同 高一64/65 的 $\bar B$ 成例）；
%   ⑦ 第 12 题新定义的运算记号 $*$ 照录（$\alpha*\beta$、$u*u$、$u*v$）；
%   ⑧ 第 17—21 题的子问原卷各自另起一行，本稿按全稿惯例用 enumerate 的 (1)(2)(3)；
%      解析中的子问标号一律用 \stepm{(1)}{…}；第 5 题解析的①②、第 21 题 (3) 解析的①②
%      按全稿惯例作 \stepm{①}{…}、\stepm{②}{…}（第 21 题 (3) 首行原卷“(3)”与“①”连排，
%      本稿合并记作 \stepm{(3)①}{…}）；
%   ⑨ 选择题作答括号原卷排半角“(　　)”（第 13、14、16 题括号内空白、第 15 题括号内印 D），
%      本稿按全稿惯例排 \mbox{（\quad）}；第 15 题括号内的答案移入 \ans{}，见上；
%   ⑩ 选项行排照原卷：第 13、15 题两行两条，第 14、16 题一行四条；
%   ⑪ 每道解答题原卷未留作答空白，本稿按全稿惯例加 \ansspace（末题用 \ansspace*）；
%      第 20、21 题题干长，各用一处 \ansneed 让题干与作答区不拆页（\ansneed 只作用于试卷版）。
%
% 原卷笔误/存疑（均照录未改，正文对应处另有 % 注）：
%   ① 第 3 题解析末句印作“则 $\bar A=(-\infty,0)\cup[2,3)$.”：按题干
%      $U=\{x\mid 1/x\leqslant 1/2\}$（即 $x<0$ 或 $x\geqslant2$）与 $A=[3,+\infty)$，
%      $\bar A$ 应为 $(-\infty,0)\cup\{2\}$；且与本解析首句自求之 $U$ 相矛盾（$3\notin U$），
%      疑为笔误；照录未改。
%   ② 第 12 题题干 $\alpha*\beta$ 定义式括号不配对：第四组“$+\cdots+x_{2020}+y_{2020}
%      -x_{2020}y_{2020}$”缺左括号、末尾连排两个右括号“))”；照录未改。
%   ③ 第 12 题解析“所以…中有 2020,个 1”一行，“$\cdots$”后脱逗号（作
%      “$x_2,\cdots x_{2020}$”，第 12 行“而…”同式样），且“2020”后衍一半角逗号
%      （作“2020,个 1”）；照录未改。
%   ④ 第 16 题解析乘积链印作“$a_1a_2<a_1a_3<a_1a_4<a_1a_5<a_2a_3<a_3a_3<a_{14}a_5$”：
%      “$a_3a_3$”同足标自乘、与单调链不合，“$a_{14}a_5$”下标 14 与 $a_1$—$a_5$ 记号
%      不合（疑为 $a_2a_4$、$a_4a_5$ 之误）；照录未改。
%   ⑤ 第 5 题题干作“$Q\subset P$”（真子集记号），解析两处均作“$Q\subseteq P$”，
%      同一题前后记号不一；照录（本稿分别排 \subset、\subseteq）。
%   ⑥ 第 17 题 (1) 解析“当 $a=1$ 时，原不等式为 $(x+1)^2<0$ 此时…”与
%      “当 $0<a<1$ 时，$-\frac1a<-1$ 此时…”两处“此时”前均脱标点；照录未改。
%   ⑦ 第 20 题 (2) 题干“若 $a=1$，$b=2k-1$，$c=k^2$，若方程 $f(x)=0$ 有两个大于 1 的根”
%      两个“若”连用，疑衍其一；照录未改。
%
% 插图：无（全卷 13 页无任何几何图形、无“如图/右图/图中”字样）。
%
% 行数核对：逐题数 \step/\stepm 条数与原卷【解析】印刷行数，18 题逐一相符
% （第 15 题无解析段，只提 \ans{}），合计 169 条 / 169 行
% （Q3—Q12：2、3、4、3、1、3、5、3、13、15；Q13—Q16：2、3、0、25；
% Q17—Q21：19、16、9、20、23）。
%
\documentclass[a4paper,11pt]{article}
\usepackage{common}

\begin{document}

\examtitlest[3]{2026--2027 年华师大一附中高一上 10 月月考}{集合与常用逻辑用语、不等式}

\bigsec{一、填空题}

\begin{enumerate}
\setcounter{enumi}{2}

\item 设全集 $U=\left\{x\mid\dfrac1x\leqslant\dfrac12\right\}$，$A=[3,+\infty)$，则 $\bar A=$\blankline[4.4em].

\ifanswer
\solbegin
\step{$\dfrac1x\leqslant\dfrac12$，解得 $x<0$ 或 $x\geqslant2$，故 $U=(-\infty,0)\cup[2,+\infty)$，}
% 原卷末句作 $\bar A=(-\infty,0)\cup[2,3)$，按题干 $U$ 应为 $(-\infty,0)\cup\{2\}$，与本解析首句矛盾，疑笔误；照录（见文件头“原卷笔误/存疑”①）
\step{$A=[3,+\infty)$，则 $\bar A=(-\infty,0)\cup[2,3)$.}
\solend
\fi

\item 已知集合 $M=\{y\mid y=x^2-2x-1\}$，$N=\{y\mid y=x^2-4x+1\}$，则 $M\cup N=$\blankline[4.4em].

\ifanswer
\solbegin
\step{$M=\{y\mid y=x^2-2x-1\}=\{y\mid y=(x-1)^2-2\}=\{y\mid y\geqslant-2\}$，}
\step{$N=\{y\mid y=x^2-4x+1\}=\{y\mid y=(x-2)^2-3\}=\{y\mid y\geqslant-3\}$，}
\step{则 $M\cup N=\{y\mid y\geqslant-3\}$.}
\solend
\fi

\item 已知 $P=\{x\mid x^2+x-6=0\}$，$Q=\{x\mid mx-1=0\}$，且 $Q\subset P$，则实数 $m$ 的值\blankline[4.4em].

\ifanswer
\solbegin
\step{$P=\{2,-3\}$，$Q=\{x\mid mx=1\}$；}
\stepm{①}{$m=0$ 时，$Q=\varnothing$，满足 $Q\subseteq P$；}
\stepm{②}{$m\neq0$ 时，要使 $Q\subseteq P$，则 $\dfrac1m=2$ 或 $-3$，所以 $m=\dfrac12$ 或 $-\dfrac13$；}
\step{所以 $m=0$ 或 $\dfrac12$ 或 $-\dfrac13$.}
\solend
\fi

\item 已知“存在正实数 $x$，使得关于 $x$ 的不等式 $ax-1>0$ 成立”为真命题，则实数 $a$ 的取值范围是\blankline[4.4em].

\ifanswer
\solbegin
\step{存在正实数 $x$，使得关于 $x$ 的不等式 $ax-1>0$ 成立为真命题，}
\step{即存在正实数 $x$，使得 $a>\dfrac1x$，因为 $\dfrac1x>0$，所以 $a>0$，}
\step{即 $a$ 的范围为 $(0,+\infty)$.}
\solend
\fi

\item 设 $\alpha:3\leqslant x\leqslant5$，$\beta:2m-1\leqslant x<m+4$，若 $\alpha$ 是 $\beta$ 的充分条件，则实数 $m$ 的取值范围是\blankline[4.4em].

\ifanswer
\solbegin
\step{若 $\alpha$ 是 $\beta$ 的充分条件，则 $\begin{cases}2m-1\leqslant3\\ m+4>5\end{cases}$，解得 $1<m\leqslant2$.}
\solend
\fi

\item 对于实数 $a$、$b$、$c$，已知关于 $x$ 的不等式 $ax^2+bx+c>0$ 的解集是 $(1,2)$，则关于 $x$ 的不等式 $cx^2-bx+a\leqslant0$ 的解集是\blankline[4.4em].

\ifanswer
\solbegin
\step{关于 $x$ 的不等式 $ax^2+bx+c>0$ 的解集是 $(1,2)$，则 $\begin{cases}a<0\\ -\dfrac ba=3\\ \dfrac ca=2\end{cases}$，则 $\begin{cases}b=-3a\\ c=2a\end{cases}$，}
\step{则关于 $x$ 的不等式 $cx^2-bx+a=2ax^2+3ax+a\leqslant0$，则 $2x^2+3x+1\geqslant0$，}
\step{则 $x\leqslant-1$ 或 $x\geqslant-\dfrac12$，则解集为 $(-\infty,-1]\cup\left[-\dfrac12,+\infty\right)$.}
\solend
\fi

\item 已知关于 $x$ 的不等式 $kx^2+kx+k+2\leqslant0$ 的解集为 $\varnothing$，则实数 $k$ 的取值范围为\blankline[4.4em].

\ifanswer
\solbegin
\step{当 $k=0$ 时，不等式化为 $2\leqslant0$，解集为 $\varnothing$，符合题意；}
\step{当 $k\neq0$ 时，若关于 $x$ 的不等式 $kx^2+kx+k+2\leqslant0$ 的解集为 $\varnothing$，}
\step{则 $\begin{cases}k>0\\ \Delta=k^2-4k(k+2)<0\end{cases}$，化简得 $\begin{cases}k>0\\ -3k^2-8k<0\end{cases}$，即 $\begin{cases}k>0\\ k(3k+8)>0\end{cases}$，}
\step{解得 $k>0$.}
\step{综上所述，$k\geqslant0$，即实数 $k$ 的取值范围为 $[0,+\infty)$.}
\solend
\fi

\item 对班级若干名学生调查对 $A$、$B$ 两事件的态度，有如下结果：赞成 $A$ 的人数是全体的五分之三，其余的不赞成；赞成 $B$ 的比赞成 $A$ 多 3 人，其余的不赞成. 已知对 $A$、$B$ 都赞成的学生有 18 人，对 $A$、$B$ 都不赞成的学生数比对 $A$、$B$ 都赞成的学生数的三分之一多 1 人，则班级总人数是\blankline[4.4em]人.

\ifanswer
\solbegin
\step{设班级总人数是 $x$ 人，则赞成 $A$ 的人数是 $\dfrac35x$，赞成 $B$ 的人数是 $\dfrac35x+3$，}
\step{$A$、$B$ 都不赞成的人数为 $18\times\dfrac13+1=7$ 人，}
\step{由容斥原理得 $x=\left(\dfrac35x+\dfrac35x+3-18\right)+7$，解得 $x=40$，即班级总人数是 40 人.}
\solend
\fi

\item 设全集 $U=\{1,2,3,4,5,6,7,8,9,10,11,12\}$，若集合 $A$ 满足：① $A\subseteq U$；②若 $x\in A$，必有 $2x\notin A$；③若 $x\in\bar A$，必有 $2x\notin\bar A$，则这样的集合 $A$ 有\blankline[4.4em]个.

\ifanswer
\solbegin
\step{设集合 $P_n=\{1,2,\cdots,n\}$,$n\in\N^*$，}
\step{任取 $x\in P_n$，将 $x$ 除以 2，若商仍为偶数，再除以 2，$\cdots$，}
\step{经过 $k$ 次以后，商必为奇数，此时记商为 $m$，}
\step{于是 $x=m\cdot2^k$，其中 $m$ 为奇数，$k\in\N^*$.}
\step{由条件得，若 $m\in A$，则 $x\in A$ 等价于 $k$ 为偶数；}
\step{若 $m\notin A$，则 $x\in A$ 等价于 $k$ 为奇数；}
\step{于是 $x$ 是否属于 $A$ 由 $m$ 是否属于 $A$ 确定；}
\step{设 $Q_n$ 是 $P_n$ 中所有奇数的集合，}
\step{设 $f(n)$ 为集合 $P_n$ 中满足题干条件的集合 $A$ 的个数，}
\step{因此 $f(n)$ 等于 $Q_n$ 的子集个数.}
\step{当 $n$ 为偶数（或奇数）时，$P_n$ 中奇数的个数是 $\dfrac n2$（或 $\dfrac{n+1}2$），}
\step{所以 $f(n)=\begin{cases}2^{\frac n2},n\text{为偶数}\\[2pt] 2^{\frac{n+1}2},n\text{为奇数}\end{cases}$，所以 $P_{12}=\{1,2,3,4,5,6,7,8,9,10,11,12\}$，}
\step{即同时满足三个条件的集合 $A$ 的个数为 $2^{\frac{12}2}=2^6=64$.}
\solend
\fi

\item 设集合 $\Omega=\{(x_1,x_2,\cdots,x_{2020})\mid x_i\in\{0,1\},i=1,2,\cdots,2020\}$，对于 $\Omega$ 中的任意两个元素 $\alpha=(x_1,x_2,\cdots,x_{2020})$，$\beta=(y_1,y_2,\cdots,y_{2020})$，记
% 原卷此式第四组缺左括号、末尾连排两个右括号“))”，括号不配对；照录（见文件头“原卷笔误/存疑”②）
$\alpha*\beta=(x_1+y_1-x_1y_1)+(x_2+y_2-x_2y_2)+\cdots+x_{2020}+y_{2020}-x_{2020}y_{2020}))$，设 $u$、$v\in\Omega$，若 $u*u+v*v=2020$，则 $u*v$ 的最大值和最小值之和为\blankline[4.4em].

\ifanswer
\solbegin
\step{设 $u=(x_1,x_2,\cdots,x_{2020})$,$v=(y_1,y_2,\cdots,y_{2020})$，}
\step{因为 $x_i\in\{0,1\}$，所以 $x_i(x_i-1)=0$，}
\step{所以 $u*u=(x_1+x_1-x_1^2)+(x_2+x_2-x_2^2)+\cdots+(x_{2020}+x_{2020}-x_{2020}^2)$}
\step{$=x_1+x_2+\cdots+x_{2020}$，}
\step{$v*v=(y_1+y_1-y_1^2)+(y_2+y_2-y_2^2)+\cdots+(y_{2020}+y_{2020}-y_{2020}^2)$}
\step{$=y_1+y_2+\cdots+y_{2020}$，}
\step{由 $u*u+v*v=2020$，得 $x_1+x_2+\cdots+x_{2020}+y_1+y_2+\cdots+y_{2020}=2020$，}
% 原卷此行“$\cdots$”后脱逗号、 “2020”后衍一半角逗号（“2020,个 1”）；照录（见文件头“原卷笔误/存疑”③）
\step{所以 $x_1,x_2,\cdots x_{2020},y_1,y_2,\cdots y_{2020}$ 中有 2020,个 1，2020 个 0，}
\step{$u*v=(x_1+y_1-x_1y_1)+(x_2+y_2-x_2y_2)+\cdots+(x_{2020}+y_{2020}-x_{2020}y_{2020})$}
\step{$=2020-(x_1y_1+x_2y_2+\cdots+x_{2020}y_{2020})$，}
\step{注意到只有 $x_i=y_i=1$ 时，$x_iy_i=1$，否则 $x_iy_i=0$，}
% 此行“$\cdots$”后同样脱逗号，与上一处同式样；照录（见文件头“原卷笔误/存疑”③）
\step{而 $x_1,x_2,\cdots x_{2020},y_1,y_2,\cdots y_{2020}$ 中有 2020 个 1，2020 个 0，}
\step{所以满足 $x_iy_i=1$ 的最多有 1010 个，所以 $u*v\geqslant2020-1010=1010$，}
\step{满足 $x_iy_i=0$ 的最多有 2020 个，所以 $u*v\leqslant2020-0=2020$，}
\step{则 $u*v$ 的最大值和最小值之和为 $2020+1010=3030$.}
\solend
\fi

\end{enumerate}

\bigsec{二、选择题}

\begin{enumerate}
\setcounter{enumi}{12}

\item 实数 $x$ 与 $y$，“$x+y>2$”的\mbox{（\quad）}条件是“$x>1$ 或 $y>1$”.

\keepwithprev
A. 充分非必要\qquad B. 必要非充分

C. 充要\qquad D. 既非充分又非必要

\ifanswer
\solbegin
\step{当 $x>1$ 或 $y>1$ 时，$x+y>2$ 不一定成立，充分性不成立；}
\step{当 $x+y>2$ 时，$x>1$ 或 $y>1$ 一定成立，必要性成立. 故选 $B$.}
\solend
\fi

\item 设 $a,b\in\R$，且 $a<0$，$-1<b<0$，那么下列不等式成立的是\mbox{（\quad）}

\keepwithprev
A. $a>ab>ab^2$\qquad B. $ab^2>ab>a$\qquad C. $ab>a>ab^2$\qquad D. $ab>ab^2>a$

\ifanswer
\solbegin
\step{因为 $-1<b<0$，$a<0$，所以 $ab>0$，$0<b^2<1$.}
\step{所以 $ab-ab^2=ab(1-b)>0$，$ab^2-a=a(b^2-1)>0$，}
\step{所以 $ab>ab^2>a$. 故选 $D$.}
\solend
\fi

% 第 15 题答案原卷印在题干括号内（“(  D  )”），且无【答案】/【解析】段；
% 按 高一37 第 13 题、高一63 第 14 题成例括号留空、答案另提 \ans{$D$}
\item 用反证法证明：“平面四边形中至少有一个内角不超过 90°”时，下列“假设”中正确的是\mbox{（\quad）}

\keepwithprev
A. 假设有两个内角超过 90°\qquad B. 假设有三个内角超过 90°

C. 假设至多有两个内角超过 90°\qquad D. 假设四个内角都超过 90°

\ifanswer
\ans{$D$}
\fi

\item 设 $A$ 为非空实数集，集合 $B=\{xy\mid x\in A,y\in A,x\neq y\}$，则对于以下两个结论说法正确的是\mbox{（\quad）}

结论①：若 $A$ 是五元集，则 $B$ 中的元素最多 10 个；

结论②：存在四元集 $A$，使得 $B=\{2,4,5,8,10,16\}$.

\keepwithprev
A. ①对②对\qquad B. ①对②错\qquad C. ①错②对\qquad D. ①错②错

\ifanswer
\solbegin
\step{对于①，想让元素个数达到最大值 10，则须避免产生符号不同的乘积，}
\step{否则正负分开计数容易导致总数受限，且容易产生相等值，}
\step{也最好避免 0 的出现，}
\step{故设 $A=\{a_1,a_2,a_3,a_4,a_5\}$，$0<a_1<a_2<a_3<a_4<a_5$.}
% 原卷此链“$a_3a_3$”同足标自乘、“$a_{14}a_5$”下标 14 均与文意不合，疑为 $a_2a_4$、$a_4a_5$ 之误；照录（见文件头“原卷笔误/存疑”④）
\step{因为 $a_1a_2<a_1a_3<a_1a_4<a_1a_5<a_2a_3<a_3a_3<a_{14}a_5$，$a_2a_3<a_2a_4<a_3a_4$，}
\step{所以 $B$ 中元素个数小于等于 10 个，}
\step{如设 $A=\{1,2,3,5,7\}$，则 $B=\{2,3,5,6,7,10,14,15,21,35\}$，}
\step{所以集合 $B$ 中元素个数的最大值为 10 个，故①正确；}
\step{对于②，由于集合 $B$ 中没有 0，所以集合 $A$ 中没有 0；}
\step{设集合 $A=\{a,b,c,d\}$，且 $a<b<c<d$，}
\step{当 $a<0<b<c<d$ 时，$a\times b<0$，$a\times c<0$，$a\times d<0$，}
\step{故集合 $B$ 至少有 3 个负数，与 $B=\{2,4,5,8,10,16\}$ 矛盾；}
\step{当 $a<b<0<c<d$ 时，$a\times b>0$，$c\times d>0$，}
\step{$a\times c<0$，$a\times d<0$，$b\times c<0$，$b\times d<0$，}
\step{故集合 $B$ 至少有 4 个负数，与 $B=\{2,4,5,8,10,16\}$ 矛盾；}
\step{当 $a<b<c<0<d$ 时，$a\times b>0$，$a\times c>0$，$b\times c>0$，}
\step{$a\times d<0$，$b\times d<0$，$c\times d<0$，}
\step{故集合 $B$ 至少有 3 个负数，与 $B=\{2,4,5,8,10,16\}$ 矛盾；}
\step{当 $0<a<b<c<d$ 时，则集合 $B=\{ab,ac,ad,bc,bd,cd\}$，}
\step{则必有 $ab=2$，$cd=16$，其 4 个正实数的乘积 $abcd=32$，}
\step{又 $ad=5$，$bc=8$ 或 $ad=8$，$bc=5$，其 4 个正实数的乘积 $abcd=40$，矛盾；}
\step{当 $a<b<c<d<0$ 时，与 $0<a<b<c<d$ 时的证明方法相同，}
\step{得集合 $A$ 也不存在；}
\step{综上不存在四元集 $A$，使其生成集 $B=\{2,4,5,8,10,16\}$，故②错误.}
\step{故选 $B$.}
\solend
\fi

\end{enumerate}

\bigsec{三、解答题}

\begin{enumerate}
\setcounter{enumi}{16}

\item 解下列不等式（组）

\begin{enumerate}[label=(\arabic*)]
\item 解不等式：$(ax+1)(x+1)<0(a\in\R)$.
\item 解不等式组：$\begin{cases}\dfrac{2x+4}{x+4}\geqslant1\\[4pt] (x+1)(2-|x|)\geqslant0\end{cases}$.
\end{enumerate}

\ifanswer
\solbegin
\stepm{(1)}{当 $a=0$ 时，原不等式为 $x+1<0$，即 $x<-1$，}
\step{所以此时不等式的解集为 $\{x\mid x<-1\}$；}
\step{当 $a\neq0$ 时，令 $(ax+1)(x+1)=0$，解得 $x_1=-\dfrac1a$,$x_2=-1$，}
\step{当 $a>1$ 时，$-\dfrac1a>-1$，此时不等式的解集为 $\left\{x\mid-1<x<-\dfrac1a\right\}$；}
% 以下两行“此时”前原卷均无标点；照录（见文件头“原卷笔误/存疑”⑥）
\step{当 $a=1$ 时，原不等式为 $(x+1)^2<0$ 此时不等式的解集为 $\varnothing$；}
\step{当 $0<a<1$ 时，$-\dfrac1a<-1$ 此时不等式的解集为 $\left\{x\mid-\dfrac1a<x<-1\right\}$；}
\step{当 $a<0$ 时，$-\dfrac1a>-1$，此时不等式的解集为 $\left\{x\mid x<-1\ \text{或}\ x>-\dfrac1a\right\}$.}
\step{综上所述，当 $a=0$ 时，$\{x\mid x<-1\}$；}
\step{当 $a>1$ 时，$\left\{x\mid-1<x<-\dfrac1a\right\}$；}
\step{当 $a=1$ 时，$\varnothing$；}
\step{当 $0<a<1$ 时，$\left\{x\mid-\dfrac1a<x<-1\right\}$；}
\step{当 $a<0$ 时，$\left\{x\mid x<-1\ \text{或}\ x>-\dfrac1a\right\}$.}
\stepm{(2)}{不等式 $\dfrac{2x+4}{x+4}\geqslant1$ 可化为 $\dfrac{2x+4}{x+4}-1\geqslant0$，即 $\dfrac x{x+4}\geqslant0$，}
\step{等价于 $\begin{cases}x(x+4)\geqslant0\\ x+4\neq0\end{cases}$，解得 $x<-4$ 或 $x\geqslant0$，}
\step{不等式 $(x+1)(2-|x|)\geqslant0$ 等价于 $\begin{cases}(x+1)(2-x)\geqslant0,x\geqslant0\\ (x+1)(2+x)\geqslant0,x<0\end{cases}$，}
\step{解不等式 $(x+1)(2-x)\geqslant0$，$x\geqslant0$，得 $0\leqslant x\leqslant2$，}
\step{解不等式 $(x+1)(2+x)\geqslant0$，$x<0$，得 $x\leqslant-2$ 或 $-1\leqslant x<0$，}
\step{所以 $(x+1)(2-|x|)\geqslant0$ 的解集为 $\{x\mid x\leqslant-2\ \text{或}\ -1\leqslant x\leqslant2\}$，}
\step{综上所述，$\begin{cases}\dfrac{2x+4}{x+4}\geqslant1\\[4pt] (x+1)(2-|x|)\geqslant0\end{cases}$ 的解集为 $\{x\mid x<-4\ \text{或}\ 0\leqslant x\leqslant2\}$.}
\solend
\fi

\ansspace[8]

% 第 18 题题干＋两小问与作答区 6 行须同页
\ansneed{20cm}
\item 已知 $a$ 为实数，设 $f(x)=x^2+2(a+1)x+1$.

\begin{enumerate}[label=(\arabic*)]
\item 已知对于任意实数 $x$，$\sqrt{f(x)}$ 均有意义，求实数 $a$ 的取值范围；
\item 设 $A=\{x\mid f(x)+a^2-2=0\}$，$B=\{-4,0\}$，$A\cup B=B$，求实数 $a$ 的取值范围.
\end{enumerate}

\ifanswer
\solbegin
\stepm{(1)}{$\sqrt{f(x)}$ 对任意实数 $x$ 都有意义，}
\step{即 $f(x)=x^2+2(a+1)x+1\geqslant0$ 对任意 $x\in\R$ 恒成立，}
\step{所以 $\Delta=4(a+1)^2-4\times1\times1\leqslant0$，即 $(a+1)^2\leqslant1$，解得 $-2\leqslant a\leqslant0$，}
\step{故 $a$ 的取值范围是 $[-2,0]$.}
\stepm{(2)}{由 $f(x)+a^2-2=0$ 得 $x^2+2(a+1)x+a^2-1=0$.}
\step{由 $A\cup B=B$ 得 $A\subseteq B=\{-4,0\}$，即该方程无实根，}
\step{或它的实根都在 $\{-4,0\}$ 中.}
\step{当方程无实根时，$\Delta=4(a+1)^2-4(a^2-1)=8a+8<0$，解得 $a<-1$；}
\step{当方程有实根时，$a\geqslant-1$，设两根为 $x_1$、$x_2$，}
\step{由根与系数的关系得 $x_1+x_2=-2(a+1)$、$x_1x_2=a^2-1$，}
\step{此时两根只能在 $\{-4,0\}$ 中取值；}
\step{若两根为 0 与 0，由 $x_1x_2=0$ 得 $a=\pm1$，}
\step{再由 $x_1+x_2=0$ 得 $a=-1$；}
\step{若两根为 0 与 $-4$，由 $x_1x_2=0$ 得 $a=\pm1$，再由 $x_1+x_2=-4$ 得 $a=1$.}
\step{经检验，$a=-1$ 时 $A=\{0\}$，$a=1$ 时 $A=\{-4,0\}$，都满足 $A\subseteq B$.}
\step{综上，实数 $a$ 的取值范围为 $(-\infty,-1]\cup\{1\}$.}
\solend
\fi

\ansspace[8]

\ansneed{18cm}
\item 共享单车给市民出行带来了诸多便利，某公司购买了一批单车投放到某地给居民使用，据市场分析，这批单车的运营累计利润 $y$（单位：元）与运营天数 $x$（单位：天，$x$ 为正整数）满足关系：$y=-\dfrac12x^2+60x-800$.

\begin{enumerate}[label=(\arabic*)]
\item 要使这批单车的运营累计利润高于 800 元，求运营天数的取值范围；
\item 判断这批单车的日平均运营利润能否超过 20 元，并说明理由.
\end{enumerate}

\ifanswer
\solbegin
\stepm{(1)}{由题意得 $y=-\dfrac12x^2+60x-800>800$，}
\step{即 $x^2-120x+3200<0$，解得 $40<x<80$，}
\step{即运营天数的取值范围为 $\{x\mid40<x<80,x\in\N^*\}$.}
\stepm{(2)}{不能，理由如下：}
\step{这批单车的日平均运营利润为 $\dfrac yx=-\dfrac12x-\dfrac{800}x+60=-\dfrac12\left(x+\dfrac{1600}x\right)+60$，}
\step{因 $x>0$，则 $x+\dfrac{1600}x\geqslant2\sqrt{x\cdot\dfrac{1600}x}=80$，}
\step{当且仅当 $x=\dfrac{1600}x$，即 $x=40$ 时取等号，}
\step{此时 $\dfrac yx=-\dfrac12\left(x+\dfrac{1600}x\right)+60\leqslant-\dfrac12\times80+60=20$，}
\step{故这批单车的日平均运营利润不能超过 20 元.}
\solend
\fi

\ansspace[8]

% 第 20 题题干＋三小问与作答区 8 行须同页
\ansneed{20cm}
\item 设 $f(x)=ax^2+bx+c$($a\neq0$，$b$、$c\in\R$).

\begin{enumerate}[label=(\arabic*)]
\item 若 $a=b=c=1$，设 $A=\{(x,y)\mid y=f(x)\}$，$B=\{(x,y)\mid y=1\}$，用列举法表示 $A\cap B$；
% 原卷两个“若”连用（“若 a=1，…，若方程 f(x)=0 有两个大于 1 的根”），疑衍其一；照录（见文件头“原卷笔误/存疑”⑦）
\item 若 $a=1$，$b=2k-1$，$c=k^2$，若方程 $f(x)=0$ 有两个大于 1 的根，求实数 $k$ 的取值范围；
\item 若 $a$、$b$、$c\in\Z$，且 $f(0)$、$f(1)$ 均为奇数，求证：方程 $f(x)=0$ 无整数解.
\end{enumerate}

\ifanswer
\solbegin
\stepm{(1)}{当 $a=b=c=1$ 时，$f(x)=x^2+x+1$，}
\step{所以 $A=\{(x,y)\mid y=x^2+x+1\}$，又 $y=1$，}
\step{由 $x^2+x+1=1$ 得 $x^2+x=0$，解得 $x=0$ 或 $x=-1$，}
\step{所以 $A\cap B=\{(0,1),(-1,1)\}$.}
\stepm{(2)}{当 $a=1$、$b=2k-1$、$c=k^2$ 时，$f(x)=x^2+(2k-1)x+k^2$.}
\step{设两根为 $x_1$、$x_2$，}
\step{则“两根都大于 1”等价于 $\Delta>0$、$f(1)>0$、$-\dfrac b{2a}>1$.}
\step{由 $\Delta=(2k-1)^2-4k^2=1-4k>0$ 得 $k<\dfrac14$；}
\step{由 $f(1)=1+(2k-1)+k^2=k^2+2k=k(k+2)>0$ 得 $k<-2$ 或 $k>0$；}
\step{由 $-\dfrac{2k-1}2=\dfrac{1-2k}2>1$ 得 $k<-\dfrac12$.}
\step{三者取交集得 $k<-2$，即 $k$ 的取值范围是 $(-\infty,-2)$.}
\stepm{(3)}{设 $x_0\in\Z$ 是方程 $f(x)=0$ 的整数解.}
\step{因为 $x_0(x_0-1)$ 是两个相邻整数之积，恒为偶数，所以 $ax_0(x_0-1)$ 为偶数.}
\step{而 $f(x_0)=ax_0^2+bx_0+c=ax_0(x_0-1)+(a+b)x_0+c$，}
\step{所以 $f(x_0)$ 与 $(a+b)x_0+c$ 同奇偶.}
\step{若 $x_0$ 为偶数，则 $(a+b)x_0$ 为偶数，$f(x_0)$ 与 $c=f(0)$ 同奇偶，为奇数；}
\step{若 $x_0$ 为奇数，则 $(a+b)x_0$ 与 $a+b$ 同奇偶，}
\step{$f(x_0)$ 与 $(a+b)+c=f(1)$ 同奇偶，也为奇数.}
\step{综上，$f(x_0)$ 恒为奇数，不可能等于 0，与 $f(x_0)=0$ 矛盾，}
\step{所以方程 $f(x)=0$ 无整数解.}
\solend
\fi

\ansspace[8]

% 压轴题（题干＋三小问）与作答区 10 行须同页
\ansneed{20cm}
\item 已知 $X$ 为有限集，且任取 $x\in X$，均有 $x\in\Z$，称 $X$ 为非空有限整数集. 定义：$X_k=\{x+k\mid x\in X\}$，$X^2=\{y\mid y=x^2,x\in X\}$.

\begin{enumerate}[label=(\arabic*)]
\item 若 $X=\{-2,2\}$，分别写出 $X_2$ 和 $X^2$（用列举法表示）；
\item 若 $X=\{-1,0,1\}$，$X_1\subseteq Y$，$(Y_{-1})^2\subseteq X$，求集合 $Y$ 中所有的元素；
\item 对于非空有限整数集 $A$、$B$，若 $A_k\subseteq B$，$(B^2)_{-2}\subseteq A$，$B$ 中的元素均为非负整数，$k$ 为整数，求证：

① $A$、$B$ 中元素的个数相同；

② 若集合 $A$ 中的元素个数为奇数，则 $k$ 为偶数.
\end{enumerate}

\ifanswer
\solbegin
\stepm{(1)}{由题意得 $X_2=\{0,4\}$，$X^2=\{4\}$；}
\stepm{(2)}{由题意得 $X_1=\{0,1,2\}$，即 $\{0,1,2\}\subseteq Y$，$Y_{-1}=\{y-1\mid y\in Y\}$，}
\step{$(Y_{-1})^2=\{(y-1)^2\mid y\in Y\}\subseteq X=\{-1,0,1\}$，}
\step{所以 $(y-1)^2=0$ 或 $(y-1)^2=1$，解得 $y=1$ 或 $y=2$ 或 $y=0$，}
\step{即 $y\in\{0,1,2\}$，即 $Y\subseteq\{0,1,2\}$，所以 $Y=\{0,1,2\}$；}
% 原卷 (3) 解析首行“(3)”与“①”连排，本稿合并记作 \stepm{(3)①}
\stepm{(3)①}{因为 $A_k=\{x+k\mid x\in A\}$，函数 $g(x)=x+k$ 在全体整数集上严格增，}
\step{所以 $\forall x_1$、$x_2\in A$，$x_1\neq x_2$，$g(x_1)\neq g(x_2)$，所以 $|A_k|=|A|$，}
\step{因为 $A_k\subseteq B$，所以 $|A|\leqslant|B|$，}
\step{设 $b\in B$，$B$ 中元素均为非负整数，$\forall b_1\neq b_2\in B$，$b_1^2\neq b_2^2$，}
\step{故函数 $f(b)=b^2-2$ 在 $B$ 上单射，所以 $|(B^2)_{-2}|=|B|$，}
\step{又因为 $(B^2)_{-2}\subseteq A$，所以 $|B|\leqslant|A|$，所以 $|A|=|B|$.}
\stepm{②}{由①得 $A_k\subseteq B$ 且 $|A_k|=|A|=|B|$，$(B^2)_{-2}\subseteq A$ 且 $|(B^2)_{-2}|=|B|=|A|$，}
\step{所以 $A_k=B$，$(B^2)_{-2}=A$，}
\step{所以 $B=A_k=\{a+k\mid a\in A\}$，$A=(B^2)_{-2}=\{b^2-2\mid b\in B\}$，}
\step{所以 $A=\{(a+k)^2-2\mid a\in A\}$，}
\step{所以映射 $h(a)=(a+k)^2-2$ 是集合 $A$ 到自身的一个双射.}
\step{假设 $k$ 是奇数，当 $a$ 为偶数时，$a+k$ 为奇数，$(a+k)^2-2$ 为奇数，}
\step{当 $a$ 为奇数时，$a+k$ 为偶数，$(a+k)^2-2$ 为偶数，}
\step{所以此时映射 $h(a)=(a+k)^2-2$，将集合 $A$ 中的奇数映射为偶数，}
\step{将集合 $A$ 中的偶数映射为奇数，即集合 $A$ 中的奇数元素和偶数元素必须配对，}
\step{此时集合 $A$ 的元素个数则为偶数，}
\step{与题设中“集合 $A$ 中的元素个数为奇数”矛盾.}
\step{故假设不成立，所以 $k$ 是偶数.}
\solend
\fi

\ansspace*[10]

\end{enumerate}

\end{document}
"
% 紧凑版：xelatex "\def\iscompact{1}% 高一67_华师大一附中_10月月考.tex
%   —— 高一上数学“2026-2027 年华师大一附中高一上 10 月月考”（讲评版）
% 试卷版：xelatex 高一67_华师大一附中_10月月考.tex
% 答案版：xelatex "\def\isanswer{1}\input{高一67_华师大一附中_10月月考.tex}"
% 紧凑版：xelatex "\def\iscompact{1}\input{高一67_华师大一附中_10月月考.tex}"
%
% 原始材料：微信公众号“魔都数学加油站”文章，作者破晓，连载“27学年高一67”；
% 原文标题“27学年高一67——2026-2027年上海市华师大一附中高一上10月月考”；
% 原文链接：https://mp.weixin.qq.com/s/Kdhx07TVjLoGTcIlnRzgcQ。
% 页面标注发布时间 2026-10-18 20:33（上海）。
% 本次转录只据 `原图/高一67_华师大一附中_10月月考/` 的 13 张 PNG 逐页读图。
% 正文为 13 张 PNG 页（`原图/高一67_华师大一附中_10月月考/01..13.png`）：
% 第 01、04、13 页 2560×3620，第 02、03、05—12 页 4762×6735
% （**同一篇各页分辨率不同**），均按页序读取。13 页全为试卷正文页（卷首标题在第 1 页页顶，
% 卷末解答题第 21 题解析止于第 13 页），无封面页、目录页、二维码或宣传页。
% 粉色斜排“魔都数学加油站”为卷面水印，与公众号页脚一并未录入。
%
% 卷首标题：照卷面印作“2026--2027 年华师大一附中高一上 10 月月考”（卷面为黑体加粗一行），
% 年份区间按全稿惯例排 en dash。原卷只有一行标题、无副标题。
%
% 篇幅与题号：原卷自第 3 题起编号（第 1、2 题不在本卷内）。填空题 3—12（10 题）、
% 选择题 13—16（4 题）、解答题 17—21（5 题），共 19 题。本稿沿用原节与原题号，
% 只改排为全稿统一的大节样式（原卷小节标题为“一、填空题”“二、选择题”“三、解答题”）。
%
% 讲评版：原卷每题下直接印【解析】。**全卷 19 题中 18 题都印【解析】（无一处只印
% 【答案】）**，按全稿口径这类题不另提 \ans{}；各题解析末句的结论（第 13、14、16 题
% 的“故选 B／D／B”、各填空与解答题的“故/综上…”）照录在解析内。
% **另有 1 题答案印在题干括号内**——第 15 题（“(  D  )”，且该题无【答案】/【解析】段），
% 本稿按 高一37 第 13 题、高一63 第 14 题成例把括号留空、答案另提为 \ans{$D$}。
%
% 副标题：原卷未印副标题。本稿按“知识模块”口径标作“集合与常用逻辑用语、不等式”：
% 19 题中集合与常用逻辑用语 12 题（第 3、4、5、7、10、11、12、13、15、16、18、21 题：
% 补集、并集、子集、充分条件、容斥计数、子集计数、新定义二元运算、条件判断、
% 反证法假设、新定义生成集、集合关系与二次方程、新定义 X_k 与 X^2），
% 不等式 7 题（第 6、8、9、14、17、19、20 题：存在量词不等式、解集反求系数、
% 解集为空、不等式性质、解不等式（组）、二次不等式与基本不等式、二次方程根的分布），
% 两类均超过三分之一，故并标。
%
% 与原卷的排版差异：
%   ① 原卷每题下直接印【解析】，本稿按全稿惯例将解析置于答案版，试卷版保持干净；
%   ② 原卷未印副标题，本稿按“知识模块”口径补出；
%   ③ 原卷填空横线约四个字宽，本稿按全稿惯例统一排 \blankline[4.4em]；
%   ④ 原卷实数集排斜体 $R$（第 14 题题干、第 17 题 (1) 题干、第 18 题解析、第 20 题题干），
%      整数集排斜体 $Z$（第 20 题 (3) 题干与解析、第 21 题题干），正整数集排粗体
%      $N^*$（第 11 题解析、第 19 题解析），本稿按全稿惯例统一改用 $\R$、$\Z$、$\N^*$
%      （记号归一，粗体／斜体不作区分）；空集原卷作 $\varnothing$，照录；
%   ⑤ 原卷在数学式内多处用半角逗号（如第 8 题 $(1,2)$ 与各 cases、第 11 题解析
%      $P_n=\{1,2,\cdots,n\}$,$n\in N^*$、第 12 题解析 $u=(\cdots)$,$v=(\cdots)$ 与
%      “2020,个 1”、第 19 题解析 $\{x|40<x<80,x\in N^*\}$），均照录未强改；
%      描述法集合的竖线照原卷排 \mid；句末句点按全稿惯例排西文句点；
%   ⑥ 原卷补集记号为上划线 $\bar A$（第 3 题题干与解析、第 11 题题干），照录
%      （同 高一64/65 的 $\bar B$ 成例）；
%   ⑦ 第 12 题新定义的运算记号 $*$ 照录（$\alpha*\beta$、$u*u$、$u*v$）；
%   ⑧ 第 17—21 题的子问原卷各自另起一行，本稿按全稿惯例用 enumerate 的 (1)(2)(3)；
%      解析中的子问标号一律用 \stepm{(1)}{…}；第 5 题解析的①②、第 21 题 (3) 解析的①②
%      按全稿惯例作 \stepm{①}{…}、\stepm{②}{…}（第 21 题 (3) 首行原卷“(3)”与“①”连排，
%      本稿合并记作 \stepm{(3)①}{…}）；
%   ⑨ 选择题作答括号原卷排半角“(　　)”（第 13、14、16 题括号内空白、第 15 题括号内印 D），
%      本稿按全稿惯例排 \mbox{（\quad）}；第 15 题括号内的答案移入 \ans{}，见上；
%   ⑩ 选项行排照原卷：第 13、15 题两行两条，第 14、16 题一行四条；
%   ⑪ 每道解答题原卷未留作答空白，本稿按全稿惯例加 \ansspace（末题用 \ansspace*）；
%      第 20、21 题题干长，各用一处 \ansneed 让题干与作答区不拆页（\ansneed 只作用于试卷版）。
%
% 原卷笔误/存疑（均照录未改，正文对应处另有 % 注）：
%   ① 第 3 题解析末句印作“则 $\bar A=(-\infty,0)\cup[2,3)$.”：按题干
%      $U=\{x\mid 1/x\leqslant 1/2\}$（即 $x<0$ 或 $x\geqslant2$）与 $A=[3,+\infty)$，
%      $\bar A$ 应为 $(-\infty,0)\cup\{2\}$；且与本解析首句自求之 $U$ 相矛盾（$3\notin U$），
%      疑为笔误；照录未改。
%   ② 第 12 题题干 $\alpha*\beta$ 定义式括号不配对：第四组“$+\cdots+x_{2020}+y_{2020}
%      -x_{2020}y_{2020}$”缺左括号、末尾连排两个右括号“))”；照录未改。
%   ③ 第 12 题解析“所以…中有 2020,个 1”一行，“$\cdots$”后脱逗号（作
%      “$x_2,\cdots x_{2020}$”，第 12 行“而…”同式样），且“2020”后衍一半角逗号
%      （作“2020,个 1”）；照录未改。
%   ④ 第 16 题解析乘积链印作“$a_1a_2<a_1a_3<a_1a_4<a_1a_5<a_2a_3<a_3a_3<a_{14}a_5$”：
%      “$a_3a_3$”同足标自乘、与单调链不合，“$a_{14}a_5$”下标 14 与 $a_1$—$a_5$ 记号
%      不合（疑为 $a_2a_4$、$a_4a_5$ 之误）；照录未改。
%   ⑤ 第 5 题题干作“$Q\subset P$”（真子集记号），解析两处均作“$Q\subseteq P$”，
%      同一题前后记号不一；照录（本稿分别排 \subset、\subseteq）。
%   ⑥ 第 17 题 (1) 解析“当 $a=1$ 时，原不等式为 $(x+1)^2<0$ 此时…”与
%      “当 $0<a<1$ 时，$-\frac1a<-1$ 此时…”两处“此时”前均脱标点；照录未改。
%   ⑦ 第 20 题 (2) 题干“若 $a=1$，$b=2k-1$，$c=k^2$，若方程 $f(x)=0$ 有两个大于 1 的根”
%      两个“若”连用，疑衍其一；照录未改。
%
% 插图：无（全卷 13 页无任何几何图形、无“如图/右图/图中”字样）。
%
% 行数核对：逐题数 \step/\stepm 条数与原卷【解析】印刷行数，18 题逐一相符
% （第 15 题无解析段，只提 \ans{}），合计 169 条 / 169 行
% （Q3—Q12：2、3、4、3、1、3、5、3、13、15；Q13—Q16：2、3、0、25；
% Q17—Q21：19、16、9、20、23）。
%
\documentclass[a4paper,11pt]{article}
\usepackage{common}

\begin{document}

\examtitlest[3]{2026--2027 年华师大一附中高一上 10 月月考}{集合与常用逻辑用语、不等式}

\bigsec{一、填空题}

\begin{enumerate}
\setcounter{enumi}{2}

\item 设全集 $U=\left\{x\mid\dfrac1x\leqslant\dfrac12\right\}$，$A=[3,+\infty)$，则 $\bar A=$\blankline[4.4em].

\ifanswer
\solbegin
\step{$\dfrac1x\leqslant\dfrac12$，解得 $x<0$ 或 $x\geqslant2$，故 $U=(-\infty,0)\cup[2,+\infty)$，}
% 原卷末句作 $\bar A=(-\infty,0)\cup[2,3)$，按题干 $U$ 应为 $(-\infty,0)\cup\{2\}$，与本解析首句矛盾，疑笔误；照录（见文件头“原卷笔误/存疑”①）
\step{$A=[3,+\infty)$，则 $\bar A=(-\infty,0)\cup[2,3)$.}
\solend
\fi

\item 已知集合 $M=\{y\mid y=x^2-2x-1\}$，$N=\{y\mid y=x^2-4x+1\}$，则 $M\cup N=$\blankline[4.4em].

\ifanswer
\solbegin
\step{$M=\{y\mid y=x^2-2x-1\}=\{y\mid y=(x-1)^2-2\}=\{y\mid y\geqslant-2\}$，}
\step{$N=\{y\mid y=x^2-4x+1\}=\{y\mid y=(x-2)^2-3\}=\{y\mid y\geqslant-3\}$，}
\step{则 $M\cup N=\{y\mid y\geqslant-3\}$.}
\solend
\fi

\item 已知 $P=\{x\mid x^2+x-6=0\}$，$Q=\{x\mid mx-1=0\}$，且 $Q\subset P$，则实数 $m$ 的值\blankline[4.4em].

\ifanswer
\solbegin
\step{$P=\{2,-3\}$，$Q=\{x\mid mx=1\}$；}
\stepm{①}{$m=0$ 时，$Q=\varnothing$，满足 $Q\subseteq P$；}
\stepm{②}{$m\neq0$ 时，要使 $Q\subseteq P$，则 $\dfrac1m=2$ 或 $-3$，所以 $m=\dfrac12$ 或 $-\dfrac13$；}
\step{所以 $m=0$ 或 $\dfrac12$ 或 $-\dfrac13$.}
\solend
\fi

\item 已知“存在正实数 $x$，使得关于 $x$ 的不等式 $ax-1>0$ 成立”为真命题，则实数 $a$ 的取值范围是\blankline[4.4em].

\ifanswer
\solbegin
\step{存在正实数 $x$，使得关于 $x$ 的不等式 $ax-1>0$ 成立为真命题，}
\step{即存在正实数 $x$，使得 $a>\dfrac1x$，因为 $\dfrac1x>0$，所以 $a>0$，}
\step{即 $a$ 的范围为 $(0,+\infty)$.}
\solend
\fi

\item 设 $\alpha:3\leqslant x\leqslant5$，$\beta:2m-1\leqslant x<m+4$，若 $\alpha$ 是 $\beta$ 的充分条件，则实数 $m$ 的取值范围是\blankline[4.4em].

\ifanswer
\solbegin
\step{若 $\alpha$ 是 $\beta$ 的充分条件，则 $\begin{cases}2m-1\leqslant3\\ m+4>5\end{cases}$，解得 $1<m\leqslant2$.}
\solend
\fi

\item 对于实数 $a$、$b$、$c$，已知关于 $x$ 的不等式 $ax^2+bx+c>0$ 的解集是 $(1,2)$，则关于 $x$ 的不等式 $cx^2-bx+a\leqslant0$ 的解集是\blankline[4.4em].

\ifanswer
\solbegin
\step{关于 $x$ 的不等式 $ax^2+bx+c>0$ 的解集是 $(1,2)$，则 $\begin{cases}a<0\\ -\dfrac ba=3\\ \dfrac ca=2\end{cases}$，则 $\begin{cases}b=-3a\\ c=2a\end{cases}$，}
\step{则关于 $x$ 的不等式 $cx^2-bx+a=2ax^2+3ax+a\leqslant0$，则 $2x^2+3x+1\geqslant0$，}
\step{则 $x\leqslant-1$ 或 $x\geqslant-\dfrac12$，则解集为 $(-\infty,-1]\cup\left[-\dfrac12,+\infty\right)$.}
\solend
\fi

\item 已知关于 $x$ 的不等式 $kx^2+kx+k+2\leqslant0$ 的解集为 $\varnothing$，则实数 $k$ 的取值范围为\blankline[4.4em].

\ifanswer
\solbegin
\step{当 $k=0$ 时，不等式化为 $2\leqslant0$，解集为 $\varnothing$，符合题意；}
\step{当 $k\neq0$ 时，若关于 $x$ 的不等式 $kx^2+kx+k+2\leqslant0$ 的解集为 $\varnothing$，}
\step{则 $\begin{cases}k>0\\ \Delta=k^2-4k(k+2)<0\end{cases}$，化简得 $\begin{cases}k>0\\ -3k^2-8k<0\end{cases}$，即 $\begin{cases}k>0\\ k(3k+8)>0\end{cases}$，}
\step{解得 $k>0$.}
\step{综上所述，$k\geqslant0$，即实数 $k$ 的取值范围为 $[0,+\infty)$.}
\solend
\fi

\item 对班级若干名学生调查对 $A$、$B$ 两事件的态度，有如下结果：赞成 $A$ 的人数是全体的五分之三，其余的不赞成；赞成 $B$ 的比赞成 $A$ 多 3 人，其余的不赞成. 已知对 $A$、$B$ 都赞成的学生有 18 人，对 $A$、$B$ 都不赞成的学生数比对 $A$、$B$ 都赞成的学生数的三分之一多 1 人，则班级总人数是\blankline[4.4em]人.

\ifanswer
\solbegin
\step{设班级总人数是 $x$ 人，则赞成 $A$ 的人数是 $\dfrac35x$，赞成 $B$ 的人数是 $\dfrac35x+3$，}
\step{$A$、$B$ 都不赞成的人数为 $18\times\dfrac13+1=7$ 人，}
\step{由容斥原理得 $x=\left(\dfrac35x+\dfrac35x+3-18\right)+7$，解得 $x=40$，即班级总人数是 40 人.}
\solend
\fi

\item 设全集 $U=\{1,2,3,4,5,6,7,8,9,10,11,12\}$，若集合 $A$ 满足：① $A\subseteq U$；②若 $x\in A$，必有 $2x\notin A$；③若 $x\in\bar A$，必有 $2x\notin\bar A$，则这样的集合 $A$ 有\blankline[4.4em]个.

\ifanswer
\solbegin
\step{设集合 $P_n=\{1,2,\cdots,n\}$,$n\in\N^*$，}
\step{任取 $x\in P_n$，将 $x$ 除以 2，若商仍为偶数，再除以 2，$\cdots$，}
\step{经过 $k$ 次以后，商必为奇数，此时记商为 $m$，}
\step{于是 $x=m\cdot2^k$，其中 $m$ 为奇数，$k\in\N^*$.}
\step{由条件得，若 $m\in A$，则 $x\in A$ 等价于 $k$ 为偶数；}
\step{若 $m\notin A$，则 $x\in A$ 等价于 $k$ 为奇数；}
\step{于是 $x$ 是否属于 $A$ 由 $m$ 是否属于 $A$ 确定；}
\step{设 $Q_n$ 是 $P_n$ 中所有奇数的集合，}
\step{设 $f(n)$ 为集合 $P_n$ 中满足题干条件的集合 $A$ 的个数，}
\step{因此 $f(n)$ 等于 $Q_n$ 的子集个数.}
\step{当 $n$ 为偶数（或奇数）时，$P_n$ 中奇数的个数是 $\dfrac n2$（或 $\dfrac{n+1}2$），}
\step{所以 $f(n)=\begin{cases}2^{\frac n2},n\text{为偶数}\\[2pt] 2^{\frac{n+1}2},n\text{为奇数}\end{cases}$，所以 $P_{12}=\{1,2,3,4,5,6,7,8,9,10,11,12\}$，}
\step{即同时满足三个条件的集合 $A$ 的个数为 $2^{\frac{12}2}=2^6=64$.}
\solend
\fi

\item 设集合 $\Omega=\{(x_1,x_2,\cdots,x_{2020})\mid x_i\in\{0,1\},i=1,2,\cdots,2020\}$，对于 $\Omega$ 中的任意两个元素 $\alpha=(x_1,x_2,\cdots,x_{2020})$，$\beta=(y_1,y_2,\cdots,y_{2020})$，记
% 原卷此式第四组缺左括号、末尾连排两个右括号“))”，括号不配对；照录（见文件头“原卷笔误/存疑”②）
$\alpha*\beta=(x_1+y_1-x_1y_1)+(x_2+y_2-x_2y_2)+\cdots+x_{2020}+y_{2020}-x_{2020}y_{2020}))$，设 $u$、$v\in\Omega$，若 $u*u+v*v=2020$，则 $u*v$ 的最大值和最小值之和为\blankline[4.4em].

\ifanswer
\solbegin
\step{设 $u=(x_1,x_2,\cdots,x_{2020})$,$v=(y_1,y_2,\cdots,y_{2020})$，}
\step{因为 $x_i\in\{0,1\}$，所以 $x_i(x_i-1)=0$，}
\step{所以 $u*u=(x_1+x_1-x_1^2)+(x_2+x_2-x_2^2)+\cdots+(x_{2020}+x_{2020}-x_{2020}^2)$}
\step{$=x_1+x_2+\cdots+x_{2020}$，}
\step{$v*v=(y_1+y_1-y_1^2)+(y_2+y_2-y_2^2)+\cdots+(y_{2020}+y_{2020}-y_{2020}^2)$}
\step{$=y_1+y_2+\cdots+y_{2020}$，}
\step{由 $u*u+v*v=2020$，得 $x_1+x_2+\cdots+x_{2020}+y_1+y_2+\cdots+y_{2020}=2020$，}
% 原卷此行“$\cdots$”后脱逗号、 “2020”后衍一半角逗号（“2020,个 1”）；照录（见文件头“原卷笔误/存疑”③）
\step{所以 $x_1,x_2,\cdots x_{2020},y_1,y_2,\cdots y_{2020}$ 中有 2020,个 1，2020 个 0，}
\step{$u*v=(x_1+y_1-x_1y_1)+(x_2+y_2-x_2y_2)+\cdots+(x_{2020}+y_{2020}-x_{2020}y_{2020})$}
\step{$=2020-(x_1y_1+x_2y_2+\cdots+x_{2020}y_{2020})$，}
\step{注意到只有 $x_i=y_i=1$ 时，$x_iy_i=1$，否则 $x_iy_i=0$，}
% 此行“$\cdots$”后同样脱逗号，与上一处同式样；照录（见文件头“原卷笔误/存疑”③）
\step{而 $x_1,x_2,\cdots x_{2020},y_1,y_2,\cdots y_{2020}$ 中有 2020 个 1，2020 个 0，}
\step{所以满足 $x_iy_i=1$ 的最多有 1010 个，所以 $u*v\geqslant2020-1010=1010$，}
\step{满足 $x_iy_i=0$ 的最多有 2020 个，所以 $u*v\leqslant2020-0=2020$，}
\step{则 $u*v$ 的最大值和最小值之和为 $2020+1010=3030$.}
\solend
\fi

\end{enumerate}

\bigsec{二、选择题}

\begin{enumerate}
\setcounter{enumi}{12}

\item 实数 $x$ 与 $y$，“$x+y>2$”的\mbox{（\quad）}条件是“$x>1$ 或 $y>1$”.

\keepwithprev
A. 充分非必要\qquad B. 必要非充分

C. 充要\qquad D. 既非充分又非必要

\ifanswer
\solbegin
\step{当 $x>1$ 或 $y>1$ 时，$x+y>2$ 不一定成立，充分性不成立；}
\step{当 $x+y>2$ 时，$x>1$ 或 $y>1$ 一定成立，必要性成立. 故选 $B$.}
\solend
\fi

\item 设 $a,b\in\R$，且 $a<0$，$-1<b<0$，那么下列不等式成立的是\mbox{（\quad）}

\keepwithprev
A. $a>ab>ab^2$\qquad B. $ab^2>ab>a$\qquad C. $ab>a>ab^2$\qquad D. $ab>ab^2>a$

\ifanswer
\solbegin
\step{因为 $-1<b<0$，$a<0$，所以 $ab>0$，$0<b^2<1$.}
\step{所以 $ab-ab^2=ab(1-b)>0$，$ab^2-a=a(b^2-1)>0$，}
\step{所以 $ab>ab^2>a$. 故选 $D$.}
\solend
\fi

% 第 15 题答案原卷印在题干括号内（“(  D  )”），且无【答案】/【解析】段；
% 按 高一37 第 13 题、高一63 第 14 题成例括号留空、答案另提 \ans{$D$}
\item 用反证法证明：“平面四边形中至少有一个内角不超过 90°”时，下列“假设”中正确的是\mbox{（\quad）}

\keepwithprev
A. 假设有两个内角超过 90°\qquad B. 假设有三个内角超过 90°

C. 假设至多有两个内角超过 90°\qquad D. 假设四个内角都超过 90°

\ifanswer
\ans{$D$}
\fi

\item 设 $A$ 为非空实数集，集合 $B=\{xy\mid x\in A,y\in A,x\neq y\}$，则对于以下两个结论说法正确的是\mbox{（\quad）}

结论①：若 $A$ 是五元集，则 $B$ 中的元素最多 10 个；

结论②：存在四元集 $A$，使得 $B=\{2,4,5,8,10,16\}$.

\keepwithprev
A. ①对②对\qquad B. ①对②错\qquad C. ①错②对\qquad D. ①错②错

\ifanswer
\solbegin
\step{对于①，想让元素个数达到最大值 10，则须避免产生符号不同的乘积，}
\step{否则正负分开计数容易导致总数受限，且容易产生相等值，}
\step{也最好避免 0 的出现，}
\step{故设 $A=\{a_1,a_2,a_3,a_4,a_5\}$，$0<a_1<a_2<a_3<a_4<a_5$.}
% 原卷此链“$a_3a_3$”同足标自乘、“$a_{14}a_5$”下标 14 均与文意不合，疑为 $a_2a_4$、$a_4a_5$ 之误；照录（见文件头“原卷笔误/存疑”④）
\step{因为 $a_1a_2<a_1a_3<a_1a_4<a_1a_5<a_2a_3<a_3a_3<a_{14}a_5$，$a_2a_3<a_2a_4<a_3a_4$，}
\step{所以 $B$ 中元素个数小于等于 10 个，}
\step{如设 $A=\{1,2,3,5,7\}$，则 $B=\{2,3,5,6,7,10,14,15,21,35\}$，}
\step{所以集合 $B$ 中元素个数的最大值为 10 个，故①正确；}
\step{对于②，由于集合 $B$ 中没有 0，所以集合 $A$ 中没有 0；}
\step{设集合 $A=\{a,b,c,d\}$，且 $a<b<c<d$，}
\step{当 $a<0<b<c<d$ 时，$a\times b<0$，$a\times c<0$，$a\times d<0$，}
\step{故集合 $B$ 至少有 3 个负数，与 $B=\{2,4,5,8,10,16\}$ 矛盾；}
\step{当 $a<b<0<c<d$ 时，$a\times b>0$，$c\times d>0$，}
\step{$a\times c<0$，$a\times d<0$，$b\times c<0$，$b\times d<0$，}
\step{故集合 $B$ 至少有 4 个负数，与 $B=\{2,4,5,8,10,16\}$ 矛盾；}
\step{当 $a<b<c<0<d$ 时，$a\times b>0$，$a\times c>0$，$b\times c>0$，}
\step{$a\times d<0$，$b\times d<0$，$c\times d<0$，}
\step{故集合 $B$ 至少有 3 个负数，与 $B=\{2,4,5,8,10,16\}$ 矛盾；}
\step{当 $0<a<b<c<d$ 时，则集合 $B=\{ab,ac,ad,bc,bd,cd\}$，}
\step{则必有 $ab=2$，$cd=16$，其 4 个正实数的乘积 $abcd=32$，}
\step{又 $ad=5$，$bc=8$ 或 $ad=8$，$bc=5$，其 4 个正实数的乘积 $abcd=40$，矛盾；}
\step{当 $a<b<c<d<0$ 时，与 $0<a<b<c<d$ 时的证明方法相同，}
\step{得集合 $A$ 也不存在；}
\step{综上不存在四元集 $A$，使其生成集 $B=\{2,4,5,8,10,16\}$，故②错误.}
\step{故选 $B$.}
\solend
\fi

\end{enumerate}

\bigsec{三、解答题}

\begin{enumerate}
\setcounter{enumi}{16}

\item 解下列不等式（组）

\begin{enumerate}[label=(\arabic*)]
\item 解不等式：$(ax+1)(x+1)<0(a\in\R)$.
\item 解不等式组：$\begin{cases}\dfrac{2x+4}{x+4}\geqslant1\\[4pt] (x+1)(2-|x|)\geqslant0\end{cases}$.
\end{enumerate}

\ifanswer
\solbegin
\stepm{(1)}{当 $a=0$ 时，原不等式为 $x+1<0$，即 $x<-1$，}
\step{所以此时不等式的解集为 $\{x\mid x<-1\}$；}
\step{当 $a\neq0$ 时，令 $(ax+1)(x+1)=0$，解得 $x_1=-\dfrac1a$,$x_2=-1$，}
\step{当 $a>1$ 时，$-\dfrac1a>-1$，此时不等式的解集为 $\left\{x\mid-1<x<-\dfrac1a\right\}$；}
% 以下两行“此时”前原卷均无标点；照录（见文件头“原卷笔误/存疑”⑥）
\step{当 $a=1$ 时，原不等式为 $(x+1)^2<0$ 此时不等式的解集为 $\varnothing$；}
\step{当 $0<a<1$ 时，$-\dfrac1a<-1$ 此时不等式的解集为 $\left\{x\mid-\dfrac1a<x<-1\right\}$；}
\step{当 $a<0$ 时，$-\dfrac1a>-1$，此时不等式的解集为 $\left\{x\mid x<-1\ \text{或}\ x>-\dfrac1a\right\}$.}
\step{综上所述，当 $a=0$ 时，$\{x\mid x<-1\}$；}
\step{当 $a>1$ 时，$\left\{x\mid-1<x<-\dfrac1a\right\}$；}
\step{当 $a=1$ 时，$\varnothing$；}
\step{当 $0<a<1$ 时，$\left\{x\mid-\dfrac1a<x<-1\right\}$；}
\step{当 $a<0$ 时，$\left\{x\mid x<-1\ \text{或}\ x>-\dfrac1a\right\}$.}
\stepm{(2)}{不等式 $\dfrac{2x+4}{x+4}\geqslant1$ 可化为 $\dfrac{2x+4}{x+4}-1\geqslant0$，即 $\dfrac x{x+4}\geqslant0$，}
\step{等价于 $\begin{cases}x(x+4)\geqslant0\\ x+4\neq0\end{cases}$，解得 $x<-4$ 或 $x\geqslant0$，}
\step{不等式 $(x+1)(2-|x|)\geqslant0$ 等价于 $\begin{cases}(x+1)(2-x)\geqslant0,x\geqslant0\\ (x+1)(2+x)\geqslant0,x<0\end{cases}$，}
\step{解不等式 $(x+1)(2-x)\geqslant0$，$x\geqslant0$，得 $0\leqslant x\leqslant2$，}
\step{解不等式 $(x+1)(2+x)\geqslant0$，$x<0$，得 $x\leqslant-2$ 或 $-1\leqslant x<0$，}
\step{所以 $(x+1)(2-|x|)\geqslant0$ 的解集为 $\{x\mid x\leqslant-2\ \text{或}\ -1\leqslant x\leqslant2\}$，}
\step{综上所述，$\begin{cases}\dfrac{2x+4}{x+4}\geqslant1\\[4pt] (x+1)(2-|x|)\geqslant0\end{cases}$ 的解集为 $\{x\mid x<-4\ \text{或}\ 0\leqslant x\leqslant2\}$.}
\solend
\fi

\ansspace[8]

% 第 18 题题干＋两小问与作答区 6 行须同页
\ansneed{20cm}
\item 已知 $a$ 为实数，设 $f(x)=x^2+2(a+1)x+1$.

\begin{enumerate}[label=(\arabic*)]
\item 已知对于任意实数 $x$，$\sqrt{f(x)}$ 均有意义，求实数 $a$ 的取值范围；
\item 设 $A=\{x\mid f(x)+a^2-2=0\}$，$B=\{-4,0\}$，$A\cup B=B$，求实数 $a$ 的取值范围.
\end{enumerate}

\ifanswer
\solbegin
\stepm{(1)}{$\sqrt{f(x)}$ 对任意实数 $x$ 都有意义，}
\step{即 $f(x)=x^2+2(a+1)x+1\geqslant0$ 对任意 $x\in\R$ 恒成立，}
\step{所以 $\Delta=4(a+1)^2-4\times1\times1\leqslant0$，即 $(a+1)^2\leqslant1$，解得 $-2\leqslant a\leqslant0$，}
\step{故 $a$ 的取值范围是 $[-2,0]$.}
\stepm{(2)}{由 $f(x)+a^2-2=0$ 得 $x^2+2(a+1)x+a^2-1=0$.}
\step{由 $A\cup B=B$ 得 $A\subseteq B=\{-4,0\}$，即该方程无实根，}
\step{或它的实根都在 $\{-4,0\}$ 中.}
\step{当方程无实根时，$\Delta=4(a+1)^2-4(a^2-1)=8a+8<0$，解得 $a<-1$；}
\step{当方程有实根时，$a\geqslant-1$，设两根为 $x_1$、$x_2$，}
\step{由根与系数的关系得 $x_1+x_2=-2(a+1)$、$x_1x_2=a^2-1$，}
\step{此时两根只能在 $\{-4,0\}$ 中取值；}
\step{若两根为 0 与 0，由 $x_1x_2=0$ 得 $a=\pm1$，}
\step{再由 $x_1+x_2=0$ 得 $a=-1$；}
\step{若两根为 0 与 $-4$，由 $x_1x_2=0$ 得 $a=\pm1$，再由 $x_1+x_2=-4$ 得 $a=1$.}
\step{经检验，$a=-1$ 时 $A=\{0\}$，$a=1$ 时 $A=\{-4,0\}$，都满足 $A\subseteq B$.}
\step{综上，实数 $a$ 的取值范围为 $(-\infty,-1]\cup\{1\}$.}
\solend
\fi

\ansspace[8]

\ansneed{18cm}
\item 共享单车给市民出行带来了诸多便利，某公司购买了一批单车投放到某地给居民使用，据市场分析，这批单车的运营累计利润 $y$（单位：元）与运营天数 $x$（单位：天，$x$ 为正整数）满足关系：$y=-\dfrac12x^2+60x-800$.

\begin{enumerate}[label=(\arabic*)]
\item 要使这批单车的运营累计利润高于 800 元，求运营天数的取值范围；
\item 判断这批单车的日平均运营利润能否超过 20 元，并说明理由.
\end{enumerate}

\ifanswer
\solbegin
\stepm{(1)}{由题意得 $y=-\dfrac12x^2+60x-800>800$，}
\step{即 $x^2-120x+3200<0$，解得 $40<x<80$，}
\step{即运营天数的取值范围为 $\{x\mid40<x<80,x\in\N^*\}$.}
\stepm{(2)}{不能，理由如下：}
\step{这批单车的日平均运营利润为 $\dfrac yx=-\dfrac12x-\dfrac{800}x+60=-\dfrac12\left(x+\dfrac{1600}x\right)+60$，}
\step{因 $x>0$，则 $x+\dfrac{1600}x\geqslant2\sqrt{x\cdot\dfrac{1600}x}=80$，}
\step{当且仅当 $x=\dfrac{1600}x$，即 $x=40$ 时取等号，}
\step{此时 $\dfrac yx=-\dfrac12\left(x+\dfrac{1600}x\right)+60\leqslant-\dfrac12\times80+60=20$，}
\step{故这批单车的日平均运营利润不能超过 20 元.}
\solend
\fi

\ansspace[8]

% 第 20 题题干＋三小问与作答区 8 行须同页
\ansneed{20cm}
\item 设 $f(x)=ax^2+bx+c$($a\neq0$，$b$、$c\in\R$).

\begin{enumerate}[label=(\arabic*)]
\item 若 $a=b=c=1$，设 $A=\{(x,y)\mid y=f(x)\}$，$B=\{(x,y)\mid y=1\}$，用列举法表示 $A\cap B$；
% 原卷两个“若”连用（“若 a=1，…，若方程 f(x)=0 有两个大于 1 的根”），疑衍其一；照录（见文件头“原卷笔误/存疑”⑦）
\item 若 $a=1$，$b=2k-1$，$c=k^2$，若方程 $f(x)=0$ 有两个大于 1 的根，求实数 $k$ 的取值范围；
\item 若 $a$、$b$、$c\in\Z$，且 $f(0)$、$f(1)$ 均为奇数，求证：方程 $f(x)=0$ 无整数解.
\end{enumerate}

\ifanswer
\solbegin
\stepm{(1)}{当 $a=b=c=1$ 时，$f(x)=x^2+x+1$，}
\step{所以 $A=\{(x,y)\mid y=x^2+x+1\}$，又 $y=1$，}
\step{由 $x^2+x+1=1$ 得 $x^2+x=0$，解得 $x=0$ 或 $x=-1$，}
\step{所以 $A\cap B=\{(0,1),(-1,1)\}$.}
\stepm{(2)}{当 $a=1$、$b=2k-1$、$c=k^2$ 时，$f(x)=x^2+(2k-1)x+k^2$.}
\step{设两根为 $x_1$、$x_2$，}
\step{则“两根都大于 1”等价于 $\Delta>0$、$f(1)>0$、$-\dfrac b{2a}>1$.}
\step{由 $\Delta=(2k-1)^2-4k^2=1-4k>0$ 得 $k<\dfrac14$；}
\step{由 $f(1)=1+(2k-1)+k^2=k^2+2k=k(k+2)>0$ 得 $k<-2$ 或 $k>0$；}
\step{由 $-\dfrac{2k-1}2=\dfrac{1-2k}2>1$ 得 $k<-\dfrac12$.}
\step{三者取交集得 $k<-2$，即 $k$ 的取值范围是 $(-\infty,-2)$.}
\stepm{(3)}{设 $x_0\in\Z$ 是方程 $f(x)=0$ 的整数解.}
\step{因为 $x_0(x_0-1)$ 是两个相邻整数之积，恒为偶数，所以 $ax_0(x_0-1)$ 为偶数.}
\step{而 $f(x_0)=ax_0^2+bx_0+c=ax_0(x_0-1)+(a+b)x_0+c$，}
\step{所以 $f(x_0)$ 与 $(a+b)x_0+c$ 同奇偶.}
\step{若 $x_0$ 为偶数，则 $(a+b)x_0$ 为偶数，$f(x_0)$ 与 $c=f(0)$ 同奇偶，为奇数；}
\step{若 $x_0$ 为奇数，则 $(a+b)x_0$ 与 $a+b$ 同奇偶，}
\step{$f(x_0)$ 与 $(a+b)+c=f(1)$ 同奇偶，也为奇数.}
\step{综上，$f(x_0)$ 恒为奇数，不可能等于 0，与 $f(x_0)=0$ 矛盾，}
\step{所以方程 $f(x)=0$ 无整数解.}
\solend
\fi

\ansspace[8]

% 压轴题（题干＋三小问）与作答区 10 行须同页
\ansneed{20cm}
\item 已知 $X$ 为有限集，且任取 $x\in X$，均有 $x\in\Z$，称 $X$ 为非空有限整数集. 定义：$X_k=\{x+k\mid x\in X\}$，$X^2=\{y\mid y=x^2,x\in X\}$.

\begin{enumerate}[label=(\arabic*)]
\item 若 $X=\{-2,2\}$，分别写出 $X_2$ 和 $X^2$（用列举法表示）；
\item 若 $X=\{-1,0,1\}$，$X_1\subseteq Y$，$(Y_{-1})^2\subseteq X$，求集合 $Y$ 中所有的元素；
\item 对于非空有限整数集 $A$、$B$，若 $A_k\subseteq B$，$(B^2)_{-2}\subseteq A$，$B$ 中的元素均为非负整数，$k$ 为整数，求证：

① $A$、$B$ 中元素的个数相同；

② 若集合 $A$ 中的元素个数为奇数，则 $k$ 为偶数.
\end{enumerate}

\ifanswer
\solbegin
\stepm{(1)}{由题意得 $X_2=\{0,4\}$，$X^2=\{4\}$；}
\stepm{(2)}{由题意得 $X_1=\{0,1,2\}$，即 $\{0,1,2\}\subseteq Y$，$Y_{-1}=\{y-1\mid y\in Y\}$，}
\step{$(Y_{-1})^2=\{(y-1)^2\mid y\in Y\}\subseteq X=\{-1,0,1\}$，}
\step{所以 $(y-1)^2=0$ 或 $(y-1)^2=1$，解得 $y=1$ 或 $y=2$ 或 $y=0$，}
\step{即 $y\in\{0,1,2\}$，即 $Y\subseteq\{0,1,2\}$，所以 $Y=\{0,1,2\}$；}
% 原卷 (3) 解析首行“(3)”与“①”连排，本稿合并记作 \stepm{(3)①}
\stepm{(3)①}{因为 $A_k=\{x+k\mid x\in A\}$，函数 $g(x)=x+k$ 在全体整数集上严格增，}
\step{所以 $\forall x_1$、$x_2\in A$，$x_1\neq x_2$，$g(x_1)\neq g(x_2)$，所以 $|A_k|=|A|$，}
\step{因为 $A_k\subseteq B$，所以 $|A|\leqslant|B|$，}
\step{设 $b\in B$，$B$ 中元素均为非负整数，$\forall b_1\neq b_2\in B$，$b_1^2\neq b_2^2$，}
\step{故函数 $f(b)=b^2-2$ 在 $B$ 上单射，所以 $|(B^2)_{-2}|=|B|$，}
\step{又因为 $(B^2)_{-2}\subseteq A$，所以 $|B|\leqslant|A|$，所以 $|A|=|B|$.}
\stepm{②}{由①得 $A_k\subseteq B$ 且 $|A_k|=|A|=|B|$，$(B^2)_{-2}\subseteq A$ 且 $|(B^2)_{-2}|=|B|=|A|$，}
\step{所以 $A_k=B$，$(B^2)_{-2}=A$，}
\step{所以 $B=A_k=\{a+k\mid a\in A\}$，$A=(B^2)_{-2}=\{b^2-2\mid b\in B\}$，}
\step{所以 $A=\{(a+k)^2-2\mid a\in A\}$，}
\step{所以映射 $h(a)=(a+k)^2-2$ 是集合 $A$ 到自身的一个双射.}
\step{假设 $k$ 是奇数，当 $a$ 为偶数时，$a+k$ 为奇数，$(a+k)^2-2$ 为奇数，}
\step{当 $a$ 为奇数时，$a+k$ 为偶数，$(a+k)^2-2$ 为偶数，}
\step{所以此时映射 $h(a)=(a+k)^2-2$，将集合 $A$ 中的奇数映射为偶数，}
\step{将集合 $A$ 中的偶数映射为奇数，即集合 $A$ 中的奇数元素和偶数元素必须配对，}
\step{此时集合 $A$ 的元素个数则为偶数，}
\step{与题设中“集合 $A$ 中的元素个数为奇数”矛盾.}
\step{故假设不成立，所以 $k$ 是偶数.}
\solend
\fi

\ansspace*[10]

\end{enumerate}

\end{document}
"
%
% 原始材料：微信公众号“魔都数学加油站”文章，作者破晓，连载“27学年高一67”；
% 原文标题“27学年高一67——2026-2027年上海市华师大一附中高一上10月月考”；
% 原文链接：https://mp.weixin.qq.com/s/Kdhx07TVjLoGTcIlnRzgcQ。
% 页面标注发布时间 2026-10-18 20:33（上海）。
% 本次转录只据 `原图/高一67_华师大一附中_10月月考/` 的 13 张 PNG 逐页读图。
% 正文为 13 张 PNG 页（`原图/高一67_华师大一附中_10月月考/01..13.png`）：
% 第 01、04、13 页 2560×3620，第 02、03、05—12 页 4762×6735
% （**同一篇各页分辨率不同**），均按页序读取。13 页全为试卷正文页（卷首标题在第 1 页页顶，
% 卷末解答题第 21 题解析止于第 13 页），无封面页、目录页、二维码或宣传页。
% 粉色斜排“魔都数学加油站”为卷面水印，与公众号页脚一并未录入。
%
% 卷首标题：照卷面印作“2026--2027 年华师大一附中高一上 10 月月考”（卷面为黑体加粗一行），
% 年份区间按全稿惯例排 en dash。原卷只有一行标题、无副标题。
%
% 篇幅与题号：原卷自第 3 题起编号（第 1、2 题不在本卷内）。填空题 3—12（10 题）、
% 选择题 13—16（4 题）、解答题 17—21（5 题），共 19 题。本稿沿用原节与原题号，
% 只改排为全稿统一的大节样式（原卷小节标题为“一、填空题”“二、选择题”“三、解答题”）。
%
% 讲评版：原卷每题下直接印【解析】。**全卷 19 题中 18 题都印【解析】（无一处只印
% 【答案】）**，按全稿口径这类题不另提 \ans{}；各题解析末句的结论（第 13、14、16 题
% 的“故选 B／D／B”、各填空与解答题的“故/综上…”）照录在解析内。
% **另有 1 题答案印在题干括号内**——第 15 题（“(  D  )”，且该题无【答案】/【解析】段），
% 本稿按 高一37 第 13 题、高一63 第 14 题成例把括号留空、答案另提为 \ans{$D$}。
%
% 副标题：原卷未印副标题。本稿按“知识模块”口径标作“集合与常用逻辑用语、不等式”：
% 19 题中集合与常用逻辑用语 12 题（第 3、4、5、7、10、11、12、13、15、16、18、21 题：
% 补集、并集、子集、充分条件、容斥计数、子集计数、新定义二元运算、条件判断、
% 反证法假设、新定义生成集、集合关系与二次方程、新定义 X_k 与 X^2），
% 不等式 7 题（第 6、8、9、14、17、19、20 题：存在量词不等式、解集反求系数、
% 解集为空、不等式性质、解不等式（组）、二次不等式与基本不等式、二次方程根的分布），
% 两类均超过三分之一，故并标。
%
% 与原卷的排版差异：
%   ① 原卷每题下直接印【解析】，本稿按全稿惯例将解析置于答案版，试卷版保持干净；
%   ② 原卷未印副标题，本稿按“知识模块”口径补出；
%   ③ 原卷填空横线约四个字宽，本稿按全稿惯例统一排 \blankline[4.4em]；
%   ④ 原卷实数集排斜体 $R$（第 14 题题干、第 17 题 (1) 题干、第 18 题解析、第 20 题题干），
%      整数集排斜体 $Z$（第 20 题 (3) 题干与解析、第 21 题题干），正整数集排粗体
%      $N^*$（第 11 题解析、第 19 题解析），本稿按全稿惯例统一改用 $\R$、$\Z$、$\N^*$
%      （记号归一，粗体／斜体不作区分）；空集原卷作 $\varnothing$，照录；
%   ⑤ 原卷在数学式内多处用半角逗号（如第 8 题 $(1,2)$ 与各 cases、第 11 题解析
%      $P_n=\{1,2,\cdots,n\}$,$n\in N^*$、第 12 题解析 $u=(\cdots)$,$v=(\cdots)$ 与
%      “2020,个 1”、第 19 题解析 $\{x|40<x<80,x\in N^*\}$），均照录未强改；
%      描述法集合的竖线照原卷排 \mid；句末句点按全稿惯例排西文句点；
%   ⑥ 原卷补集记号为上划线 $\bar A$（第 3 题题干与解析、第 11 题题干），照录
%      （同 高一64/65 的 $\bar B$ 成例）；
%   ⑦ 第 12 题新定义的运算记号 $*$ 照录（$\alpha*\beta$、$u*u$、$u*v$）；
%   ⑧ 第 17—21 题的子问原卷各自另起一行，本稿按全稿惯例用 enumerate 的 (1)(2)(3)；
%      解析中的子问标号一律用 \stepm{(1)}{…}；第 5 题解析的①②、第 21 题 (3) 解析的①②
%      按全稿惯例作 \stepm{①}{…}、\stepm{②}{…}（第 21 题 (3) 首行原卷“(3)”与“①”连排，
%      本稿合并记作 \stepm{(3)①}{…}）；
%   ⑨ 选择题作答括号原卷排半角“(　　)”（第 13、14、16 题括号内空白、第 15 题括号内印 D），
%      本稿按全稿惯例排 \mbox{（\quad）}；第 15 题括号内的答案移入 \ans{}，见上；
%   ⑩ 选项行排照原卷：第 13、15 题两行两条，第 14、16 题一行四条；
%   ⑪ 每道解答题原卷未留作答空白，本稿按全稿惯例加 \ansspace（末题用 \ansspace*）；
%      第 20、21 题题干长，各用一处 \ansneed 让题干与作答区不拆页（\ansneed 只作用于试卷版）。
%
% 原卷笔误/存疑（均照录未改，正文对应处另有 % 注）：
%   ① 第 3 题解析末句印作“则 $\bar A=(-\infty,0)\cup[2,3)$.”：按题干
%      $U=\{x\mid 1/x\leqslant 1/2\}$（即 $x<0$ 或 $x\geqslant2$）与 $A=[3,+\infty)$，
%      $\bar A$ 应为 $(-\infty,0)\cup\{2\}$；且与本解析首句自求之 $U$ 相矛盾（$3\notin U$），
%      疑为笔误；照录未改。
%   ② 第 12 题题干 $\alpha*\beta$ 定义式括号不配对：第四组“$+\cdots+x_{2020}+y_{2020}
%      -x_{2020}y_{2020}$”缺左括号、末尾连排两个右括号“))”；照录未改。
%   ③ 第 12 题解析“所以…中有 2020,个 1”一行，“$\cdots$”后脱逗号（作
%      “$x_2,\cdots x_{2020}$”，第 12 行“而…”同式样），且“2020”后衍一半角逗号
%      （作“2020,个 1”）；照录未改。
%   ④ 第 16 题解析乘积链印作“$a_1a_2<a_1a_3<a_1a_4<a_1a_5<a_2a_3<a_3a_3<a_{14}a_5$”：
%      “$a_3a_3$”同足标自乘、与单调链不合，“$a_{14}a_5$”下标 14 与 $a_1$—$a_5$ 记号
%      不合（疑为 $a_2a_4$、$a_4a_5$ 之误）；照录未改。
%   ⑤ 第 5 题题干作“$Q\subset P$”（真子集记号），解析两处均作“$Q\subseteq P$”，
%      同一题前后记号不一；照录（本稿分别排 \subset、\subseteq）。
%   ⑥ 第 17 题 (1) 解析“当 $a=1$ 时，原不等式为 $(x+1)^2<0$ 此时…”与
%      “当 $0<a<1$ 时，$-\frac1a<-1$ 此时…”两处“此时”前均脱标点；照录未改。
%   ⑦ 第 20 题 (2) 题干“若 $a=1$，$b=2k-1$，$c=k^2$，若方程 $f(x)=0$ 有两个大于 1 的根”
%      两个“若”连用，疑衍其一；照录未改。
%
% 插图：无（全卷 13 页无任何几何图形、无“如图/右图/图中”字样）。
%
% 行数核对：逐题数 \step/\stepm 条数与原卷【解析】印刷行数，18 题逐一相符
% （第 15 题无解析段，只提 \ans{}），合计 169 条 / 169 行
% （Q3—Q12：2、3、4、3、1、3、5、3、13、15；Q13—Q16：2、3、0、25；
% Q17—Q21：19、16、9、20、23）。
%
\documentclass[a4paper,11pt]{article}
\usepackage{common}

\begin{document}

\examtitlest[3]{2026--2027 年华师大一附中高一上 10 月月考}{集合与常用逻辑用语、不等式}

\bigsec{一、填空题}

\begin{enumerate}
\setcounter{enumi}{2}

\item 设全集 $U=\left\{x\mid\dfrac1x\leqslant\dfrac12\right\}$，$A=[3,+\infty)$，则 $\bar A=$\blankline[4.4em].

\ifanswer
\solbegin
\step{$\dfrac1x\leqslant\dfrac12$，解得 $x<0$ 或 $x\geqslant2$，故 $U=(-\infty,0)\cup[2,+\infty)$，}
% 原卷末句作 $\bar A=(-\infty,0)\cup[2,3)$，按题干 $U$ 应为 $(-\infty,0)\cup\{2\}$，与本解析首句矛盾，疑笔误；照录（见文件头“原卷笔误/存疑”①）
\step{$A=[3,+\infty)$，则 $\bar A=(-\infty,0)\cup[2,3)$.}
\solend
\fi

\item 已知集合 $M=\{y\mid y=x^2-2x-1\}$，$N=\{y\mid y=x^2-4x+1\}$，则 $M\cup N=$\blankline[4.4em].

\ifanswer
\solbegin
\step{$M=\{y\mid y=x^2-2x-1\}=\{y\mid y=(x-1)^2-2\}=\{y\mid y\geqslant-2\}$，}
\step{$N=\{y\mid y=x^2-4x+1\}=\{y\mid y=(x-2)^2-3\}=\{y\mid y\geqslant-3\}$，}
\step{则 $M\cup N=\{y\mid y\geqslant-3\}$.}
\solend
\fi

\item 已知 $P=\{x\mid x^2+x-6=0\}$，$Q=\{x\mid mx-1=0\}$，且 $Q\subset P$，则实数 $m$ 的值\blankline[4.4em].

\ifanswer
\solbegin
\step{$P=\{2,-3\}$，$Q=\{x\mid mx=1\}$；}
\stepm{①}{$m=0$ 时，$Q=\varnothing$，满足 $Q\subseteq P$；}
\stepm{②}{$m\neq0$ 时，要使 $Q\subseteq P$，则 $\dfrac1m=2$ 或 $-3$，所以 $m=\dfrac12$ 或 $-\dfrac13$；}
\step{所以 $m=0$ 或 $\dfrac12$ 或 $-\dfrac13$.}
\solend
\fi

\item 已知“存在正实数 $x$，使得关于 $x$ 的不等式 $ax-1>0$ 成立”为真命题，则实数 $a$ 的取值范围是\blankline[4.4em].

\ifanswer
\solbegin
\step{存在正实数 $x$，使得关于 $x$ 的不等式 $ax-1>0$ 成立为真命题，}
\step{即存在正实数 $x$，使得 $a>\dfrac1x$，因为 $\dfrac1x>0$，所以 $a>0$，}
\step{即 $a$ 的范围为 $(0,+\infty)$.}
\solend
\fi

\item 设 $\alpha:3\leqslant x\leqslant5$，$\beta:2m-1\leqslant x<m+4$，若 $\alpha$ 是 $\beta$ 的充分条件，则实数 $m$ 的取值范围是\blankline[4.4em].

\ifanswer
\solbegin
\step{若 $\alpha$ 是 $\beta$ 的充分条件，则 $\begin{cases}2m-1\leqslant3\\ m+4>5\end{cases}$，解得 $1<m\leqslant2$.}
\solend
\fi

\item 对于实数 $a$、$b$、$c$，已知关于 $x$ 的不等式 $ax^2+bx+c>0$ 的解集是 $(1,2)$，则关于 $x$ 的不等式 $cx^2-bx+a\leqslant0$ 的解集是\blankline[4.4em].

\ifanswer
\solbegin
\step{关于 $x$ 的不等式 $ax^2+bx+c>0$ 的解集是 $(1,2)$，则 $\begin{cases}a<0\\ -\dfrac ba=3\\ \dfrac ca=2\end{cases}$，则 $\begin{cases}b=-3a\\ c=2a\end{cases}$，}
\step{则关于 $x$ 的不等式 $cx^2-bx+a=2ax^2+3ax+a\leqslant0$，则 $2x^2+3x+1\geqslant0$，}
\step{则 $x\leqslant-1$ 或 $x\geqslant-\dfrac12$，则解集为 $(-\infty,-1]\cup\left[-\dfrac12,+\infty\right)$.}
\solend
\fi

\item 已知关于 $x$ 的不等式 $kx^2+kx+k+2\leqslant0$ 的解集为 $\varnothing$，则实数 $k$ 的取值范围为\blankline[4.4em].

\ifanswer
\solbegin
\step{当 $k=0$ 时，不等式化为 $2\leqslant0$，解集为 $\varnothing$，符合题意；}
\step{当 $k\neq0$ 时，若关于 $x$ 的不等式 $kx^2+kx+k+2\leqslant0$ 的解集为 $\varnothing$，}
\step{则 $\begin{cases}k>0\\ \Delta=k^2-4k(k+2)<0\end{cases}$，化简得 $\begin{cases}k>0\\ -3k^2-8k<0\end{cases}$，即 $\begin{cases}k>0\\ k(3k+8)>0\end{cases}$，}
\step{解得 $k>0$.}
\step{综上所述，$k\geqslant0$，即实数 $k$ 的取值范围为 $[0,+\infty)$.}
\solend
\fi

\item 对班级若干名学生调查对 $A$、$B$ 两事件的态度，有如下结果：赞成 $A$ 的人数是全体的五分之三，其余的不赞成；赞成 $B$ 的比赞成 $A$ 多 3 人，其余的不赞成. 已知对 $A$、$B$ 都赞成的学生有 18 人，对 $A$、$B$ 都不赞成的学生数比对 $A$、$B$ 都赞成的学生数的三分之一多 1 人，则班级总人数是\blankline[4.4em]人.

\ifanswer
\solbegin
\step{设班级总人数是 $x$ 人，则赞成 $A$ 的人数是 $\dfrac35x$，赞成 $B$ 的人数是 $\dfrac35x+3$，}
\step{$A$、$B$ 都不赞成的人数为 $18\times\dfrac13+1=7$ 人，}
\step{由容斥原理得 $x=\left(\dfrac35x+\dfrac35x+3-18\right)+7$，解得 $x=40$，即班级总人数是 40 人.}
\solend
\fi

\item 设全集 $U=\{1,2,3,4,5,6,7,8,9,10,11,12\}$，若集合 $A$ 满足：① $A\subseteq U$；②若 $x\in A$，必有 $2x\notin A$；③若 $x\in\bar A$，必有 $2x\notin\bar A$，则这样的集合 $A$ 有\blankline[4.4em]个.

\ifanswer
\solbegin
\step{设集合 $P_n=\{1,2,\cdots,n\}$,$n\in\N^*$，}
\step{任取 $x\in P_n$，将 $x$ 除以 2，若商仍为偶数，再除以 2，$\cdots$，}
\step{经过 $k$ 次以后，商必为奇数，此时记商为 $m$，}
\step{于是 $x=m\cdot2^k$，其中 $m$ 为奇数，$k\in\N^*$.}
\step{由条件得，若 $m\in A$，则 $x\in A$ 等价于 $k$ 为偶数；}
\step{若 $m\notin A$，则 $x\in A$ 等价于 $k$ 为奇数；}
\step{于是 $x$ 是否属于 $A$ 由 $m$ 是否属于 $A$ 确定；}
\step{设 $Q_n$ 是 $P_n$ 中所有奇数的集合，}
\step{设 $f(n)$ 为集合 $P_n$ 中满足题干条件的集合 $A$ 的个数，}
\step{因此 $f(n)$ 等于 $Q_n$ 的子集个数.}
\step{当 $n$ 为偶数（或奇数）时，$P_n$ 中奇数的个数是 $\dfrac n2$（或 $\dfrac{n+1}2$），}
\step{所以 $f(n)=\begin{cases}2^{\frac n2},n\text{为偶数}\\[2pt] 2^{\frac{n+1}2},n\text{为奇数}\end{cases}$，所以 $P_{12}=\{1,2,3,4,5,6,7,8,9,10,11,12\}$，}
\step{即同时满足三个条件的集合 $A$ 的个数为 $2^{\frac{12}2}=2^6=64$.}
\solend
\fi

\item 设集合 $\Omega=\{(x_1,x_2,\cdots,x_{2020})\mid x_i\in\{0,1\},i=1,2,\cdots,2020\}$，对于 $\Omega$ 中的任意两个元素 $\alpha=(x_1,x_2,\cdots,x_{2020})$，$\beta=(y_1,y_2,\cdots,y_{2020})$，记
% 原卷此式第四组缺左括号、末尾连排两个右括号“))”，括号不配对；照录（见文件头“原卷笔误/存疑”②）
$\alpha*\beta=(x_1+y_1-x_1y_1)+(x_2+y_2-x_2y_2)+\cdots+x_{2020}+y_{2020}-x_{2020}y_{2020}))$，设 $u$、$v\in\Omega$，若 $u*u+v*v=2020$，则 $u*v$ 的最大值和最小值之和为\blankline[4.4em].

\ifanswer
\solbegin
\step{设 $u=(x_1,x_2,\cdots,x_{2020})$,$v=(y_1,y_2,\cdots,y_{2020})$，}
\step{因为 $x_i\in\{0,1\}$，所以 $x_i(x_i-1)=0$，}
\step{所以 $u*u=(x_1+x_1-x_1^2)+(x_2+x_2-x_2^2)+\cdots+(x_{2020}+x_{2020}-x_{2020}^2)$}
\step{$=x_1+x_2+\cdots+x_{2020}$，}
\step{$v*v=(y_1+y_1-y_1^2)+(y_2+y_2-y_2^2)+\cdots+(y_{2020}+y_{2020}-y_{2020}^2)$}
\step{$=y_1+y_2+\cdots+y_{2020}$，}
\step{由 $u*u+v*v=2020$，得 $x_1+x_2+\cdots+x_{2020}+y_1+y_2+\cdots+y_{2020}=2020$，}
% 原卷此行“$\cdots$”后脱逗号、 “2020”后衍一半角逗号（“2020,个 1”）；照录（见文件头“原卷笔误/存疑”③）
\step{所以 $x_1,x_2,\cdots x_{2020},y_1,y_2,\cdots y_{2020}$ 中有 2020,个 1，2020 个 0，}
\step{$u*v=(x_1+y_1-x_1y_1)+(x_2+y_2-x_2y_2)+\cdots+(x_{2020}+y_{2020}-x_{2020}y_{2020})$}
\step{$=2020-(x_1y_1+x_2y_2+\cdots+x_{2020}y_{2020})$，}
\step{注意到只有 $x_i=y_i=1$ 时，$x_iy_i=1$，否则 $x_iy_i=0$，}
% 此行“$\cdots$”后同样脱逗号，与上一处同式样；照录（见文件头“原卷笔误/存疑”③）
\step{而 $x_1,x_2,\cdots x_{2020},y_1,y_2,\cdots y_{2020}$ 中有 2020 个 1，2020 个 0，}
\step{所以满足 $x_iy_i=1$ 的最多有 1010 个，所以 $u*v\geqslant2020-1010=1010$，}
\step{满足 $x_iy_i=0$ 的最多有 2020 个，所以 $u*v\leqslant2020-0=2020$，}
\step{则 $u*v$ 的最大值和最小值之和为 $2020+1010=3030$.}
\solend
\fi

\end{enumerate}

\bigsec{二、选择题}

\begin{enumerate}
\setcounter{enumi}{12}

\item 实数 $x$ 与 $y$，“$x+y>2$”的\mbox{（\quad）}条件是“$x>1$ 或 $y>1$”.

\keepwithprev
A. 充分非必要\qquad B. 必要非充分

C. 充要\qquad D. 既非充分又非必要

\ifanswer
\solbegin
\step{当 $x>1$ 或 $y>1$ 时，$x+y>2$ 不一定成立，充分性不成立；}
\step{当 $x+y>2$ 时，$x>1$ 或 $y>1$ 一定成立，必要性成立. 故选 $B$.}
\solend
\fi

\item 设 $a,b\in\R$，且 $a<0$，$-1<b<0$，那么下列不等式成立的是\mbox{（\quad）}

\keepwithprev
A. $a>ab>ab^2$\qquad B. $ab^2>ab>a$\qquad C. $ab>a>ab^2$\qquad D. $ab>ab^2>a$

\ifanswer
\solbegin
\step{因为 $-1<b<0$，$a<0$，所以 $ab>0$，$0<b^2<1$.}
\step{所以 $ab-ab^2=ab(1-b)>0$，$ab^2-a=a(b^2-1)>0$，}
\step{所以 $ab>ab^2>a$. 故选 $D$.}
\solend
\fi

% 第 15 题答案原卷印在题干括号内（“(  D  )”），且无【答案】/【解析】段；
% 按 高一37 第 13 题、高一63 第 14 题成例括号留空、答案另提 \ans{$D$}
\item 用反证法证明：“平面四边形中至少有一个内角不超过 90°”时，下列“假设”中正确的是\mbox{（\quad）}

\keepwithprev
A. 假设有两个内角超过 90°\qquad B. 假设有三个内角超过 90°

C. 假设至多有两个内角超过 90°\qquad D. 假设四个内角都超过 90°

\ifanswer
\ans{$D$}
\fi

\item 设 $A$ 为非空实数集，集合 $B=\{xy\mid x\in A,y\in A,x\neq y\}$，则对于以下两个结论说法正确的是\mbox{（\quad）}

结论①：若 $A$ 是五元集，则 $B$ 中的元素最多 10 个；

结论②：存在四元集 $A$，使得 $B=\{2,4,5,8,10,16\}$.

\keepwithprev
A. ①对②对\qquad B. ①对②错\qquad C. ①错②对\qquad D. ①错②错

\ifanswer
\solbegin
\step{对于①，想让元素个数达到最大值 10，则须避免产生符号不同的乘积，}
\step{否则正负分开计数容易导致总数受限，且容易产生相等值，}
\step{也最好避免 0 的出现，}
\step{故设 $A=\{a_1,a_2,a_3,a_4,a_5\}$，$0<a_1<a_2<a_3<a_4<a_5$.}
% 原卷此链“$a_3a_3$”同足标自乘、“$a_{14}a_5$”下标 14 均与文意不合，疑为 $a_2a_4$、$a_4a_5$ 之误；照录（见文件头“原卷笔误/存疑”④）
\step{因为 $a_1a_2<a_1a_3<a_1a_4<a_1a_5<a_2a_3<a_3a_3<a_{14}a_5$，$a_2a_3<a_2a_4<a_3a_4$，}
\step{所以 $B$ 中元素个数小于等于 10 个，}
\step{如设 $A=\{1,2,3,5,7\}$，则 $B=\{2,3,5,6,7,10,14,15,21,35\}$，}
\step{所以集合 $B$ 中元素个数的最大值为 10 个，故①正确；}
\step{对于②，由于集合 $B$ 中没有 0，所以集合 $A$ 中没有 0；}
\step{设集合 $A=\{a,b,c,d\}$，且 $a<b<c<d$，}
\step{当 $a<0<b<c<d$ 时，$a\times b<0$，$a\times c<0$，$a\times d<0$，}
\step{故集合 $B$ 至少有 3 个负数，与 $B=\{2,4,5,8,10,16\}$ 矛盾；}
\step{当 $a<b<0<c<d$ 时，$a\times b>0$，$c\times d>0$，}
\step{$a\times c<0$，$a\times d<0$，$b\times c<0$，$b\times d<0$，}
\step{故集合 $B$ 至少有 4 个负数，与 $B=\{2,4,5,8,10,16\}$ 矛盾；}
\step{当 $a<b<c<0<d$ 时，$a\times b>0$，$a\times c>0$，$b\times c>0$，}
\step{$a\times d<0$，$b\times d<0$，$c\times d<0$，}
\step{故集合 $B$ 至少有 3 个负数，与 $B=\{2,4,5,8,10,16\}$ 矛盾；}
\step{当 $0<a<b<c<d$ 时，则集合 $B=\{ab,ac,ad,bc,bd,cd\}$，}
\step{则必有 $ab=2$，$cd=16$，其 4 个正实数的乘积 $abcd=32$，}
\step{又 $ad=5$，$bc=8$ 或 $ad=8$，$bc=5$，其 4 个正实数的乘积 $abcd=40$，矛盾；}
\step{当 $a<b<c<d<0$ 时，与 $0<a<b<c<d$ 时的证明方法相同，}
\step{得集合 $A$ 也不存在；}
\step{综上不存在四元集 $A$，使其生成集 $B=\{2,4,5,8,10,16\}$，故②错误.}
\step{故选 $B$.}
\solend
\fi

\end{enumerate}

\bigsec{三、解答题}

\begin{enumerate}
\setcounter{enumi}{16}

\item 解下列不等式（组）

\begin{enumerate}[label=(\arabic*)]
\item 解不等式：$(ax+1)(x+1)<0(a\in\R)$.
\item 解不等式组：$\begin{cases}\dfrac{2x+4}{x+4}\geqslant1\\[4pt] (x+1)(2-|x|)\geqslant0\end{cases}$.
\end{enumerate}

\ifanswer
\solbegin
\stepm{(1)}{当 $a=0$ 时，原不等式为 $x+1<0$，即 $x<-1$，}
\step{所以此时不等式的解集为 $\{x\mid x<-1\}$；}
\step{当 $a\neq0$ 时，令 $(ax+1)(x+1)=0$，解得 $x_1=-\dfrac1a$,$x_2=-1$，}
\step{当 $a>1$ 时，$-\dfrac1a>-1$，此时不等式的解集为 $\left\{x\mid-1<x<-\dfrac1a\right\}$；}
% 以下两行“此时”前原卷均无标点；照录（见文件头“原卷笔误/存疑”⑥）
\step{当 $a=1$ 时，原不等式为 $(x+1)^2<0$ 此时不等式的解集为 $\varnothing$；}
\step{当 $0<a<1$ 时，$-\dfrac1a<-1$ 此时不等式的解集为 $\left\{x\mid-\dfrac1a<x<-1\right\}$；}
\step{当 $a<0$ 时，$-\dfrac1a>-1$，此时不等式的解集为 $\left\{x\mid x<-1\ \text{或}\ x>-\dfrac1a\right\}$.}
\step{综上所述，当 $a=0$ 时，$\{x\mid x<-1\}$；}
\step{当 $a>1$ 时，$\left\{x\mid-1<x<-\dfrac1a\right\}$；}
\step{当 $a=1$ 时，$\varnothing$；}
\step{当 $0<a<1$ 时，$\left\{x\mid-\dfrac1a<x<-1\right\}$；}
\step{当 $a<0$ 时，$\left\{x\mid x<-1\ \text{或}\ x>-\dfrac1a\right\}$.}
\stepm{(2)}{不等式 $\dfrac{2x+4}{x+4}\geqslant1$ 可化为 $\dfrac{2x+4}{x+4}-1\geqslant0$，即 $\dfrac x{x+4}\geqslant0$，}
\step{等价于 $\begin{cases}x(x+4)\geqslant0\\ x+4\neq0\end{cases}$，解得 $x<-4$ 或 $x\geqslant0$，}
\step{不等式 $(x+1)(2-|x|)\geqslant0$ 等价于 $\begin{cases}(x+1)(2-x)\geqslant0,x\geqslant0\\ (x+1)(2+x)\geqslant0,x<0\end{cases}$，}
\step{解不等式 $(x+1)(2-x)\geqslant0$，$x\geqslant0$，得 $0\leqslant x\leqslant2$，}
\step{解不等式 $(x+1)(2+x)\geqslant0$，$x<0$，得 $x\leqslant-2$ 或 $-1\leqslant x<0$，}
\step{所以 $(x+1)(2-|x|)\geqslant0$ 的解集为 $\{x\mid x\leqslant-2\ \text{或}\ -1\leqslant x\leqslant2\}$，}
\step{综上所述，$\begin{cases}\dfrac{2x+4}{x+4}\geqslant1\\[4pt] (x+1)(2-|x|)\geqslant0\end{cases}$ 的解集为 $\{x\mid x<-4\ \text{或}\ 0\leqslant x\leqslant2\}$.}
\solend
\fi

\ansspace[8]

% 第 18 题题干＋两小问与作答区 6 行须同页
\ansneed{20cm}
\item 已知 $a$ 为实数，设 $f(x)=x^2+2(a+1)x+1$.

\begin{enumerate}[label=(\arabic*)]
\item 已知对于任意实数 $x$，$\sqrt{f(x)}$ 均有意义，求实数 $a$ 的取值范围；
\item 设 $A=\{x\mid f(x)+a^2-2=0\}$，$B=\{-4,0\}$，$A\cup B=B$，求实数 $a$ 的取值范围.
\end{enumerate}

\ifanswer
\solbegin
\stepm{(1)}{$\sqrt{f(x)}$ 对任意实数 $x$ 都有意义，}
\step{即 $f(x)=x^2+2(a+1)x+1\geqslant0$ 对任意 $x\in\R$ 恒成立，}
\step{所以 $\Delta=4(a+1)^2-4\times1\times1\leqslant0$，即 $(a+1)^2\leqslant1$，解得 $-2\leqslant a\leqslant0$，}
\step{故 $a$ 的取值范围是 $[-2,0]$.}
\stepm{(2)}{由 $f(x)+a^2-2=0$ 得 $x^2+2(a+1)x+a^2-1=0$.}
\step{由 $A\cup B=B$ 得 $A\subseteq B=\{-4,0\}$，即该方程无实根，}
\step{或它的实根都在 $\{-4,0\}$ 中.}
\step{当方程无实根时，$\Delta=4(a+1)^2-4(a^2-1)=8a+8<0$，解得 $a<-1$；}
\step{当方程有实根时，$a\geqslant-1$，设两根为 $x_1$、$x_2$，}
\step{由根与系数的关系得 $x_1+x_2=-2(a+1)$、$x_1x_2=a^2-1$，}
\step{此时两根只能在 $\{-4,0\}$ 中取值；}
\step{若两根为 0 与 0，由 $x_1x_2=0$ 得 $a=\pm1$，}
\step{再由 $x_1+x_2=0$ 得 $a=-1$；}
\step{若两根为 0 与 $-4$，由 $x_1x_2=0$ 得 $a=\pm1$，再由 $x_1+x_2=-4$ 得 $a=1$.}
\step{经检验，$a=-1$ 时 $A=\{0\}$，$a=1$ 时 $A=\{-4,0\}$，都满足 $A\subseteq B$.}
\step{综上，实数 $a$ 的取值范围为 $(-\infty,-1]\cup\{1\}$.}
\solend
\fi

\ansspace[8]

\ansneed{18cm}
\item 共享单车给市民出行带来了诸多便利，某公司购买了一批单车投放到某地给居民使用，据市场分析，这批单车的运营累计利润 $y$（单位：元）与运营天数 $x$（单位：天，$x$ 为正整数）满足关系：$y=-\dfrac12x^2+60x-800$.

\begin{enumerate}[label=(\arabic*)]
\item 要使这批单车的运营累计利润高于 800 元，求运营天数的取值范围；
\item 判断这批单车的日平均运营利润能否超过 20 元，并说明理由.
\end{enumerate}

\ifanswer
\solbegin
\stepm{(1)}{由题意得 $y=-\dfrac12x^2+60x-800>800$，}
\step{即 $x^2-120x+3200<0$，解得 $40<x<80$，}
\step{即运营天数的取值范围为 $\{x\mid40<x<80,x\in\N^*\}$.}
\stepm{(2)}{不能，理由如下：}
\step{这批单车的日平均运营利润为 $\dfrac yx=-\dfrac12x-\dfrac{800}x+60=-\dfrac12\left(x+\dfrac{1600}x\right)+60$，}
\step{因 $x>0$，则 $x+\dfrac{1600}x\geqslant2\sqrt{x\cdot\dfrac{1600}x}=80$，}
\step{当且仅当 $x=\dfrac{1600}x$，即 $x=40$ 时取等号，}
\step{此时 $\dfrac yx=-\dfrac12\left(x+\dfrac{1600}x\right)+60\leqslant-\dfrac12\times80+60=20$，}
\step{故这批单车的日平均运营利润不能超过 20 元.}
\solend
\fi

\ansspace[8]

% 第 20 题题干＋三小问与作答区 8 行须同页
\ansneed{20cm}
\item 设 $f(x)=ax^2+bx+c$($a\neq0$，$b$、$c\in\R$).

\begin{enumerate}[label=(\arabic*)]
\item 若 $a=b=c=1$，设 $A=\{(x,y)\mid y=f(x)\}$，$B=\{(x,y)\mid y=1\}$，用列举法表示 $A\cap B$；
% 原卷两个“若”连用（“若 a=1，…，若方程 f(x)=0 有两个大于 1 的根”），疑衍其一；照录（见文件头“原卷笔误/存疑”⑦）
\item 若 $a=1$，$b=2k-1$，$c=k^2$，若方程 $f(x)=0$ 有两个大于 1 的根，求实数 $k$ 的取值范围；
\item 若 $a$、$b$、$c\in\Z$，且 $f(0)$、$f(1)$ 均为奇数，求证：方程 $f(x)=0$ 无整数解.
\end{enumerate}

\ifanswer
\solbegin
\stepm{(1)}{当 $a=b=c=1$ 时，$f(x)=x^2+x+1$，}
\step{所以 $A=\{(x,y)\mid y=x^2+x+1\}$，又 $y=1$，}
\step{由 $x^2+x+1=1$ 得 $x^2+x=0$，解得 $x=0$ 或 $x=-1$，}
\step{所以 $A\cap B=\{(0,1),(-1,1)\}$.}
\stepm{(2)}{当 $a=1$、$b=2k-1$、$c=k^2$ 时，$f(x)=x^2+(2k-1)x+k^2$.}
\step{设两根为 $x_1$、$x_2$，}
\step{则“两根都大于 1”等价于 $\Delta>0$、$f(1)>0$、$-\dfrac b{2a}>1$.}
\step{由 $\Delta=(2k-1)^2-4k^2=1-4k>0$ 得 $k<\dfrac14$；}
\step{由 $f(1)=1+(2k-1)+k^2=k^2+2k=k(k+2)>0$ 得 $k<-2$ 或 $k>0$；}
\step{由 $-\dfrac{2k-1}2=\dfrac{1-2k}2>1$ 得 $k<-\dfrac12$.}
\step{三者取交集得 $k<-2$，即 $k$ 的取值范围是 $(-\infty,-2)$.}
\stepm{(3)}{设 $x_0\in\Z$ 是方程 $f(x)=0$ 的整数解.}
\step{因为 $x_0(x_0-1)$ 是两个相邻整数之积，恒为偶数，所以 $ax_0(x_0-1)$ 为偶数.}
\step{而 $f(x_0)=ax_0^2+bx_0+c=ax_0(x_0-1)+(a+b)x_0+c$，}
\step{所以 $f(x_0)$ 与 $(a+b)x_0+c$ 同奇偶.}
\step{若 $x_0$ 为偶数，则 $(a+b)x_0$ 为偶数，$f(x_0)$ 与 $c=f(0)$ 同奇偶，为奇数；}
\step{若 $x_0$ 为奇数，则 $(a+b)x_0$ 与 $a+b$ 同奇偶，}
\step{$f(x_0)$ 与 $(a+b)+c=f(1)$ 同奇偶，也为奇数.}
\step{综上，$f(x_0)$ 恒为奇数，不可能等于 0，与 $f(x_0)=0$ 矛盾，}
\step{所以方程 $f(x)=0$ 无整数解.}
\solend
\fi

\ansspace[8]

% 压轴题（题干＋三小问）与作答区 10 行须同页
\ansneed{20cm}
\item 已知 $X$ 为有限集，且任取 $x\in X$，均有 $x\in\Z$，称 $X$ 为非空有限整数集. 定义：$X_k=\{x+k\mid x\in X\}$，$X^2=\{y\mid y=x^2,x\in X\}$.

\begin{enumerate}[label=(\arabic*)]
\item 若 $X=\{-2,2\}$，分别写出 $X_2$ 和 $X^2$（用列举法表示）；
\item 若 $X=\{-1,0,1\}$，$X_1\subseteq Y$，$(Y_{-1})^2\subseteq X$，求集合 $Y$ 中所有的元素；
\item 对于非空有限整数集 $A$、$B$，若 $A_k\subseteq B$，$(B^2)_{-2}\subseteq A$，$B$ 中的元素均为非负整数，$k$ 为整数，求证：

① $A$、$B$ 中元素的个数相同；

② 若集合 $A$ 中的元素个数为奇数，则 $k$ 为偶数.
\end{enumerate}

\ifanswer
\solbegin
\stepm{(1)}{由题意得 $X_2=\{0,4\}$，$X^2=\{4\}$；}
\stepm{(2)}{由题意得 $X_1=\{0,1,2\}$，即 $\{0,1,2\}\subseteq Y$，$Y_{-1}=\{y-1\mid y\in Y\}$，}
\step{$(Y_{-1})^2=\{(y-1)^2\mid y\in Y\}\subseteq X=\{-1,0,1\}$，}
\step{所以 $(y-1)^2=0$ 或 $(y-1)^2=1$，解得 $y=1$ 或 $y=2$ 或 $y=0$，}
\step{即 $y\in\{0,1,2\}$，即 $Y\subseteq\{0,1,2\}$，所以 $Y=\{0,1,2\}$；}
% 原卷 (3) 解析首行“(3)”与“①”连排，本稿合并记作 \stepm{(3)①}
\stepm{(3)①}{因为 $A_k=\{x+k\mid x\in A\}$，函数 $g(x)=x+k$ 在全体整数集上严格增，}
\step{所以 $\forall x_1$、$x_2\in A$，$x_1\neq x_2$，$g(x_1)\neq g(x_2)$，所以 $|A_k|=|A|$，}
\step{因为 $A_k\subseteq B$，所以 $|A|\leqslant|B|$，}
\step{设 $b\in B$，$B$ 中元素均为非负整数，$\forall b_1\neq b_2\in B$，$b_1^2\neq b_2^2$，}
\step{故函数 $f(b)=b^2-2$ 在 $B$ 上单射，所以 $|(B^2)_{-2}|=|B|$，}
\step{又因为 $(B^2)_{-2}\subseteq A$，所以 $|B|\leqslant|A|$，所以 $|A|=|B|$.}
\stepm{②}{由①得 $A_k\subseteq B$ 且 $|A_k|=|A|=|B|$，$(B^2)_{-2}\subseteq A$ 且 $|(B^2)_{-2}|=|B|=|A|$，}
\step{所以 $A_k=B$，$(B^2)_{-2}=A$，}
\step{所以 $B=A_k=\{a+k\mid a\in A\}$，$A=(B^2)_{-2}=\{b^2-2\mid b\in B\}$，}
\step{所以 $A=\{(a+k)^2-2\mid a\in A\}$，}
\step{所以映射 $h(a)=(a+k)^2-2$ 是集合 $A$ 到自身的一个双射.}
\step{假设 $k$ 是奇数，当 $a$ 为偶数时，$a+k$ 为奇数，$(a+k)^2-2$ 为奇数，}
\step{当 $a$ 为奇数时，$a+k$ 为偶数，$(a+k)^2-2$ 为偶数，}
\step{所以此时映射 $h(a)=(a+k)^2-2$，将集合 $A$ 中的奇数映射为偶数，}
\step{将集合 $A$ 中的偶数映射为奇数，即集合 $A$ 中的奇数元素和偶数元素必须配对，}
\step{此时集合 $A$ 的元素个数则为偶数，}
\step{与题设中“集合 $A$ 中的元素个数为奇数”矛盾.}
\step{故假设不成立，所以 $k$ 是偶数.}
\solend
\fi

\ansspace*[10]

\end{enumerate}

\end{document}
"
%
% 原始材料：微信公众号“魔都数学加油站”文章，作者破晓，连载“27学年高一67”；
% 原文标题“27学年高一67——2026-2027年上海市华师大一附中高一上10月月考”；
% 原文链接：https://mp.weixin.qq.com/s/Kdhx07TVjLoGTcIlnRzgcQ。
% 页面标注发布时间 2026-10-18 20:33（上海）。
% 本次转录只据 `原图/高一67_华师大一附中_10月月考/` 的 13 张 PNG 逐页读图。
% 正文为 13 张 PNG 页（`原图/高一67_华师大一附中_10月月考/01..13.png`）：
% 第 01、04、13 页 2560×3620，第 02、03、05—12 页 4762×6735
% （**同一篇各页分辨率不同**），均按页序读取。13 页全为试卷正文页（卷首标题在第 1 页页顶，
% 卷末解答题第 21 题解析止于第 13 页），无封面页、目录页、二维码或宣传页。
% 粉色斜排“魔都数学加油站”为卷面水印，与公众号页脚一并未录入。
%
% 卷首标题：照卷面印作“2026--2027 年华师大一附中高一上 10 月月考”（卷面为黑体加粗一行），
% 年份区间按全稿惯例排 en dash。原卷只有一行标题、无副标题。
%
% 篇幅与题号：原卷自第 3 题起编号（第 1、2 题不在本卷内）。填空题 3—12（10 题）、
% 选择题 13—16（4 题）、解答题 17—21（5 题），共 19 题。本稿沿用原节与原题号，
% 只改排为全稿统一的大节样式（原卷小节标题为“一、填空题”“二、选择题”“三、解答题”）。
%
% 讲评版：原卷每题下直接印【解析】。**全卷 19 题中 18 题都印【解析】（无一处只印
% 【答案】）**，按全稿口径这类题不另提 \ans{}；各题解析末句的结论（第 13、14、16 题
% 的“故选 B／D／B”、各填空与解答题的“故/综上…”）照录在解析内。
% **另有 1 题答案印在题干括号内**——第 15 题（“(  D  )”，且该题无【答案】/【解析】段），
% 本稿按 高一37 第 13 题、高一63 第 14 题成例把括号留空、答案另提为 \ans{$D$}。
%
% 副标题：原卷未印副标题。本稿按“知识模块”口径标作“集合与常用逻辑用语、不等式”：
% 19 题中集合与常用逻辑用语 12 题（第 3、4、5、7、10、11、12、13、15、16、18、21 题：
% 补集、并集、子集、充分条件、容斥计数、子集计数、新定义二元运算、条件判断、
% 反证法假设、新定义生成集、集合关系与二次方程、新定义 X_k 与 X^2），
% 不等式 7 题（第 6、8、9、14、17、19、20 题：存在量词不等式、解集反求系数、
% 解集为空、不等式性质、解不等式（组）、二次不等式与基本不等式、二次方程根的分布），
% 两类均超过三分之一，故并标。
%
% 与原卷的排版差异：
%   ① 原卷每题下直接印【解析】，本稿按全稿惯例将解析置于答案版，试卷版保持干净；
%   ② 原卷未印副标题，本稿按“知识模块”口径补出；
%   ③ 原卷填空横线约四个字宽，本稿按全稿惯例统一排 \blankline[4.4em]；
%   ④ 原卷实数集排斜体 $R$（第 14 题题干、第 17 题 (1) 题干、第 18 题解析、第 20 题题干），
%      整数集排斜体 $Z$（第 20 题 (3) 题干与解析、第 21 题题干），正整数集排粗体
%      $N^*$（第 11 题解析、第 19 题解析），本稿按全稿惯例统一改用 $\R$、$\Z$、$\N^*$
%      （记号归一，粗体／斜体不作区分）；空集原卷作 $\varnothing$，照录；
%   ⑤ 原卷在数学式内多处用半角逗号（如第 8 题 $(1,2)$ 与各 cases、第 11 题解析
%      $P_n=\{1,2,\cdots,n\}$,$n\in N^*$、第 12 题解析 $u=(\cdots)$,$v=(\cdots)$ 与
%      “2020,个 1”、第 19 题解析 $\{x|40<x<80,x\in N^*\}$），均照录未强改；
%      描述法集合的竖线照原卷排 \mid；句末句点按全稿惯例排西文句点；
%   ⑥ 原卷补集记号为上划线 $\bar A$（第 3 题题干与解析、第 11 题题干），照录
%      （同 高一64/65 的 $\bar B$ 成例）；
%   ⑦ 第 12 题新定义的运算记号 $*$ 照录（$\alpha*\beta$、$u*u$、$u*v$）；
%   ⑧ 第 17—21 题的子问原卷各自另起一行，本稿按全稿惯例用 enumerate 的 (1)(2)(3)；
%      解析中的子问标号一律用 \stepm{(1)}{…}；第 5 题解析的①②、第 21 题 (3) 解析的①②
%      按全稿惯例作 \stepm{①}{…}、\stepm{②}{…}（第 21 题 (3) 首行原卷“(3)”与“①”连排，
%      本稿合并记作 \stepm{(3)①}{…}）；
%   ⑨ 选择题作答括号原卷排半角“(　　)”（第 13、14、16 题括号内空白、第 15 题括号内印 D），
%      本稿按全稿惯例排 \mbox{（\quad）}；第 15 题括号内的答案移入 \ans{}，见上；
%   ⑩ 选项行排照原卷：第 13、15 题两行两条，第 14、16 题一行四条；
%   ⑪ 每道解答题原卷未留作答空白，本稿按全稿惯例加 \ansspace（末题用 \ansspace*）；
%      第 20、21 题题干长，各用一处 \ansneed 让题干与作答区不拆页（\ansneed 只作用于试卷版）。
%
% 原卷笔误/存疑（均照录未改，正文对应处另有 % 注）：
%   ① 第 3 题解析末句印作“则 $\bar A=(-\infty,0)\cup[2,3)$.”：按题干
%      $U=\{x\mid 1/x\leqslant 1/2\}$（即 $x<0$ 或 $x\geqslant2$）与 $A=[3,+\infty)$，
%      $\bar A$ 应为 $(-\infty,0)\cup\{2\}$；且与本解析首句自求之 $U$ 相矛盾（$3\notin U$），
%      疑为笔误；照录未改。
%   ② 第 12 题题干 $\alpha*\beta$ 定义式括号不配对：第四组“$+\cdots+x_{2020}+y_{2020}
%      -x_{2020}y_{2020}$”缺左括号、末尾连排两个右括号“))”；照录未改。
%   ③ 第 12 题解析“所以…中有 2020,个 1”一行，“$\cdots$”后脱逗号（作
%      “$x_2,\cdots x_{2020}$”，第 12 行“而…”同式样），且“2020”后衍一半角逗号
%      （作“2020,个 1”）；照录未改。
%   ④ 第 16 题解析乘积链印作“$a_1a_2<a_1a_3<a_1a_4<a_1a_5<a_2a_3<a_3a_3<a_{14}a_5$”：
%      “$a_3a_3$”同足标自乘、与单调链不合，“$a_{14}a_5$”下标 14 与 $a_1$—$a_5$ 记号
%      不合（疑为 $a_2a_4$、$a_4a_5$ 之误）；照录未改。
%   ⑤ 第 5 题题干作“$Q\subset P$”（真子集记号），解析两处均作“$Q\subseteq P$”，
%      同一题前后记号不一；照录（本稿分别排 \subset、\subseteq）。
%   ⑥ 第 17 题 (1) 解析“当 $a=1$ 时，原不等式为 $(x+1)^2<0$ 此时…”与
%      “当 $0<a<1$ 时，$-\frac1a<-1$ 此时…”两处“此时”前均脱标点；照录未改。
%   ⑦ 第 20 题 (2) 题干“若 $a=1$，$b=2k-1$，$c=k^2$，若方程 $f(x)=0$ 有两个大于 1 的根”
%      两个“若”连用，疑衍其一；照录未改。
%
% 插图：无（全卷 13 页无任何几何图形、无“如图/右图/图中”字样）。
%
% 行数核对：逐题数 \step/\stepm 条数与原卷【解析】印刷行数，18 题逐一相符
% （第 15 题无解析段，只提 \ans{}），合计 169 条 / 169 行
% （Q3—Q12：2、3、4、3、1、3、5、3、13、15；Q13—Q16：2、3、0、25；
% Q17—Q21：19、16、9、20、23）。
%
\documentclass[a4paper,11pt]{article}
\usepackage{common}

\begin{document}

\examtitlest[3]{2026--2027 年华师大一附中高一上 10 月月考}{集合与常用逻辑用语、不等式}

\bigsec{一、填空题}

\begin{enumerate}
\setcounter{enumi}{2}

\item 设全集 $U=\left\{x\mid\dfrac1x\leqslant\dfrac12\right\}$，$A=[3,+\infty)$，则 $\bar A=$\blankline[4.4em].

\ifanswer
\solbegin
\step{$\dfrac1x\leqslant\dfrac12$，解得 $x<0$ 或 $x\geqslant2$，故 $U=(-\infty,0)\cup[2,+\infty)$，}
% 原卷末句作 $\bar A=(-\infty,0)\cup[2,3)$，按题干 $U$ 应为 $(-\infty,0)\cup\{2\}$，与本解析首句矛盾，疑笔误；照录（见文件头“原卷笔误/存疑”①）
\step{$A=[3,+\infty)$，则 $\bar A=(-\infty,0)\cup[2,3)$.}
\solend
\fi

\item 已知集合 $M=\{y\mid y=x^2-2x-1\}$，$N=\{y\mid y=x^2-4x+1\}$，则 $M\cup N=$\blankline[4.4em].

\ifanswer
\solbegin
\step{$M=\{y\mid y=x^2-2x-1\}=\{y\mid y=(x-1)^2-2\}=\{y\mid y\geqslant-2\}$，}
\step{$N=\{y\mid y=x^2-4x+1\}=\{y\mid y=(x-2)^2-3\}=\{y\mid y\geqslant-3\}$，}
\step{则 $M\cup N=\{y\mid y\geqslant-3\}$.}
\solend
\fi

\item 已知 $P=\{x\mid x^2+x-6=0\}$，$Q=\{x\mid mx-1=0\}$，且 $Q\subset P$，则实数 $m$ 的值\blankline[4.4em].

\ifanswer
\solbegin
\step{$P=\{2,-3\}$，$Q=\{x\mid mx=1\}$；}
\stepm{①}{$m=0$ 时，$Q=\varnothing$，满足 $Q\subseteq P$；}
\stepm{②}{$m\neq0$ 时，要使 $Q\subseteq P$，则 $\dfrac1m=2$ 或 $-3$，所以 $m=\dfrac12$ 或 $-\dfrac13$；}
\step{所以 $m=0$ 或 $\dfrac12$ 或 $-\dfrac13$.}
\solend
\fi

\item 已知“存在正实数 $x$，使得关于 $x$ 的不等式 $ax-1>0$ 成立”为真命题，则实数 $a$ 的取值范围是\blankline[4.4em].

\ifanswer
\solbegin
\step{存在正实数 $x$，使得关于 $x$ 的不等式 $ax-1>0$ 成立为真命题，}
\step{即存在正实数 $x$，使得 $a>\dfrac1x$，因为 $\dfrac1x>0$，所以 $a>0$，}
\step{即 $a$ 的范围为 $(0,+\infty)$.}
\solend
\fi

\item 设 $\alpha:3\leqslant x\leqslant5$，$\beta:2m-1\leqslant x<m+4$，若 $\alpha$ 是 $\beta$ 的充分条件，则实数 $m$ 的取值范围是\blankline[4.4em].

\ifanswer
\solbegin
\step{若 $\alpha$ 是 $\beta$ 的充分条件，则 $\begin{cases}2m-1\leqslant3\\ m+4>5\end{cases}$，解得 $1<m\leqslant2$.}
\solend
\fi

\item 对于实数 $a$、$b$、$c$，已知关于 $x$ 的不等式 $ax^2+bx+c>0$ 的解集是 $(1,2)$，则关于 $x$ 的不等式 $cx^2-bx+a\leqslant0$ 的解集是\blankline[4.4em].

\ifanswer
\solbegin
\step{关于 $x$ 的不等式 $ax^2+bx+c>0$ 的解集是 $(1,2)$，则 $\begin{cases}a<0\\ -\dfrac ba=3\\ \dfrac ca=2\end{cases}$，则 $\begin{cases}b=-3a\\ c=2a\end{cases}$，}
\step{则关于 $x$ 的不等式 $cx^2-bx+a=2ax^2+3ax+a\leqslant0$，则 $2x^2+3x+1\geqslant0$，}
\step{则 $x\leqslant-1$ 或 $x\geqslant-\dfrac12$，则解集为 $(-\infty,-1]\cup\left[-\dfrac12,+\infty\right)$.}
\solend
\fi

\item 已知关于 $x$ 的不等式 $kx^2+kx+k+2\leqslant0$ 的解集为 $\varnothing$，则实数 $k$ 的取值范围为\blankline[4.4em].

\ifanswer
\solbegin
\step{当 $k=0$ 时，不等式化为 $2\leqslant0$，解集为 $\varnothing$，符合题意；}
\step{当 $k\neq0$ 时，若关于 $x$ 的不等式 $kx^2+kx+k+2\leqslant0$ 的解集为 $\varnothing$，}
\step{则 $\begin{cases}k>0\\ \Delta=k^2-4k(k+2)<0\end{cases}$，化简得 $\begin{cases}k>0\\ -3k^2-8k<0\end{cases}$，即 $\begin{cases}k>0\\ k(3k+8)>0\end{cases}$，}
\step{解得 $k>0$.}
\step{综上所述，$k\geqslant0$，即实数 $k$ 的取值范围为 $[0,+\infty)$.}
\solend
\fi

\item 对班级若干名学生调查对 $A$、$B$ 两事件的态度，有如下结果：赞成 $A$ 的人数是全体的五分之三，其余的不赞成；赞成 $B$ 的比赞成 $A$ 多 3 人，其余的不赞成. 已知对 $A$、$B$ 都赞成的学生有 18 人，对 $A$、$B$ 都不赞成的学生数比对 $A$、$B$ 都赞成的学生数的三分之一多 1 人，则班级总人数是\blankline[4.4em]人.

\ifanswer
\solbegin
\step{设班级总人数是 $x$ 人，则赞成 $A$ 的人数是 $\dfrac35x$，赞成 $B$ 的人数是 $\dfrac35x+3$，}
\step{$A$、$B$ 都不赞成的人数为 $18\times\dfrac13+1=7$ 人，}
\step{由容斥原理得 $x=\left(\dfrac35x+\dfrac35x+3-18\right)+7$，解得 $x=40$，即班级总人数是 40 人.}
\solend
\fi

\item 设全集 $U=\{1,2,3,4,5,6,7,8,9,10,11,12\}$，若集合 $A$ 满足：① $A\subseteq U$；②若 $x\in A$，必有 $2x\notin A$；③若 $x\in\bar A$，必有 $2x\notin\bar A$，则这样的集合 $A$ 有\blankline[4.4em]个.

\ifanswer
\solbegin
\step{设集合 $P_n=\{1,2,\cdots,n\}$,$n\in\N^*$，}
\step{任取 $x\in P_n$，将 $x$ 除以 2，若商仍为偶数，再除以 2，$\cdots$，}
\step{经过 $k$ 次以后，商必为奇数，此时记商为 $m$，}
\step{于是 $x=m\cdot2^k$，其中 $m$ 为奇数，$k\in\N^*$.}
\step{由条件得，若 $m\in A$，则 $x\in A$ 等价于 $k$ 为偶数；}
\step{若 $m\notin A$，则 $x\in A$ 等价于 $k$ 为奇数；}
\step{于是 $x$ 是否属于 $A$ 由 $m$ 是否属于 $A$ 确定；}
\step{设 $Q_n$ 是 $P_n$ 中所有奇数的集合，}
\step{设 $f(n)$ 为集合 $P_n$ 中满足题干条件的集合 $A$ 的个数，}
\step{因此 $f(n)$ 等于 $Q_n$ 的子集个数.}
\step{当 $n$ 为偶数（或奇数）时，$P_n$ 中奇数的个数是 $\dfrac n2$（或 $\dfrac{n+1}2$），}
\step{所以 $f(n)=\begin{cases}2^{\frac n2},n\text{为偶数}\\[2pt] 2^{\frac{n+1}2},n\text{为奇数}\end{cases}$，所以 $P_{12}=\{1,2,3,4,5,6,7,8,9,10,11,12\}$，}
\step{即同时满足三个条件的集合 $A$ 的个数为 $2^{\frac{12}2}=2^6=64$.}
\solend
\fi

\item 设集合 $\Omega=\{(x_1,x_2,\cdots,x_{2020})\mid x_i\in\{0,1\},i=1,2,\cdots,2020\}$，对于 $\Omega$ 中的任意两个元素 $\alpha=(x_1,x_2,\cdots,x_{2020})$，$\beta=(y_1,y_2,\cdots,y_{2020})$，记
% 原卷此式第四组缺左括号、末尾连排两个右括号“))”，括号不配对；照录（见文件头“原卷笔误/存疑”②）
$\alpha*\beta=(x_1+y_1-x_1y_1)+(x_2+y_2-x_2y_2)+\cdots+x_{2020}+y_{2020}-x_{2020}y_{2020}))$，设 $u$、$v\in\Omega$，若 $u*u+v*v=2020$，则 $u*v$ 的最大值和最小值之和为\blankline[4.4em].

\ifanswer
\solbegin
\step{设 $u=(x_1,x_2,\cdots,x_{2020})$,$v=(y_1,y_2,\cdots,y_{2020})$，}
\step{因为 $x_i\in\{0,1\}$，所以 $x_i(x_i-1)=0$，}
\step{所以 $u*u=(x_1+x_1-x_1^2)+(x_2+x_2-x_2^2)+\cdots+(x_{2020}+x_{2020}-x_{2020}^2)$}
\step{$=x_1+x_2+\cdots+x_{2020}$，}
\step{$v*v=(y_1+y_1-y_1^2)+(y_2+y_2-y_2^2)+\cdots+(y_{2020}+y_{2020}-y_{2020}^2)$}
\step{$=y_1+y_2+\cdots+y_{2020}$，}
\step{由 $u*u+v*v=2020$，得 $x_1+x_2+\cdots+x_{2020}+y_1+y_2+\cdots+y_{2020}=2020$，}
% 原卷此行“$\cdots$”后脱逗号、 “2020”后衍一半角逗号（“2020,个 1”）；照录（见文件头“原卷笔误/存疑”③）
\step{所以 $x_1,x_2,\cdots x_{2020},y_1,y_2,\cdots y_{2020}$ 中有 2020,个 1，2020 个 0，}
\step{$u*v=(x_1+y_1-x_1y_1)+(x_2+y_2-x_2y_2)+\cdots+(x_{2020}+y_{2020}-x_{2020}y_{2020})$}
\step{$=2020-(x_1y_1+x_2y_2+\cdots+x_{2020}y_{2020})$，}
\step{注意到只有 $x_i=y_i=1$ 时，$x_iy_i=1$，否则 $x_iy_i=0$，}
% 此行“$\cdots$”后同样脱逗号，与上一处同式样；照录（见文件头“原卷笔误/存疑”③）
\step{而 $x_1,x_2,\cdots x_{2020},y_1,y_2,\cdots y_{2020}$ 中有 2020 个 1，2020 个 0，}
\step{所以满足 $x_iy_i=1$ 的最多有 1010 个，所以 $u*v\geqslant2020-1010=1010$，}
\step{满足 $x_iy_i=0$ 的最多有 2020 个，所以 $u*v\leqslant2020-0=2020$，}
\step{则 $u*v$ 的最大值和最小值之和为 $2020+1010=3030$.}
\solend
\fi

\end{enumerate}

\bigsec{二、选择题}

\begin{enumerate}
\setcounter{enumi}{12}

\item 实数 $x$ 与 $y$，“$x+y>2$”的\mbox{（\quad）}条件是“$x>1$ 或 $y>1$”.

\keepwithprev
A. 充分非必要\qquad B. 必要非充分

C. 充要\qquad D. 既非充分又非必要

\ifanswer
\solbegin
\step{当 $x>1$ 或 $y>1$ 时，$x+y>2$ 不一定成立，充分性不成立；}
\step{当 $x+y>2$ 时，$x>1$ 或 $y>1$ 一定成立，必要性成立. 故选 $B$.}
\solend
\fi

\item 设 $a,b\in\R$，且 $a<0$，$-1<b<0$，那么下列不等式成立的是\mbox{（\quad）}

\keepwithprev
A. $a>ab>ab^2$\qquad B. $ab^2>ab>a$\qquad C. $ab>a>ab^2$\qquad D. $ab>ab^2>a$

\ifanswer
\solbegin
\step{因为 $-1<b<0$，$a<0$，所以 $ab>0$，$0<b^2<1$.}
\step{所以 $ab-ab^2=ab(1-b)>0$，$ab^2-a=a(b^2-1)>0$，}
\step{所以 $ab>ab^2>a$. 故选 $D$.}
\solend
\fi

% 第 15 题答案原卷印在题干括号内（“(  D  )”），且无【答案】/【解析】段；
% 按 高一37 第 13 题、高一63 第 14 题成例括号留空、答案另提 \ans{$D$}
\item 用反证法证明：“平面四边形中至少有一个内角不超过 90°”时，下列“假设”中正确的是\mbox{（\quad）}

\keepwithprev
A. 假设有两个内角超过 90°\qquad B. 假设有三个内角超过 90°

C. 假设至多有两个内角超过 90°\qquad D. 假设四个内角都超过 90°

\ifanswer
\ans{$D$}
\fi

\item 设 $A$ 为非空实数集，集合 $B=\{xy\mid x\in A,y\in A,x\neq y\}$，则对于以下两个结论说法正确的是\mbox{（\quad）}

结论①：若 $A$ 是五元集，则 $B$ 中的元素最多 10 个；

结论②：存在四元集 $A$，使得 $B=\{2,4,5,8,10,16\}$.

\keepwithprev
A. ①对②对\qquad B. ①对②错\qquad C. ①错②对\qquad D. ①错②错

\ifanswer
\solbegin
\step{对于①，想让元素个数达到最大值 10，则须避免产生符号不同的乘积，}
\step{否则正负分开计数容易导致总数受限，且容易产生相等值，}
\step{也最好避免 0 的出现，}
\step{故设 $A=\{a_1,a_2,a_3,a_4,a_5\}$，$0<a_1<a_2<a_3<a_4<a_5$.}
% 原卷此链“$a_3a_3$”同足标自乘、“$a_{14}a_5$”下标 14 均与文意不合，疑为 $a_2a_4$、$a_4a_5$ 之误；照录（见文件头“原卷笔误/存疑”④）
\step{因为 $a_1a_2<a_1a_3<a_1a_4<a_1a_5<a_2a_3<a_3a_3<a_{14}a_5$，$a_2a_3<a_2a_4<a_3a_4$，}
\step{所以 $B$ 中元素个数小于等于 10 个，}
\step{如设 $A=\{1,2,3,5,7\}$，则 $B=\{2,3,5,6,7,10,14,15,21,35\}$，}
\step{所以集合 $B$ 中元素个数的最大值为 10 个，故①正确；}
\step{对于②，由于集合 $B$ 中没有 0，所以集合 $A$ 中没有 0；}
\step{设集合 $A=\{a,b,c,d\}$，且 $a<b<c<d$，}
\step{当 $a<0<b<c<d$ 时，$a\times b<0$，$a\times c<0$，$a\times d<0$，}
\step{故集合 $B$ 至少有 3 个负数，与 $B=\{2,4,5,8,10,16\}$ 矛盾；}
\step{当 $a<b<0<c<d$ 时，$a\times b>0$，$c\times d>0$，}
\step{$a\times c<0$，$a\times d<0$，$b\times c<0$，$b\times d<0$，}
\step{故集合 $B$ 至少有 4 个负数，与 $B=\{2,4,5,8,10,16\}$ 矛盾；}
\step{当 $a<b<c<0<d$ 时，$a\times b>0$，$a\times c>0$，$b\times c>0$，}
\step{$a\times d<0$，$b\times d<0$，$c\times d<0$，}
\step{故集合 $B$ 至少有 3 个负数，与 $B=\{2,4,5,8,10,16\}$ 矛盾；}
\step{当 $0<a<b<c<d$ 时，则集合 $B=\{ab,ac,ad,bc,bd,cd\}$，}
\step{则必有 $ab=2$，$cd=16$，其 4 个正实数的乘积 $abcd=32$，}
\step{又 $ad=5$，$bc=8$ 或 $ad=8$，$bc=5$，其 4 个正实数的乘积 $abcd=40$，矛盾；}
\step{当 $a<b<c<d<0$ 时，与 $0<a<b<c<d$ 时的证明方法相同，}
\step{得集合 $A$ 也不存在；}
\step{综上不存在四元集 $A$，使其生成集 $B=\{2,4,5,8,10,16\}$，故②错误.}
\step{故选 $B$.}
\solend
\fi

\end{enumerate}

\bigsec{三、解答题}

\begin{enumerate}
\setcounter{enumi}{16}

\item 解下列不等式（组）

\begin{enumerate}[label=(\arabic*)]
\item 解不等式：$(ax+1)(x+1)<0(a\in\R)$.
\item 解不等式组：$\begin{cases}\dfrac{2x+4}{x+4}\geqslant1\\[4pt] (x+1)(2-|x|)\geqslant0\end{cases}$.
\end{enumerate}

\ifanswer
\solbegin
\stepm{(1)}{当 $a=0$ 时，原不等式为 $x+1<0$，即 $x<-1$，}
\step{所以此时不等式的解集为 $\{x\mid x<-1\}$；}
\step{当 $a\neq0$ 时，令 $(ax+1)(x+1)=0$，解得 $x_1=-\dfrac1a$,$x_2=-1$，}
\step{当 $a>1$ 时，$-\dfrac1a>-1$，此时不等式的解集为 $\left\{x\mid-1<x<-\dfrac1a\right\}$；}
% 以下两行“此时”前原卷均无标点；照录（见文件头“原卷笔误/存疑”⑥）
\step{当 $a=1$ 时，原不等式为 $(x+1)^2<0$ 此时不等式的解集为 $\varnothing$；}
\step{当 $0<a<1$ 时，$-\dfrac1a<-1$ 此时不等式的解集为 $\left\{x\mid-\dfrac1a<x<-1\right\}$；}
\step{当 $a<0$ 时，$-\dfrac1a>-1$，此时不等式的解集为 $\left\{x\mid x<-1\ \text{或}\ x>-\dfrac1a\right\}$.}
\step{综上所述，当 $a=0$ 时，$\{x\mid x<-1\}$；}
\step{当 $a>1$ 时，$\left\{x\mid-1<x<-\dfrac1a\right\}$；}
\step{当 $a=1$ 时，$\varnothing$；}
\step{当 $0<a<1$ 时，$\left\{x\mid-\dfrac1a<x<-1\right\}$；}
\step{当 $a<0$ 时，$\left\{x\mid x<-1\ \text{或}\ x>-\dfrac1a\right\}$.}
\stepm{(2)}{不等式 $\dfrac{2x+4}{x+4}\geqslant1$ 可化为 $\dfrac{2x+4}{x+4}-1\geqslant0$，即 $\dfrac x{x+4}\geqslant0$，}
\step{等价于 $\begin{cases}x(x+4)\geqslant0\\ x+4\neq0\end{cases}$，解得 $x<-4$ 或 $x\geqslant0$，}
\step{不等式 $(x+1)(2-|x|)\geqslant0$ 等价于 $\begin{cases}(x+1)(2-x)\geqslant0,x\geqslant0\\ (x+1)(2+x)\geqslant0,x<0\end{cases}$，}
\step{解不等式 $(x+1)(2-x)\geqslant0$，$x\geqslant0$，得 $0\leqslant x\leqslant2$，}
\step{解不等式 $(x+1)(2+x)\geqslant0$，$x<0$，得 $x\leqslant-2$ 或 $-1\leqslant x<0$，}
\step{所以 $(x+1)(2-|x|)\geqslant0$ 的解集为 $\{x\mid x\leqslant-2\ \text{或}\ -1\leqslant x\leqslant2\}$，}
\step{综上所述，$\begin{cases}\dfrac{2x+4}{x+4}\geqslant1\\[4pt] (x+1)(2-|x|)\geqslant0\end{cases}$ 的解集为 $\{x\mid x<-4\ \text{或}\ 0\leqslant x\leqslant2\}$.}
\solend
\fi

\ansspace[8]

% 第 18 题题干＋两小问与作答区 6 行须同页
\ansneed{20cm}
\item 已知 $a$ 为实数，设 $f(x)=x^2+2(a+1)x+1$.

\begin{enumerate}[label=(\arabic*)]
\item 已知对于任意实数 $x$，$\sqrt{f(x)}$ 均有意义，求实数 $a$ 的取值范围；
\item 设 $A=\{x\mid f(x)+a^2-2=0\}$，$B=\{-4,0\}$，$A\cup B=B$，求实数 $a$ 的取值范围.
\end{enumerate}

\ifanswer
\solbegin
\stepm{(1)}{$\sqrt{f(x)}$ 对任意实数 $x$ 都有意义，}
\step{即 $f(x)=x^2+2(a+1)x+1\geqslant0$ 对任意 $x\in\R$ 恒成立，}
\step{所以 $\Delta=4(a+1)^2-4\times1\times1\leqslant0$，即 $(a+1)^2\leqslant1$，解得 $-2\leqslant a\leqslant0$，}
\step{故 $a$ 的取值范围是 $[-2,0]$.}
\stepm{(2)}{由 $f(x)+a^2-2=0$ 得 $x^2+2(a+1)x+a^2-1=0$.}
\step{由 $A\cup B=B$ 得 $A\subseteq B=\{-4,0\}$，即该方程无实根，}
\step{或它的实根都在 $\{-4,0\}$ 中.}
\step{当方程无实根时，$\Delta=4(a+1)^2-4(a^2-1)=8a+8<0$，解得 $a<-1$；}
\step{当方程有实根时，$a\geqslant-1$，设两根为 $x_1$、$x_2$，}
\step{由根与系数的关系得 $x_1+x_2=-2(a+1)$、$x_1x_2=a^2-1$，}
\step{此时两根只能在 $\{-4,0\}$ 中取值；}
\step{若两根为 0 与 0，由 $x_1x_2=0$ 得 $a=\pm1$，}
\step{再由 $x_1+x_2=0$ 得 $a=-1$；}
\step{若两根为 0 与 $-4$，由 $x_1x_2=0$ 得 $a=\pm1$，再由 $x_1+x_2=-4$ 得 $a=1$.}
\step{经检验，$a=-1$ 时 $A=\{0\}$，$a=1$ 时 $A=\{-4,0\}$，都满足 $A\subseteq B$.}
\step{综上，实数 $a$ 的取值范围为 $(-\infty,-1]\cup\{1\}$.}
\solend
\fi

\ansspace[8]

\ansneed{18cm}
\item 共享单车给市民出行带来了诸多便利，某公司购买了一批单车投放到某地给居民使用，据市场分析，这批单车的运营累计利润 $y$（单位：元）与运营天数 $x$（单位：天，$x$ 为正整数）满足关系：$y=-\dfrac12x^2+60x-800$.

\begin{enumerate}[label=(\arabic*)]
\item 要使这批单车的运营累计利润高于 800 元，求运营天数的取值范围；
\item 判断这批单车的日平均运营利润能否超过 20 元，并说明理由.
\end{enumerate}

\ifanswer
\solbegin
\stepm{(1)}{由题意得 $y=-\dfrac12x^2+60x-800>800$，}
\step{即 $x^2-120x+3200<0$，解得 $40<x<80$，}
\step{即运营天数的取值范围为 $\{x\mid40<x<80,x\in\N^*\}$.}
\stepm{(2)}{不能，理由如下：}
\step{这批单车的日平均运营利润为 $\dfrac yx=-\dfrac12x-\dfrac{800}x+60=-\dfrac12\left(x+\dfrac{1600}x\right)+60$，}
\step{因 $x>0$，则 $x+\dfrac{1600}x\geqslant2\sqrt{x\cdot\dfrac{1600}x}=80$，}
\step{当且仅当 $x=\dfrac{1600}x$，即 $x=40$ 时取等号，}
\step{此时 $\dfrac yx=-\dfrac12\left(x+\dfrac{1600}x\right)+60\leqslant-\dfrac12\times80+60=20$，}
\step{故这批单车的日平均运营利润不能超过 20 元.}
\solend
\fi

\ansspace[8]

% 第 20 题题干＋三小问与作答区 8 行须同页
\ansneed{20cm}
\item 设 $f(x)=ax^2+bx+c$($a\neq0$，$b$、$c\in\R$).

\begin{enumerate}[label=(\arabic*)]
\item 若 $a=b=c=1$，设 $A=\{(x,y)\mid y=f(x)\}$，$B=\{(x,y)\mid y=1\}$，用列举法表示 $A\cap B$；
% 原卷两个“若”连用（“若 a=1，…，若方程 f(x)=0 有两个大于 1 的根”），疑衍其一；照录（见文件头“原卷笔误/存疑”⑦）
\item 若 $a=1$，$b=2k-1$，$c=k^2$，若方程 $f(x)=0$ 有两个大于 1 的根，求实数 $k$ 的取值范围；
\item 若 $a$、$b$、$c\in\Z$，且 $f(0)$、$f(1)$ 均为奇数，求证：方程 $f(x)=0$ 无整数解.
\end{enumerate}

\ifanswer
\solbegin
\stepm{(1)}{当 $a=b=c=1$ 时，$f(x)=x^2+x+1$，}
\step{所以 $A=\{(x,y)\mid y=x^2+x+1\}$，又 $y=1$，}
\step{由 $x^2+x+1=1$ 得 $x^2+x=0$，解得 $x=0$ 或 $x=-1$，}
\step{所以 $A\cap B=\{(0,1),(-1,1)\}$.}
\stepm{(2)}{当 $a=1$、$b=2k-1$、$c=k^2$ 时，$f(x)=x^2+(2k-1)x+k^2$.}
\step{设两根为 $x_1$、$x_2$，}
\step{则“两根都大于 1”等价于 $\Delta>0$、$f(1)>0$、$-\dfrac b{2a}>1$.}
\step{由 $\Delta=(2k-1)^2-4k^2=1-4k>0$ 得 $k<\dfrac14$；}
\step{由 $f(1)=1+(2k-1)+k^2=k^2+2k=k(k+2)>0$ 得 $k<-2$ 或 $k>0$；}
\step{由 $-\dfrac{2k-1}2=\dfrac{1-2k}2>1$ 得 $k<-\dfrac12$.}
\step{三者取交集得 $k<-2$，即 $k$ 的取值范围是 $(-\infty,-2)$.}
\stepm{(3)}{设 $x_0\in\Z$ 是方程 $f(x)=0$ 的整数解.}
\step{因为 $x_0(x_0-1)$ 是两个相邻整数之积，恒为偶数，所以 $ax_0(x_0-1)$ 为偶数.}
\step{而 $f(x_0)=ax_0^2+bx_0+c=ax_0(x_0-1)+(a+b)x_0+c$，}
\step{所以 $f(x_0)$ 与 $(a+b)x_0+c$ 同奇偶.}
\step{若 $x_0$ 为偶数，则 $(a+b)x_0$ 为偶数，$f(x_0)$ 与 $c=f(0)$ 同奇偶，为奇数；}
\step{若 $x_0$ 为奇数，则 $(a+b)x_0$ 与 $a+b$ 同奇偶，}
\step{$f(x_0)$ 与 $(a+b)+c=f(1)$ 同奇偶，也为奇数.}
\step{综上，$f(x_0)$ 恒为奇数，不可能等于 0，与 $f(x_0)=0$ 矛盾，}
\step{所以方程 $f(x)=0$ 无整数解.}
\solend
\fi

\ansspace[8]

% 压轴题（题干＋三小问）与作答区 10 行须同页
\ansneed{20cm}
\item 已知 $X$ 为有限集，且任取 $x\in X$，均有 $x\in\Z$，称 $X$ 为非空有限整数集. 定义：$X_k=\{x+k\mid x\in X\}$，$X^2=\{y\mid y=x^2,x\in X\}$.

\begin{enumerate}[label=(\arabic*)]
\item 若 $X=\{-2,2\}$，分别写出 $X_2$ 和 $X^2$（用列举法表示）；
\item 若 $X=\{-1,0,1\}$，$X_1\subseteq Y$，$(Y_{-1})^2\subseteq X$，求集合 $Y$ 中所有的元素；
\item 对于非空有限整数集 $A$、$B$，若 $A_k\subseteq B$，$(B^2)_{-2}\subseteq A$，$B$ 中的元素均为非负整数，$k$ 为整数，求证：

① $A$、$B$ 中元素的个数相同；

② 若集合 $A$ 中的元素个数为奇数，则 $k$ 为偶数.
\end{enumerate}

\ifanswer
\solbegin
\stepm{(1)}{由题意得 $X_2=\{0,4\}$，$X^2=\{4\}$；}
\stepm{(2)}{由题意得 $X_1=\{0,1,2\}$，即 $\{0,1,2\}\subseteq Y$，$Y_{-1}=\{y-1\mid y\in Y\}$，}
\step{$(Y_{-1})^2=\{(y-1)^2\mid y\in Y\}\subseteq X=\{-1,0,1\}$，}
\step{所以 $(y-1)^2=0$ 或 $(y-1)^2=1$，解得 $y=1$ 或 $y=2$ 或 $y=0$，}
\step{即 $y\in\{0,1,2\}$，即 $Y\subseteq\{0,1,2\}$，所以 $Y=\{0,1,2\}$；}
% 原卷 (3) 解析首行“(3)”与“①”连排，本稿合并记作 \stepm{(3)①}
\stepm{(3)①}{因为 $A_k=\{x+k\mid x\in A\}$，函数 $g(x)=x+k$ 在全体整数集上严格增，}
\step{所以 $\forall x_1$、$x_2\in A$，$x_1\neq x_2$，$g(x_1)\neq g(x_2)$，所以 $|A_k|=|A|$，}
\step{因为 $A_k\subseteq B$，所以 $|A|\leqslant|B|$，}
\step{设 $b\in B$，$B$ 中元素均为非负整数，$\forall b_1\neq b_2\in B$，$b_1^2\neq b_2^2$，}
\step{故函数 $f(b)=b^2-2$ 在 $B$ 上单射，所以 $|(B^2)_{-2}|=|B|$，}
\step{又因为 $(B^2)_{-2}\subseteq A$，所以 $|B|\leqslant|A|$，所以 $|A|=|B|$.}
\stepm{②}{由①得 $A_k\subseteq B$ 且 $|A_k|=|A|=|B|$，$(B^2)_{-2}\subseteq A$ 且 $|(B^2)_{-2}|=|B|=|A|$，}
\step{所以 $A_k=B$，$(B^2)_{-2}=A$，}
\step{所以 $B=A_k=\{a+k\mid a\in A\}$，$A=(B^2)_{-2}=\{b^2-2\mid b\in B\}$，}
\step{所以 $A=\{(a+k)^2-2\mid a\in A\}$，}
\step{所以映射 $h(a)=(a+k)^2-2$ 是集合 $A$ 到自身的一个双射.}
\step{假设 $k$ 是奇数，当 $a$ 为偶数时，$a+k$ 为奇数，$(a+k)^2-2$ 为奇数，}
\step{当 $a$ 为奇数时，$a+k$ 为偶数，$(a+k)^2-2$ 为偶数，}
\step{所以此时映射 $h(a)=(a+k)^2-2$，将集合 $A$ 中的奇数映射为偶数，}
\step{将集合 $A$ 中的偶数映射为奇数，即集合 $A$ 中的奇数元素和偶数元素必须配对，}
\step{此时集合 $A$ 的元素个数则为偶数，}
\step{与题设中“集合 $A$ 中的元素个数为奇数”矛盾.}
\step{故假设不成立，所以 $k$ 是偶数.}
\solend
\fi

\ansspace*[10]

\end{enumerate}

\end{document}
"
%
% 原始材料：微信公众号“魔都数学加油站”文章，作者破晓，连载“27学年高一67”；
% 原文标题“27学年高一67——2026-2027年上海市华师大一附中高一上10月月考”；
% 原文链接：https://mp.weixin.qq.com/s/Kdhx07TVjLoGTcIlnRzgcQ。
% 页面标注发布时间 2026-10-18 20:33（上海）。
% 本次转录只据 `原图/高一67_华师大一附中_10月月考/` 的 13 张 PNG 逐页读图。
% 正文为 13 张 PNG 页（`原图/高一67_华师大一附中_10月月考/01..13.png`）：
% 第 01、04、13 页 2560×3620，第 02、03、05—12 页 4762×6735
% （**同一篇各页分辨率不同**），均按页序读取。13 页全为试卷正文页（卷首标题在第 1 页页顶，
% 卷末解答题第 21 题解析止于第 13 页），无封面页、目录页、二维码或宣传页。
% 粉色斜排“魔都数学加油站”为卷面水印，与公众号页脚一并未录入。
%
% 卷首标题：照卷面印作“2026--2027 年华师大一附中高一上 10 月月考”（卷面为黑体加粗一行），
% 年份区间按全稿惯例排 en dash。原卷只有一行标题、无副标题。
%
% 篇幅与题号：原卷自第 3 题起编号（第 1、2 题不在本卷内）。填空题 3—12（10 题）、
% 选择题 13—16（4 题）、解答题 17—21（5 题），共 19 题。本稿沿用原节与原题号，
% 只改排为全稿统一的大节样式（原卷小节标题为“一、填空题”“二、选择题”“三、解答题”）。
%
% 讲评版：原卷每题下直接印【解析】。**全卷 19 题中 18 题都印【解析】（无一处只印
% 【答案】）**，按全稿口径这类题不另提 \ans{}；各题解析末句的结论（第 13、14、16 题
% 的“故选 B／D／B”、各填空与解答题的“故/综上…”）照录在解析内。
% **另有 1 题答案印在题干括号内**——第 15 题（“(  D  )”，且该题无【答案】/【解析】段），
% 本稿按 高一37 第 13 题、高一63 第 14 题成例把括号留空、答案另提为 \ans{$D$}。
%
% 副标题：原卷未印副标题。本稿按“知识模块”口径标作“集合与常用逻辑用语、不等式”：
% 19 题中集合与常用逻辑用语 12 题（第 3、4、5、7、10、11、12、13、15、16、18、21 题：
% 补集、并集、子集、充分条件、容斥计数、子集计数、新定义二元运算、条件判断、
% 反证法假设、新定义生成集、集合关系与二次方程、新定义 X_k 与 X^2），
% 不等式 7 题（第 6、8、9、14、17、19、20 题：存在量词不等式、解集反求系数、
% 解集为空、不等式性质、解不等式（组）、二次不等式与基本不等式、二次方程根的分布），
% 两类均超过三分之一，故并标。
%
% 与原卷的排版差异：
%   ① 原卷每题下直接印【解析】，本稿按全稿惯例将解析置于答案版，试卷版保持干净；
%   ② 原卷未印副标题，本稿按“知识模块”口径补出；
%   ③ 原卷填空横线约四个字宽，本稿按全稿惯例统一排 \blankline[4.4em]；
%   ④ 原卷实数集排斜体 $R$（第 14 题题干、第 17 题 (1) 题干、第 18 题解析、第 20 题题干），
%      整数集排斜体 $Z$（第 20 题 (3) 题干与解析、第 21 题题干），正整数集排粗体
%      $N^*$（第 11 题解析、第 19 题解析），本稿按全稿惯例统一改用 $\R$、$\Z$、$\N^*$
%      （记号归一，粗体／斜体不作区分）；空集原卷作 $\varnothing$，照录；
%   ⑤ 原卷在数学式内多处用半角逗号（如第 8 题 $(1,2)$ 与各 cases、第 11 题解析
%      $P_n=\{1,2,\cdots,n\}$,$n\in N^*$、第 12 题解析 $u=(\cdots)$,$v=(\cdots)$ 与
%      “2020,个 1”、第 19 题解析 $\{x|40<x<80,x\in N^*\}$），均照录未强改；
%      描述法集合的竖线照原卷排 \mid；句末句点按全稿惯例排西文句点；
%   ⑥ 原卷补集记号为上划线 $\bar A$（第 3 题题干与解析、第 11 题题干），照录
%      （同 高一64/65 的 $\bar B$ 成例）；
%   ⑦ 第 12 题新定义的运算记号 $*$ 照录（$\alpha*\beta$、$u*u$、$u*v$）；
%   ⑧ 第 17—21 题的子问原卷各自另起一行，本稿按全稿惯例用 enumerate 的 (1)(2)(3)；
%      解析中的子问标号一律用 \stepm{(1)}{…}；第 5 题解析的①②、第 21 题 (3) 解析的①②
%      按全稿惯例作 \stepm{①}{…}、\stepm{②}{…}（第 21 题 (3) 首行原卷“(3)”与“①”连排，
%      本稿合并记作 \stepm{(3)①}{…}）；
%   ⑨ 选择题作答括号原卷排半角“(　　)”（第 13、14、16 题括号内空白、第 15 题括号内印 D），
%      本稿按全稿惯例排 \mbox{（\quad）}；第 15 题括号内的答案移入 \ans{}，见上；
%   ⑩ 选项行排照原卷：第 13、15 题两行两条，第 14、16 题一行四条；
%   ⑪ 每道解答题原卷未留作答空白，本稿按全稿惯例加 \ansspace（末题用 \ansspace*）；
%      第 20、21 题题干长，各用一处 \ansneed 让题干与作答区不拆页（\ansneed 只作用于试卷版）。
%
% 原卷笔误/存疑（均照录未改，正文对应处另有 % 注）：
%   ① 第 3 题解析末句印作“则 $\bar A=(-\infty,0)\cup[2,3)$.”：按题干
%      $U=\{x\mid 1/x\leqslant 1/2\}$（即 $x<0$ 或 $x\geqslant2$）与 $A=[3,+\infty)$，
%      $\bar A$ 应为 $(-\infty,0)\cup\{2\}$；且与本解析首句自求之 $U$ 相矛盾（$3\notin U$），
%      疑为笔误；照录未改。
%   ② 第 12 题题干 $\alpha*\beta$ 定义式括号不配对：第四组“$+\cdots+x_{2020}+y_{2020}
%      -x_{2020}y_{2020}$”缺左括号、末尾连排两个右括号“))”；照录未改。
%   ③ 第 12 题解析“所以…中有 2020,个 1”一行，“$\cdots$”后脱逗号（作
%      “$x_2,\cdots x_{2020}$”，第 12 行“而…”同式样），且“2020”后衍一半角逗号
%      （作“2020,个 1”）；照录未改。
%   ④ 第 16 题解析乘积链印作“$a_1a_2<a_1a_3<a_1a_4<a_1a_5<a_2a_3<a_3a_3<a_{14}a_5$”：
%      “$a_3a_3$”同足标自乘、与单调链不合，“$a_{14}a_5$”下标 14 与 $a_1$—$a_5$ 记号
%      不合（疑为 $a_2a_4$、$a_4a_5$ 之误）；照录未改。
%   ⑤ 第 5 题题干作“$Q\subset P$”（真子集记号），解析两处均作“$Q\subseteq P$”，
%      同一题前后记号不一；照录（本稿分别排 \subset、\subseteq）。
%   ⑥ 第 17 题 (1) 解析“当 $a=1$ 时，原不等式为 $(x+1)^2<0$ 此时…”与
%      “当 $0<a<1$ 时，$-\frac1a<-1$ 此时…”两处“此时”前均脱标点；照录未改。
%   ⑦ 第 20 题 (2) 题干“若 $a=1$，$b=2k-1$，$c=k^2$，若方程 $f(x)=0$ 有两个大于 1 的根”
%      两个“若”连用，疑衍其一；照录未改。
%
% 插图：无（全卷 13 页无任何几何图形、无“如图/右图/图中”字样）。
%
% 行数核对：逐题数 \step/\stepm 条数与原卷【解析】印刷行数，18 题逐一相符
% （第 15 题无解析段，只提 \ans{}），合计 169 条 / 169 行
% （Q3—Q12：2、3、4、3、1、3、5、3、13、15；Q13—Q16：2、3、0、25；
% Q17—Q21：19、16、9、20、23）。
%
\documentclass[a4paper,11pt]{article}
\usepackage{common}

\begin{document}

\examtitlest[3]{2026--2027 年华师大一附中高一上 10 月月考}{集合与常用逻辑用语、不等式}

\bigsec{一、填空题}

\begin{enumerate}
\setcounter{enumi}{2}

\item 设全集 $U=\left\{x\mid\dfrac1x\leqslant\dfrac12\right\}$，$A=[3,+\infty)$，则 $\bar A=$\blankline[4.4em].

\ifanswer
\solbegin
\step{$\dfrac1x\leqslant\dfrac12$，解得 $x<0$ 或 $x\geqslant2$，故 $U=(-\infty,0)\cup[2,+\infty)$，}
% 原卷末句作 $\bar A=(-\infty,0)\cup[2,3)$，按题干 $U$ 应为 $(-\infty,0)\cup\{2\}$，与本解析首句矛盾，疑笔误；照录（见文件头“原卷笔误/存疑”①）
\step{$A=[3,+\infty)$，则 $\bar A=(-\infty,0)\cup[2,3)$.}
\solend
\fi

\item 已知集合 $M=\{y\mid y=x^2-2x-1\}$，$N=\{y\mid y=x^2-4x+1\}$，则 $M\cup N=$\blankline[4.4em].

\ifanswer
\solbegin
\step{$M=\{y\mid y=x^2-2x-1\}=\{y\mid y=(x-1)^2-2\}=\{y\mid y\geqslant-2\}$，}
\step{$N=\{y\mid y=x^2-4x+1\}=\{y\mid y=(x-2)^2-3\}=\{y\mid y\geqslant-3\}$，}
\step{则 $M\cup N=\{y\mid y\geqslant-3\}$.}
\solend
\fi

\item 已知 $P=\{x\mid x^2+x-6=0\}$，$Q=\{x\mid mx-1=0\}$，且 $Q\subset P$，则实数 $m$ 的值\blankline[4.4em].

\ifanswer
\solbegin
\step{$P=\{2,-3\}$，$Q=\{x\mid mx=1\}$；}
\stepm{①}{$m=0$ 时，$Q=\varnothing$，满足 $Q\subseteq P$；}
\stepm{②}{$m\neq0$ 时，要使 $Q\subseteq P$，则 $\dfrac1m=2$ 或 $-3$，所以 $m=\dfrac12$ 或 $-\dfrac13$；}
\step{所以 $m=0$ 或 $\dfrac12$ 或 $-\dfrac13$.}
\solend
\fi

\item 已知“存在正实数 $x$，使得关于 $x$ 的不等式 $ax-1>0$ 成立”为真命题，则实数 $a$ 的取值范围是\blankline[4.4em].

\ifanswer
\solbegin
\step{存在正实数 $x$，使得关于 $x$ 的不等式 $ax-1>0$ 成立为真命题，}
\step{即存在正实数 $x$，使得 $a>\dfrac1x$，因为 $\dfrac1x>0$，所以 $a>0$，}
\step{即 $a$ 的范围为 $(0,+\infty)$.}
\solend
\fi

\item 设 $\alpha:3\leqslant x\leqslant5$，$\beta:2m-1\leqslant x<m+4$，若 $\alpha$ 是 $\beta$ 的充分条件，则实数 $m$ 的取值范围是\blankline[4.4em].

\ifanswer
\solbegin
\step{若 $\alpha$ 是 $\beta$ 的充分条件，则 $\begin{cases}2m-1\leqslant3\\ m+4>5\end{cases}$，解得 $1<m\leqslant2$.}
\solend
\fi

\item 对于实数 $a$、$b$、$c$，已知关于 $x$ 的不等式 $ax^2+bx+c>0$ 的解集是 $(1,2)$，则关于 $x$ 的不等式 $cx^2-bx+a\leqslant0$ 的解集是\blankline[4.4em].

\ifanswer
\solbegin
\step{关于 $x$ 的不等式 $ax^2+bx+c>0$ 的解集是 $(1,2)$，则 $\begin{cases}a<0\\ -\dfrac ba=3\\ \dfrac ca=2\end{cases}$，则 $\begin{cases}b=-3a\\ c=2a\end{cases}$，}
\step{则关于 $x$ 的不等式 $cx^2-bx+a=2ax^2+3ax+a\leqslant0$，则 $2x^2+3x+1\geqslant0$，}
\step{则 $x\leqslant-1$ 或 $x\geqslant-\dfrac12$，则解集为 $(-\infty,-1]\cup\left[-\dfrac12,+\infty\right)$.}
\solend
\fi

\item 已知关于 $x$ 的不等式 $kx^2+kx+k+2\leqslant0$ 的解集为 $\varnothing$，则实数 $k$ 的取值范围为\blankline[4.4em].

\ifanswer
\solbegin
\step{当 $k=0$ 时，不等式化为 $2\leqslant0$，解集为 $\varnothing$，符合题意；}
\step{当 $k\neq0$ 时，若关于 $x$ 的不等式 $kx^2+kx+k+2\leqslant0$ 的解集为 $\varnothing$，}
\step{则 $\begin{cases}k>0\\ \Delta=k^2-4k(k+2)<0\end{cases}$，化简得 $\begin{cases}k>0\\ -3k^2-8k<0\end{cases}$，即 $\begin{cases}k>0\\ k(3k+8)>0\end{cases}$，}
\step{解得 $k>0$.}
\step{综上所述，$k\geqslant0$，即实数 $k$ 的取值范围为 $[0,+\infty)$.}
\solend
\fi

\item 对班级若干名学生调查对 $A$、$B$ 两事件的态度，有如下结果：赞成 $A$ 的人数是全体的五分之三，其余的不赞成；赞成 $B$ 的比赞成 $A$ 多 3 人，其余的不赞成. 已知对 $A$、$B$ 都赞成的学生有 18 人，对 $A$、$B$ 都不赞成的学生数比对 $A$、$B$ 都赞成的学生数的三分之一多 1 人，则班级总人数是\blankline[4.4em]人.

\ifanswer
\solbegin
\step{设班级总人数是 $x$ 人，则赞成 $A$ 的人数是 $\dfrac35x$，赞成 $B$ 的人数是 $\dfrac35x+3$，}
\step{$A$、$B$ 都不赞成的人数为 $18\times\dfrac13+1=7$ 人，}
\step{由容斥原理得 $x=\left(\dfrac35x+\dfrac35x+3-18\right)+7$，解得 $x=40$，即班级总人数是 40 人.}
\solend
\fi

\item 设全集 $U=\{1,2,3,4,5,6,7,8,9,10,11,12\}$，若集合 $A$ 满足：① $A\subseteq U$；②若 $x\in A$，必有 $2x\notin A$；③若 $x\in\bar A$，必有 $2x\notin\bar A$，则这样的集合 $A$ 有\blankline[4.4em]个.

\ifanswer
\solbegin
\step{设集合 $P_n=\{1,2,\cdots,n\}$,$n\in\N^*$，}
\step{任取 $x\in P_n$，将 $x$ 除以 2，若商仍为偶数，再除以 2，$\cdots$，}
\step{经过 $k$ 次以后，商必为奇数，此时记商为 $m$，}
\step{于是 $x=m\cdot2^k$，其中 $m$ 为奇数，$k\in\N^*$.}
\step{由条件得，若 $m\in A$，则 $x\in A$ 等价于 $k$ 为偶数；}
\step{若 $m\notin A$，则 $x\in A$ 等价于 $k$ 为奇数；}
\step{于是 $x$ 是否属于 $A$ 由 $m$ 是否属于 $A$ 确定；}
\step{设 $Q_n$ 是 $P_n$ 中所有奇数的集合，}
\step{设 $f(n)$ 为集合 $P_n$ 中满足题干条件的集合 $A$ 的个数，}
\step{因此 $f(n)$ 等于 $Q_n$ 的子集个数.}
\step{当 $n$ 为偶数（或奇数）时，$P_n$ 中奇数的个数是 $\dfrac n2$（或 $\dfrac{n+1}2$），}
\step{所以 $f(n)=\begin{cases}2^{\frac n2},n\text{为偶数}\\[2pt] 2^{\frac{n+1}2},n\text{为奇数}\end{cases}$，所以 $P_{12}=\{1,2,3,4,5,6,7,8,9,10,11,12\}$，}
\step{即同时满足三个条件的集合 $A$ 的个数为 $2^{\frac{12}2}=2^6=64$.}
\solend
\fi

\item 设集合 $\Omega=\{(x_1,x_2,\cdots,x_{2020})\mid x_i\in\{0,1\},i=1,2,\cdots,2020\}$，对于 $\Omega$ 中的任意两个元素 $\alpha=(x_1,x_2,\cdots,x_{2020})$，$\beta=(y_1,y_2,\cdots,y_{2020})$，记
% 原卷此式第四组缺左括号、末尾连排两个右括号“))”，括号不配对；照录（见文件头“原卷笔误/存疑”②）
$\alpha*\beta=(x_1+y_1-x_1y_1)+(x_2+y_2-x_2y_2)+\cdots+x_{2020}+y_{2020}-x_{2020}y_{2020}))$，设 $u$、$v\in\Omega$，若 $u*u+v*v=2020$，则 $u*v$ 的最大值和最小值之和为\blankline[4.4em].

\ifanswer
\solbegin
\step{设 $u=(x_1,x_2,\cdots,x_{2020})$,$v=(y_1,y_2,\cdots,y_{2020})$，}
\step{因为 $x_i\in\{0,1\}$，所以 $x_i(x_i-1)=0$，}
\step{所以 $u*u=(x_1+x_1-x_1^2)+(x_2+x_2-x_2^2)+\cdots+(x_{2020}+x_{2020}-x_{2020}^2)$}
\step{$=x_1+x_2+\cdots+x_{2020}$，}
\step{$v*v=(y_1+y_1-y_1^2)+(y_2+y_2-y_2^2)+\cdots+(y_{2020}+y_{2020}-y_{2020}^2)$}
\step{$=y_1+y_2+\cdots+y_{2020}$，}
\step{由 $u*u+v*v=2020$，得 $x_1+x_2+\cdots+x_{2020}+y_1+y_2+\cdots+y_{2020}=2020$，}
% 原卷此行“$\cdots$”后脱逗号、 “2020”后衍一半角逗号（“2020,个 1”）；照录（见文件头“原卷笔误/存疑”③）
\step{所以 $x_1,x_2,\cdots x_{2020},y_1,y_2,\cdots y_{2020}$ 中有 2020,个 1，2020 个 0，}
\step{$u*v=(x_1+y_1-x_1y_1)+(x_2+y_2-x_2y_2)+\cdots+(x_{2020}+y_{2020}-x_{2020}y_{2020})$}
\step{$=2020-(x_1y_1+x_2y_2+\cdots+x_{2020}y_{2020})$，}
\step{注意到只有 $x_i=y_i=1$ 时，$x_iy_i=1$，否则 $x_iy_i=0$，}
% 此行“$\cdots$”后同样脱逗号，与上一处同式样；照录（见文件头“原卷笔误/存疑”③）
\step{而 $x_1,x_2,\cdots x_{2020},y_1,y_2,\cdots y_{2020}$ 中有 2020 个 1，2020 个 0，}
\step{所以满足 $x_iy_i=1$ 的最多有 1010 个，所以 $u*v\geqslant2020-1010=1010$，}
\step{满足 $x_iy_i=0$ 的最多有 2020 个，所以 $u*v\leqslant2020-0=2020$，}
\step{则 $u*v$ 的最大值和最小值之和为 $2020+1010=3030$.}
\solend
\fi

\end{enumerate}

\bigsec{二、选择题}

\begin{enumerate}
\setcounter{enumi}{12}

\item 实数 $x$ 与 $y$，“$x+y>2$”的\mbox{（\quad）}条件是“$x>1$ 或 $y>1$”.

\keepwithprev
A. 充分非必要\qquad B. 必要非充分

C. 充要\qquad D. 既非充分又非必要

\ifanswer
\solbegin
\step{当 $x>1$ 或 $y>1$ 时，$x+y>2$ 不一定成立，充分性不成立；}
\step{当 $x+y>2$ 时，$x>1$ 或 $y>1$ 一定成立，必要性成立. 故选 $B$.}
\solend
\fi

\item 设 $a,b\in\R$，且 $a<0$，$-1<b<0$，那么下列不等式成立的是\mbox{（\quad）}

\keepwithprev
A. $a>ab>ab^2$\qquad B. $ab^2>ab>a$\qquad C. $ab>a>ab^2$\qquad D. $ab>ab^2>a$

\ifanswer
\solbegin
\step{因为 $-1<b<0$，$a<0$，所以 $ab>0$，$0<b^2<1$.}
\step{所以 $ab-ab^2=ab(1-b)>0$，$ab^2-a=a(b^2-1)>0$，}
\step{所以 $ab>ab^2>a$. 故选 $D$.}
\solend
\fi

% 第 15 题答案原卷印在题干括号内（“(  D  )”），且无【答案】/【解析】段；
% 按 高一37 第 13 题、高一63 第 14 题成例括号留空、答案另提 \ans{$D$}
\item 用反证法证明：“平面四边形中至少有一个内角不超过 90°”时，下列“假设”中正确的是\mbox{（\quad）}

\keepwithprev
A. 假设有两个内角超过 90°\qquad B. 假设有三个内角超过 90°

C. 假设至多有两个内角超过 90°\qquad D. 假设四个内角都超过 90°

\ifanswer
\ans{$D$}
\fi

\item 设 $A$ 为非空实数集，集合 $B=\{xy\mid x\in A,y\in A,x\neq y\}$，则对于以下两个结论说法正确的是\mbox{（\quad）}

结论①：若 $A$ 是五元集，则 $B$ 中的元素最多 10 个；

结论②：存在四元集 $A$，使得 $B=\{2,4,5,8,10,16\}$.

\keepwithprev
A. ①对②对\qquad B. ①对②错\qquad C. ①错②对\qquad D. ①错②错

\ifanswer
\solbegin
\step{对于①，想让元素个数达到最大值 10，则须避免产生符号不同的乘积，}
\step{否则正负分开计数容易导致总数受限，且容易产生相等值，}
\step{也最好避免 0 的出现，}
\step{故设 $A=\{a_1,a_2,a_3,a_4,a_5\}$，$0<a_1<a_2<a_3<a_4<a_5$.}
% 原卷此链“$a_3a_3$”同足标自乘、“$a_{14}a_5$”下标 14 均与文意不合，疑为 $a_2a_4$、$a_4a_5$ 之误；照录（见文件头“原卷笔误/存疑”④）
\step{因为 $a_1a_2<a_1a_3<a_1a_4<a_1a_5<a_2a_3<a_3a_3<a_{14}a_5$，$a_2a_3<a_2a_4<a_3a_4$，}
\step{所以 $B$ 中元素个数小于等于 10 个，}
\step{如设 $A=\{1,2,3,5,7\}$，则 $B=\{2,3,5,6,7,10,14,15,21,35\}$，}
\step{所以集合 $B$ 中元素个数的最大值为 10 个，故①正确；}
\step{对于②，由于集合 $B$ 中没有 0，所以集合 $A$ 中没有 0；}
\step{设集合 $A=\{a,b,c,d\}$，且 $a<b<c<d$，}
\step{当 $a<0<b<c<d$ 时，$a\times b<0$，$a\times c<0$，$a\times d<0$，}
\step{故集合 $B$ 至少有 3 个负数，与 $B=\{2,4,5,8,10,16\}$ 矛盾；}
\step{当 $a<b<0<c<d$ 时，$a\times b>0$，$c\times d>0$，}
\step{$a\times c<0$，$a\times d<0$，$b\times c<0$，$b\times d<0$，}
\step{故集合 $B$ 至少有 4 个负数，与 $B=\{2,4,5,8,10,16\}$ 矛盾；}
\step{当 $a<b<c<0<d$ 时，$a\times b>0$，$a\times c>0$，$b\times c>0$，}
\step{$a\times d<0$，$b\times d<0$，$c\times d<0$，}
\step{故集合 $B$ 至少有 3 个负数，与 $B=\{2,4,5,8,10,16\}$ 矛盾；}
\step{当 $0<a<b<c<d$ 时，则集合 $B=\{ab,ac,ad,bc,bd,cd\}$，}
\step{则必有 $ab=2$，$cd=16$，其 4 个正实数的乘积 $abcd=32$，}
\step{又 $ad=5$，$bc=8$ 或 $ad=8$，$bc=5$，其 4 个正实数的乘积 $abcd=40$，矛盾；}
\step{当 $a<b<c<d<0$ 时，与 $0<a<b<c<d$ 时的证明方法相同，}
\step{得集合 $A$ 也不存在；}
\step{综上不存在四元集 $A$，使其生成集 $B=\{2,4,5,8,10,16\}$，故②错误.}
\step{故选 $B$.}
\solend
\fi

\end{enumerate}

\bigsec{三、解答题}

\begin{enumerate}
\setcounter{enumi}{16}

\item 解下列不等式（组）

\begin{enumerate}[label=(\arabic*)]
\item 解不等式：$(ax+1)(x+1)<0(a\in\R)$.
\item 解不等式组：$\begin{cases}\dfrac{2x+4}{x+4}\geqslant1\\[4pt] (x+1)(2-|x|)\geqslant0\end{cases}$.
\end{enumerate}

\ifanswer
\solbegin
\stepm{(1)}{当 $a=0$ 时，原不等式为 $x+1<0$，即 $x<-1$，}
\step{所以此时不等式的解集为 $\{x\mid x<-1\}$；}
\step{当 $a\neq0$ 时，令 $(ax+1)(x+1)=0$，解得 $x_1=-\dfrac1a$,$x_2=-1$，}
\step{当 $a>1$ 时，$-\dfrac1a>-1$，此时不等式的解集为 $\left\{x\mid-1<x<-\dfrac1a\right\}$；}
% 以下两行“此时”前原卷均无标点；照录（见文件头“原卷笔误/存疑”⑥）
\step{当 $a=1$ 时，原不等式为 $(x+1)^2<0$ 此时不等式的解集为 $\varnothing$；}
\step{当 $0<a<1$ 时，$-\dfrac1a<-1$ 此时不等式的解集为 $\left\{x\mid-\dfrac1a<x<-1\right\}$；}
\step{当 $a<0$ 时，$-\dfrac1a>-1$，此时不等式的解集为 $\left\{x\mid x<-1\ \text{或}\ x>-\dfrac1a\right\}$.}
\step{综上所述，当 $a=0$ 时，$\{x\mid x<-1\}$；}
\step{当 $a>1$ 时，$\left\{x\mid-1<x<-\dfrac1a\right\}$；}
\step{当 $a=1$ 时，$\varnothing$；}
\step{当 $0<a<1$ 时，$\left\{x\mid-\dfrac1a<x<-1\right\}$；}
\step{当 $a<0$ 时，$\left\{x\mid x<-1\ \text{或}\ x>-\dfrac1a\right\}$.}
\stepm{(2)}{不等式 $\dfrac{2x+4}{x+4}\geqslant1$ 可化为 $\dfrac{2x+4}{x+4}-1\geqslant0$，即 $\dfrac x{x+4}\geqslant0$，}
\step{等价于 $\begin{cases}x(x+4)\geqslant0\\ x+4\neq0\end{cases}$，解得 $x<-4$ 或 $x\geqslant0$，}
\step{不等式 $(x+1)(2-|x|)\geqslant0$ 等价于 $\begin{cases}(x+1)(2-x)\geqslant0,x\geqslant0\\ (x+1)(2+x)\geqslant0,x<0\end{cases}$，}
\step{解不等式 $(x+1)(2-x)\geqslant0$，$x\geqslant0$，得 $0\leqslant x\leqslant2$，}
\step{解不等式 $(x+1)(2+x)\geqslant0$，$x<0$，得 $x\leqslant-2$ 或 $-1\leqslant x<0$，}
\step{所以 $(x+1)(2-|x|)\geqslant0$ 的解集为 $\{x\mid x\leqslant-2\ \text{或}\ -1\leqslant x\leqslant2\}$，}
\step{综上所述，$\begin{cases}\dfrac{2x+4}{x+4}\geqslant1\\[4pt] (x+1)(2-|x|)\geqslant0\end{cases}$ 的解集为 $\{x\mid x<-4\ \text{或}\ 0\leqslant x\leqslant2\}$.}
\solend
\fi

\ansspace[8]

% 第 18 题题干＋两小问与作答区 6 行须同页
\ansneed{20cm}
\item 已知 $a$ 为实数，设 $f(x)=x^2+2(a+1)x+1$.

\begin{enumerate}[label=(\arabic*)]
\item 已知对于任意实数 $x$，$\sqrt{f(x)}$ 均有意义，求实数 $a$ 的取值范围；
\item 设 $A=\{x\mid f(x)+a^2-2=0\}$，$B=\{-4,0\}$，$A\cup B=B$，求实数 $a$ 的取值范围.
\end{enumerate}

\ifanswer
\solbegin
\stepm{(1)}{$\sqrt{f(x)}$ 对任意实数 $x$ 都有意义，}
\step{即 $f(x)=x^2+2(a+1)x+1\geqslant0$ 对任意 $x\in\R$ 恒成立，}
\step{所以 $\Delta=4(a+1)^2-4\times1\times1\leqslant0$，即 $(a+1)^2\leqslant1$，解得 $-2\leqslant a\leqslant0$，}
\step{故 $a$ 的取值范围是 $[-2,0]$.}
\stepm{(2)}{由 $f(x)+a^2-2=0$ 得 $x^2+2(a+1)x+a^2-1=0$.}
\step{由 $A\cup B=B$ 得 $A\subseteq B=\{-4,0\}$，即该方程无实根，}
\step{或它的实根都在 $\{-4,0\}$ 中.}
\step{当方程无实根时，$\Delta=4(a+1)^2-4(a^2-1)=8a+8<0$，解得 $a<-1$；}
\step{当方程有实根时，$a\geqslant-1$，设两根为 $x_1$、$x_2$，}
\step{由根与系数的关系得 $x_1+x_2=-2(a+1)$、$x_1x_2=a^2-1$，}
\step{此时两根只能在 $\{-4,0\}$ 中取值；}
\step{若两根为 0 与 0，由 $x_1x_2=0$ 得 $a=\pm1$，}
\step{再由 $x_1+x_2=0$ 得 $a=-1$；}
\step{若两根为 0 与 $-4$，由 $x_1x_2=0$ 得 $a=\pm1$，再由 $x_1+x_2=-4$ 得 $a=1$.}
\step{经检验，$a=-1$ 时 $A=\{0\}$，$a=1$ 时 $A=\{-4,0\}$，都满足 $A\subseteq B$.}
\step{综上，实数 $a$ 的取值范围为 $(-\infty,-1]\cup\{1\}$.}
\solend
\fi

\ansspace[8]

\ansneed{18cm}
\item 共享单车给市民出行带来了诸多便利，某公司购买了一批单车投放到某地给居民使用，据市场分析，这批单车的运营累计利润 $y$（单位：元）与运营天数 $x$（单位：天，$x$ 为正整数）满足关系：$y=-\dfrac12x^2+60x-800$.

\begin{enumerate}[label=(\arabic*)]
\item 要使这批单车的运营累计利润高于 800 元，求运营天数的取值范围；
\item 判断这批单车的日平均运营利润能否超过 20 元，并说明理由.
\end{enumerate}

\ifanswer
\solbegin
\stepm{(1)}{由题意得 $y=-\dfrac12x^2+60x-800>800$，}
\step{即 $x^2-120x+3200<0$，解得 $40<x<80$，}
\step{即运营天数的取值范围为 $\{x\mid40<x<80,x\in\N^*\}$.}
\stepm{(2)}{不能，理由如下：}
\step{这批单车的日平均运营利润为 $\dfrac yx=-\dfrac12x-\dfrac{800}x+60=-\dfrac12\left(x+\dfrac{1600}x\right)+60$，}
\step{因 $x>0$，则 $x+\dfrac{1600}x\geqslant2\sqrt{x\cdot\dfrac{1600}x}=80$，}
\step{当且仅当 $x=\dfrac{1600}x$，即 $x=40$ 时取等号，}
\step{此时 $\dfrac yx=-\dfrac12\left(x+\dfrac{1600}x\right)+60\leqslant-\dfrac12\times80+60=20$，}
\step{故这批单车的日平均运营利润不能超过 20 元.}
\solend
\fi

\ansspace[8]

% 第 20 题题干＋三小问与作答区 8 行须同页
\ansneed{20cm}
\item 设 $f(x)=ax^2+bx+c$($a\neq0$，$b$、$c\in\R$).

\begin{enumerate}[label=(\arabic*)]
\item 若 $a=b=c=1$，设 $A=\{(x,y)\mid y=f(x)\}$，$B=\{(x,y)\mid y=1\}$，用列举法表示 $A\cap B$；
% 原卷两个“若”连用（“若 a=1，…，若方程 f(x)=0 有两个大于 1 的根”），疑衍其一；照录（见文件头“原卷笔误/存疑”⑦）
\item 若 $a=1$，$b=2k-1$，$c=k^2$，若方程 $f(x)=0$ 有两个大于 1 的根，求实数 $k$ 的取值范围；
\item 若 $a$、$b$、$c\in\Z$，且 $f(0)$、$f(1)$ 均为奇数，求证：方程 $f(x)=0$ 无整数解.
\end{enumerate}

\ifanswer
\solbegin
\stepm{(1)}{当 $a=b=c=1$ 时，$f(x)=x^2+x+1$，}
\step{所以 $A=\{(x,y)\mid y=x^2+x+1\}$，又 $y=1$，}
\step{由 $x^2+x+1=1$ 得 $x^2+x=0$，解得 $x=0$ 或 $x=-1$，}
\step{所以 $A\cap B=\{(0,1),(-1,1)\}$.}
\stepm{(2)}{当 $a=1$、$b=2k-1$、$c=k^2$ 时，$f(x)=x^2+(2k-1)x+k^2$.}
\step{设两根为 $x_1$、$x_2$，}
\step{则“两根都大于 1”等价于 $\Delta>0$、$f(1)>0$、$-\dfrac b{2a}>1$.}
\step{由 $\Delta=(2k-1)^2-4k^2=1-4k>0$ 得 $k<\dfrac14$；}
\step{由 $f(1)=1+(2k-1)+k^2=k^2+2k=k(k+2)>0$ 得 $k<-2$ 或 $k>0$；}
\step{由 $-\dfrac{2k-1}2=\dfrac{1-2k}2>1$ 得 $k<-\dfrac12$.}
\step{三者取交集得 $k<-2$，即 $k$ 的取值范围是 $(-\infty,-2)$.}
\stepm{(3)}{设 $x_0\in\Z$ 是方程 $f(x)=0$ 的整数解.}
\step{因为 $x_0(x_0-1)$ 是两个相邻整数之积，恒为偶数，所以 $ax_0(x_0-1)$ 为偶数.}
\step{而 $f(x_0)=ax_0^2+bx_0+c=ax_0(x_0-1)+(a+b)x_0+c$，}
\step{所以 $f(x_0)$ 与 $(a+b)x_0+c$ 同奇偶.}
\step{若 $x_0$ 为偶数，则 $(a+b)x_0$ 为偶数，$f(x_0)$ 与 $c=f(0)$ 同奇偶，为奇数；}
\step{若 $x_0$ 为奇数，则 $(a+b)x_0$ 与 $a+b$ 同奇偶，}
\step{$f(x_0)$ 与 $(a+b)+c=f(1)$ 同奇偶，也为奇数.}
\step{综上，$f(x_0)$ 恒为奇数，不可能等于 0，与 $f(x_0)=0$ 矛盾，}
\step{所以方程 $f(x)=0$ 无整数解.}
\solend
\fi

\ansspace[8]

% 压轴题（题干＋三小问）与作答区 10 行须同页
\ansneed{20cm}
\item 已知 $X$ 为有限集，且任取 $x\in X$，均有 $x\in\Z$，称 $X$ 为非空有限整数集. 定义：$X_k=\{x+k\mid x\in X\}$，$X^2=\{y\mid y=x^2,x\in X\}$.

\begin{enumerate}[label=(\arabic*)]
\item 若 $X=\{-2,2\}$，分别写出 $X_2$ 和 $X^2$（用列举法表示）；
\item 若 $X=\{-1,0,1\}$，$X_1\subseteq Y$，$(Y_{-1})^2\subseteq X$，求集合 $Y$ 中所有的元素；
\item 对于非空有限整数集 $A$、$B$，若 $A_k\subseteq B$，$(B^2)_{-2}\subseteq A$，$B$ 中的元素均为非负整数，$k$ 为整数，求证：

① $A$、$B$ 中元素的个数相同；

② 若集合 $A$ 中的元素个数为奇数，则 $k$ 为偶数.
\end{enumerate}

\ifanswer
\solbegin
\stepm{(1)}{由题意得 $X_2=\{0,4\}$，$X^2=\{4\}$；}
\stepm{(2)}{由题意得 $X_1=\{0,1,2\}$，即 $\{0,1,2\}\subseteq Y$，$Y_{-1}=\{y-1\mid y\in Y\}$，}
\step{$(Y_{-1})^2=\{(y-1)^2\mid y\in Y\}\subseteq X=\{-1,0,1\}$，}
\step{所以 $(y-1)^2=0$ 或 $(y-1)^2=1$，解得 $y=1$ 或 $y=2$ 或 $y=0$，}
\step{即 $y\in\{0,1,2\}$，即 $Y\subseteq\{0,1,2\}$，所以 $Y=\{0,1,2\}$；}
% 原卷 (3) 解析首行“(3)”与“①”连排，本稿合并记作 \stepm{(3)①}
\stepm{(3)①}{因为 $A_k=\{x+k\mid x\in A\}$，函数 $g(x)=x+k$ 在全体整数集上严格增，}
\step{所以 $\forall x_1$、$x_2\in A$，$x_1\neq x_2$，$g(x_1)\neq g(x_2)$，所以 $|A_k|=|A|$，}
\step{因为 $A_k\subseteq B$，所以 $|A|\leqslant|B|$，}
\step{设 $b\in B$，$B$ 中元素均为非负整数，$\forall b_1\neq b_2\in B$，$b_1^2\neq b_2^2$，}
\step{故函数 $f(b)=b^2-2$ 在 $B$ 上单射，所以 $|(B^2)_{-2}|=|B|$，}
\step{又因为 $(B^2)_{-2}\subseteq A$，所以 $|B|\leqslant|A|$，所以 $|A|=|B|$.}
\stepm{②}{由①得 $A_k\subseteq B$ 且 $|A_k|=|A|=|B|$，$(B^2)_{-2}\subseteq A$ 且 $|(B^2)_{-2}|=|B|=|A|$，}
\step{所以 $A_k=B$，$(B^2)_{-2}=A$，}
\step{所以 $B=A_k=\{a+k\mid a\in A\}$，$A=(B^2)_{-2}=\{b^2-2\mid b\in B\}$，}
\step{所以 $A=\{(a+k)^2-2\mid a\in A\}$，}
\step{所以映射 $h(a)=(a+k)^2-2$ 是集合 $A$ 到自身的一个双射.}
\step{假设 $k$ 是奇数，当 $a$ 为偶数时，$a+k$ 为奇数，$(a+k)^2-2$ 为奇数，}
\step{当 $a$ 为奇数时，$a+k$ 为偶数，$(a+k)^2-2$ 为偶数，}
\step{所以此时映射 $h(a)=(a+k)^2-2$，将集合 $A$ 中的奇数映射为偶数，}
\step{将集合 $A$ 中的偶数映射为奇数，即集合 $A$ 中的奇数元素和偶数元素必须配对，}
\step{此时集合 $A$ 的元素个数则为偶数，}
\step{与题设中“集合 $A$ 中的元素个数为奇数”矛盾.}
\step{故假设不成立，所以 $k$ 是偶数.}
\solend
\fi

\ansspace*[10]

\end{enumerate}

\end{document}
